\documentclass[11pt]{article}
\usepackage[T1]{fontenc}
\usepackage[utf8]{inputenc}
\usepackage{newtxtext,newtxmath}

\setcounter{tocdepth}{1}
\usepackage[a4paper,margin=35mm]{geometry}
\linespread{0.90}
\setlength{\parindent}{0pt}
\setlength{\parskip}{6pt}
\usepackage{float}
\usepackage{cmap}
\usepackage{amsfonts}
\let\openbox\relax
\usepackage{amsthm}
\usepackage{amsmath}
\usepackage{mathtools}
\allowdisplaybreaks
\usepackage{bm}
\usepackage{graphicx}
\usepackage{tikz}
\usepackage{pgfplots}
\pgfplotsset{compat=1.18}

\hyphenation{superselection}
\usepackage{microtype}
\usepackage[numbers,sort&compress]{natbib}
\usepackage{xcolor}
\usepackage{booktabs}
\usepackage{enumitem}
\usepackage{array}
\usepackage{tabularx}
\usepackage{setspace}
\usepackage[explicit]{titlesec}
\usepackage[labelfont=bf]{caption}
\DeclareCaptionFont{captionthirtysmaller}{\fontsize{7.7}{9.2}\selectfont}
\captionsetup[figure]{font=captionthirtysmaller}
\captionsetup[table]{font=captionthirtysmaller}
\usepackage{mdframed}
\usepackage{needspace}
\usepackage{tocloft}
\usepackage{placeins}
\newcolumntype{Y}{>{\raggedright\arraybackslash}X}
\setlength{\tabcolsep}{4.5pt}
\renewcommand{\arraystretch}{1.1}
\renewcommand{\qedsymbol}{}
\AtBeginEnvironment{tabular}{\small}
\AtBeginEnvironment{tabularx}{\small}

\numberwithin{equation}{section}
\mathtoolsset{showonlyrefs=false}
\DeclarePairedDelimiter{\abs}{\lvert}{\rvert}
\DeclarePairedDelimiter{\norm}{\lVert}{\rVert}
\DeclareMathOperator{\Tr}{Tr}
\DeclareMathOperator{\diag}{diag}
\DeclareMathOperator{\rank}{rank}
\DeclareMathOperator{\sr}{sr}
\DeclareMathOperator{\logdet}{logdet}
\DeclareMathOperator{\pdet}{pdet}
\DeclareMathOperator{\supp}{supp}
\newcommand{\Sym}{\mathrm{Sym}}
\newcommand{\dd}{\mathrm{d}}
\newcommand{\ii}{\mathrm{i}}
\newcommand{\SPD}{\mathrm{SPD}}
\newcommand{\eps}{\varepsilon}
\newcommand{\Loew}{\mathcal{L}}

\setlist{nosep,leftmargin=1.25em}

\titleformat{\section}
{\normalfont\normalsize\bfseries\raggedright}
{\thesection}{0.75em}{#1}
[\vspace{0.12em}\titlerule]

\newcommand{\subaccent}{\vspace{0.01em}\noindent\rule{7ex}{0.2pt}}
\titleformat{\subsection}
{\normalfont\normalsize\bfseries}{\thesubsection}{0.5em}{#1}
[\subaccent]
\titlespacing*{\subsection}{0pt}{2.2ex plus 0.6ex}{1.5ex}

\titleformat{\subsubsection}
{\normalfont\normalsize\bfseries}{\thesubsubsection}{0.5em}{#1}
\titlespacing*{\subsubsection}{0pt}{1.6ex plus 0.4ex}{0.6ex}

\usepackage[
  pdfusetitle,
  hidelinks,
  bookmarksopen=true,
  bookmarksnumbered=true
]{hyperref}

\newtheoremstyle{thmcompact}
{0.5ex}{0.5ex}
{\itshape}
{0pt}
{\bfseries}
{.}
{0.6em}
{\thmname{#1}\ \thmnumber{#2}\ \thmnote{(\textit{#3})}}

\theoremstyle{thmcompact}
\newtheorem{theorem}{Theorem}[section]
\newtheorem{proposition}[theorem]{Proposition}
\newtheorem{lemma}[theorem]{Lemma}
\newtheorem{corollary}[theorem]{Corollary}
\newtheorem{definition}[theorem]{Definition}

\theoremstyle{remark}
\newtheorem*{remarkplain}{Remark}

\mdfdefinestyle{thmframe}{%
	leftline=true, rightline=false, topline=false, bottomline=false,
	linewidth=0.8pt, linecolor=black,
	innerleftmargin=16pt, innerrightmargin=8pt,
	innertopmargin=6pt, innerbottommargin=6pt,
	skipabove=6pt, skipbelow=6pt,
	nobreak=true
}
\surroundwithmdframed[style=thmframe]{theorem}
\surroundwithmdframed[style=thmframe]{proposition}
\surroundwithmdframed[style=thmframe]{lemma}
\surroundwithmdframed[style=thmframe]{corollary}
\surroundwithmdframed[style=thmframe]{definition}

\newmdenv[
linewidth=0.6pt,
topline=false,
bottomline=false,
rightline=false,
skipabove=\baselineskip,
skipbelow=\baselineskip,
backgroundcolor=gray!3,
leftmargin=0pt,
innerleftmargin=8pt,
innerrightmargin=6pt,
frametitleaboveskip=0.4em,
frametitlebelowskip=0.2em,
frametitlefont=\bfseries\small,
]{remarkbox}

\mdfdefinestyle{eqbox}{%
	leftline=true, rightline=true, topline=true, bottomline=true,
	linewidth=0.8pt, linecolor=black,
	innerleftmargin=14pt, innerrightmargin=10pt,
	innertopmargin=8pt, innerbottommargin=8pt,
	skipabove=8pt, skipbelow=8pt,
	backgroundcolor=gray!3,
	nobreak=true
}
\usepackage{titling}
\newenvironment{remark}[1][]%
{\begin{remarkbox}\begin{remarkplain}[#1]\normalfont}
		{\end{remarkplain}\end{remarkbox}}

\renewcommand{\contentsname}{Contents}

\titleformat{name=\section,numberless}
{\normalfont\normalsize\bfseries\raggedright}
{}{0em}{#1}
[\vspace{0.12em}\titlerule]
\makeatletter
\renewcommand{\l@section}[2]{\@dottedtocline{1}{0em}{2.3em}{\raggedright #1}{#2}}
\renewcommand{\l@subsection}[2]{\@dottedtocline{2}{1.8em}{3.2em}{\raggedright #1}{#2}}
\makeatother

\emergencystretch=3em
\hfuzz=1pt
\tolerance=1800
\hbadness=2500

\newcommand{\thmneed}{\Needspace{6\baselineskip}}
\AtBeginEnvironment{theorem}{\thmneed}
\AtBeginEnvironment{proposition}{\thmneed}
\AtBeginEnvironment{lemma}{\thmneed}
\AtBeginEnvironment{corollary}{\thmneed}

\AtBeginEnvironment{thebibliography}{\small}

\widowpenalty=10000
\clubpenalty=10000

\title{\textbf{Quotient Observation}\\[0.2em]
{\large Hidden Locality, Schur Mediation, and Visible Precision Geometry}\\[0.2em]
{\large \textbf{DRAFT PREPRINT}}}
\author{J.\,R.~Dunkley}
\date{21st March 2026}

\hypersetup{
  pdftitle={Quotient Observation: Hidden Locality, Schur Mediation, and Visible Precision Geometry},
  pdfauthor={J. R. Dunkley},
  pdfsubject={Nonequilibrium Markov chains, quotient observation, visible precision geometry}
}

\begin{document}
\maketitle

\begin{abstract}
We study what can be recovered under partial observation and argue that the intrinsic visible object is quotient-visible precision, denoted Phi, rather than the Donsker-Varadhan Hessian itself. The paper develops a finite, compositional theory of this visible object on the positive cone of precision operators. We prove that quotient observation composes along towers of observation, identify the visible algebraic descendant of the bridge developed in Finite, and show that in the scalar visible case the resulting structure closes into a closed system of precision, clock, transport, and conservation laws. This yields a picture of hidden mediation, including support rigidity, a hidden-load parametrisation, a unique additive interior clock, and a sharp distinction between quotient observation and naive compression. We then turn to nonequilibrium systems and analyse how the visible object is realised in practice. At a symmetry point we obtain a canonical hidden response carrier and a scalar collapse. Away from that point we identify a fast-memory response layer with a two-state homogenised limit, a first-correction Schur carrier, and a source law with an explicit finite-epsilon obstruction. In a separate strong-bias corridor regime we derive a different asymptotic sector and show that these regimes are both real and not presently unified by a single higher operator architecture. The theory is tested against a sequence of benchmark systems used to keep the developing framework tied to a real signal while algebraic closure was being pinned down. These include kinesin, proofreading, transformed semi-Markov, and hidden-feedback ladder tests. Together they support the central claim that quotient-visible precision is the renormalisation-natural visible invariant, while also making the present transfer boundaries explicit. The result is a finished quotient-visible core theory together with a sharply delimited correction-layer package, a falsification battery, and a clear map of what remains open.
\end{abstract}

\newpage

\tableofcontents

\newpage

\section*{Foreword}
This work records an ongoing research programme whose full development is not yet complete. The surrounding theory has grown substantially in scope and density, and a proper treatment would require a much longer document than is presently practical.

We therefore release the current manuscript in the interest of dissemination. The main structures are already in view, much of the territory has been mapped, and the work has been subjected throughout to direct falsification attempts, engagement with real data, and operational testing. Thus far, these checks have supported the theory and, in doing so, have opened further lines of investigation.

Readers primarily interested in the extracted geometry may consult \href{https://doi.org/10.5281/zenodo.19361501}{\emph{Geometric Observation}}, where the geometric core toward which this paper moves is presented in a shorter and more concentrated form.

The present manuscript will continue to be refined, but we expect its main structure to remain stable, with finalisation in the coming weeks.
\newpage

\section{Introduction}
\label{sec:01_introduction}

\subsection{The observation problem}

Given that you only see part of a system, what can you recover about its internal structure?

A Markov process runs on $n$ states; an observer sees only the image under a linear surjection $C : \mathbb{R}^n \to \mathbb{R}^m$ with $m < n$. The question is, what quadratic structure on the visible sector is canonically determined by the full precision and the observation map, with no further choices?

Here we push the full precision through the observation constraint and extract the unique visible quadratic form that records the minimum energy required to realise each visible displacement. The result is a definite positive operator on the visible sector, the quotient-visible precision $\Phi$, and it is this operator that is the natural visible invariant.

The setting is the Donsker-Varadhan large-deviation geometry of empirical flows in finite-state Markov chains, \cite{DonskerVaradhan1975a,DonskerVaradhan1975b,DonskerVaradhan1976,Ellis1985,Touchette2009,BertiniFaggionatoGabrielli2015} following our earlier work \cite{Dunkley2026FiniteObservation} which established a finite quadratic bridge between the stationary backbone and the nonequilibrium Hessian. It also sits naturally alongside the network and stochastic-thermodynamic description of nonequilibrium observables \cite{Schnakenberg1976,LebowitzSpohn1999,Seifert2012,Esposito2012,Maes2020}. Now we ask, what does that bridge produce once hidden degrees of freedom are eliminated by quotient observation?

\subsection{Main claim and conceptual picture}

The central claim is that $\Phi$ is the intrinsic visible variable. It is defined variationally, it composes along towers of observation, it admits a stable first-order calculus, and it is provably different from compression. No other candidate satisfies all properties simultaneously:

\[
\begingroup
\setlength{\fboxsep}{8pt}%
\setlength{\fboxrule}{1.15pt}%
\boxed{\scalebox{1.08}{$\displaystyle \Phi_T^{\circ}=\left\{X=T^{1/2}\Pi T^{1/2}:\; 0\prec \Pi\le I\ \text{ on }\ S=\operatorname{Ran}(T)\right\}$}}%
\endgroup
\]

\[
\begingroup
\setlength{\fboxsep}{8pt}%
\setlength{\fboxrule}{1.15pt}%
\boxed{\scalebox{1.08}{$\displaystyle \overline{\Phi}_T=\left\{X=T^{1/2}\Pi T^{1/2}:\; 0\le \Pi\le I\ \text{ on }\ S=\operatorname{Ran}(T)\right\}=T^{1/2}[0,I]_S T^{1/2}$}}%
\endgroup
\]

The Donsker-Varadhan bridge \cite{Dunkley2026FiniteObservation} identifies a finite visible algebraic object, the tangent descendant of $\Phi$, but it is $\Phi$ itself that closes under observation. The Hessian is a bridge object: it matters because it leads to $\Phi$. More precisely, once the tangent ceiling has been fixed on the part of visible space where it acts nontrivially, the exact observed class consists of all visible laws obtained by contracting within that same active sector, without ever creating new visible directions. Finite hidden sectors realise the strict interior of that class, while allowing limits adds the boundary cases. Section~\ref{sec:04_hidden_mediation} makes this precise.

The paper contains two species of results. The first is an algebraic theory of the visible object: its variational characterisation, composition law, conservation structure, hidden-load parametrisation, and graph-local decay. These hold for any quotient observation on any positive precision operator.

The second evaluates this theory in specific nonequilibrium systems: symmetry theorems, fast-memory homogenisation, corridor asymptotics, and thermodynamic constraints. The boundary between the two is maintained throughout.

\subsection{Main results}

The paper establishes the following results, listed in the order they appear.

\begin{enumerate}[label=(\roman*)]
\item Quotient observation on the positive cone is variational, composes along observation towers, and produces a coisometric tangent transport with monotone Ky Fan detectability envelopes.

\item The Donsker-Varadhan bridge produces a finite visible algebraic object via quotient reduction. The first genuinely new quotient effect beyond the tangent law is quartic and sign-definite; smooth divergence functionals built from the visible shadow lose only at sextic order.

\item In scalar visible rank, the visible precision admits a generator-level Schur simplification and does not depend on the hidden completion at the matched visible-law level.

\item The visible precision splits as $\Phi = F_0 + W$, with $F_0$ the reversible backbone contribution and $W$ the irreversible correction. This conservation identity holds without remainder.

\item Hidden mediation is parametrised by a positive hidden-load operator on the tangent support. Sequential transport, a unique additive determinant clock, and support rigidity follow.

\item Graph-local decay of hidden contributions: the visible effect of a hidden perturbation decays with the graph distance between the perturbation site and the visible sector.

\item On the Bell square, common-cause visible Gaussian gluing is a law-level compatibility problem whose frontier has two independent coordinates.

\item At the symmetric point of the two-state visible ladder, a symmetry theorem identifies a canonical hidden response carrier, and the visible algebra collapses to one free parameter: $F_0 = \alpha^2/(T+\alpha)$ and $W = \alpha T/(T+\alpha)$.

\item Fast-memory homogenisation yields a source law and an explicit obstruction to hidden-structure recovery.

\item Corridor asymptotics separate three regimes (fast memory, intermediate, slow memory) with sharp crossover diagnostics.

\item A thermodynamic ceiling constrains visible precision: $\Phi \leq F_0 + \sigma/2$, where $\sigma$ is the entropy production rate.

\item Timescale duality, the EPR-Fisher bound, and tightness at both timescale limits constrain the visible object within a closed region of parameter space.
\end{enumerate}

\subsection{Roadmap}

Section~\ref{sec:quotient-precision} formulates quotient observation on the positive cone, proves the variational characterisation and the Fr\'echet calculus, establishes tower composition, identifies the Schur gap, and names $\Phi$ as the intrinsic visible invariant. Section~\ref{sec:03_dv_bridge} connects the quotient theory to the Donsker-Varadhan bridge from \cite{Dunkley2026FiniteObservation}, establishes the visible algebraic object, and proves the quartic defect and sextic robustness theorems.

Sections~\ref{sec:04_hidden_mediation}-\ref{sec:06_bell_frontier} develop the internal structure of $\Phi$: hidden mediation and the conservation split, hidden locality and graph-local decay, and the Bell-square gluing frontier. Sections~\ref{sec:07_exact_symmetry}-\ref{sec:09_corridor_regime} evaluate the algebraic theory in nonequilibrium systems: symmetry at the symmetric point, fast-memory homogenisation with source law and obstruction, and corridor asymptotics with regime separation.

Section~\ref{sec:10_benchmark_synthesis} synthesises the benchmark evidence, Section~\ref{sec:11_falsification} presents falsification diagnostics, Section~\ref{sec:12_scope} states the proved boundary, and Section~\ref{sec:13_discussion} discusses what remains open and what geometric structure lies at the frontier.

\subsection{Notation and glossary}

The \emph{upstairs} space $\mathbb{R}^n$ carries the full precision operator $H$; the \emph{downstairs} space $\mathbb{R}^m$ carries the visible precision $\Phi$. An \emph{observation} is a surjective linear map $C : \mathbb{R}^n \to \mathbb{R}^m$. An \emph{observation tower} is a chain of surjections $\mathbb{R}^n \to \mathbb{R}^{m_1} \to \mathbb{R}^{m_2}$; the composition theorem says that observing once or in two steps gives the same visible precision.

The \emph{backbone} $H_0$ is the time-symmetric part of the full precision. The \emph{irreversible correction} $W$ is the contribution from detailed-balance-breaking currents. The \emph{conservation split} is $\Phi = F_0 + W$, where $F_0 := \Phi_C(H_0)$ is the quotient-visible backbone.

We write $\Sym_{++}(n)$ for the cone of $n \times n$ symmetric positive definite matrices, $\Phi_C(H)$ or $F_C(H)$ for the quotient map, and $\succeq$ for the Loewner partial order.


\section{Quotient observation and the intrinsic visible object}
\label{sec:quotient-precision}

For a positive precision operator $H \in \Sym_{++}(n)$ and a surjective observation map $C \in \mathbb{R}^{m \times n}$, the visible operator is
\begin{equation}
\Phi_C(H)
:=
F_C(H)
:=
\bigl(C H^{-1} C^{\mathsf T}\bigr)^{-1}.
\label{eq:quotient-observation-map}
\end{equation}
This is the visible precision canonically induced by $(H,C)$. It is the positive-definite shorted or Schur-complement construction associated with the observation map in the operator-theoretic sense of \cite{Schur1917,Anderson1971,AndersonDuffin1969,AndersonTrapp1975,Bhatia2007}. It is not chosen or assumed, instead forced by the observation problem itself: for each visible displacement, it records the least full-space energy needed to realise that displacement.

It has a variational meaning, it admits a clean differential calculus, it composes along towers of observation, and it differs in a controlled way from naive visible compression. Those facts identify $\Phi$ as the visible invariant carried by quotient observation.

\subsection{Setting and quotient observation on the positive cone}

Fix a surjective linear map $C : \mathbb{R}^n \to \mathbb{R}^m$. For each visible displacement $y \in \mathbb{R}^m$, consider the constrained quadratic energy
\begin{equation}
\mathcal E_H(y)
:=
\inf\!\left\{ x^{\mathsf T} H x : Cx = y \right\}.
\label{eq:full-space-energy}
\end{equation}
The positive cone is forced from the start: once the ambient precision is positive definite, the minimisation problem is strictly convex, the minimiser is unique, and the visible quadratic form is again positive definite. The quotient-visible precision is characterised by
\begin{equation}
\mathcal E_H(y)=y^{\mathsf T} \Phi_C(H) y.
\label{eq:variational-characterisation}
\end{equation}
So the visible operator is not an auxiliary construction, it is the quadratic form selected by the problem.

The same minimisation admits a matching dual formulation. The constrained primal energy upstairs and the dual visible energy downstairs agree, and that agreement is the algebra behind the quotient map. From it one obtains well posedness on $\Sym_{++}(n)$, positive homogeneity, monotonicity in Loewner order, and concavity, in the standard positive-operator framework of \cite{Loewner1934,Ando1979,HansenPedersen1982,Bhatia1997,Bhatia2007}. One also obtains a canonical $H_0$-horizontal chart, which separates genuine visible motion from gauge motion inside the hidden kernel and makes the later tangent calculus transparent.

\begin{theorem}[Variational characterisation]
\label{thm:variational}
For every $H \in \Sym_{++}(n)$ and every surjective $C$, the quotient-visible precision $\Phi_C(H)$ is the unique operator in $\Sym_{++}(m)$ satisfying \eqref{eq:variational-characterisation}. Equivalently,
\begin{equation}
y^{\mathsf T}\Phi_C(H)y
=
\inf_{Cx=y} x^{\mathsf T} H x
\qquad
(y \in \mathbb{R}^m).
\end{equation}
\end{theorem}

\begin{corollary}[Primal-dual decomposition]
The visible energy admits matching primal and dual representations. In particular, the minimising lift and the maximising visible multiplier encode the same quadratic form.
\end{corollary}

In this language quotient precision is the minimum-energy visible form induced by the upstairs law, and it sits naturally in the shorted-operator and parallel-addition side of positive-operator theory \cite{AndersonDuffin1969,Anderson1971,KuboAndo1980,PuszWoronowicz1975}.

\begin{theorem}[Basic positive-cone calculus]
\label{thm:basic-F}
The map $H \mapsto \Phi_C(H)$ is well posed on $\Sym_{++}(n)$, positively homogeneous, monotone, and concave. Relative to a fixed base point $H_0$, it admits a canonical horizontal chart adapted to quotient observation.
\end{theorem}

What has been gained already is decisive. Observation no longer looks like a loss of structure followed by arbitrary repair. It produces a definite positive operator with a canonical variational meaning.

\subsection{Fr\'echet calculus and the tangent ceiling}

Once the quotient map is fixed, its first-order calculus is immediate. For a symmetric perturbation $\dot H$, the derivative is the visible shadow of that perturbation transported through $H^{-1}$ and returned by $C$. The main point is not the formula itself, but what it bounds: positive perturbations upstairs induce a maximal first-order visible response, and later hidden-mediation effects can only sit beneath that ceiling.

\begin{theorem}[Fr\'echet derivative]
\label{thm:frechet-derivative}
The quotient map is Fr\'echet differentiable on $\Sym_{++}(n)$. Its derivative at $H$ in direction $\dot H$ is
\begin{equation}
D\Phi_C[H](\dot H)
=
\Phi_C(H)\, C H^{-1} \dot H H^{-1} C^{\mathsf T}\, \Phi_C(H).
\label{eq:derivative-schur}
\end{equation}
\end{theorem}

\begin{proposition}[Tangent ceiling]
If $\dot H \succeq 0$, then $D\Phi_C[H](\dot H) \succeq 0$. More sharply, $D\Phi_C[H](\dot H)$ is the maximal first-order visible response compatible with the positive perturbation $\dot H$.
\end{proposition}

The tangent ceiling is infrastructure rather than endpoint. It identifies the maximal first-order visible response compatible with a given positive perturbation. Later hidden-mediation results will show that hidden structure redistributes or attenuates visible strength beneath this ceiling, but does not create a new visible direction from nothing.

\subsection{Exact composition along observation towers}

The decisive structural fact is composition. Suppose observation occurs in stages,
\begin{equation}
\mathbb{R}^n \xrightarrow{\ C_1\ } \mathbb{R}^{m_1} \xrightarrow{\ C_2\ } \mathbb{R}^{m_2},
\qquad
C = C_2 C_1.
\end{equation}
Then one may observe once or observe in two steps. The result is the same:
\begin{equation}
\Phi_C(H)
=
\Phi_{C_2}\!\bigl(\Phi_{C_1}(H)\bigr).
\label{eq:tower-composition}
\end{equation}
So quotient observation closes on itself. The visible object keeps its type under repeated elimination of hidden structure.

\begin{theorem}[Composition law]
\label{thm:composition-law}
For surjective observation maps $C_1$ and $C_2$ with $C=C_2C_1$,
\begin{equation}
\Phi_C(H)
=
\Phi_{C_2}\!\bigl(\Phi_{C_1}(H)\bigr).
\end{equation}
\end{theorem}

This is what makes $\Phi$ a renormalisation-natural object rather than an ad hoc one, in the precise sense that repeated elimination stays inside one visible variable rather than generating a new effective object at each stage \cite{Kadanoff1966,WilsonKogut1974,Polchinski1984}. No correction term is needed. No auxiliary state has to be remembered. Repeated observation stays inside the same category and returns the same kind of operator.

The tangent theory composes as well. The visible derivative shadow transports down an observation tower by coisometric compression. At the spectral level this yields monotone Ky Fan detectability envelopes: as observation becomes coarser, the visible singular spectrum can only contract \cite{Fan1951,HornJohnson2013}.

\begin{theorem}[Tangent shadow transport]
\label{thm:shadow-transport}
Under tower composition, the derivative shadow of quotient observation transports by a coisometric map $W$ satisfying $WW^{\mathsf T} = I$. Specifically, if $C = C_2 C_1$ and $\Phi_1 := \Phi_{C_1}(H)$, the derivative at the coarser level factors through the derivative at the finer level via coisometric compression. In particular, the Ky Fan singular values of the visible derivative shadow contract monotonically along any observation tower.
\end{theorem}

With composition known, the role of $\Phi$ changes. It is no longer merely a good definition. It is the object that survives observation in the only way a visible invariant should.

\subsection{Compression versus quotient: the Schur gap}

Quotient observation is not naive visible compression. Fix a visible-hidden splitting and write
\begin{equation}
H
=
\begin{pmatrix}
H_{vv} & H_{vh} \\
H_{hv} & H_{hh}
\end{pmatrix}.
\end{equation}
Compression keeps the principal visible block $H_{vv}$. Quotient observation eliminates the hidden sector through the inverse and returns the Schur complement \cite{Zhang2005,Haynsworth1968,CarlsonHaynsworthMarkham1974}
\begin{equation}
\Phi(H)=H_{vv}-H_{vh}H_{hh}^{-1}H_{hv}.
\label{eq:schur-complement}
\end{equation}
The difference is the Schur gap
\begin{equation}
G(H)
:=
H_{vv}-\Phi(H)
=
H_{vh}H_{hh}^{-1}H_{hv}
\succeq 0.
\label{eq:schur-gap}
\end{equation}
So compression and quotient coincide only when the hidden sector does not feed back into the visible one.

\begin{proposition}[Schur gap formula]
For every block decomposition adapted to a visible-hidden splitting, the difference between naive compression and quotient observation is the positive semidefinite Schur gap.
\end{proposition}

\begin{remark}[Scalar diagnostics]
In scalar visible rank, useful summaries include $\operatorname{tr} G$, $\|G\|_{\mathrm{op}}$, and the corresponding log-determinant defect. These are diagnostics of hidden contribution.
\end{remark}

\subsection{The intrinsic visible invariant \texorpdfstring{$\Phi$}{Phi}}

The story is now fixed and quotient observation produces a definite visible precision. That operator has a variational characterisation, a stable first-order calculus, a composition law along observation towers, and a precise distinction from naive compression. Those properties identify $\Phi$ as the intrinsic visible invariant.

The Donsker-Varadhan Hessian has a different role. It is a bridge object upstairs, and it will matter in the next section because it leads to a finite visible algebraic descendant. But the object that closes under observation is $\Phi$, with the Donsker-Varadhan law surviving as its tangent ceiling.

That is the operator one can transport, compare, and compose without leaving the observation category. The paired-spectrum obstruction (\S\ref{sec:04_hidden_mediation}) makes this precise: at fixed visible dimension, the Donsker-Varadhan Hessian generically does not close, and the quotient precision is the correct replacement.

A second structural theme begins here. Later results will show that, in scalar visible rank, the visible precision splits as
\begin{equation}
\Phi = F_0 + W,
\label{eq:conservation-split-preview}
\end{equation}
with $F_0$ the reversible backbone and $W$ the irreversible correction. The point here is only conceptual: $\Phi$ is large enough to carry both contributions while remaining stable under observation.


\section{The Donsker-Varadhan bridge and the visible algebraic object}
\label{sec:03_dv_bridge}

The Donsker-Varadhan bridge from \cite{Dunkley2026FiniteObservation} produces a finite visible algebraic object, the tangent descendant of $\Phi$. The quartic and sextic defect theorems below measure how much structure survives the transition from tangent to full covariance level.

\subsection{The bridge and its quotient reduction}

Section~3 of \cite{Dunkley2026FiniteObservation} proved the exact quadratic Donsker-Varadhan bridge
\begin{equation}
H_{\mathrm{DV}} = H_0 + \Delta_{\mathrm{DV}},
\qquad
\Delta_{\mathrm{DV}} = \tfrac{1}{4}\, K\, H_0^{-1}\, K^{\mathsf T} \succeq 0,
\label{eq:dv-bridge}
\end{equation}
where $H_0$ is the time-symmetric backbone precision and $K$ is the skew current operator. We read the output of this bridge through quotient observation.

Introduce a perturbation parameter by $K = \eps K_1$, so that
\begin{equation}
H(\eps) = H_0 + \eps^2 \Delta_2 + O(\eps^4).
\label{eq:eps-bookkeeping}
\end{equation}
In particular,
\begin{equation}
\Delta_2 = \tfrac{1}{4}\, K_1\, H_0^{-1}\, K_1^{\mathsf T}.
\label{eq:Delta2-definition}
\end{equation}

\begin{definition}[$H_0$-adapted quotient sector]
\label{def:h0-adapted}
A decomposition $V = U \oplus U^{\perp}$ is called $H_0$-adapted if
\begin{equation}
H_0 =
\begin{pmatrix}
H_{0,S} & 0 \\
0 & H_{0,F}
\end{pmatrix}
\label{eq:h0-adapted}
\end{equation}
in the associated block coordinates.
\end{definition}

The algebraic cascade from bridge to visible object proceeds in three steps:
\begin{equation}
\begin{gathered}
H_{\mathrm{DV}} \;=\; H_0 + \tfrac{1}{4}\, K\, H_0^{-1}\, K^{\mathsf T}
\qquad [n \times n \text{ bridge}] \\[4pt]
\big\downarrow \quad \text{quotient observation } F_C \\[4pt]
\Phi_b \;=\; \operatorname{Schur}_v(H_b)
\qquad [m \times m \text{ visible algebraic object}] \\[4pt]
\big\downarrow \quad \text{composition (Theorem~\ref{thm:composition-law})} \\[4pt]
\Phi \;=\; F_C(H)
\qquad [\text{intrinsic visible invariant, closes}]
\end{gathered}
\label{eq:algebraic-cascade}
\end{equation}

\begin{remark}[Source, propagator, field structure]
The bridge \eqref{eq:dv-bridge} has a physical reading. The current operator $K$ is the source, the resolvent $H_0^{-1}$ is the propagator, and the product $K H_0^{-1} K^{\mathsf T}$ is the induced field correction. The visible algebraic object is the gauge-invariant visible projection of this chain.
\end{remark}

\subsection{The visible algebraic object}

Write the finite reduced block precision as
\begin{equation}
H_b
\;=\;
H_{0,b} + \tfrac{1}{4}\, K_b\, H_{0,b}^{-1}\, K_b^{\mathsf T}
\;=\;
\begin{pmatrix}
A & B \\
B^{\mathsf T} & D
\end{pmatrix},
\label{eq:hb-block}
\end{equation}
where the first coordinate is the visible direction and the remaining coordinates are hidden reduced directions.

\begin{theorem}[Visible algebraic object]
\label{thm:visible-algebraic-object}
The visible algebraic descendant of the bridge output is the Schur complement
\begin{equation}
\Phi_b \;:=\; \operatorname{Schur}_v(H_b) \;=\; A - B\, D^{-1}\, B^{\mathsf T}.
\label{eq:phi-schur}
\end{equation}
The Donsker-Varadhan correction is quadratic in the skew channel, and quotient elimination of the hidden reduced directions turns the bridge output into a visible Schur complement. By the composition theorem (Theorem~\ref{thm:composition-law}), $\Phi_b$ is a tangent descendant of the intrinsic visible precision $\Phi$.
\end{theorem}

\subsection{Scalar visible Schur identity and completion-independence package}
\label{subsec:scalar-schur}

In scalar visible rank ($m = 1$), the Schur complement reduces to a generator-level identity that does not require explicit diagonalisation of the full precision. The visible precision is a rational function of the generator entries, computable from the stationary distribution and transition rates without spectral decomposition.

\begin{theorem}[Completion independence at $M=1$]
\label{thm:completion-independence}
At the matched visible-law level, the quotient-visible precision $\Phi$ does not depend on the hidden completion. Two hidden completions that produce the same visible law yield the same visible precision.
\end{theorem}

The completion-independence hierarchy:

\begin{itemize}
\item Exact for $M = 1$ (proved).
\item Verified numerically for tested ladder families at $M = 2$.
\item Open for general $M$ and general graph topologies.
\end{itemize}

\subsection{Intrinsic quartic shadow defect and sextic divergence robustness}

The quartic order is where quotient observation first departs from the tangent law.

\begin{theorem}[Signed quartic bridge]
\label{thm:quartic-bridge}
Assume the Donsker-Varadhan expansion \eqref{eq:dv-bridge} with bookkeeping \eqref{eq:eps-bookkeeping}, and let $V = U \oplus U^{\perp}$ be $H_0$-adapted. Write $\Delta_{SF}^{(2)}$ for the $(S,F)$ block of $\Delta_2$ from \eqref{eq:Delta2-definition}. Then:
\begin{enumerate}
\item compression recovers the finite-observation visible quadratic law on $U$;
\item quotient and compression agree through order $\eps^2$;
\item the first genuinely new quotient effect is quartic and sign-definite:
\begin{equation}
H_U^{\mathrm{comp}} - H_U^{\mathrm{quot}}
=
\eps^4\, \Delta_{SF}^{(2)}\, H_{0,F}^{-1}\, \Delta_{FS}^{(2)} + O(\eps^6) \succeq 0.
\label{eq:quartic-gap}
\end{equation}
\end{enumerate}
\end{theorem}

\begin{proof}
Because the splitting is $H_0$-adapted, the mixed backbone blocks vanish. Hence
\[
H_{SF}(\eps) = \eps^2 \Delta_{SF}^{(2)} + O(\eps^4),
\qquad
H_{FF}(\eps) = H_{0,F} + O(\eps^2).
\]
Compression gives the visible block $H_{SS} = H_{0,S} + \eps^2 \Delta_{SS}^{(2)} + O(\eps^4)$, the finite-observation visible quadratic law from \cite{Dunkley2026FiniteObservation}. The Schur gap formula (\S\ref{sec:quotient-precision}) gives
\[
H_U^{\mathrm{comp}} - H_U^{\mathrm{quot}} = H_{SF}\, H_{FF}^{-1}\, H_{FS}.
\]
Substituting the block expansions yields \eqref{eq:quartic-gap}.
\end{proof}

\begin{corollary}[Tangent closure of visible Donsker-Varadhan geometry]
\label{cor:tangent-closure}
In the $H_0$-adapted setting,
\begin{equation}
H_U^{\mathrm{quot}} = H_{0,S} + (\Delta_{\mathrm{DV}})_{SS} + O(\norm{K}^4).
\label{eq:tangent-closure}
\end{equation}
The finite-observation visible law is the quadratic tangent law of quotient observation in the Donsker-Varadhan sector.
\end{corollary}

\begin{proof}
Combine Theorem~\ref{thm:quartic-bridge} with \eqref{eq:dv-bridge}.
\end{proof}

The signed quartic bridge identifies where quotient observation departs from compression. At the algebraic level it is a square-completion defect with definite sign, in the same broad spirit as energy-splitting identities familiar from other positive-definite variational settings \cite{Bogomolnyi1976}. The next theorem shows that this departure is intrinsic to the quotient geometry.

\begin{theorem}[Intrinsic quartic shadow defect and sextic divergence robustness]
\label{thm:shadow-sextic}
Fix a surjection $C : V \to Y$ and a backbone $H_0 \in \SPD(V)$. Let
\begin{equation}
H(\eps) = H_0 + \eps^2 \Delta_2 + O(\eps^4),
\qquad
\Delta_2 \succeq 0,
\label{eq:dv-general-family}
\end{equation}
and set $F_0 := F_C(H_0)$. Define the tangent visible correction
\begin{equation}
T_0 := DF_C(H_0)[\Delta_2],
\label{eq:tangent-visible-correction}
\end{equation}
the visible shadow
\begin{equation}
\Sigma_{\mathrm{exact}}(\eps)
:=
F_0^{-1/2}\bigl(F_C(H(\eps)) - F_0\bigr)F_0^{-1/2},
\label{eq:exact-shadow}
\end{equation}
and the tangent shadow
\begin{equation}
\Sigma_{\mathrm{tan}}(\eps)
:=
\eps^2\, F_0^{-1/2}\, T_0\, F_0^{-1/2}.
\label{eq:tangent-shadow}
\end{equation}
Then:
\begin{enumerate}
\item the quotient-visible expansion has the intrinsic form
\begin{equation}
F_C(H(\eps))
=
F_0 + \eps^2 T_0 - \eps^4 Q_0(\Delta_2) + O(\eps^6),
\label{eq:hessian-defect-expansion}
\end{equation}
where
\begin{equation}
Q_0(\Delta_2) := -\tfrac{1}{2}\, D^2 F_C(H_0)[\Delta_2, \Delta_2] \succeq 0;
\label{eq:hessian-defect}
\end{equation}
\item the visible shadow loses first at quartic order,
\begin{equation}
\Sigma_{\mathrm{exact}}(\eps)
=
\Sigma_{\mathrm{tan}}(\eps) - \eps^4 M_0 + O(\eps^6),
\qquad
M_0 := F_0^{-1/2}\, Q_0(\Delta_2)\, F_0^{-1/2} \succeq 0;
\label{eq:quartic-shadow-defect}
\end{equation}
\item for every smooth spectral functional $\Psi$ on symmetric matrices with $\Psi(0) = 0$ and $D\Psi(0) = 0$,
\begin{equation}
\Psi\!\bigl(\Sigma_{\mathrm{tan}}(\eps)\bigr) - \Psi\!\bigl(\Sigma_{\mathrm{exact}}(\eps)\bigr) = O(\eps^6).
\label{eq:sextic-functional-gap}
\end{equation}
In particular, reverse Gaussian KL, forward Gaussian KL, and squared Hellinger first lose only at sextic order.
\end{enumerate}
\end{theorem}

\begin{proof}
The quotient map $F_C$ is smooth on the positive cone, so Taylor expansion at $H_0$ along the curve $H(\eps)=H_0+\eps^2\Delta_2$ gives
\[
F_C(H(\eps))
=
F_C(H_0)+\eps^2 DF_C(H_0)[\Delta_2]+\tfrac12\eps^4 D^2F_C(H_0)[\Delta_2,\Delta_2]+O(\eps^6).
\]
Using the definitions of $F_0$, $T_0$, and $Q_0(\Delta_2):=-\tfrac12 D^2F_C(H_0)[\Delta_2,\Delta_2]$ yields \eqref{eq:hessian-defect-expansion}. Conjugating by $F_0^{-1/2}$ gives \eqref{eq:quartic-shadow-defect}.

For the spectral-functional claim, both $\Sigma_{\mathrm{tan}}(\eps)$ and $\Sigma_{\mathrm{exact}}(\eps)$ are $O(\eps^2)$, and their difference is $O(\eps^4)$ by \eqref{eq:quartic-shadow-defect}. Since $\Psi(0)=0$ and $D\Psi(0)=0$, the Taylor expansion of $\Psi$ at the origin begins quadratically. Applying that expansion to the pair $(\Sigma_{\mathrm{tan}}(\eps),\Sigma_{\mathrm{exact}}(\eps))$ shows that the first possible difference appears at order $\eps^6$, which proves \eqref{eq:sextic-functional-gap}.
\end{proof}

\begin{remark}
In the canonical $H_0$-horizontal chart of \S\ref{sec:quotient-precision}, the quartic Hessian defect $Q_0(\Delta_2)$ reduces to the signed quartic bridge of Theorem~\ref{thm:quartic-bridge}. The adapted block formula is the coordinate representation of an intrinsic quotient-geometric defect.
\end{remark}

\subsection{What the bridge still contributes conceptually}

\begin{remark}[Computation versus interpretation]
In the scalar visible case ($m = 1$), the bridge may be algebraically bypassed: the visible precision can be computed directly from the generator. The bridge nonetheless matters as the systematic path from the Donsker-Varadhan Hessian to the quotient-visible precision. The distinction is between computation (which the scalar case streamlines) and interpretation (which the bridge provides at any rank).
\end{remark}

\subsection{Boundaries of the scalar simplification}

The scalar simplification of \S\ref{subsec:scalar-schur} does not extend to general visible rank. When $m > 1$, the visible precision $\Phi$ is a matrix whose internal structure (eigenvalues, eigenvectors, parametric dependence) is not determined by scalar diagnostics.

Visible-law encoding of $\Phi$ is $M = 1$-specific: one can read $\Phi$ from the visible switching rates. For general $M$, the visible law constrains but does not determine $\Phi$. What additional structure is needed at higher rank remains open.

The algebraic theory (\S\ref{sec:quotient-precision}) and the bridge connection are in hand. The next sections develop the internal structure of $\Phi$: hidden mediation, conservation, locality, and the gluing frontier.


% Sections 4 onwards intentionally reduced to headers only

\section{Hidden mediation, scalar clocks, and sequential transport}
\label{sec:04_hidden_mediation}

Section~\ref{sec:03_dv_bridge} identified the visible algebraic object. The next question is internal: what hidden mechanism produces a visible law beneath the tangent ceiling, how does that mechanism compose, and which scalar quantity measures its cumulative action? The answer is hidden mediation. Beneath a fixed tangent ceiling, every support-preserving visible law is classified by a positive hidden load, sequential composition transports that load by a rigid conjugation law, and the only continuous orthogonally invariant additive clock on the interior is the determinant clock. The section closes by isolating the obstruction that prevents same-dimension Donsker-Varadhan closure from being the correct observable class.

\subsection{Exact latent-contraction formula and the tangent ceiling}

Work in an $H_0$-adapted quotient chart and whiten by the backbone so that the upstairs precision takes the form
\begin{equation}\label{eq:whitened-LLt}
\widetilde H = I + L L^{\mathsf T},
\qquad
L=
\begin{pmatrix}
B\\
M
\end{pmatrix},
\end{equation}
where $B$ contains the visible rows and $M$ the hidden rows. In this representation the tangent ceiling is explicit: quotient observation keeps the same visible support ceiling as compression, but contracts beneath it by a positive hidden contribution. At the level of Gaussian graphical models this is hidden-variable precision geometry, not a heuristic latent-variable ansatz \cite{Dempster1972,Lauritzen1996,ChandrasekaranParriloWillsky2012}.

\begin{theorem}[Exact latent-contraction formula]\label{thm:latent-contraction}
In the setting \eqref{eq:whitened-LLt}, the visible quotient correction is
\begin{equation}\label{eq:exact-visible-correction}
X = B(I+M^{\mathsf T} M)^{-1} B^{\mathsf T}.
\end{equation}
The tangent ceiling is
\begin{equation}\label{eq:tangent-ceiling}
T:=BB^{\mathsf T},
\end{equation}
and the gap is
\begin{equation}\label{eq:exact-gap}
T-X = B\,M^{\mathsf T}(I+MM^{\mathsf T})^{-1}M\,B^{\mathsf T} \succeq 0.
\end{equation}
Equivalently,
\begin{equation}\label{eq:X-as-BRBt}
X = B R B^{\mathsf T},
\qquad
R:=(I+M^{\mathsf T} M)^{-1},
\qquad 0\prec R\le I.
\end{equation}
\end{theorem}

\begin{proof}
In block form,
\[
\widetilde H=
\begin{pmatrix}
I+BB^{\mathsf T} & BM^{\mathsf T}\\
MB^{\mathsf T} & I+MM^{\mathsf T}
\end{pmatrix}.
\]
The visible quotient precision is its Schur complement:
\[
I+BB^{\mathsf T} - BM^{\mathsf T}(I+MM^{\mathsf T})^{-1}MB^{\mathsf T}.
\]
Using the identity
\[
I - M^{\mathsf T}(I+MM^{\mathsf T})^{-1}M = (I+M^{\mathsf T}M)^{-1},
\]
we obtain
\[
I + B(I+M^{\mathsf T}M)^{-1}B^{\mathsf T}.
\]
Subtracting the whitened backbone $I$ gives \eqref{eq:exact-visible-correction}. The formulae \eqref{eq:exact-gap} and \eqref{eq:X-as-BRBt} follow immediately.
\end{proof}

\begin{corollary}[Exact support theorem]\label{cor:support-theorem}
For every finite hidden sector,
\begin{equation}\label{eq:support-theorem}
\operatorname{Ran}(X)=\operatorname{Ran}(T),
\qquad
\rank(X)=\rank(T).
\end{equation}
Thus finite hidden mediation preserves the visible correction support. In particular, the hidden sector attenuates inside the existing visible support rather than creating new visible directions, which is also the natural completion geometry for latent Gaussian models \cite{GroneJohnsonSaWolkowicz1984,ChandrasekaranParriloWillsky2012}.
\end{corollary}

\begin{proof}
Since $R\succ 0$ on the latent space, $BRB^{\mathsf T}$ and $BB^{\mathsf T}$ have the same range and rank.
\end{proof}

\begin{remark}[Blind directions remain blind]
Corollary~\ref{cor:support-theorem} is the support rigidity statement in ordinary language. If a visible direction receives no signal from the tangent ceiling $T$, then no finite hidden architecture can make that direction informative later, because the actual visible correction $X$ has the same support. Hidden complexity can attenuate or redistribute strength inside the tangent support, but it cannot create a new visible direction from nothing.
\end{remark}

\subsection{Support-preserving visible cone, its closure, and hidden-load parametrisation}

The latent-contraction formula says that hidden mediation acts by attenuation inside the support of the tangent ceiling. The next result identifies the full image of that action.

\begin{proposition}[Support-preserving interval and closure theorem]\label{prop:interval}
Let $T=BB^{\mathsf T}$ and let $S:=\operatorname{Ran}(T)$. A matrix $X$ arises from a finite hidden sector if and only if
\begin{equation}\label{eq:interval-form}
X=T^{1/2}\Pi T^{1/2},
\qquad 0\prec \Pi\le I \text{ on } S.
\end{equation}
Taking closure adds the boundary $0\le \Pi\le I$ on $S$. Equivalently, the closure of the visible image is the Loewner interval beneath $T$, whose interior is support-preserving.
\end{proposition}

\begin{proof}
From Theorem~\ref{thm:latent-contraction},
\[
X = BRB^{\mathsf T},
\qquad 0\prec R\le I.
\]
By polar decomposition, $B=T^{1/2}U$ for a partial isometry $U$ with initial space the latent support of $B$ and final space $S$. Therefore
\[
X=T^{1/2}(URU^{\mathsf T})T^{1/2},
\]
so \eqref{eq:interval-form} holds with $\Pi:=URU^{\mathsf T}$. Since $R\succ 0$ and $R\le I$, one has $0\prec\Pi\le I$ on $S$.

Conversely, let $0\prec\Pi\le I$ on $S$. Choose the latent space to be $S$, set $B:=T^{1/2}$, and define $R:=\Pi$. Since $R\succ 0$ and $R\le I$, the matrix $G:=R^{-1}-I$ is positive semidefinite. Pick $M$ with $M^{\mathsf T}M=G$. Then $(I+M^{\mathsf T}M)^{-1}=R$, so $X=BRB^{\mathsf T}$. The boundary case follows by limit.
\end{proof}

\noindent\textbf{Exact visible class and its closure.} With $S:=\operatorname{Ran}(T)$,

\[
\begingroup
\setlength{\fboxsep}{8pt}%
\setlength{\fboxrule}{1.05pt}%
\boxed{\displaystyle \Phi_T^{\circ}=\left\{X=T^{1/2}\Pi T^{1/2}:\; 0\prec \Pi\le I\ \text{ on }\ S=\operatorname{Ran}(T)\right\}}%
\endgroup
\]

\[
\begingroup
\setlength{\fboxsep}{8pt}%
\setlength{\fboxrule}{1.05pt}%
\boxed{\displaystyle \overline{\Phi}_T=\left\{X=T^{1/2}\Pi T^{1/2}:\; 0\le \Pi\le I\ \text{ on }\ S=\operatorname{Ran}(T)\right\}=T^{1/2}[0,I]_S T^{1/2}}%
\endgroup
\]

Here $[0,I]_S:=\{\Pi\in\Sym(S):0\le \Pi\le I_S\}$. Thus finite hidden sectors realise the support-preserving interior $\Phi_T^{\circ}$, while taking closure adds the boundary $\overline{\Phi}_T$. The hidden-load coordinate of Theorem~\ref{thm:hidden-load} is defined on $\Phi_T^{\circ}$, and the next results identify its rank and determinant invariants.

\begin{theorem}[Intrinsic hidden-load parametrisation]\label{thm:hidden-load}
Let $T\succeq 0$ and let $S:=\operatorname{Ran}(T)$. A visible law $X$ in the support-preserving interior beneath $T$ is equivalently an operator of the form
\[
X=T^{1/2}\Pi T^{1/2},
\qquad
0\prec \Pi\le I \text{ on } S.
\]
For such $X$, define the intrinsic hidden-load operator on $S$ by
\begin{equation}\label{eq:hidden-load}
\Lambda
:=
(T|_S)^{1/2}(X|_S)^{-1}(T|_S)^{1/2}-I_S.
\end{equation}
Then $\Lambda\succeq 0$ on $S$, and
\begin{equation}\label{eq:X-from-load}
X=T^{1/2}(I_S+\Lambda)^{-1}T^{1/2}.
\end{equation}
Conversely, every $\Lambda\succeq 0$ on $S$ determines a visible law by \eqref{eq:X-from-load}. Hence the support-preserving visible interior beneath $T$ is in bijection with the positive cone on $S$.
\end{theorem}

\begin{proof}
Write
\[
\Pi:=(T|_S)^{-1/2}(X|_S)(T|_S)^{-1/2}.
\]
Then $0\prec \Pi\le I_S$ on $S$. Therefore
\[
\Lambda=\Pi^{-1}-I_S \succeq 0,
\]
and
\[
\Pi=(I_S+\Lambda)^{-1}.
\]
Substituting this into $X=T^{1/2}\Pi T^{1/2}$ gives \eqref{eq:X-from-load}.

Conversely, if $\Lambda\succeq 0$ on $S$, then $\Pi:=(I_S+\Lambda)^{-1}$ satisfies
\[
0\prec \Pi\le I_S,
\]
so Proposition~\ref{prop:interval} yields a visible realisation
\[
X=T^{1/2}\Pi T^{1/2}=T^{1/2}(I_S+\Lambda)^{-1}T^{1/2}.
\]
This proves the bijection.
\end{proof}

\begin{corollary}[Hidden-load monotonicity]\label{cor:hidden-load-monotone}
Let
\[
X_i=T^{1/2}(I_S+\Lambda_i)^{-1}T^{1/2},
\qquad \Lambda_i\succeq 0 \text{ on } S,
\]
for $i=0,1$. Then
\begin{equation}\label{eq:hidden-load-order}
\Lambda_1\succeq \Lambda_0
\quad\Longleftrightarrow\quad
X_1\preceq X_0.
\end{equation}
Thus increasing hidden load can only attenuate the visible correction.
\end{corollary}

\begin{proof}
Restrict to $S$, where $T|_S\succ 0$. Congruence by $(T|_S)^{-1/2}$ gives
\[
X_1\preceq X_0
\quad\Longleftrightarrow\quad
(I_S+\Lambda_1)^{-1}\preceq (I_S+\Lambda_0)^{-1}.
\]
Since inversion reverses Loewner order on positive definite operators,
\[
(I_S+\Lambda_1)^{-1}\preceq (I_S+\Lambda_0)^{-1}
\quad\Longleftrightarrow\quad
I_S+\Lambda_1 \succeq I_S+\Lambda_0,
\]
which is \eqref{eq:hidden-load-order}.
\end{proof}

\begin{remark}[One positive-cone geometry seen from two sides]
The quotient programme now contains two positive-cone geometries. First, the global quotient map
\[
H \mapsto F_C(H)=(CH^{-1}C^{\mathsf T})^{-1}
\]
is monotone and concave on the upstairs positive cone. Second, beneath a fixed tangent ceiling $T$ on the active support $S=\operatorname{Ran}(T)$, every support-preserving visible law in the interior has the form
\[
X=T^{1/2}(I_S+\Lambda)^{-1}T^{1/2},
\qquad \Lambda\succeq 0.
\]
The first is the global quotient-side cone, the second is the intrinsic hidden-load cone after the ceiling has been fixed. They are the global and support-adapted faces of the same attenuation geometry.
\end{remark}

\begin{theorem}[Minimal hidden latent dimension]\label{thm:min-hidden-dim}
Let $X=T^{1/2}\Pi T^{1/2}$ with $0\prec\Pi\le I$ on $S=\operatorname{Ran}(T)$. Then the minimal hidden latent dimension needed to realise $X$ is
\begin{equation}\label{eq:min-hidden-dim}
d_{\min} = \rank(I-\Pi) = \rank(T-X).
\end{equation}
\end{theorem}

\begin{proof}
By Theorem~\ref{thm:hidden-load}, write
\[
X=T^{1/2}(I_S+\Lambda)^{-1}T^{1/2},
\qquad
\Lambda\succeq 0.
\]
Any realisation requires a Gram factorisation $\Lambda=M^{\mathsf T}M$, so the hidden latent dimension must be at least $\rank(\Lambda)$. Choosing a Gram factorisation of $\Lambda$ shows that $\rank(\Lambda)$ hidden dimensions suffice. Since
\[
T-X = T^{1/2}\Bigl(I_S-(I_S+\Lambda)^{-1}\Bigr)T^{1/2},
\]
and $I_S-(I_S+\Lambda)^{-1}=\Lambda(I_S+\Lambda)^{-1}$ has the same rank as $\Lambda$, we obtain
\[
\rank(\Lambda)=\rank(T-X)=\rank(I-\Pi).
\]
This proves the claim.
\end{proof}

\begin{corollary}[Rank and volume of hidden load]\label{cor:hidden-load-rank-volume}
With $\Lambda$ as in Theorem~\ref{thm:hidden-load},
\begin{equation}\label{eq:hidden-load-rank}
\rank(\Lambda)=\rank(T-X),
\end{equation}
and
\begin{equation}\label{eq:hidden-load-volume}
\log\pdet(T)-\log\pdet(X)=\log\det(I_S+\Lambda).
\end{equation}
In particular, the minimal hidden latent dimension is $\rank(\Lambda)=\rank(T-X)$.
\end{corollary}

\begin{proof}
From \eqref{eq:X-from-load},
\[
T-X
=
T^{1/2}\Bigl(I_S-(I_S+\Lambda)^{-1}\Bigr)T^{1/2}
=
T^{1/2}\Lambda(I_S+\Lambda)^{-1}T^{1/2}.
\]
Since $(I_S+\Lambda)^{-1}$ is invertible on $S$, this implies
\[
\rank(T-X)=\rank(\Lambda).
\]
Also, on $S$,
\[
X|_S=(T|_S)^{1/2}(I_S+\Lambda)^{-1}(T|_S)^{1/2},
\]
hence
\[
\det(X|_S)=\det(T|_S)\det(I_S+\Lambda)^{-1}.
\]
Taking logarithms gives \eqref{eq:hidden-load-volume}.
\end{proof}

\begin{remark}[Operator-interval picture]
Beneath a fixed tangent ceiling $T$, finite hidden sectors realise the support-preserving interior
\[
X=T^{1/2}\Pi T^{1/2},
\qquad 0\prec \Pi\le I \text{ on } S,
\]
and taking closure adds the boundary $0\le \Pi\le I$ on $S$. On the interior, the hidden-load coordinate $\Lambda=\Pi^{-1}-I$ turns the support-preserving class into the positive cone on the active support. The scalar quantity
\[
\log\det(I_S+\Lambda)=\log\pdet(T)-\log\pdet(X)
\]
measures latent contraction volume. Same-dimension skew-square closure will appear below as a special paired-spectrum slice inside this larger cone.
\end{remark}

\subsection{Sequential composition, hidden-load transport, and scalar clocks}
\label{sec:support_towers}

All statements in this subsection are made on one fixed active support stratum $S$. No claim is made here about a finite additive scalar clock across genuine rank-loss events.

\begin{remark}
The coupled field layer leaves the transport algebra of this subsection unchanged on support-stable strata. What changes there is only the infinitesimal generator: the one-step increment becomes $M=dt\,A_{\mathrm{cpl}}+o(dt)$ with
\[
A_{\mathrm{cpl}}:=-\tfrac12 V_S^{-1/2}(D_t^\alpha V)_S V_S^{-1/2},
\]
so the present hidden-load transport law applies stratumwise in the coupled regime as well.
\end{remark}

\begin{proposition}[Support-preserving sequential effect law]\label{prop:sequential-effect-law}
Let
\[
\mathcal E^\circ(S):=\{\Pi\in\mathrm{Sym}(S):0\prec \Pi\le I_S\}.
\]
For $\Pi,\Xi\in\mathcal E^\circ(S)$ define
\begin{equation}\label{eq:sequential-effect-law}
\Pi\circ\Xi:=\Pi^{1/2}\Xi\Pi^{1/2}.
\end{equation}
Then $\Pi\circ\Xi\in\mathcal E^\circ(S)$.
\end{proposition}

\begin{proof}
Since $\Xi\succ 0$, congruence by $\Pi^{1/2}$ preserves positivity and gives $\Pi\circ\Xi\succ 0$. Also,
\[
0\prec \Xi\le I_S
\quad\Longrightarrow\quad
0\prec \Pi^{1/2}\Xi\Pi^{1/2}\le \Pi\le I_S.
\]
Hence $\Pi\circ\Xi\in\mathcal E^\circ(S)$.
\end{proof}

\begin{remark}[The visible law is not the associative object]
The visible product $\circ$ is not associative in general. The true associative object is the underlying contraction composition downstairs. The visible interval is the Gram-shadow image of that contraction composition.
\end{remark}

\begin{theorem}[Exact hidden-load transport law]\label{thm:hidden-load-transport}
Let
\[
\Pi=(I_S+\Lambda)^{-1},
\qquad
\Xi=(I_S+M)^{-1},
\qquad
\Lambda,M\succeq 0 \text{ on } S.
\]
Then
\begin{equation}\label{eq:load-transport-product}
\Pi\circ\Xi=(I_S+\Lambda_{\mathrm{tot}})^{-1},
\end{equation}
with
\begin{equation}\label{eq:load-transport-law}
\Lambda_{\mathrm{tot}}
=
\Lambda + (I_S+\Lambda)^{1/2} M (I_S+\Lambda)^{1/2}.
\end{equation}
\end{theorem}

\begin{proof}
Because
\[
\Pi^{-1/2}=(I_S+\Lambda)^{1/2},
\qquad
\Xi^{-1}=I_S+M,
\]
we have
\[
(\Pi\circ\Xi)^{-1}
=
\Pi^{-1/2}\Xi^{-1}\Pi^{-1/2}
=
(I_S+\Lambda)^{1/2}(I_S+M)(I_S+\Lambda)^{1/2}.
\]
Expanding the right-hand side gives
\[
(\Pi\circ\Xi)^{-1}
=
I_S + \Lambda + (I_S+\Lambda)^{1/2} M (I_S+\Lambda)^{1/2}.
\]
This is \eqref{eq:load-transport-law}.
\end{proof}

\begin{corollary}[Survivor transport law]\label{cor:survivor-transport}
Define the survivor complement
\[
\Sigma(\Pi):=I_S-\Pi.
\]
Then
\begin{equation}\label{eq:survivor-transport}
\Sigma(\Pi\circ\Xi)
=
\Sigma(\Pi)+\Pi^{1/2}\Sigma(\Xi)\Pi^{1/2}.
\end{equation}
\end{corollary}

\begin{proof}
Use
\[
I_S-\Pi^{1/2}\Xi\Pi^{1/2}
=
(I_S-\Pi)+\Pi^{1/2}(I_S-\Xi)\Pi^{1/2}.
\]
\end{proof}

The transport law isolates the correct associative structure, and the scalar clock theorem identifies the only continuous additive quantity compatible with it.

\begin{corollary}[Additive determinant clock and barrier identity]\label{cor:det-clock}
Define
\begin{equation}\label{eq:det-clock}
\tau(\Pi):=-\log\det\Pi
\end{equation}
on $\mathcal E^\circ(S)$. Then
\begin{equation}\label{eq:det-clock-additive}
\tau(\Pi\circ\Xi)=\tau(\Pi)+\tau(\Xi).
\end{equation}
In hidden-load coordinates,
\begin{equation}\label{eq:clock-load-barrier}
\tau(\Pi)=\log\det(I_S+\Lambda).
\end{equation}
If the unwhitened visible law is
\[
X=T^{1/2}\Pi T^{1/2},
\]
then on the active support,
\begin{equation}\label{eq:pdet-barrier}
\log\pdet(T)-\log\pdet(X)=\log\det(I_S+\Lambda).
\end{equation}
\end{corollary}

\begin{proof}
Since
\[
\det(\Pi^{1/2}\Xi\Pi^{1/2})=\det(\Pi)\det(\Xi),
\]
we obtain \eqref{eq:det-clock-additive}. Because $\Pi=(I_S+\Lambda)^{-1}$, this is equivalent to \eqref{eq:clock-load-barrier}. The pseudo-determinant identity is Corollary~\ref{cor:hidden-load-rank-volume}.
\end{proof}

\begin{theorem}[Uniqueness of the additive interior clock]\label{thm:clock-uniqueness}
Let $F:\mathcal E^\circ(S)\to\mathbb R$ be continuous, orthogonally invariant, and additive under $\circ$:
\[
F(\Pi\circ\Xi)=F(\Pi)+F(\Xi),
\qquad
F(I_S)=0.
\]
Then there exists $c\in\mathbb R$ such that
\begin{equation}\label{eq:clock-uniqueness}
F(\Pi)=c\,\log\det\Pi
\end{equation}
for all $\Pi\in\mathcal E^\circ(S)$.
\end{theorem}

\begin{proof}
By orthogonal invariance, it suffices to evaluate $F$ on diagonal effects. For diagonal matrices,
\[
\operatorname{diag}(a_1,\dots,a_r)\circ \operatorname{diag}(b_1,\dots,b_r)
=
\operatorname{diag}(a_1b_1,\dots,a_rb_r),
\]
with each $a_i,b_i\in(0,1]$. Hence continuity and additivity under coordinatewise multiplication imply
\[
F(\operatorname{diag}(a_1,\dots,a_r)) = c\sum_i \log a_i
\]
for some constant $c$. Orthogonal invariance extends this to all $\Pi\in\mathcal E^\circ(S)$.
\end{proof}

\begin{corollary}[No finite continuous additive clock on the closed interval]\label{cor:no-global-finite-clock}
Let $F:[0,I_S]\to\mathbb R$ be continuous, orthogonally invariant, additive under $\circ$, and satisfy $F(I_S)=0$. Then
\[
F\equiv 0.
\]
\end{corollary}

\begin{proof}
Restrict to the interior $0\prec \Pi\le I_S$. By Theorem~\ref{thm:clock-uniqueness},
\[
F(\Pi)=c\,\log\det\Pi
\]
there. Consider
\[
\Pi_\varepsilon=\operatorname{diag}(1,\dots,1,\varepsilon),
\qquad \varepsilon\downarrow 0.
\]
If $c\neq 0$, then $\log\det\Pi_\varepsilon=\log\varepsilon\to -\infty$, contradicting finite continuity on $[0,I_S]$. Thus $c=0$, hence $F\equiv 0$.
\end{proof}

\begin{corollary}[Extended-real completion of the clock]\label{cor:extended-clock}
Define
\[
\tau_{\mathrm{ext}}(\Pi)
=
\begin{cases}
-\log\det\Pi, & 0\prec \Pi\le I_S,\\
+\infty, & \Pi \text{ singular.}
\end{cases}
\]
Then $\tau_{\mathrm{ext}}$ is additive under $\circ$ in the extended-real sense.
\end{corollary}

\begin{proof}
Interior additivity is Corollary~\ref{cor:det-clock}. If either factor is singular, then $\Pi\circ\Xi$ is singular, so $+\infty$ is stable under composition.
\end{proof}

\begin{proposition}[Scalar mediation invariants]\label{prop:scalar-mediation}
Let
\[
G:=I_S-\Pi=\Lambda(I_S+\Lambda)^{-1}.
\]
If $\lambda_i$ are the eigenvalues of $\Lambda$, then the eigenvalues of $G$ are
\[
g_i=\frac{\lambda_i}{1+\lambda_i}.
\]
Hence
\begin{equation}\label{eq:trace-gap-load}
\operatorname{tr}(G)=\sum_i \frac{\lambda_i}{1+\lambda_i},
\end{equation}
\begin{equation}\label{eq:op-gap-load}
\|G\|_{\mathrm{op}}=\max_i \frac{\lambda_i}{1+\lambda_i},
\end{equation}
and
\begin{equation}\label{eq:logdet-gap-load}
\tau(\Pi)=\log\det(I_S+\Lambda)=\sum_i \log(1+\lambda_i).
\end{equation}
Thus trace, operator norm, and log-determinant provide complementary scalar diagnostics of total, peak, and volumetric hidden mediation.
\end{proposition}

\begin{proof}
Since $\Pi=(I_S+\Lambda)^{-1}$ is a spectral function of $\Lambda$, the operators $\Pi$, $G$, and $\Lambda$ commute and are simultaneously diagonalisable. On each eigenvector of $\Lambda$ with eigenvalue $\lambda_i$, the corresponding eigenvalue of $G$ is $\lambda_i/(1+\lambda_i)$. The displayed trace, norm, and log-determinant formulas follow immediately.
\end{proof}

\begin{proposition}[Exact superadditive ceiling-KL identities]\label{prop:ceiling-kl-superadditive}
Let
\[
A:=I_S+\Lambda,
\qquad
B:=I_S+M,
\qquad
A_{\mathrm{tot}}:=A^{1/2}BA^{1/2},
\]
and write
\[
P_0:=\mathcal N(0,I_S),
\qquad
P_\Lambda:=\mathcal N(0,A),
\qquad
P_M:=\mathcal N(0,B),
\qquad
P_{\Lambda_{\mathrm{tot}}}:=\mathcal N(0,A_{\mathrm{tot}}).
\]
Then
\begin{equation}\label{eq:reverse-kl-superadditive}
D_{KL}(P_{\Lambda_{\mathrm{tot}}}\|P_0)
=
D_{KL}(P_\Lambda\|P_0)+D_{KL}(P_M\|P_0)+\frac12\operatorname{tr}(\Lambda M),
\end{equation}
and, with
\[
G_\Lambda:=I_S-(I_S+\Lambda)^{-1},
\qquad
G_M:=I_S-(I_S+M)^{-1},
\]
one also has
\begin{equation}\label{eq:forward-kl-superadditive}
D_{KL}(P_0\|P_{\Lambda_{\mathrm{tot}}})
=
D_{KL}(P_0\|P_\Lambda)+D_{KL}(P_0\|P_M)+\frac12\operatorname{tr}(G_\Lambda G_M).
\end{equation}
In particular, both Gaussian KL directions relative to the ceiling are superadditive under support-preserving sequential quotienting.
\end{proposition}

\begin{proof}
For the reverse direction,
\[
D_{KL}(P_\Lambda\|P_0)=\frac12\bigl(\operatorname{tr}(A)-r-\log\det A\bigr),
\qquad r:=\dim S.
\]
Because
\[
\operatorname{tr}(A_{\mathrm{tot}})=\operatorname{tr}(AB)=r+\operatorname{tr}\Lambda+\operatorname{tr}M+\operatorname{tr}(\Lambda M)
\]
and
\[
\log\det A_{\mathrm{tot}}=\log\det A+\log\det B,
\]
substitution gives \eqref{eq:reverse-kl-superadditive}. For the forward direction,
\[
D_{KL}(P_0\|P_\Lambda)=\frac12\bigl(\log\det A+\operatorname{tr}(A^{-1})-r\bigr).
\]
Now
\[
A_{\mathrm{tot}}^{-1}=A^{-1/2}B^{-1}A^{-1/2},
\]
so cyclicity of trace gives
\[
\operatorname{tr}(A_{\mathrm{tot}}^{-1})=\operatorname{tr}(A^{-1}B^{-1}).
\]
Since
\[
A^{-1}=I_S-G_\Lambda,
\qquad
B^{-1}=I_S-G_M,
\]
we have
\[
\operatorname{tr}(A^{-1}B^{-1})=r-\operatorname{tr}G_\Lambda-\operatorname{tr}G_M+\operatorname{tr}(G_\Lambda G_M).
\]
Substituting into the forward Gaussian KL formula yields \eqref{eq:forward-kl-superadditive}.
\end{proof}

\begin{remark}[Hidden-load interaction measure]
The cross-term $\tfrac12\operatorname{tr}(\Lambda M)$ in \eqref{eq:reverse-kl-superadditive} is not bookkeeping noise. It is the interaction measure for sequential hidden mediations. It vanishes only when the two mediations are orthogonal in hidden-load space, and otherwise quantifies how much the second mediation is transported by the first before the visible law is read out.
\end{remark}

\subsection{Conservation split: $F_0 + W = \Phi$}

The determinant clock identifies the only additive scalar content of sequential mediation. The next statement isolates the additive operator content of the visible precision itself.

\begin{theorem}[Conservation split]\label{thm:conservation-split}
Fix a completion and write
\begin{equation}\label{eq:conservation-defs}
F_0 := F_{\mathrm{vis}}(H_0),
\qquad
\Phi := F_{\mathrm{vis}}(H_{\mathrm{alg}}),
\qquad
W := \Phi - F_0,
\end{equation}
where $H_0$ is the reversible backbone and $H_{\mathrm{alg}}$ is the algebraic precision carrying the full visible response. Then
\begin{equation}\label{eq:conservation-split}
\Phi = F_0 + W.
\end{equation}
In scalar visible rank, and more generally on every class of completions for which $\Phi$ is known to be completion-independent, the completion-dependent variations of $F_0$ and $W$ compensate so that the total visible precision remains fixed.
\end{theorem}

\begin{proof}
The identity \eqref{eq:conservation-split} is the definition of $W$ rewritten. The nontrivial statement is the compensation claim. In scalar visible rank this follows from the completion-independence package of Section~\ref{subsec:scalar-schur}: when two completions produce the same visible law, they produce the same $\Phi$, hence any change in $F_0$ must be cancelled by the opposite change in $W$. The same implication holds on any broader class of completions for which completion-independence of $\Phi$ has been established.
\end{proof}

\begin{remark}[Ward-identity reading]
The split \eqref{eq:conservation-split} has the structure of a Ward identity. The total visible precision $\Phi$ is the invariant object, while $F_0$ and $W$ depend on the chosen completion but are locked in sum. The paper does not need gauge language for the mathematics, but it is the right interpretation of the compensation mechanism.
\end{remark}

\begin{remark}[Reversible and irreversible sectors]
In the Donsker-Varadhan bridge, $F_0$ is the visible contribution of the time-symmetric backbone $H_0$ and $W$ is the visible increment carried by the quadratic skew correction $\Delta$. The split therefore separates the reversible backbone from the irreversible correction at the level of the intrinsic visible object. It is this separation, not the fixed-dimension Hessian itself, that survives quotient observation.
\end{remark}

\subsection{Same-dimension closure criterion and its failure boundary}

The hidden-load cone is the class closed by quotient elimination. The same observed-dimension Donsker-Varadhan skew-square class is smaller, and its failure boundary is sharp.

\begin{proposition}[Paired-spectrum criterion]\label{prop:paired-spectrum}
Let $H_{0,\mathrm{vis}}\succ0$ and let $\Delta_{\mathrm{vis}}\succeq0$ be a visible correction. Set
\begin{equation}\label{eq:whitened-visible}
W:=H_{0,\mathrm{vis}}^{-1/2}\,\Delta_{\mathrm{vis}}\,H_{0,\mathrm{vis}}^{-1/2}.
\end{equation}
There exists a visible skew matrix $J_{\mathrm{vis}}$ of the same observed dimension such that
\[
\Delta_{\mathrm{vis}}=\tfrac14 J_{\mathrm{vis}}H_{0,\mathrm{vis}}^{-1}J_{\mathrm{vis}}^{\mathsf T}
\]
if and only if the positive spectrum of $W$ is pairwise degenerate. In particular, odd visible correction rank is impossible for same-dimension skew-square closure.
\end{proposition}

This is the same real skew-symmetric pairing obstruction that underlies Williamson-type normal forms and the canonical even-multiplicity structure of skew squares \cite{Williamson1936,HornJohnson1991}.

\begin{proof}
If such $J_{\mathrm{vis}}$ exists, define
\[
K:=H_{0,\mathrm{vis}}^{-1/2}J_{\mathrm{vis}}H_{0,\mathrm{vis}}^{-1/2}.
\]
Then $K$ is real skew-symmetric and
\[
W=\tfrac14 K K^{\mathsf T}.
\]
The nonzero singular values of a real skew matrix come in equal pairs, hence the positive eigenvalues of $W$ do as well.

Conversely, if the positive spectrum of $W$ is pairwise degenerate, choose an orthogonal basis in which
\[
W=Q\,\diag(\mu_1,\mu_1,\mu_2,\mu_2,\dots,0,\dots,0)\,Q^{\mathsf T},
\qquad \mu_i>0.
\]
Let
\[
K := 2Q\,\diag\!\bigl(\sqrt{\mu_1}J_2,\sqrt{\mu_2}J_2,\dots,0\bigr)Q^{\mathsf T},
\]
where $J_2=\begin{psmallmatrix}0&1\\-1&0\end{psmallmatrix}$. Then $K$ is real skew-symmetric and $W=\tfrac14 K K^{\mathsf T}$. Setting
\[
J_{\mathrm{vis}}:=H_{0,\mathrm{vis}}^{1/2} K H_{0,\mathrm{vis}}^{1/2}
\]
gives the required representation.
\end{proof}

\begin{corollary}[Paired-spectrum defect as exact distance]\label{cor:paired-defect}
Let $W$ be as in \eqref{eq:whitened-visible}, with ordered eigenvalues
\[
\lambda_1\ge \lambda_2\ge \cdots \ge \lambda_n\ge 0.
\]
Write $m:=\lfloor n/2\rfloor$ and define
\[
\bar\lambda_i:=\frac{\lambda_{2i-1}+\lambda_{2i}}{2},
\qquad i=1,\dots,m.
\]
Let
\begin{equation}\label{eq:pair-projection}
P_{\mathrm{pair}}(W)
:=
Q\,\diag(\bar\lambda_1,\bar\lambda_1,\dots,\bar\lambda_m,\bar\lambda_m,\,0\text{ if }n\text{ is odd})\,Q^{\mathsf T},
\end{equation}
where $W=Q\diag(\lambda_1,\dots,\lambda_n)Q^{\mathsf T}$ is an orthogonal spectral decomposition. Then $P_{\mathrm{pair}}(W)$ belongs to the same observed-dimension skew-square slice from Proposition~\ref{prop:paired-spectrum}, and it is the unique Frobenius-nearest point in that slice. Equivalently,
\begin{equation}\label{eq:paired-defect-distance}
\inf\Bigl\{\norm{W-Z}_F: Z=\tfrac14 K K^{\mathsf T},\ K^{\mathsf T}=-K,\ K\in\mathbb R^{n\times n}\Bigr\}
=
\norm{W-P_{\mathrm{pair}}(W)}_F.
\end{equation}
Moreover,
\begin{equation}\label{eq:paired-defect-formula}
\norm{W-P_{\mathrm{pair}}(W)}_F^2
=
\frac12\sum_{i=1}^{m}(\lambda_{2i-1}-\lambda_{2i})^2
\,+\,
\begin{cases}
0, & n\text{ even},\\
\lambda_n^2, & n\text{ odd}.
\end{cases}
\end{equation}
In particular, the defect vanishes if and only if same observed-dimension skew-square closure holds.
\end{corollary}

\begin{proof}
By Proposition~\ref{prop:paired-spectrum}, the target class consists of matrices of the form $\tfrac14 K K^{\mathsf T}$ with $K^{\mathsf T}=-K$, equivalently positive semidefinite matrices whose positive eigenvalues occur in equal adjacent pairs, with one forced zero when $n$ is odd. Fix such a target matrix $Z$. For fixed eigenvalues of $Z$, the Frobenius distance
\[
\norm{W-Z}_F^2=\operatorname{tr}(W^2)+\operatorname{tr}(Z^2)-2\operatorname{tr}(WZ)
\]
is minimised when $W$ and $Z$ commute, by the symmetric form of von Neumann's trace inequality. Hence the nearest point may be taken in the eigenbasis of $W$, and the matrix problem reduces to Euclidean projection of $(\lambda_1,\dots,\lambda_n)$ onto the cone of vectors with equal adjacent pairs and final zero when $n$ is odd.

The orthogonal projection onto that linear constraint set is obtained by averaging each adjacent pair, which gives \eqref{eq:pair-projection}. Because the eigenvalues are ordered, these pair averages remain ordered and nonnegative, so the projected vector already lies in the admissible cone. This proves \eqref{eq:paired-defect-distance}. Expanding the squared norm pair by pair gives
\[
(\lambda_{2i-1}-\bar\lambda_i)^2+(\lambda_{2i}-\bar\lambda_i)^2
=\frac12(\lambda_{2i-1}-\lambda_{2i})^2,
\]
and when $n$ is odd the final coordinate is projected to zero, contributing $\lambda_n^2$. This yields \eqref{eq:paired-defect-formula}. The final claim is immediate from Proposition~\ref{prop:paired-spectrum}.\qedhere
\end{proof}

\begin{remark}[Robust empirical certificate]
Define the paired-spectrum defect by
\[
D_{\mathrm{pair}}(W):=\norm{W-P_{\mathrm{pair}}(W)}_F.
\]
Distance to a closed set is $1$-Lipschitz in Frobenius norm, so for any estimate $\widehat W$ one has
\[
\abs{D_{\mathrm{pair}}(\widehat W)-D_{\mathrm{pair}}(W)}\le \norm{\widehat W-W}_F.
\]
Hence any certified lower bound on $D_{\mathrm{pair}}(\widehat W)$ beyond estimation error is a robust witness that the visible correction does not lie in the same observed-dimension skew-square class. The defect is therefore the quantitative version of the binary obstruction in Proposition~\ref{prop:paired-spectrum}.
\end{remark}

\begin{remark}[What is and is not closed]
Theorems~\ref{thm:latent-contraction} and \ref{thm:hidden-load}, together with Proposition~\ref{prop:paired-spectrum}, show the correct closure picture. Exact quotienting preserves Gaussian precision geometry and the support-preserving attenuation class beneath a fixed tangent ceiling. It does not generically preserve the same observed-dimension Donsker-Varadhan skew-square class. The microscopic Donsker-Varadhan Hessian should therefore be regarded as the tangent ceiling and microscopic cone inside the larger quotient-visible class, not as the class closed by Schur elimination itself.
\end{remark}

\section{Hidden locality and graph-local structure}
\label{sec:05_hidden_locality}

Section~\ref{sec:04_hidden_mediation} showed that the quotient-visible object is produced by hidden mediation beneath a fixed tangent ceiling. The next question is spatial: if the upstairs precision is graph-local on the hidden sector, does the induced visible law respect that locality, or can Schur elimination generate arbitrary long-range coupling? The answer is no. Visible fill-in is organised by hidden connected components and is exponentially suppressed by hidden graph distance, with a universal rate determined by the hidden spectral window. The section closes by identifying corridor extremisers for this long-range transfer problem. Those extremisers belong to the abstract positive-cone locality theory developed here. Section~\ref{sec:09_corridor_regime} later studies a Markov corridor regime in the hidden-feedback ladder; it is inspired by the same geometry but is not the same statement.

\subsection{Componentwise hidden mediation}
\label{subsec:componentwise-hidden-mediation}

Let $V=V_{\mathrm{vis}}\oplus V_{\mathrm{hid}}$ and let
\begin{equation}\label{eq:local-block}
K_X=
\begin{pmatrix}
K_{VV} & K_{VH}\\
K_{HV} & K_{HH}
\end{pmatrix}
\in\SPD(V)
\end{equation}
be a graph-local precision operator. In this Gaussian graphical-model setting, precision sparsity encodes conditional independence and Schur elimination creates visible fill-in \cite{Dempster1972,Lauritzen1996,GeorgeLiu1981}. The point is that this fill-in is not arbitrary. It is first split by hidden connected component, and only then can one ask how large the resulting visible coupling can be.

\begin{theorem}[Hidden-component mediation]\label{thm:hidden-components}
The visible effective precision after hidden elimination is
\begin{equation}\label{eq:keff}
K_{\mathrm{eff}}=K_{VV}-K_{VH}K_{HH}^{-1}K_{HV}.
\end{equation}
If the hidden graph decomposes into connected components $H=\bigsqcup_{\alpha} H_\alpha$ respected by $K_{HH}$, then
\begin{equation}\label{eq:component-sum}
K_{VH}K_{HH}^{-1}K_{HV}
=
\sum_{\alpha} K_{V H_\alpha}K_{H_\alpha H_\alpha}^{-1}K_{H_\alpha V}.
\end{equation}
Consequently a visible pair can acquire a new coupling only if both visible coordinates connect to the same hidden component.
\end{theorem}

The decomposition in \eqref{eq:component-sum} is the Green-kernel or walk-sum mediation picture for hidden elimination \cite{DoyleSnell1984,MalioutovJohnsonWillsky2006}. It is the componentwise version of the hidden-mediation mechanism from Section~\ref{sec:04_hidden_mediation}: locality is not lost at the level of support, even though fill-in becomes dense inside the mediated visible block.

\begin{proof}
Equation~\eqref{eq:keff} is the Schur complement formula. If the hidden graph splits into connected components respected by $K_{HH}$, then after permuting indices one has
\[
K_{HH}=\diag(K_{H_1H_1},\dots,K_{H_mH_m}),
\]
so the inverse is block diagonal with the same decomposition. Substituting this into the Schur term yields \eqref{eq:component-sum}. The last claim is immediate from the support structure of the sum.
\end{proof}

\subsection{Intrinsic graph-local decay}
\label{subsec:intrinsic-graph-local-decay}

Componentwise mediation is only the first locality statement. The sharper question is quantitative: once two visible channels are forced to communicate through a distant hidden region, how strongly can they couple? The next theorem gives the universal answer. It extends classical decay-of-inverse results for sparse positive matrices to the hidden-mediation setting relevant for quotient observation \cite{DemkoMossSmith1984,BenziGolub1999,BenziRazouk2007}.

\begin{proposition}[Graph-distance decay via polynomial approximation]\label{thm:chebyshev-decay}
Let $D\succ0$ be a hidden precision matrix supported on a hidden graph $G_H$, and suppose
\[
\operatorname{spec}(D)\subset[\alpha,\beta],
\qquad 0<\alpha<\beta<\infty.
\]
Let $S,T$ be hidden vertex sets with graph distance at least $\delta$, and define
\begin{equation}\label{eq:best-poly-error}
E_\delta(\alpha,\beta):=
\inf_{\deg p\le \delta-1}
\max_{x\in[\alpha,\beta]}\abs{\frac1x-p(x)}.
\end{equation}
Then
\begin{equation}\label{eq:chebyshev-decay}
\|P_S D^{-1} P_T\|_{\mathrm{op}}\le E_\delta(\alpha,\beta).
\end{equation}
Moreover, classical Chebyshev approximation yields constants $C_{\mathrm{dec}}(\alpha,\beta)>0$ and
\[
\rho:=\frac{\sqrt{\kappa}-1}{\sqrt{\kappa}+1},
\qquad \kappa:=\beta/\alpha,
\]
for which
\begin{equation}\label{eq:chebyshev-decay-exp}
E_\delta(\alpha,\beta)\le C_{\mathrm{dec}}(\alpha,\beta)\,\rho^\delta.
\end{equation}
\end{proposition}

\begin{proof}
Let $p$ be any polynomial of degree at most $\delta-1$. Since $D^j$ propagates only along hidden walks of length at most $j$, one has
\[
P_S p(D) P_T = 0.
\]
Hence
\[
P_S D^{-1} P_T = P_S\bigl(D^{-1}-p(D)\bigr)P_T,
\]
so
\[
\|P_S D^{-1} P_T\|_{\mathrm{op}}
\le
\|D^{-1}-p(D)\|_{\mathrm{op}}
\le
\max_{x\in[\alpha,\beta]} \abs{\frac1x-p(x)}.
\]
Taking the infimum over all such $p$ gives \eqref{eq:chebyshev-decay}. The exponential estimate \eqref{eq:chebyshev-decay-exp} is the standard Chebyshev-approximation consequence for $1/x$ on $[\alpha,\beta]$ \cite{DemkoMossSmith1984,Rivlin1974,Trefethen2013}.
\end{proof}

\begin{corollary}[Visible coupling decay]\label{cor:visible-decay}
Let $b_u,b_v$ be visible-to-hidden coupling vectors supported on hidden sets at graph distance at least $\delta$. Then
\begin{equation}\label{eq:visible-decay}
\abs{b_u^T D^{-1} b_v}
\le
\norm{b_u}_2\,\norm{b_v}_2\,E_\delta(\alpha,\beta)
\le
\norm{b_u}_2\,\norm{b_v}_2\,C_{\mathrm{dec}}(\alpha,\beta)\,\rho^\delta.
\end{equation}
\end{corollary}

\begin{proof}
Apply Proposition~\ref{thm:chebyshev-decay} and Cauchy-Schwarz.
\end{proof}

\begin{remark}[Locality as a visible prediction]
Corollary~\ref{cor:visible-decay} is the headline locality statement of the programme-core theory. Once the hidden spectral window is controlled, long-range visible coupling is exponentially suppressed by hidden graph distance. This is a prediction about the visible object, not merely a structural statement about the upstairs operator.
\end{remark}

\subsection{Corridor extremisers and bundle structure}
\label{subsec:corridor-extremisers}

The Chebyshev bound gives a universal decay rate. The next theorem identifies the asymptotic graph shape that realises that rate. This is still an abstract statement about graph-local hidden precision geometry. The strong-driving corridor analysis in Section~\ref{sec:09_corridor_regime} concerns a specific hidden-feedback ladder and should not be conflated with the present extremiser theorem.

\begin{proposition}[Corridor realisation of the universal exponential rate]\label{thm:corridor}
Let $P_{\delta+1}$ be the path on $\delta+1$ hidden vertices and define the constant-coefficient corridor precision
\begin{equation}\label{eq:corridor-precision}
D_\delta^{\mathrm{corr}} := cI-dA_{P_{\delta+1}},
\qquad
c:=\frac{\alpha+\beta}{2},
\qquad
d:=\frac{\beta-\alpha}{4}.
\end{equation}
Then $\operatorname{spec}(D_\delta^{\mathrm{corr}})\subset[\alpha,\beta]$, the endpoint distance is $\delta$, and the endpoint transfer obeys
\begin{equation}\label{eq:corridor-rate}
\bigl(D_\delta^{\mathrm{corr}}\bigr)^{-1}_{1,\delta+1}
=
C_{\mathrm{corr}}(\alpha,\beta)\,\rho^\delta\,(1+o(1)),
\qquad \delta\to\infty,
\end{equation}
for an explicit positive constant $C_{\mathrm{corr}}(\alpha,\beta)$ and the same decay rate $\rho=(\sqrt\kappa-1)/(\sqrt\kappa+1)$ as in Proposition~\ref{thm:chebyshev-decay}. Consequently the one-dimensional corridor realises the universal exponential decay rate at the level of the exponent.
\end{proposition}

\begin{proof}
The spectral inclusion for \eqref{eq:corridor-precision} follows from the path adjacency spectrum. The inverse of a constant-coefficient Jacobi chain is explicit, and its endpoint entry has the stated asymptotic form \eqref{eq:corridor-rate}. Proposition~\ref{thm:chebyshev-decay} shows that no graph-local hidden precision with spectral window $[\alpha,\beta]$ can decay more slowly than exponentially at rate $\rho^\delta$. Hence the corridor family realises the universal exponent.
\end{proof}

\begin{remark}[Rate optimality versus extremiser questions]
Proposition~\ref{thm:corridor} is a rate-level statement. It identifies the exponential decay exponent and a corridor family that realises it. It does not prove a sharp finite-$\delta$ minimax constant, and it does not prove the stronger many-channel threshold-count extremiser statement. Those questions remain open in the abstract locality problem. The corridor regime in Section~\ref{sec:09_corridor_regime} should therefore be read as a separate dynamical investigation built on the same rate-level geometry, not as a proof of full finite-distance optimality.
\end{remark}

\subsection{Scope of the current locality theory}
\label{subsec:scope-locality-theory}

The locality theory proved here has two parts. First, hidden mediation splits by connected component. Second, once the hidden spectral window is fixed, off-mask visible coupling is exponentially suppressed with graph distance, and the corridor family realises the universal decay exponent asymptotically. What is not proved here is a sharp finite-distance minimax constant or a full classification of graph extremisers beyond the corridor bundle. Adjacency-loss diagnostics and fill-in summaries are deferred to Appendix~C.

\section{Bell and temporal frontiers as law-level evidence}
\label{sec:06_bell_frontier}

Sections~\ref{sec:04_hidden_mediation} and \ref{sec:05_hidden_locality} developed the internal geometry of $\Phi$: hidden mediation, transport, conservation, and graph-local structure. The next question is evidential. Does the quotient-visible viewpoint detect genuine law-level structure that simpler summaries erase? On the Bell square the answer is yes. The visible datum is not merely a correlator table. It is a family of visible pair laws, and the right compatibility question is whether those pair laws arise as quotient marginals of one common visible Gaussian law. On that square the resulting frontier has two independent coordinates: variance consistency and the arcsine CHSH obstruction. Correlator-only summaries retain only the second.

\subsection{Common visible gluing as a law-level compatibility problem}

The Bell square is the first nontrivial test of the programme-core viewpoint because it asks for compatibility of several visible quotient outputs at once. The question is not whether each pair law is individually admissible. It is whether the four pair laws can be realised simultaneously as quotient marginals of one common law upstairs on the visible variables. That is the kind of compatibility problem the quotient-visible formalism is designed to express.

\begin{definition}[Common Schur gluing of context families]\label{def:common-schur-gluing}
Let
\[
Q_{xy}\in\SPD(2),
\qquad x,y\in\{0,1\},
\]
be four visible pair precisions, and let
\[
P_{xy}:\mathbb R^4\to\mathbb R^2
\]
select the coordinate pair $(A_x,B_y)$ from $(A_0,A_1,B_0,B_1)$. We say that the family $(Q_{xy})$ is \emph{commonly Schur-glued} if there exists $Q_\ast\in\SPD(4)$ such that
\begin{equation}\label{eq:common-schur-gluing}
Q_{xy}=F_{P_{xy}}(Q_\ast)
\qquad\text{for all }x,y\in\{0,1\}.
\end{equation}
\end{definition}

Definition~\ref{def:common-schur-gluing} is the Bell-square instance of quotient compatibility. It asks whether four visible pair laws sit on one common point of the positive cone after quotienting by the four coordinate projections. The next theorem shows that, on the Bell square, this compatibility problem collapses to two law-level conditions.

\begin{theorem}[Bell-square collapse theorem]\label{thm:bell-collapse}
Let
\[
Q_{xy}=\begin{pmatrix} a_{xy} & b_{xy}\\ b_{xy} & d_{xy}\end{pmatrix}\in\SPD(2),
\qquad
\Sigma_{xy}:=Q_{xy}^{-1}=\begin{pmatrix} v^A_{xy} & c_{xy}\\ c_{xy} & v^B_{xy}\end{pmatrix},
\]
and define
\begin{equation}\label{eq:bell-rho-def}
\rho_{xy}:=\frac{c_{xy}}{\sqrt{v^A_{xy}v^B_{xy}}}=\frac{-b_{xy}}{\sqrt{a_{xy}d_{xy}}},
\qquad
\phi_{xy}:=\arcsin(\rho_{xy}).
\end{equation}
Then the following are equivalent.
\begin{enumerate}
\item There exists a centred Gaussian law on $(A_0,A_1,B_0,B_1)$ whose $(A_x,B_y)$ marginals are exactly the pair laws with covariances $\Sigma_{xy}$. In the positive definite case this is equivalent to common Schur gluing in the sense of Definition~\ref{def:common-schur-gluing}.
\item The family satisfies remote-setting variance consistency,
\begin{equation}\label{eq:variance-consistency-body}
v^A_{x0}=v^A_{x1}\quad (x=0,1),
\qquad
v^B_{0y}=v^B_{1y}\quad (y=0,1),
\end{equation}
and the four arcsine CHSH inequalities
\begin{equation}\label{eq:arcsine-chsh-body-a}
|\phi_{00}+\phi_{01}+\phi_{10}-\phi_{11}|\le \pi,
\end{equation}
\begin{equation}\label{eq:arcsine-chsh-body-b}
|\phi_{00}+\phi_{01}-\phi_{10}+\phi_{11}|\le \pi,
\end{equation}
\begin{equation}\label{eq:arcsine-chsh-body-c}
|\phi_{00}-\phi_{01}+\phi_{10}+\phi_{11}|\le \pi,
\end{equation}
\begin{equation}\label{eq:arcsine-chsh-body-d}
|-\phi_{00}+\phi_{01}+\phi_{10}+\phi_{11}|\le \pi.
\end{equation}
Equivalently, with
\begin{equation}\label{eq:bell-lifted-correlators}
E_{xy}:=\frac{2}{\pi}\phi_{xy}=\frac{2}{\pi}\arcsin\!\Bigl(\frac{-b_{xy}}{\sqrt{a_{xy}d_{xy}}}\Bigr),
\end{equation}
common gluing is equivalent to variance consistency together with Bell-locality of the lifted correlator table $E=(E_{xy})$.
\end{enumerate}
\end{theorem}

\begin{proof}
Appendix~A isolates the Gaussian gluing criterion and then specialises the Bell square using the standard arcsine lift. The new point for the present paper is structural: the quotient-visible compatibility problem is not exhausted by correlators. On the Bell square the full visible law contributes the diagonal variance-gluing constraints \eqref{eq:variance-consistency-body} and the lifted CHSH constraints \eqref{eq:arcsine-chsh-body-a}-\eqref{eq:arcsine-chsh-body-d}.
\end{proof}

The theorem should be read as support for the programme-core theory, not as an isolated Bell observation. Sections~\ref{sec:quotient-precision}-\ref{sec:04_hidden_mediation} say that visible objects must be compared at law level, because quotient observation lives on the positive cone and composes by Schur transport. The Bell square tests that claim in the first genuinely constrained multi-context setting. The answer is rigid: the compatibility frontier is two-coordinate from the outset.

\subsection{Bell-square frontier and strict refinement beyond Bell-locality}

The Bell-local shadow detects only one part of the frontier from Theorem~\ref{thm:bell-collapse}. The variance sector is a distinct obstruction, and it is invisible to correlator-only summaries.

\begin{proposition}[Strict refinement over Bell-locality]\label{prop:bell-strict-refinement}
Bell-locality of the lifted correlator table $E$ is necessary but not sufficient for common visible Gaussian gluing of the full pair laws.
\end{proposition}

\begin{proof}
Consider
\[
Q_{00}=\begin{pmatrix}1&0\\0&1\end{pmatrix},
\qquad
Q_{01}=\begin{pmatrix}1/4&0\\0&1\end{pmatrix},
\qquad
Q_{10}=Q_{11}=\begin{pmatrix}1&0\\0&1\end{pmatrix}.
\]
Then every off-diagonal entry vanishes, so \eqref{eq:bell-lifted-correlators} gives $E_{xy}=0$ for all $x,y$, hence the lifted correlator table is Bell-local. But
\[
Q_{00}^{-1}=\begin{pmatrix}1&0\\0&1\end{pmatrix},
\qquad
Q_{01}^{-1}=\begin{pmatrix}4&0\\0&1\end{pmatrix},
\]
so the variance of $A_0$ depends on Bob's setting. The variance-consistency conditions \eqref{eq:variance-consistency-body} therefore fail, and no common visible Gaussian gluing exists.
\end{proof}

\begin{remark}[Two Bell-square obstructions]\label{rem:two-bell-obstructions}
For a Bell-square family, define
\begin{equation}\label{eq:epsilon-var}
\varepsilon_{\mathrm{var}}:=\max\!\bigl\{|v^A_{00}-v^A_{01}|,|v^A_{10}-v^A_{11}|,|v^B_{00}-v^B_{10}|,|v^B_{01}-v^B_{11}|\bigr\},
\end{equation}
and let $\varepsilon_{\mathrm{CHSH}}$ be the excess of the largest left-hand side of \eqref{eq:arcsine-chsh-body-a}-\eqref{eq:arcsine-chsh-body-d} above $\pi$, truncated below at zero. Then common gluing on the Bell square is equivalent to the simultaneous vanishing of both obstructions,
\[
\varepsilon_{\mathrm{var}}=0,
\qquad
\varepsilon_{\mathrm{CHSH}}=0.
\]
This packages the Bell-square frontier as a two-coordinate compatibility problem: one coordinate is invisible to Bell correlators, while the other is the usual arcsine CHSH obstruction.
\end{remark}

\begin{corollary}[Bell cube normal form]\label{cor:bell-cube-body}
Let $S$ be the $4\times 4$ matrix whose rows are the four CHSH sign vectors
\[
(1,1,1,-1),\qquad (1,1,-1,1),\qquad (1,-1,1,1),\qquad (-1,1,1,1),
\]
and define lifted Bell coordinates
\[
y:=\frac{1}{\pi}S\phi,
\qquad
\phi:=(\phi_{00},\phi_{01},\phi_{10},\phi_{11})^\top.
\]
Then the Bell-compatible region in arcsine coordinates is exactly the cube
\[
|y_i|\le 1,
\qquad i=1,\dots,4.
\]
Equivalently, the four CHSH inequalities are the coordinate inequalities of one global hypercubic normal form. In particular, Euclidean nearest-point repair in $y$ is coordinatewise clipping, and the Bell support function is
\[
h_{\mathcal B}(u)=\frac{\pi}{4}\,\|Su\|_1.
\]
Appendix~A records the derivation and the exact repair formulas.
\end{corollary}

\begin{corollary}[Single-readout CHSH locality]\label{cor:chsh-guardrail}
Suppose a Bell-context family $(Q_{xy})$ is commonly Schur-glued, and let
\[
Z=(Z_{A_0},Z_{A_1},Z_{B_0},Z_{B_1})\sim N(0,Q_\ast^{-1})
\]
for some $Q_\ast$ satisfying \eqref{eq:common-schur-gluing}. Let
\[
A_x=f_x(Z_{A_x}),\qquad B_y=g_y(Z_{B_y}),
\qquad |f_x|\le 1,\quad |g_y|\le 1.
\]
Then the correlators
\[
\widetilde E_{xy}:=\mathbb E[A_xB_y]
\]
obey the CHSH bound
\begin{equation}\label{eq:chsh-guardrail}
\bigl|\widetilde E_{00}+\widetilde E_{01}+\widetilde E_{10}-\widetilde E_{11}\bigr|\le 2,
\end{equation}
and likewise for the other three sign choices with one minus sign.
\end{corollary}

\begin{proof}
All four outputs are measurable functions of the single common random vector $Z$. Hence for every realisation one has
\[
\bigl|A_0(B_0+B_1)+A_1(B_0-B_1)\bigr|
\le |B_0+B_1|+|B_0-B_1|\le 2,
\]
because $|A_x|\le 1$ and $|B_y|\le 1$. Taking expectations gives \eqref{eq:chsh-guardrail}, and the other sign choices follow by relabelling.
\end{proof}

\subsection{Why the quotient-visible perspective sees more than correlators}

The Bell-facing conclusion is now sharper than a CHSH guardrail. Common visible Gaussian gluing on the Bell square is a law-level compatibility problem. Bell correlators detect only the arcsine shadow of that problem, whereas the quotient-visible formalism keeps the full pair law and therefore keeps the variance sector as well. Proposition~\ref{prop:bell-strict-refinement} proves that this extra sector is not cosmetic. Two families can have identical Bell-local correlator shadow and still differ on the existence of a common quotient-visible law.

This is the kind of evidence the paper needs at this stage. Theorem~\ref{thm:bell-collapse} does not merely add another reformulation of Bell-locality. It shows that the intrinsic visible variable from Section~\ref{sec:quotient-precision} resolves frontier structure that correlator compression destroys. The Bell square is therefore a genuine law-level stress test for $\Phi$. The passage from pair laws to correlators is a further readout layer, and it loses one of the two compatibility coordinates. The quotient-visible perspective sees more because it works before that readout collapse.

Corollary~\ref{cor:chsh-guardrail} then recovers the operational statement inside the fuller law-level picture. Once common gluing exists, every single bounded local readout sits under the classical CHSH bound. The usual guardrail is therefore recovered as a downstream shadow of common visible gluing, not as the primary object.

\subsection{Temporal frontier}

The same compatibility question can be asked on an observation graph of times rather than Bell settings. One specifies visible pair laws on the edges of a time-indexed graph and asks whether they admit one common Gaussian law whose edge marginals realise the prescribed family. Trees have only repeated-marginal consistency, whereas the first genuinely temporal clique already introduces a nontrivial law-level obstruction. Appendix~A records the first fully observed temporal case and its sign-shadow interpretation. We do not develop a full temporal theory in the main text. The point here is only structural: Bell and temporal frontiers are two graph-indexed instances of one quotient-visible gluing problem.
\section{Exact symmetry and the hidden response carrier}
\label{sec:07_exact_symmetry}

We now turn from the abstract quotient-visible theory to the controlled hidden-feedback ladder family. The visible variable is the activity $a\in\{0,1\}$, the hidden variable is the memory level $m\in\{0,\dots,M\}$, and the observer sees only the activity projection. The question is the first one worth asking once the abstract theory is in place: does the intrinsic visible object $\Phi$ ever close in a completely rigid way inside a genuine nonequilibrium model? At $h=0$ and $\lambda=0$ the answer is yes. A discrete symmetry forces an exact eigenmode, the visible line becomes an exact eigenline of the reduced conjugated operator, and the scalar visible algebra collapses to one parameter.

\subsection{The exact eigenmode at $\lambda = 0$}

Write the ladder generator in the split form
\begin{equation}\label{eq:ladder-generator-split}
Q = \alpha\bigl(Q_{\mathrm{vis}} + \eps Q_{\mathrm{mem}}\bigr),
\end{equation}
where $Q_{\mathrm{vis}}$ changes activity and $Q_{\mathrm{mem}}$ changes memory only. At $h=0$ and $\lambda=0$ the visible flip rates are symmetric, so the lumped visible chain has rates $\beta_0=\beta_1=\alpha$. Define the activity sign vector
\begin{equation}\label{eq:activity-sign-vector}
u(a,m) := \mathbf 1_{\{a=0\}} - \mathbf 1_{\{a=1\}}.
\end{equation}
This vector is constant along each memory fibre and records only the visible activity parity.

\begin{theorem}[Exact symmetry at the symmetric point]\label{thm:exact-symmetry-point}
Consider the hidden-feedback ladder at $h=0$ and $\lambda=0$, with scalar visible observation given by activity. Then the following statements hold exactly for every finite $M$ and every finite $\eps$.
\begin{enumerate}[label=(\roman*)]
\item The activity sign vector $u$ satisfies
\begin{equation}\label{eq:Qu-eigenmode}
Q u = -2\alpha u.
\end{equation}
\item In the visible-first reduced tangent basis, the visible line is an exact right eigenline of the reduced conjugated generator with eigenvalue $-2\alpha$.
\item The quotient-visible algebraic precision satisfies
\begin{equation}\label{eq:phi-equals-alpha}
\Phi = \alpha.
\end{equation}
\item Writing $F_0$ for the reversible visible backbone, $W:=\Phi-F_0$ for the irreversible correction, and $T$ for the scalar tangent load on the visible line, one has
\begin{equation}\label{eq:exact-symmetry-scalars}
v^{\mathsf T} H_0 v = \alpha,
\qquad
W\alpha = F_0 T.
\end{equation}
\end{enumerate}
\end{theorem}

\begin{proof}
Memory moves preserve activity and therefore preserve the sign vector, so $Q_{\mathrm{mem}}u=0$. At $h=0$ and $\lambda=0$, every visible flip changes the sign and occurs at rate $\alpha$, hence $Q_{\mathrm{vis}}u=-2u$. Multiplying by the prefactor in \eqref{eq:ladder-generator-split} gives \eqref{eq:Qu-eigenmode}. The reduction to the visible-first tangent basis preserves the visible line generated by the activity mode, so the induced reduced operator has the same visible eigenvalue. In scalar visible rank, the generator-level Schur identity from Section~\ref{subsec:scalar-schur} identifies $2\Phi$ with the visible Schur complement of $-L_{\mathrm{red}}$. Because the visible line is exact and carries eigenvalue $-2\alpha$, that Schur complement equals $2\alpha$, which gives \eqref{eq:phi-equals-alpha}. The identities in \eqref{eq:exact-symmetry-scalars} are then the scalar Schur relations for the decomposition of the reduced conjugated operator into reversible backbone plus skew correction.
\end{proof}

The theorem is the centre of Part III. Nothing asymptotic has happened yet. The visible clock is already pinned to a constant, and the remaining freedom is only in how that fixed total precision is split between reversible and irreversible sectors.

\begin{corollary}[Single-parameter reduction and frozen clock]\label{cor:single-parameter-reduction}
On the slice $h=0$ and $\lambda=0$, the scalar visible algebra is determined by the single parameter $T(\eps,M)$. More precisely,
\begin{equation}\label{eq:single-parameter-reduction}
F_0 = \frac{\alpha^2}{T+\alpha},
\qquad
W = \frac{\alpha T}{T+\alpha},
\qquad
F_0 + W = \alpha.
\end{equation}
Consequently $\Phi=\alpha$ is frozen across all $\eps$ and $M$, while $F_0$ decreases and $W$ increases monotonically as functions of $T$.
\end{corollary}

\begin{proof}
By Theorem~\ref{thm:exact-symmetry-point}, one has $\Phi=\alpha$ and $W\alpha=F_0T$. Together with the exact conservation split from Theorem~\ref{thm:conservation-split}, namely $F_0+W=\Phi$, this gives
\[
(\alpha-F_0)\alpha = F_0T.
\]
Solving for $F_0$ and then substituting into $W=\alpha-F_0$ yields \eqref{eq:single-parameter-reduction}. The monotonicity statements are immediate from
\[
\frac{dF_0}{dT} = -\frac{\alpha^2}{(T+\alpha)^2} < 0,
\qquad
\frac{dW}{dT} = \frac{\alpha^2}{(T+\alpha)^2} > 0.
\]
\end{proof}

\begin{remarkplain}[Trading segment and corridor endpoint]
Corollary~\ref{cor:single-parameter-reduction} puts every point on the symmetric $\lambda=0$ slice onto the line segment $F_0+W=\alpha$. The equilibrium endpoint is $(F_0,W)=(\alpha,0)$. Section~\ref{sec:09_corridor_regime} shows that the strong-corridor endpoint is $(\alpha/(M+1),\alpha M/(M+1))$. The total visible precision does not move. Only its internal partition moves.
\end{remarkplain}

\begin{remarkplain}[Why the eigenmode exists]
The mechanism behind Theorem~\ref{thm:exact-symmetry-point} is a $\mathbb Z_2$ activity-exchange symmetry. At $h=0$ and $\lambda=0$, exchanging the two activity classes leaves the generator invariant and makes the sign representation one-dimensional. The mode $u$ is therefore forced, not guessed. This is also the reason the phenomenon transfers to any binary-visible model whose generator has the same activity-exchange symmetry.
\end{remarkplain}

\subsection{The visible-first Schur reduction}

The previous theorem identifies the correct visible line. Once that line is fixed, the correct reduced basis is forced: first the true visible tangent vector, then an orthonormal basis of its hidden complement. The point is structural, not cosmetic. The first-correction analysis lives in this basis because only in this basis does the visible eigenline become a block constraint.

Let $L_{\mathrm{red}}(\lambda)$ denote the reduced conjugated tangent operator in a visible-first orthonormal basis, and write its block form as
\begin{equation}\label{eq:Lred-block-general}
L_{\mathrm{red}}(\lambda)=
\begin{pmatrix}
a(\lambda) & r(\lambda)^{\mathsf T}\\
c(\lambda) & \mathcal H(\lambda)
\end{pmatrix},
\end{equation}
where the first coordinate is the visible line and the remaining coordinates form the hidden complement.

\begin{proposition}[Visible-first Schur reduction at $\lambda=0$]\label{prop:visible-first-schur}
For arbitrary field $h$, provided $\lambda=0$, the visible line is the exact tangent line of the lumped two-state visible chain and is an exact right eigenline of the reduced conjugated generator. Equivalently,
\begin{equation}\label{eq:Lred-block-triangular}
L_{\mathrm{red}}(0)=
\begin{pmatrix}
a_0 & r_0^{\mathsf T}\\
0 & \mathcal H_0
\end{pmatrix},
\qquad
a_0 = -(\beta_0+\beta_1).
\end{equation}
In particular, the lower-left block vanishes exactly, so all hidden-to-visible correction passes through the single Schur channel determined by $r_0$, $\mathcal H_0$, and the derivative of the lower-left block.
\end{proposition}

\begin{proof}
Exact lumpability at $\lambda=0$ identifies the true visible tangent direction with the activity mode of the two-state visible chain. Writing the reduced conjugated operator in a basis whose first vector spans that line, the eigenline condition means exactly that the image of the first basis vector has no hidden-complement component. This is equivalent to $c(0)=0$, which yields the block upper-triangular form \eqref{eq:Lred-block-triangular}. The visible scalar block is the negative total visible escape rate of the lumped chain, namely $-(\beta_0+\beta_1)$.
\end{proof}

Proposition~\ref{prop:visible-first-schur} is the basis theorem for the correction layer. It explains why the naive odd-pair basis is not the right global coordinate system for the fast-memory expansion. The visible object responds through the true visible tangent line, not through an arbitrary parity chart. Section~\ref{sec:08_fast_memory} uses this exact block form as the starting point for the first-correction expansion.

\subsection{The canonical hidden response carrier}

Once the visible-first block form is exact, the first hidden derivative that can influence the visible Schur coefficient is forced. There is only one scalar contraction left.

\begin{theorem}[Canonical hidden response carrier]\label{thm:canonical-hidden-carrier}
Assume $\lambda=0$ and write the visible-first reduced operator as in \eqref{eq:Lred-block-general}, with $c(0)=0$ as in Proposition~\ref{prop:visible-first-schur}. Then the odd hidden contribution to the first derivative of the scalar visible Schur coefficient is carried by the single scalar
\begin{equation}\label{eq:Jhid-definition}
J_{\mathrm{hid}} := -\frac{\eps}{2}\, r_0^{\mathsf T}\mathcal H_0^{-1} c'(0).
\end{equation}
Equivalently, all competing hidden derivative channels are killed by the exact upper-triangular reduction at $\lambda=0$.

Moreover, once the visible line is fixed, $J_{\mathrm{hid}}$ is invariant under arbitrary orthogonal changes of hidden complement. Thus $J_{\mathrm{hid}}$ is canonical inside the visible-first Schur setting.
\end{theorem}

\begin{proof}
For scalar visible rank, the visible Schur coefficient of \eqref{eq:Lred-block-general} is
\begin{equation}\label{eq:scalar-schur-lambda}
S(\lambda) = a(\lambda) - r(\lambda)^{\mathsf T}\mathcal H(\lambda)^{-1}c(\lambda).
\end{equation}
Differentiating at $\lambda=0$ and using $c(0)=0$ gives
\begin{equation}\label{eq:scalar-schur-derivative}
S'(0) = a'(0) - r_0^{\mathsf T}\mathcal H_0^{-1}c'(0).
\end{equation}
Every term involving $r'(0)$, $\mathcal H'(0)$, or derivatives of the inverse multiplied by $c(0)$ vanishes. The only hidden contraction that survives is therefore $r_0^{\mathsf T}\mathcal H_0^{-1}c'(0)$. The normalization in \eqref{eq:Jhid-definition} is the one natural for the ladder scaling. For canonicality, let $O$ be an orthogonal change of basis on the hidden complement. Then
\[
r_0^{\mathsf T}\mapsto r_0^{\mathsf T}O^{\mathsf T},
\qquad
\mathcal H_0\mapsto O\mathcal H_0 O^{\mathsf T},
\qquad
c'(0)\mapsto O c'(0),
\]
so the scalar contraction is unchanged:
\[
r_0^{\mathsf T}O^{\mathsf T}(O\mathcal H_0 O^{\mathsf T})^{-1}Oc'(0)
= r_0^{\mathsf T}\mathcal H_0^{-1}c'(0).
\]
This proves both claims.
\end{proof}

\begin{remarkplain}[Meaning of canonical]
The word ``canonical'' is used here in a precise and limited sense. The scalar $J_{\mathrm{hid}}$ is not a basis-free universal observable on every reduced model. It is canonical once the visible line has been fixed and the reduction is written in the visible-first Schur basis forced by Proposition~\ref{prop:visible-first-schur}. Inside that setting there is exactly one hidden derivative contraction, and it is invariant under all residual orthogonal freedom.
\end{remarkplain}

\subsection{What this does and does not solve}

Section~\ref{sec:07_exact_symmetry} solves the exact algebraic centre of the correction analysis at $\lambda=0$. It identifies the correct reduction basis, proves that the visible precision freezes to $\Phi=\alpha$ on the symmetric slice, and collapses the scalar visible algebra to the single function $T(\eps,M)$. What it does not solve is equally important. It does not show that the eigenmode or the identity $\Phi=\alpha$ survive when $h\neq0$ or $\lambda\neq0$, it does not determine the fast-memory obstruction, and it does not explain the corridor regime. Those are the jobs of Sections~\ref{sec:08_fast_memory} and \ref{sec:09_corridor_regime}. The mechanistic Woodbury termination behind the multiplicative identity remains open even though the algebraic identity itself is proved. Finally, the canonical carrier identified here is shown in Section~\ref{sec:08_fast_memory} to be the unique survivor among the plausible first-correction reduction routes.

\section{Fast-memory response and thermodynamic constraints}
\label{sec:08_fast_memory}

The exact carrier from Section~\ref{sec:07_exact_symmetry} answers the symmetry question at the algebraic centre. The next question is asymptotic. When hidden memory equilibrates much faster than the visible activity flips, what visible law survives, and how does hidden nonequilibrium re-enter at the first nontrivial order? The answer has two parts. First, the visible law homogenises to a two-state chain, so the leading visible precision is completely explicit. Second, the first correction already remembers hidden geometry through one Schur channel, and that channel comes with both a source law and a sharp list of failures of the obvious simplifications.

\subsection{Fast-memory homogenisation}

Write $k_{a\to 1-a}(m)$ for the visible flip rate out of activity $a$ at memory level $m$, and let $\nu_a$ denote the stationary law of the pure memory dynamics on the fibre with activity fixed to $a$. The corresponding homogenised visible exit rates are
\begin{equation}\label{eq:beta-hom-def}
\beta_a:=\sum_{m=0}^M \nu_a(m)\,k_{a\to 1-a}(m),
\qquad a\in\{0,1\}.
\end{equation}
These are the rates seen after the memory sector has fully mixed before the next visible jump.

\begin{proposition}[Fast-memory homogenisation]\label{prop:fast-memory-homogenisation}
Fix finite $M$ and fixed parameters $(\mu,\lambda,h)$. As $\eps\to\infty$, the visible activity process homogenises to a two-state continuous-time Markov chain with rates $\beta_0$ and $\beta_1$ from \eqref{eq:beta-hom-def}. In particular, the leading visible dwell law is exponential in each activity sector.
\end{proposition}

\begin{proof}
This is the standard fast-memory timescale-separation mechanism, used here as a verified asymptotic input rather than as a standalone homogenisation proof. At timescale $\eps^{-1}$ the memory generator relaxes inside each activity fibre before a visible flip typically occurs. The visible jump hazard therefore averages against the fibre equilibrium $\nu_a$, which yields the effective exit rates $\beta_a$. Once those rates are fixed, the homogenised visible law is the two-state chain with generator
\[
Q_{\mathrm{hom}}=
\begin{pmatrix}
-\beta_0 & \beta_0\\
\beta_1 & -\beta_1
\end{pmatrix}.
\]
The exponential visible dwell law is then immediate.
\end{proof}

The first point of contact with the quotient-visible theory is that the homogenised limit is already scalar-visible, so the visible precision can be read off directly.

\begin{corollary}[Homogenised visible precision]\label{cor:phi-hom}
For the homogenised two-state chain from Proposition~\ref{prop:fast-memory-homogenisation},
\begin{equation}\label{eq:phi-hom}
\Phi_{\mathrm{hom}}=\frac{\beta_0+\beta_1}{2}.
\end{equation}
Moreover, along the ladder family,
\begin{equation}\label{eq:phi-fast-limit}
\Phi\to \Phi_{\mathrm{hom}},
\qquad
W\to 0,
\qquad \eps\to\infty.
\end{equation}
Thus the irreversible contribution is subleading in the fast-memory limit.
\end{corollary}

\begin{proof}
For a two-state visible chain with rates $a$ and $b$, the scalar visible Schur computation gives $\Phi=(a+b)/2$. Applying this to $Q_{\mathrm{hom}}$ proves \eqref{eq:phi-hom}. The convergence statement is the visible consequence of Proposition~\ref{prop:fast-memory-homogenisation}: once the visible law converges to the homogenised two-state chain, the scalar visible precision converges to its exact two-state value, while the hidden irreversible correction disappears at leading order.
\end{proof}

The basis theorem from Section~\ref{sec:07_exact_symmetry} is already active here. The correct asymptotic basis is not an arbitrary parity chart. It is the visible-first basis forced by Proposition~\ref{prop:visible-first-schur}, because that basis keeps the true visible tangent line fixed while the hidden block diverges.

\subsection{First-correction structure}

The homogenised limit is too coarse to explain how hidden nonequilibrium first re-enters the visible precision. That information appears one order later. The visible-first basis from Section~\ref{sec:07_exact_symmetry} is the right coordinate system because it separates scalar visible drift from hidden-mediated feedback before any asymptotic truncation is taken.

The first object that enters is the correct entry law for a visible dwell. If the process is about to enter activity $a$, it does not enter according to the within-activity equilibrium $\nu_a$. It enters according to jump-weighted mass coming from the opposite activity.

\begin{proposition}[Leading fast-memory entry law]\label{prop:fast-memory-entry-law}
Define
\begin{equation}\label{eq:eta-entry-law}
\eta_a(m):=\frac{\nu_{1-a}(m)\,k_{1-a\to a}(m)}{\beta_{1-a}},
\qquad a\in\{0,1\}.
\end{equation}
Then $\eta_a$ is the leading fast-memory entry law for dwells in activity $a$ in the fast-memory asymptotic regime. In particular, the exact entry law differs from $\eta_a$ by $O(\eps^{-1})$, and at fixed $(M,\mu,\lambda)$ the law $\eta_a$ is independent of the uniform field $h$.
\end{proposition}

\begin{proof}
A dwell in activity $a$ begins with a jump from the opposite activity, so the entering memory distribution is weighted by the source equilibrium $\nu_{1-a}$ multiplied by the jump hazard into $a$. Normalisation by the total inflow rate $\beta_a$ gives \eqref{eq:eta-entry-law}. The $O(\eps^{-1})$ correction is the usual first transient defect before complete fibre mixing. The field-independence is a normalisation cancellation: the uniform field multiplies both numerator and denominator in the same way and therefore drops out of the leading ratio.
\end{proof}

The next input is spectral. Let $T_{a,\eps}=\eps G_a-\diag(k_{a\to 1-a})$ denote the killed memory operator in activity sector $a$.

\begin{proposition}[Principal eigenvalue shift]\label{prop:fast-memory-eigenvalue}
For each activity sector,
\begin{equation}\label{eq:lambda-max-fast-memory}
\lambda_{\max}(T_{a,\eps})=-\beta_a+\frac{c_{\mathrm{eig},a}}{\eps}+o(\eps^{-1}),
\qquad \eps\to\infty.
\end{equation}
The first correction to the visible mean dwell is therefore not purely an effective-rate effect: the leading eigenvalue shift is only one contribution to the full first-order visible correction.
\end{proposition}

\begin{proof}
At leading order the operator $\eps G_a$ mixes to the fibre equilibrium, so the killed generator sees only the averaged hazard $\beta_a$. Standard first-order spectral perturbation around that averaged hazard gives \eqref{eq:lambda-max-fast-memory}. The final sentence records that the entry-law and transient-relaxation terms survive at the same order, so the eigenvalue shift alone does not close the visible correction.
\end{proof}

The visible-first Schur basis now exposes the full first coefficient. Write the reduced conjugated operator of $-L$ in that basis as
\begin{equation}\label{eq:fast-memory-block-expansion}
-L_{\mathrm{red},\eps}=
\begin{pmatrix}
2\Phi_{\mathrm{hom}}+a_1\eps^{-1}+o(\eps^{-1}) & r_0^{\mathsf T}+o(1)\\
c_0+o(1) & \eps\mathcal H_1+o(\eps)
\end{pmatrix}.
\end{equation}
The hidden block diverges like $\eps$, but the visible line stays fixed. This is precisely why the Schur feedback survives at order $\eps^{-1}$.

\begin{theorem}[First fast-memory correction in visible-first form]\label{thm:fast-memory-first-correction}
Assume the expansion \eqref{eq:fast-memory-block-expansion}. Then the scalar visible precision satisfies
\begin{equation}\label{eq:phi-first-correction}
\Phi_{\mathrm{hom}}-\Phi
=
\frac{1}{2\eps}\Bigl[-a_1+r_0^{\mathsf T}\mathcal H_1^{-1}c_0\Bigr]+o(\eps^{-1}).
\end{equation}
Thus the first visible correction splits into a scalar drift term and a hidden-mediated Schur term. At $\lambda=0$, the second term is exactly the fast-memory manifestation of the canonical hidden carrier from Theorem~\ref{thm:canonical-hidden-carrier}.
\end{theorem}

\begin{proof}
For scalar visible rank, Section~\ref{subsec:scalar-schur} identifies $2\Phi$ with the visible Schur complement of $-L_{\mathrm{red},\eps}$. Applying the scalar Schur formula to \eqref{eq:fast-memory-block-expansion} gives
\[
2\Phi
=
2\Phi_{\mathrm{hom}}+a_1\eps^{-1}-r_0^{\mathsf T}(\eps\mathcal H_1)^{-1}c_0+o(\eps^{-1}).
\]
Since $(\eps\mathcal H_1)^{-1}=\eps^{-1}\mathcal H_1^{-1}$, rearranging yields \eqref{eq:phi-first-correction}. The identification with the carrier from Section~\ref{sec:07_exact_symmetry} is exactly the statement that the only hidden-mediated channel surviving in visible-first coordinates is the canonical Schur contraction.
\end{proof}

Equation~\eqref{eq:phi-first-correction} is the technical centre of the section. It shows why the first correction cannot be read from homogenised rates alone. The visible coefficient remembers hidden resolvent geometry through the single Schur channel $r_0^{\mathsf T}\mathcal H_1^{-1}c_0$.

\subsection{Source law and obstruction}

We now specialise to the symmetric coupling slice $\lambda=0$, where the activity-swap involution is exact. Section~\ref{sec:07_exact_symmetry} isolated the canonical hidden carrier $J_{\mathrm{hid}}$. The present question is how that carrier varies with the field $h$ in the fast-memory regime.

\begin{theorem}[Parity and exact perturbation decomposition]\label{thm:field-parity-decomposition}
Assume $\lambda=0$ and fix the visible-first basis from Proposition~\ref{prop:visible-first-schur}. Write
\[
r(h)=r_0+\Delta r(h),
\qquad
c'_h=c'(0)+\Delta c(h),
\qquad
\mathcal H(h)=\mathcal H_0+\Delta\mathcal H(h),
\]
and let
\[
R(h):=\mathcal H(h)^{-1}-\mathcal H_0^{-1}.
\]
Then the canonical hidden carrier
\[
J_{\mathrm{hid}}(h):=-\frac{\eps}{2}\,r(h)^{\mathsf T}\mathcal H(h)^{-1}c'_h
\]
is an even function of $h$. The normalised datum
\begin{equation}\label{eq:Jstar-def}
J_*(h):=\frac{J_{\mathrm{hid}}(h)}{1+\cosh(h)}
\end{equation}
is also even. Moreover, the field defect
\[
\Delta J_{\mathrm{hid}}:=J_{\mathrm{hid}}(h)-J_{\mathrm{hid}}(0)
\]
admits the exact decomposition
\begin{equation}\label{eq:Jh-defect-decomp}
\Delta J_{\mathrm{hid}}=T_r+T_c+T_H+T_{\mathrm{mix}},
\end{equation}
where
\begin{align}
T_r&:=-\frac{\eps}{2}\,(\Delta r)^{\mathsf T}\mathcal H_0^{-1}c'(0),\label{eq:Tr-def}\\
T_c&:=-\frac{\eps}{2}\,r_0^{\mathsf T}\mathcal H_0^{-1}\Delta c,\label{eq:Tc-def}\\
T_H&:=-\frac{\eps}{2}\,r_0^{\mathsf T}R(h)c'(0),\label{eq:TH-def}
\end{align}
and
\begin{align}\label{eq:Tmix-def}
T_{\mathrm{mix}}:={}&-\frac{\eps}{2}\Bigl[(\Delta r)^{\mathsf T}\mathcal H_0^{-1}\Delta c
+(\Delta r)^{\mathsf T}R(h)c'(0)
+r_0^{\mathsf T}R(h)\Delta c
+(\Delta r)^{\mathsf T}R(h)\Delta c\Bigr].
\end{align}
Here $T_r$ is row-tail drift, $T_c$ is input-channel drift, $T_H$ is hidden-resolvent drift, and $T_{\mathrm{mix}}$ collects the mixed interaction terms.
\end{theorem}

\begin{proof}
The evenness is the activity-swap symmetry from Section~\ref{sec:07_exact_symmetry}: replacing $h$ by $-h$ exchanges the two activity classes while leaving the visible-first scalar contraction invariant, so $J_{\mathrm{hid}}(h)=J_{\mathrm{hid}}(-h)$. Since $1+\cosh(h)$ is itself even, the same holds for $J_*(h)$.

For the decomposition, expand
\[
r(h)^{\mathsf T}\mathcal H(h)^{-1}c'_h
=
\bigl(r_0+\Delta r\bigr)^{\mathsf T}\bigl(\mathcal H_0^{-1}+R(h)\bigr)\bigl(c'(0)+\Delta c\bigr)
\]
and subtract the $h=0$ contribution $r_0^{\mathsf T}\mathcal H_0^{-1}c'(0)$. Grouping the surviving terms by whether the perturbation enters through the row, the input channel, the hidden resolvent, or a mixed product yields \eqref{eq:Jh-defect-decomp} with \eqref{eq:Tr-def}-\eqref{eq:Tmix-def}.
\end{proof}

The point of Theorem~\ref{thm:field-parity-decomposition} is structural. There is a clean source law, but there is also a genuine obstruction. After normalisation by $1+\cosh(h)$, the carrier becomes asymptotically field-blind at leading fast-memory order, yet not exactly field-blind at finite $\eps$.

\begin{proposition}[Asymptotic source law and obstruction class]\label{prop:source-law-obstruction}
Assume $\lambda=0$. Then in the fast-memory regime,
\begin{equation}\label{eq:Jstar-field-blind}
J_*(h)=J_*(0)+O\!\left(\frac{\cosh(h)-1}{\eps}\right),
\end{equation}
and hence
\begin{equation}\label{eq:Jstar-obstruction-h2}
J_*(h)-J_*(0)=O\!\left(\frac{h^2}{\eps}\right)
\qquad (h\to 0).
\end{equation}
In the same regime the odd first-correction coefficients obey the asymptotic redistribution laws
\begin{equation}\label{eq:odd-redistribution-scalar}
-\frac{c_{\mathrm{scalar}}(h)}{\cosh(h)}\sim J_*(0),
\qquad
\frac{c_{\mathrm{schur}}(h)}{1+\cosh(h)}\sim J_*(0),
\qquad \eps\to\infty.
\end{equation}
Thus the source law is asymptotically clean, while the residual defect sits in an even obstruction class of order $O(h^2\eps^{-1})$.
\end{proposition}

\begin{proof}
Because $J_*$ is even, its small-field defect must begin at quadratic order in $h$. Fast-memory scaling pushes the coefficient of that even defect down by one power of $\eps^{-1}$, which yields \eqref{eq:Jstar-field-blind} and \eqref{eq:Jstar-obstruction-h2}. The redistribution laws follow from the visible-first Schur decomposition of the odd first-correction slope: the scalar part carries a $\cosh(h)$ factor, while the hidden Schur part carries a $(1+\cosh(h))$ factor. After those factors are divided out, both coefficients approach the same source datum $J_*(0)$.
\end{proof}

At larger depth the same mechanism persists in a parity-resolved form. The fast-memory correction separates into a dominant odd part and a smaller negative even defect; the odd response is linear to first order in small $\lambda$; and, after parity decomposition, both parity sectors are extensive in depth. That is the sense in which the hidden carrier from Section~\ref{sec:07_exact_symmetry} really does organise the fast-memory layer. It does not solve the whole correction problem, but it isolates the unique dominant odd channel.

\subsection{What the first correction is not}

The first-correction formula is informative partly because of what it excludes. Every obvious scalar shortcut was tested against the ladder family, and each one fails for a structural reason: the visible coefficient remembers hidden resolvent geometry, not just a few coarse summaries.

First, the carrier itself is not recovered by replacing the exact stationary law with the naive product-limit stationary law. The product-limit picture gets the scaling intuition right, but it does not reproduce the exact hidden contraction at finite $\eps$.

Second, the coefficient
\begin{equation}\label{eq:Cphi-def}
C_\Phi:=\eps\bigl(\Phi_{\mathrm{hom}}-\Phi\bigr)
\end{equation}
is not determined by the homogenised rate pair $(\beta_0,\beta_1)$ alone. Nor is it captured by simple summaries of the entry laws $\eta_a$ together with those rates. The reason is visible in Theorem~\ref{thm:fast-memory-first-correction}: the term $r_0^{\mathsf T}\mathcal H_1^{-1}c_0$ depends on hidden resolvent data that no low-dimensional rate summary can see.

Third, the first-order visible dwell correction is not determined solely by the principal eigenvalue shift \eqref{eq:lambda-max-fast-memory}. Entry-law and transient-relaxation contributions survive at the same order. Likewise, in the large-$M$ fast-memory regime the odd correction density does not collapse to simple Poisson-energy scalars such as $\nu(u_1^2)$, $\nu(k u_1^2)$, $\eta(u_1)$, or $\eta(u_1^2)$. Those quantities do not encode the hidden resolvent geometry of the Schur term.

Fourth, the field defect from Proposition~\ref{prop:source-law-obstruction} does not admit a one-slot reduction. The exact decomposition \eqref{eq:Jh-defect-decomp} proves that row-tail drift alone, input-channel drift alone, and hidden-resolvent drift alone are each insufficient. The mixed channel is not optional. For the same reason, the negative even defect is not generated by continuing the odd carrier $J_{\mathrm{hid}}$ to second order in $\lambda$.

Finally, the finite-forcing odd redistribution law is not exact on the full parameter set. The asymptotic source law is real, but away from the fast-memory limit higher-order finite-forcing terms spoil exact factorisation. This is why the source law is a leading-order theorem rather than a full finite-parameter identity.

Taken together, these exclusions are not side remarks. They are the rigidity part of the correction theory. They show that the first visible correction cannot be reduced to homogenised rates, entry-law moments, product-limit stationarity, one-slot field drift, or a continuation of the odd carrier. The visible object is already genuinely higher-order at first subleading order.

\subsection{Thermodynamic constraints on the visible object}

The correction layer is not unconstrained. The same ladder family that reveals the source law also puts a thermodynamic ceiling on how large the visible object can become. The cleanest summary is one line.

\begin{corollary}[Thermodynamic ceiling]\label{cor:thermodynamic-ceiling}
Assume the verified half-bound $W\le \sigma/2$, where $\sigma$ is the total entropy production rate. Then
\begin{equation}\label{eq:thermo-ceiling}
\Phi\le F_0+\frac{\sigma}{2}.
\end{equation}
\end{corollary}

\begin{proof}
By the exact conservation split from Theorem~\ref{thm:conservation-split},
\[
\Phi=F_0+W.
\]
Combining this with the verified half-bound $W\le \sigma/2$ yields \eqref{eq:thermo-ceiling}.
\end{proof}

The ceiling inherits the status of the bound behind it. Within the ladder family the stronger verified statement is
\begin{equation}\label{eq:epr-fisher-sum}
W+T\le \sigma,
\end{equation}
with the individual half-bounds
\begin{equation}\label{eq:epr-fisher-half}
W\le \frac{\sigma}{2},
\qquad
T\le \frac{\sigma}{2}.
\end{equation}
These inequalities were verified across the full tested parameter grid, including nonzero $\lambda$. They are programme-core constraints, not a peculiarity of one asymptotic slice.

There is also an exact symmetry statement on the simplest slice. At $M=1$, $\lambda=0$, and $h=0$, all visible Fisher quantities are invariant under the timescale duality
\begin{equation}\label{eq:eps-duality}
\eps\longmapsto \eps^{-1}.
\end{equation}
In particular,
\begin{equation}\label{eq:eps-duality-components}
F_0(\eps)=F_0(\eps^{-1}),
\qquad
T(\eps)=T(\eps^{-1}),
\qquad
W(\eps)=W(\eps^{-1}),
\qquad
\Phi(\eps)=\Phi(\eps^{-1}).
\end{equation}
Since the unique fixed point of the duality is $\eps=1$, the irreversible correction is maximised there on that slice.

\begin{proposition}[Complete $\eps$-profile at $M=1$]\label{prop:eps-profile-M1}
At $M=1$, $\lambda=0$, and $h=0$, the quantities $F_0$, $T$, $W$, and $\Phi$ are symmetric in $\log\eps$ around $\eps=1$. Moreover,
\begin{equation}\label{eq:T-peak-eps1}
T(\eps)\le T(1)
\end{equation}
for all $\eps>0$, and the maximum of $W$ occurs at $\eps=1$.
\end{proposition}

\begin{proof}
The symmetry in $\log\eps$ is exactly the duality \eqref{eq:eps-duality-components}. The fixed point of the duality is $\eps=1$, so any unique extremum must occur there. The complete profile follows once this is combined with the verified tightness at both timescale extremes recorded below.
\end{proof}

The duality is exceptional. It fails cleanly once either of its structural hypotheses is removed: at $M\ge 2$ the odd sector is no longer two-dimensional, and at $\lambda\neq0$ the activity-exchange symmetry is broken. Those failure modes matter because they show that the duality is a rigid symmetry, not a numerical accident.

The asymptotic tightness statements are the two endpoint companions to the ceiling. In the slow-memory regime,
\begin{equation}\label{eq:slow-memory-tightness}
1-\frac{W}{T}=O(\eps),
\qquad \eps\to 0,
\end{equation}
and in the fast-memory regime,
\begin{equation}\label{eq:fast-memory-tightness}
1-\frac{W}{T}=O(\eps^{-1}),
\qquad \eps\to \infty.
\end{equation}
So the tangent envelope becomes tight in both timescale extremes, from opposite sides.

Finally, the correction layer has a clean near-equilibrium and small-driving structure. Near detailed balance, with plaquette affinity $A_\square\to 0$,
\begin{equation}\label{eq:near-equilibrium-scaling}
W\asymp (A_\square^2)^2,
\qquad
T\asymp (A_\square^2)^2,
\qquad
\sigma\asymp (A_\square^2)^2.
\end{equation}
Thus the ratios $W/\sigma$ and $T/\sigma$ remain finite in the linear-response regime. But $\Phi$ is not determined by $A_\square$ alone. The full Schur object depends on more than the scalar thermodynamic force. Likewise, at small driving with $\lambda=0$ one has
\begin{equation}\label{eq:small-mu-symmetry}
W=T+O(\mu^4),
\qquad \mu\to 0,
\end{equation}
so the tangent and irreversible corrections coincide to leading order before separating at quartic order.

\subsection{Rigor boundary}

Section~\ref{sec:08_fast_memory} contains four different epistemic layers, and they should be kept distinct. The homogenised two-state visible precision \eqref{eq:phi-hom} and the visible-first Schur expansion \eqref{eq:phi-first-correction} are exact once the asymptotic block expansion is granted. The parity statement for $J_{\mathrm{hid}}$ and the perturbation decomposition \eqref{eq:Jh-defect-decomp} are exact on the slice $\lambda=0$. The field-obstruction class, the EPR-Fisher bounds, the $M=1$ duality, the endpoint tightness laws, and the scaling claims are verified asymptotic or verified structural statements in the frozen ladder programme. The open pieces are the sharp obstruction constants, the operator that generates the even defect, and analytic proofs of the EPR-Fisher bound and the $M=1$ timescale duality.

Two proof routes are visible already. First, an analytic proof of $W+T\le \sigma$ would likely follow if the tangent quantity $T$ can be identified as a current-type precision in the sense required by the thermodynamic uncertainty relation. Second, the corridor decay transition discussed in Section~\ref{sec:09_corridor_regime} looks accessible by a direct spectral-gap analysis of the path-kernel restricted generator. Both are genuine research directions rather than present results.

\section{Corridor regime and the regime boundary}
\label{sec:09_corridor_regime}

Section~\ref{sec:08_fast_memory} analysed the fast-memory direction $\eps\to\infty$. We now turn a different knob. Fix $h=0$ and $\lambda=0$, keep the exact symmetry backbone from Section~\ref{sec:07_exact_symmetry}, and drive the bias $\mu$ large. A second asymptotic mechanism appears. Its natural carrier is not the visible-first Schur channel of the fast-memory layer, but an odd-sector path kernel with an endpoint Schur complement. The section has two jobs. First, it records the corridor asymptotics cleanly. Second, it proves that this corridor sector is genuinely different from the fast-memory response sector. That separation is a result, not a matter of taste. \par

\subsection{The strong-driving limit}
\label{subsec:strong-driving-limit}

At $h=0$ and $\lambda=0$, Section~\ref{sec:07_exact_symmetry} gave the exact scalar backbone $\Phi=\alpha$. The corridor regime asks how the decomposition of that fixed visible precision changes when $\mu\to\infty$. The answer is that the reversible part collapses to an endpoint law of order $O(M^{-1})$, while the tangent load grows linearly in depth.

\begin{proposition}[Corridor endpoint law and refined asymptotics]\label{thm:corridor-path-kernel}
Fix $h=0$, $\lambda=0$, and finite depth $M$. In the strong-driving limit $\mu\to\infty$, the odd sector is governed by a path-kernel reduction whose visible output is an endpoint Schur complement. The endpoint law is
\begin{equation}\label{eq:corridor-F0-limit}
F_0(M,\mu)\to \frac{\alpha}{M+1}.
\end{equation}
Consequently, using the exact scalar identity $\Phi=\alpha$ on the same slice,
\begin{equation}\label{eq:corridor-T-limit}
T(M,\mu)\to \alpha M,
\end{equation}
\begin{equation}\label{eq:corridor-WT-limit}
\frac{W}{T}(M,\mu)\to \frac{1}{M+1},
\end{equation}
and
\begin{equation}\label{eq:corridor-W-limit}
W(M,\mu)\to \frac{\alpha M}{M+1}.
\end{equation}
In the archived ladder computations, this convergence is consistent with the refined asymptotic form
\begin{equation}\label{eq:corridor-T-refined}
T(M,\mu)=\alpha M\tanh^2\!\left(\frac{\mu}{4}\right)+O\!\left(e^{-\gamma_M\mu}\right),
\qquad
\gamma_M=
\begin{cases}
1, & M=2,\\
\tfrac12, & M\ge 3.
\end{cases}
\end{equation}
\end{proposition}

\begin{proof}
The endpoint law \eqref{eq:corridor-F0-limit} is the corridor endpoint Schur computation on the odd-sector path kernel. Section~\ref{sec:07_exact_symmetry} gives the exact identity $\Phi=\alpha$ on the present slice, so
\[
W=\Phi-F_0=\alpha-F_0.
\]
Taking the limit in \eqref{eq:corridor-F0-limit} gives \eqref{eq:corridor-W-limit}. The ratio limit \eqref{eq:corridor-WT-limit} and the tangent limit \eqref{eq:corridor-T-limit} are then equivalent through \eqref{eq:corridor-W-limit}. The refined form \eqref{eq:corridor-T-refined} is the verified corridor spectral pattern from the archived ladder computations.
\end{proof}

\begin{corollary}[Trading relation and corridor endpoint]\label{cor:corridor-trading-endpoint}
On the slice $h=0$ and $\lambda=0$, the pair $(F_0,W)$ lies on the exact line
\begin{equation}\label{eq:corridor-trading-segment}
F_0+W=\alpha.
\end{equation}
Its equilibrium endpoint is $(\alpha,0)$ and its corridor endpoint is
\begin{equation}\label{eq:corridor-endpoint-pair}
\left(\frac{\alpha}{M+1},\frac{\alpha M}{M+1}\right).
\end{equation}
No monotonicity or no-overshoot statement in $\mu$ is claimed here.
\end{corollary}

\begin{proof}
The exact identity $F_0+W=\Phi=\alpha$ holds on the whole $h=0$, $\lambda=0$ slice by Section~\ref{sec:07_exact_symmetry}. The equilibrium endpoint is the small-driving limit $\mu\to 0$, and the corridor endpoint is given by Proposition~\ref{thm:corridor-path-kernel}.
\end{proof}

\begin{remark}[Different asymptotic direction, different carrier]
The fast-memory limit and the corridor limit are independent asymptotic directions in parameter space. One sends $\eps\to\infty$ at fixed $(M,\mu,h,\lambda)$ and keeps the visible-first Schur carrier. The other sends $\mu\to\infty$ on the exact-symmetry slice and replaces that carrier by an endpoint Schur law on the odd-sector path kernel. The asymptotic formulas may both involve the same scalar visible objects $(F_0,T,W,\Phi)$, but they are produced by different internal reductions.
\end{remark}

\subsection{Why this is not the same regime}
\label{subsec:not-the-same-regime}

The point of the corridor analysis is not merely that one gets a second formula. The point is that the second formula comes from a different operator sector. The next statement is the load-bearing boundary theorem for Part III.

\begin{proposition}[Verified regime separation in the current ladder family]\label{thm:regime-separation}
Within the current hidden-feedback ladder analysis, the fast-memory response sector and the large-$\mu$ corridor sector do not behave as one operator regime in disguise. The separation has three independent components.

First, the two sectors use different natural bases: the fast-memory expansion is carried by the true visible tangent vector completed to a visible-first orthonormal basis, whereas the corridor reduction is carried by the odd pair basis.

Second, the two sectors use different carrier geometries: the fast-memory odd response is organised by the canonical hidden Schur carrier
\[
J_{\mathrm{hid}}=-\frac{\eps}{2}r_0^{\mathsf T}\mathcal H_0^{-1}c'(0),
\]
while the corridor sector is organised by an endpoint path-kernel carrier on the odd chain.

Third, the two sectors have different scaling classes on strong-bias slices: the fast-memory odd hidden density remains $O(1)$ in $M$, whereas the corridor visible reversible law is $O(M^{-1})$.

In addition, for $M\ge 2$ at $\lambda=0$ and $\eps=1$, the tangent correction does not factorise as
\begin{equation}\label{eq:no-separable-cMg}
T(M,\mu)=c(M)g(\mu).
\end{equation}
So even inside the corridor sector, depth and driving remain genuinely entangled.
\end{proposition}

\begin{proof}
The basis statement is the contrast between Section~\ref{sec:08_fast_memory}, where the visible-first Schur basis is forced by the true visible tangent line, and Theorem~\ref{thm:corridor-path-kernel}, where the strong-driving reduction lives on the odd pair sector. The carrier statement is immediate from the constructions: Section~\ref{sec:08_fast_memory} is driven by the canonical hidden Schur scalar $J_{\mathrm{hid}}$, while the corridor reduction is an endpoint Schur complement of a path Green kernel. The scaling-class statement follows by comparing the verified large-$M$ odd-density class from the fast-memory package with \eqref{eq:corridor-F0-limit}, which is of order $O(M^{-1})$ at fixed large $\mu$. The non-factorisation claim is exactly the observed corridor negative boundary: the dependence of $T(M,\mu)/\tanh^2(\mu/4)$ on $\mu$ varies with $M$ for every tested $M\ge2$, so no decomposition of the form \eqref{eq:no-separable-cMg} can hold on that domain.
\end{proof}

\begin{remark}[What the theorem does not say]
Theorem~\ref{thm:regime-separation} does not deny that a higher containing architecture may exist. It says something narrower and stronger: the present fast-memory and corridor sectors are already provably distinct at the level of basis, carrier, and scaling class. Any future programme would therefore have to contain both sectors without collapsing these differences.
\end{remark}

\subsection{The regime boundary}
\label{subsec:regime-boundary}

The present theory is therefore genuinely multi-regime. Section~\ref{sec:07_exact_symmetry} gives the exact scalar backbone at $h=0$ and $\lambda=0$. Section~\ref{sec:08_fast_memory} gives the visible-first Schur carrier in the fast-memory direction. The present section gives the endpoint path-kernel carrier in the strong-driving direction. These are all real pieces of the same correction programme, but they are not yet one unified operator sector.

\begin{proposition}[Current regime boundary]\label{prop:current-regime-boundary}
The current correction theory consists of at least two distinct asymptotic sectors beyond the exact-symmetry backbone: the fast-memory sector of Section~\ref{sec:08_fast_memory} and the corridor sector of Theorem~\ref{thm:corridor-path-kernel}. A higher operator architecture containing both sectors without collapsing their distinctions remains open.
\end{proposition}

\begin{proof}
The existence of the two sectors follows from Section~\ref{sec:08_fast_memory} and Theorem~\ref{thm:corridor-path-kernel}. Their non-collapse is exactly Theorem~\ref{thm:regime-separation}. The remaining statement is therefore the open containing-architecture problem.
\end{proof}

The strength of Proposition~\ref{prop:current-regime-boundary} is that it says exactly enough. It does not pretend that the present paper has a single master kernel. It proves the boundary instead. That is the honest state of the theory.

\subsection{Large-$M$ scaling}
\label{subsec:large-M-scaling}

The strong-driving corridor endpoint is only one part of the depth story. The fast-memory package also revealed a depth-normalised odd hidden density whose behaviour across $M$ is remarkably stable. That object is not the corridor carrier, but it is the main large-$M$ datum that survives the present correction analysis, so it belongs here at the regime boundary.

\begin{proposition}[Verified large-$M$ odd-density class]\label{prop:large-M-odd-density}
At $h=0$ and on tested $\mu$ slices, the normalised odd hidden density is consistent with the asymptotic class
\begin{equation}\label{eq:largeM-density-class}
\frac{J_*(M,\mu)}{M-1}=L(\mu)+\frac{A(\mu)}{M}+O\!\left(\frac{1}{M^2}\right).
\end{equation}
This is finite-limit evidence in a stable asymptotic class, not a proved limit theorem.
\end{proposition}

\begin{proof}
The statement summarises the verified large-depth fits: after the extensive factor $(M-1)$ is removed, the data are consistent with a stable linear-in-$1/M$ correction and a smaller residual. The point here is not to claim a proof of limit existence, but to record the stable asymptotic form supported by the current computations.
\end{proof}

\begin{proposition}[Crossover at $\mu=1$ is not a singular threshold]\label{prop:mu-one-crossover}
Across the tested large-$M$ window, the approach direction in \eqref{eq:largeM-density-class} changes across $\mu=1$: it is from below for $\mu<1$, from above for $\mu>1$, and nearly flat near $\mu=1$. At present there is no numerical evidence that $\mu=1$ is a singular threshold in the odd hidden density itself.
\end{proposition}

\begin{proof}
The sign of the leading $1/M$ correction changes across $\mu=1$, which forces the approach direction to switch. At the same time, the underlying scans remain smooth in a neighbourhood of $\mu=1$, so the available evidence supports a crossover centre rather than a singular carrier threshold.
\end{proof}

Two open problems remain attached to this large-$M$ sector. First, the actual limit function $L(\mu)$ is not yet identified analytically. Second, the special role of $\mu=1$ is only phenomenological. We know it marks the centre of the observed crossover in approach direction, but we do not yet know whether it carries a deeper operator meaning. Those are the natural questions left by the present data.


\section{Benchmark synthesis and evidence}
\label{sec:10_benchmark_synthesis}

Parts~I-III built the quotient-visible theory and then opened the hidden-feedback ladder family in enough detail to expose its multi-regime correction architecture. The remaining question is methodological as well as external: while algebraic closure was still being pinned down, did the theory stay tied to a real signal under deliberately varied tests? The four archived studies were the sequence of tests used to answer that question. They were not run after the final closure layer was in hand, and the paper should say so openly. Their role was to keep the developing theory sharply grounded while the algebraic picture was still being fixed. That is also why the main text does not pretend they already use the final closure method. Read at their true historical strength, they do exactly what Part~IV needs. One study supplies an exact transformed semi-Markov anchor, one supplies same-visible-law completion tests, one supplies a decoupling benchmark in which visible precision survives while observed entropy production vanishes, and one supplies the exact visible-law family that eventually made the completion picture and closure push possible.

\subsection{Why these four tests matter}

These four tests probe four different structural predictions, and they do so in roughly the order the theory itself was being stabilised. The Kapustin-Ghosal-Bisker transformed semi-Markov test asks whether exact transformed visible structure can exist in a finite partially observed system and whether the exact compression mechanism is genuinely structural rather than numerical \citep{KapustinGhosalBisker2024}. The matched-completion tests ask whether the total algebraic visible quantity survives when the hidden completion is changed without changing the visible law; in the archive this includes both abstract matched completions and a real-model sanity check built from the Liepelt-Lipowsky six-state kinesin chemomechanical network \citep{LiepeltLipowsky2007}. The Banerjee-Kolomeisky-Igoshin proofreading test asks whether visible precision can remain nonzero after coarse observation kills the observed entropy-production signal on a biologically motivated kinetic proofreading model \citep{BanerjeeKolomeiskyIgoshin2017}. The hidden-feedback ladder test family asks whether visible memory, hidden dissipation, and visible algebraic correction can be separated inside one exact visible-law family with genuine two-way hidden feedback. Together they test exact structure, invariance, decoupling, and frontier behaviour while also documenting the route by which the theory was kept empirically disciplined before final algebraic closure was achieved.


\subsection{Exact transformed semi-Markov anchor}

The Kapustin-Ghosal-Bisker transformed semi-Markov benchmark \citep{KapustinGhosalBisker2024} remains the cleanest exact transformed anchor in the archive. On the canonical four-state benchmark, the transformed observer has an exact sequence-level gain, the transformed waiting-time contribution vanishes exactly, and the exact compact observer is the four-symbol second-order parity lift recorded in the archived theorem spine. That exact compression mechanism is not merely approximate: it retains the transformed quantity at machine precision on the benchmark, while the threshold-2 observer remains near-exact but not exact. The same archive also carries the right negative contrast. In the flashing ratchet test there is no comparable compact exact lift, and the best compact retentions remain far below the exact four-state parity benchmark.

The algebraic layer extracted from the same test is also scientifically useful. On the benchmark four-state system it reports
\[
\begin{aligned}
F_0 &= 5.4104119541, & \Phi &= 5.4286720632,\\
W &= 0.0182601092, & T &= 4.7306689027,\\
\frac{W}{T} &= 0.003860.
\end{aligned}
\]
So the exact transformed anchor sits in an overwhelmingly backbone-dominated regime. This test therefore does two jobs at once. At the semi-Markov level it provides the exact transformed benchmark promised in Appendix~D. At the algebraic level it gives a small but nonzero visible correction on a benchmark whose transformed structure is known exactly.

\subsection{Same-visible-law completion tests}

These matched-completion tests probe the claim that the total scalar visible quantity survives same-visible-law completion changes even when the hidden architecture is radically altered. Its strongest archived result is the four-state versus fourteen-state same-law comparison:
\[
\Phi_{4\text{-state}}=5.4286720632,
\qquad
\Phi_{14\text{-state}}=5.4286720632,
\qquad
|\Delta\Phi|=2.7\times 10^{-15}.
\]
This is the exact kind of test Section~3 asks for. The visible law is held fixed while the hidden phase representation is enlarged, and the total visible algebraic precision survives to machine precision. The decomposition beneath that total does \emph{not} stay fixed. In the archived outputs the four-state system has
\[
F_0=5.4104119541,
\qquad W=0.0182601092,
\qquad T=4.7306689027,
\]
whereas the fourteen-state completion has
\[
F_0=4.1993218800,
\qquad W=1.2293501833,
\qquad T=5.1336332970.
\]
So this test demonstrates the right invariant and the right non-invariants in the same comparison.

The archive also contains a second independent same-law check at the renewal level. Duplicating one hidden state to produce an augmented five-state system preserves the visible semi-Markov law exactly and leaves the visible correction unchanged to machine precision,
\[
W_{\mathrm{orig}}=0.0230962835,
\qquad
W_{\mathrm{dup}}=0.0230962835,
\qquad
|\Delta W|=8.88\times 10^{-16}.
\]
That second check matters because it rules out the possibility that the fourteen-state match is a one-off feature of a particular phase-type construction.


The same archive also contains a useful real-model sanity check based on the six-state Liepelt-Lipowsky kinesin chemomechanical network, observed through a two-class D-versus-T readout \citep{LiepeltLipowsky2007}. For that test the archived values are
\[
W=7.3411775678,
\qquad
T=9.9969400245,
\qquad
W/T=0.734342.
\]
This is not a same-law completion test. It is a reminder that once the hidden topology really carries a cycle current visible to the chosen partition, the quotient-visible correction can be large.

\subsection{Banerjee-Kolomeisky-Igoshin proofreading test}

The Banerjee-Kolomeisky-Igoshin proofreading test probes a different prediction, namely that visible precision need not track observed entropy production after coarse observation on a biological error-correction network \citep{BanerjeeKolomeiskyIgoshin2017}. Under the default closure, all three coarse observers in the archived Banerjee benchmark retain exactly zero observed entropy production, but the quotient analysis reports nonzero visible precision for all four presets. On the scalar E-only observer the archived values are
\[
17.5036,\quad 16.2444,\quad 15.7350,\quad 387.9618,
\]
for the WT, HYP, ERR, and T7 presets respectively. On the correct-versus-wrong observer the corresponding traces are
\[
31.3115,\quad 22.1071,\quad 37.3374,\quad 811.0236.
\]
So in this test zero observed entropy production emphatically does not imply zero visible precision.

The test also tells us what kind of decoupling this is. The archived irreversibility fractions remain modest,
\[
W/T\in[0.070,0.258],
\]
across the four presets. The visible algebraic object is therefore still backbone-dominated. What survives the coarse observation is not a large visible irreversible sector but a finite local precision object supported by hidden structure that the coarse entropy-production observer cannot see.

This test archive also contains the right negative control. When the closure is changed in a way that changes the visible semi-Markov law, completion independence fails exactly as it should. The archive compares a five-state turnover closure to a seven-state product closure and finds visibly different $F_{\mathrm{vis}}$ values for the same coarse observer. That failure is a feature, not a bug. It shows the test is probing the correct invariance statement: same visible law, not arbitrary closure change.


\subsection{Hidden-feedback ladder test family}

The hidden-feedback ladder test family has a dual status in the paper. It is a benchmark family in its own right, and it is also the system through which the completion picture and later closure push became possible. It therefore has to be framed more carefully than the other three archives. The archived ladder runs predate the final algebraic-closure stage, but they already give an exact visible-law family of the right kind: under the activity-only observer the visible process is an exact alternating semi-Markov process with phase-type dwell laws. The promoted atlas then separates baseline Markovity, recorder-only hidden dissipation, modulator-only visible memory, two-way adaptive feedback, and equilibrium-feedback reversible controls. Historically, this is the family that kept the theory anchored while the algebraic closure mechanism was still being isolated.

The equilibrium-feedback line is the cleanest control in the archive. The points labelled \texttt{eq\_feedback\_weak} and \texttt{eq\_feedback\_strong} both have $A_\square=0$ and entropy production at numerical zero, but their visible dwell coefficients of variation are $1.094$ and $1.294$ respectively. So the archive contains memory-rich visible laws with effectively zero plaquette-level dissipation. At the algebraic layer both points satisfy
\[
W=0,
\qquad
F_0=\Phi,
\]
and therefore act as reversible memory controls for the correction theory.

The recorder and modulator points then separate hidden cost from visible memory. The archived \texttt{pure\_recorder} point has
\[
\begin{aligned}
\mathrm{EPR} &= 0.3618, & cv &= 1,\\
F_0 &= 0.86768, & \Phi &= 1.00000,\\
W &= 0.13232, & \frac{W}{T} &= 0.86768,
\end{aligned}
\]
so it carries substantial algebraic visible correction without visible non-Markovity. By contrast the archived \texttt{pure\_modulator} point has
\[
\begin{aligned}
\mathrm{EPR} &= 0.0210, & cv &= 1.0754,\\
F_0 &= 0.49433, & \Phi &= 0.49925,\\
W &= 0.00492, & \frac{W}{T} &= 0.54210,
\end{aligned}
\]
so it carries visible memory with only a very small algebraic correction. That single pair is already enough to show that visible memory, hidden cost, and visible correction are three different coordinates in this family.

The adaptive and depth-sensitive points push the family to the current boundary. The archived \texttt{adaptive\_strong} point has $cv\approx 1.280$ and $W/T\approx 0.216$; \texttt{adaptive\_slow} has smaller dissipation but remains visibly non-Markov; and the $M=8$ \texttt{depth\_sensitive\_hidden\_resolution} point changes the visible law only moderately while remaining noticeably nonlinear and numerically stiff. The frozen nonlinear report is already scientifically informative, but it explicitly classifies several promoted points as optimisation-sensitive rather than fully bankable closure constants. That is the correct boundary to state in the main text, especially because these runs were part of the route toward closure rather than post-closure victory laps. The exact backbone and the promoted algebraic layer are bankable now; the nonlinear layer is qualitatively supportive and quantitatively selective.


\subsection{Cross-test synthesis}

\begin{enumerate}[label=(\roman*)]
\item The Kapustin-Ghosal-Bisker transformed semi-Markov test shows that exact transformed visible structure can exist in a finite benchmark and can admit an exact compact observer, but the ratchet contrast shows that this is a structural mechanism, not a generic compression miracle.
\item The matched-completion tests show that at scalar visible rank the total visible algebraic precision can survive drastic same-law completion changes even while the backbone-correction split moves substantially underneath it.
\item The Banerjee-Kolomeisky-Igoshin proofreading test shows that visible precision and observed entropy production can decouple sharply under coarse observation, exactly in the direction predicted by the quotient-visible theory.
\item The hidden-feedback ladder test family shows that visible memory, hidden dissipation, and visible algebraic correction are genuinely distinct coordinates and can be separated inside one exact visible-law family.
\item Across all four archives, the honest reading is methodological as well as mathematical: these were the tests used to keep the developing theory tied to a real signal while algebraic closure was still being isolated. That is why they matter historically, and it is why the main text reads them at their archived strength rather than pretending they were post-closure reruns.
\item The current boundary is not absence of structure but controlled selectivity: the exact and algebraic layers are already strong, whereas the nonlinear branch is informative only where the archived optimisation diagnostics say it is stable enough to trust.
\end{enumerate}



\section{Synthetic falsification battery and blind benchmark}
\label{sec:11_falsification}

Section~\ref{sec:10_benchmark_synthesis} showed that the theory remains coherent across independent benchmark tests grounded in transformed semi-Markov inference, biological proofreading, and molecular-motor coarse observation \cite{KapustinGhosalBisker2024,BanerjeeKolomeiskyIgoshin2017,LiepeltLipowsky2007}. This section applies a stricter standard. The point is not further illustration, but deliberate adversarial pressure on the sharpest structural claims. We therefore separate two forms of stress. First, a synthetic falsification battery targets the exposed theorem-level predictions directly. Second, a blind digital benchmark asks whether locked predictions made from restricted visible data survive reveal of the hidden world.

\subsection{Falsification principles and synthetic battery}

The synthetic battery was designed to attack the claims that would most clearly fail if the theory were tracking numerical artefacts rather than genuine structure: quartic bridge scaling, support preservation beneath quotient observation, the paired-spectrum obstruction to fixed-dimension closure, graph-local decay, and the corridor-versus-bundle extremiser picture. The tests were not tuning exercises. Each family was built to try to break a definite prediction.

\begin{table}[H]
\centering
\begin{tabularx}{\textwidth}{>{\raggedright\arraybackslash}p{3.1cm} >{\raggedright\arraybackslash}p{4.0cm} Y Y}
\toprule
Diagnostic & Tested regime & Quantitative outcome & Interpretation\\
\midrule
Quartic bridge slope & $H_0$-adapted bridge sweeps with visible and hidden dimensions $2$ to $4$ & Grand mean slope $3.9957$; nine dimension-pair means from $3.9936$ to $3.9977$ & Consistent with the quartic hidden-mediation bridge and stable across the tested dimension range.\\
Worst bridge-gap eigenvalue & Same adapted bridge sweeps & Worst reported bridge-gap eigenvalue $-1.17\times 10^{-15}$ & Numerically nonnegative up to floating-point tolerance, as required by the sign-definite quartic bridge law.\\
Support preservation & Same adapted bridge sweeps & Rank-preservation rate $100\%$; maximum reported principal-angle sine $5.58\times 10^{-8}$ & Finite hidden sectors preserve tangent support to numerical precision, matching the support-rigidity theorem.\\
Same-dimension closure & Generic quotient Gaussian ensembles with visible dimensions $3$ to $8$ & Failure fraction $1.0$ throughout; pairable fraction $0.0$ throughout & Generic same-dimension skew-square closure fails completely in the tested ensemble, matching the paired-spectrum obstruction.\\
Hidden-local decay bound & Random local hidden graphs of size $24$, $36$, and $48$ & Bound-compliance rate $100\%$ across all sampled size-distance cells & Visible long-range transfer behaves as a controlled hidden-local tail, consistent with the graph-local decay theory.\\
Topology and bundle tests & Paths, cycles, ladders, trees, and disjoint corridor bundles & Corridor geometry realises the sharp decay class; disjoint bundles give flat singular spectra; under the tested finite-distance normalisation the best one-channel performer was the cycle & Supports the rate-level corridor picture and the flat-spectrum many-channel extremiser law while refining the finite-distance prefactor statement.\\
\bottomrule
\end{tabularx}
\caption{Synthetic falsification highlights. Each row is a deliberate attempt to break a sharp structural prediction.}
\label{tab:falsification_battery}
\end{table}

The outcomes are clear. Across the adapted bridge sweeps the quartic slope stays pinned near $4$, the bridge-gap eigenvalue remains at floating-point scale, and support preservation holds to numerical precision. Across visible dimensions $3$ to $8$, generic same-dimension skew-square closure fails in every sampled case. Across the hidden-local graph families, every tested transfer entry obeys the decay bound. The one substantive refinement forced by the synthetic battery is finite-distance rather than structural: the one-channel topology statement should be made at the level of universal exponential rate, not universal finite-distance constant, because looped topologies can improve prefactors under the tested spectral-window normalisation.

\subsection{Blind digital benchmark}

The blind benchmark was designed as a lock-then-reveal protocol. The hidden world was fixed first, only a restricted visible sector was exposed, predictions were committed before reveal, and only then was the concealed world released for scoring. The point of this protocol was not to enlarge the theorem domain of the paper. It was to check whether the existing finite tangent law, quotient-visible attenuation geometry, and Bell-frontier compatibility theory can function as an instrument when the hidden world is genuinely withheld.

Each digital world coupled three exact layers. First, a finite-state nonequilibrium backbone fixed the reduced detailed-balance geometry $H_0$, the skew object $\widehat J$, and the finite Donsker-Varadhan signal $\Delta_{DV}=\tfrac14\widehat J H_0^{-1}\widehat J^\top$ on the observer space. Second, an adapted visible sector $U$ fixed the tangent ceiling $T=P_U\Delta_{DV}P_U|_U$, and the actual quotient-visible correction was generated inside the proved attenuation class
\[
X=T^{1/2}(I+\Lambda)^{-1}T^{1/2},\qquad \Lambda\succeq 0.
\]
Third, Bell-context families were attached in three exact types: commonly glueable families, Bell-local but variance-inconsistent families, and variance-consistent families with arcsine-CHSH violation. Hidden truth and visible challenge artefacts were separated on disk, and the analysis stage read only the visible manifest before reveal.

On the exact-mode hero world, the relative errors were $2.61\times 10^{-2}$ for the retained-signal curve $K_d$ and $1.67\times 10^{-3}$ for the detectability curve $\Theta_d$. On the quotient side, the minimum hidden dimension and contraction-volume scores were recovered exactly in the recorded score files, with identity residuals at machine scale. The stronger test is the frozen 256-world exact ensemble. Across that ensemble the mean relative error was $1.461\times 10^{-1}$ for the finite $K$-curve and $2.541\times 10^{-2}$ for the detectability curve $\Theta_d$. Quotient scores again remained exact on the recorded ensemble, Bell verdict agreement was $1.0$, locality compatibility matched in $92.58\%$ of worlds, and motif agreement was $45.70\%$. We therefore treat locality only as a class-level diagnostic reader, not as an exact hidden-graph reconstruction claim.

Two cautions are essential. First, the sampled track is a robustness layer, not the primary claim. On the sampled hero world the finite module remained stable and quotient scores again remained exact, but the sampled Bell layer missed glueable families, so the sampled benchmark should be read as interface support rather than as the cleanest statement of the law-level frontier. Second, adaptive reveal is already informative but should still be described cautiously. On a 64-world exact panel, oracle theorem-target reveal improved the mean second-step raw revealed signal over a random baseline by $6.39\times 10^{-2}$ and the corresponding detectability score by $1.64\times 10^{-2}$. We therefore treat adaptive reveal as a diagnostic of what the observation ladder can target next, not yet as a deployed blind acquisition policy. Detailed benchmark diagnostics are collected in Appendix~\ref{app:blind-benchmark}.

\subsection{What survived and what did not}

The falsification pressure was productive because it did not leave the theory unchanged. What survived is the structural core. The quartic bridge law, support preservation, the paired-spectrum obstruction, and the hidden-local decay bound all survived direct adversarial testing. The exact-track blind benchmark also survived in the sense that the finite tangent law, quotient-visible attenuation geometry, and Bell-frontier verdicts remained operational under a real lock-then-reveal protocol.

What did not survive unchanged are the stronger informal readings that were never entitled to theorem status. The one-channel topology claim had to be weakened from a universal finite-distance statement to a universal rate statement. Motif reconstruction did not become an exact hidden-graph reader. The sampled Bell layer is not yet a clean benchmark instrument. And adaptive reveal remains a promising diagnostic policy rather than a finished acquisition procedure. That is the right outcome for this section. The theory survives where it claimed structural necessity, and it bends precisely where the data says only a weaker statement is justified.

\section{Scope, guardrails, and transfer boundaries}
\label{sec:12_scope}

This section states the status boundaries of the paper in one place. Its purpose is not rhetorical. It is to tell the reader exactly which statements are proved, which are asymptotic theorems, which are verified but not yet proved, and which questions remain open. The division below is the contract under which the rest of the paper should be read.

\subsection{What is proved: the abstract algebraic theory}

The abstract quotient-visible theory developed in Sections~\ref{sec:quotient-precision}-\ref{sec:06_bell_frontier} is proved at the level stated in the theorem bodies. For any quotient observation on the positive cone, the paper proves the variational characterisation of quotient precision, its homogeneity, monotonicity, and concavity, and its exact composition along towers of surjections. In that sense the central visible object $\Phi$ is not an ad hoc diagnostic but the intrinsic visible variable of quotient observation.

The same abstract theory also proves the internal algebraic structure of the visible object once a tangent ceiling has been fixed. Hidden mediation preserves support, every support-preserving visible law in the interior beneath the ceiling is parametrised by a positive hidden-load operator, sequential mediation transports that load by an exact Schur law, the determinant clock is the unique continuous orthogonally invariant additive interior clock, and the visible precision splits exactly into the reversible contribution and the irreversible correction through the conservation identity $\Phi=F_0+W$. These statements are exact structural results, not asymptotic approximations.

The locality and compatibility results proved in Sections~\ref{sec:05_hidden_locality} and \ref{sec:06_bell_frontier} belong to the same exact quotient-visible theory. Graph-local hidden precision yields componentwise mediation and exponential graph-distance suppression of long-range visible coupling, with corridor architectures realising the universal asymptotic rate class. On the Bell square, common visible Gaussian gluing is proved to be equivalent to remote-setting variance consistency together with Bell-locality of the arcsine lift, and this yields the single-readout CHSH guardrail as a downstream consequence of common gluing. These results hold at the level of quotient-visible law geometry itself; they do not depend on the hidden-feedback ladder analysis of Sections~\ref{sec:07_exact_symmetry}-\ref{sec:09_corridor_regime}.

\subsection{What is proved and verified: the nonequilibrium analysis}

The nonequilibrium analysis in Sections~\ref{sec:07_exact_symmetry}-\ref{sec:09_corridor_regime} is carried out for the hidden-feedback ladder family. Within that family, several statements are exact. At $h=0$ and $\lambda=0$, the activity sign mode is an exact generator eigenmode, the visible tangent line is an exact eigenline of the reduced conjugated generator, the scalar visible precision satisfies $\Phi=\alpha$ exactly, and the symmetric-point identities linking $\Phi$, $F_0$, $W$, and $T$ hold exactly. At $\lambda=0$ for arbitrary field, the true visible line remains exact, the visible-first reduced basis is exact, and the odd first-correction slope is carried by the canonical hidden scalar $J_{\mathrm{hid}}$, which is invariant under orthogonal changes of hidden complement.

The same ladder analysis also contains proved asymptotic statements. In the strong-bias corridor regime at $h=0$ and $\lambda=0$, the reversible visible law, the tangent correction, and the ratio $W/T$ obey the explicit corridor asymptotics stated in Section~\ref{sec:09_corridor_regime}. Those are genuine asymptotic theorems for that regime. They are not merely numerical trends, but they do not imply that the corridor carrier and the fast-memory carrier are one and the same object.

A larger part of the ladder analysis is verified rather than proved. This includes fast-memory homogenisation to a two-state visible chain, convergence of $\Phi$ to the homogenised value, decay of $W$ in the fast-memory tail, the first-correction structure beyond the leading homogenised rates, the field-normalised source law and its obstruction class, the large-$M$ normalised odd sector, the timescale-duality profile on the $M=1$, $\lambda=0$, $h=0$ slice, the EPR-Fisher bounds, endpoint tightness in the slow-memory and fast-memory limits, near-equilibrium scaling, and the small-driving symmetry of $W$ and $T$. Throughout the paper these statements are presented as verified asymptotic or verified structural facts, not as already-promoted theorems.

The benchmark studies of Sections~\ref{sec:10_benchmark_synthesis} and \ref{sec:11_falsification} do not enlarge the theorem boundary. Their role is methodological. They show that the developing algebraic picture was kept tied to a real signal while closure was being pinned down, and they apply adversarial pressure to the visible object and its diagnostics. They support transfer and robustness, but they are not used to convert verified statements into proved ones.

\subsection{What remains open}

Completion independence of $\Phi$ beyond the proved $M=1$ case remains open, and this matters because a full visible-law descent theorem would identify the exact minimal visible carrier of the algebraic quantity beyond the scalar benchmark. The multiplicative mechanism behind the $\lambda=0$ identities is identified but not yet proved analytically in full Schur-termination form. The field-law package does not yet have sharp operator constants beyond the present obstruction class $O((\cosh h-1)/\eps)$.

The large-$M$ odd sector has a stable asymptotic class in the tested range, but a proved limit theorem and an explicit identification of the limit function remain open. The status of the crossover near $\mu=1$ is also open at the operator level: the numerical evidence identifies a change in approach direction, but not yet a deeper structural interpretation. The operator generating the negative even defect has not yet been identified, and its relation, if any, to the odd hidden carrier is not known.

At the level of regime architecture, it remains open whether there exists a higher operator framework that contains both the fast-memory and corridor sectors without erasing their distinction. The EPR-Fisher bound and the $M=1$ timescale duality are currently verified rather than analytically proved. The full nonlinear Markov re-embedding problem after quotient observation also remains open. Finally, the transfer from the present discrete quotient-visible architecture to any genuine continuum or dynamical gravity theory is open; the paper proves structural parallels only, not a continuum field theory.

\subsection{Transfer question}

The abstract quotient-visible theory of Sections~\ref{sec:quotient-precision}-\ref{sec:06_bell_frontier} is independent of the hidden-feedback ladder. Its statements are formulated for quotient observation on positive precision operators and therefore transfer at that level without reference to any one benchmark family. The ladder enters only when the paper asks what the visible object looks like inside a nonequilibrium system and how the first correction organises itself.

The analysis of Sections~\ref{sec:07_exact_symmetry}-\ref{sec:09_corridor_regime} is ladder-specific in its present form. The exact symmetry results, the fast-memory source law, the obstruction package, the corridor asymptotics, and the regime-separation statements are established for that family. Sections~\ref{sec:10_benchmark_synthesis} and \ref{sec:11_falsification} then test partial transfer: the archived kinesin and proofreading studies show that the algebraic observables remain meaningful under controlled biochemical observers, the transformed semi-Markov benchmark shows that the visible programme survives a nontrivial observer change, and the falsification battery shows that the diagnostics are not tuned to one toy topology. What those tests do not yet prove is that the ladder-specific correction architecture transfers unchanged to every binary-visible nonequilibrium system.

The paper therefore draws a sharp line. The quotient-visible theory transfers as abstract geometry. The correction architecture is presently established in one benchmark family and partially pressure-tested beyond it. The transfer question is real, and the paper treats it as open where it has not been proved.

\subsection{What the present paper does not claim}

The paper does \textit{not} claim a full ontological unification of the fast-memory and corridor regimes.

It does not claim a continuum field theory, a dynamical gravity theory, or an Einstein-equation analogue.

It does not claim a nonlinear Markov re-embedding theorem after quotient observation.

It does not claim a general visible-law descent theorem beyond the proved $M=1$ scalar case.

\section{Discussion}
\label{sec:13_discussion}

\subsection{Why $\Phi$ is the intrinsic visible variable}

The central interpretive claim of the paper is now sharp. The intrinsic visible variable is quotient precision $\Phi$, not the Donsker-Varadhan Hessian taken at fixed visible dimension. The reason is structural, not stylistic. Section~\ref{sec:quotient-precision} proves that quotient observation composes exactly along towers, and Theorem~\ref{thm:composition-law} makes that closure explicit. Section~\ref{sec:03_dv_bridge} then shows that the Donsker-Varadhan bridge remains essential, but as a bridge variable: it is the route by which the microscopic skew-square structure descends to a visible algebraic object. What closes under observation is the quotient-visible precision itself.

That distinction is not cosmetic. The tangent closure statement in Corollary~\ref{cor:tangent-closure} identifies the Donsker-Varadhan skew-square sector as the tangent microscopic cone inside a larger quotient-visible attenuation class. Section~\ref{sec:04_hidden_mediation} explains the geometry of that larger class by hidden load, support rigidity, additive clock, and conservation. Once those structures are in place, the right reading is forced: the Donsker-Varadhan Hessian is the bridge variable that generates the correct visible object, while $\Phi$ is the visible invariant that survives composition, hidden elimination, and law-level comparison.

At scalar visible rank the conclusion is even stronger. Within that domain, ceiling, clock, precision, clock additivity, and conservation form a closed scalar system. The tangent ceiling fixes where visible correction can live, the determinant clock measures hidden contraction volume, and the conservation split from Section~\ref{sec:04_hidden_mediation} separates reversible and irreversible visible contributions without leaving the quotient-visible framework. In that precise sense, the scalar visible theory is complete: it is not merely a collection of formulas, but a closed algebraic system for the visible object. Conceptually this places the visible side of the theory closer to an information-geometric coarse description than to a fixed-dimension microscopic parametrisation \cite{Rao1945,AmariNagaoka2000,AyJostLeSchwachhofer2017}.

\subsection{What the DV bridge means after quotient reduction}

The bridge from \emph{Finite} is not demoted by the quotient theory. It is repositioned. Before quotient reduction, the bridge organises the microscopic correction as a skew-square law built from the reversible backbone. After quotient reduction, that same bridge is understood as the doorway to the right visible object. It tells us how the microscopic correction enters the visible sector, but it does not itself define the renormalisation-natural invariant. The visible algebraic object of Section~\ref{sec:03_dv_bridge} is therefore best read as the bridge's finite descendant, while $\Phi$ is the quotient-theoretic object that remains well defined after observation towers, hidden mediation, and law-level comparison have been taken seriously.

\subsection{Apparent nonlinearity as coordinate artefact}

Within the paper's proved domain, apparent nonlinearity repeatedly resolves into one of two things. Either it disappears after the correct coordinate choice, or it reappears as a finite-dimensional algebraic feature with explicit operational meaning. Section~\ref{sec:04_hidden_mediation} is the clearest example. Hidden-load composition looks nonlinear in the raw visible variable $X$, but becomes additive transport in the determinant clock and linear order on the hidden-load cone. The same phenomenon appears in Section~\ref{sec:06_bell_frontier}: Bell-square compatibility looks nonlinear in correlator language, but after the arcsine lift the compatible region becomes a hypercubic coordinate constraint with exact repair by clipping. In both cases the apparent nonlinearity is not false, but misplaced. It belongs to the choice of coordinates rather than to the underlying compatibility law.

This claim must be stated with its boundary. The paper does not assert that all nonlinearity in nonequilibrium observation theory is a coordinate artefact. It asserts a narrower and stronger statement: within the quotient-visible theory developed here, every nonlinear structure that has been brought under exact control either linearises in canonical coordinates or reduces to a finite-dimensional algebraic feature with explicit physical meaning. Hidden-load composition, Bell-square compatibility, and the scalar visible Schur laws are the proved examples. Beyond that domain, the paper makes no general ontological claim.

\subsection{The multi-regime architecture}

One of the paper's mature conclusions is that the nonequilibrium visible theory is genuinely multi-regime. The abstract quotient-visible theory of Sections~\ref{sec:quotient-precision}-\ref{sec:06_bell_frontier} is exact and regime-independent. The hidden-feedback analysis of Sections~\ref{sec:07_exact_symmetry}-\ref{sec:09_corridor_regime} is not. It splits into an exact symmetry backbone, a fast-memory response sector, and a strong-driving corridor sector. The paper does not hide that separation. It proves it.

That separation is a strength. Section~\ref{sec:08_fast_memory} shows that fast-memory response is organised by the visible-first Schur carrier, the canonical odd hidden response scalar, and a source law with explicit obstruction. Section~\ref{sec:09_corridor_regime} shows that the strong-driving corridor regime is controlled instead by path-kernel endpoint Schur laws and a distinct scaling architecture. The regime-separation theorem, Theorem~\ref{thm:regime-separation}, and the current boundary statement in Proposition~\ref{prop:current-regime-boundary} make the point sharply: these are not two parametrisations of one already unified operator sector. They are two real asymptotic sectors with different carriers, different dominant mechanisms, and a currently open containing architecture.

The hidden-feedback ladder therefore teaches two lessons at once. First, the visible object is rigid enough to support exact symmetry laws, asymptotic source laws, regime boundaries, and thermodynamic constraints. Second, those laws do not collapse into one prematurely simplified picture. The theory has become strong enough to prove the separation, and that is a more valuable outcome than a forced unity would have been.

\subsection{Geometric structure of quotient observation}

The quotient map is also geometric in a precise sense.

\begin{remark}[Affine-invariant geometric viewpoint]
Equip $\Sym_{++}(n)$ with the affine-invariant metric
\[
g_H(U,V)=\operatorname{tr}(H^{-1}UH^{-1}V).
\]
In the metric-normalised shadow coordinates of Section~\ref{sec:quotient-precision}, the differential of
\[
F_C(H)=(CH^{-1}C^{\mathsf T})^{-1}
\]
acts by the coisometric compression of Theorem~\ref{thm:shadow-transport}. This is the tangent-level isometry condition one would need for a Riemannian-submersion interpretation of quotient observation in the affine-invariant geometry of the positive cone.

That geometric reading is enough for the present paper. A full Riemannian-submersion theorem would additionally require an explicit horizontal distribution, surjectivity of the differential onto the visible tangent space, and a complete metric-preservation argument on horizontal directions. We do not supply that full proof here. What the paper uses is the tangent coisometry and the resulting structural analogy.

It is equally important to state what this does \emph{not} mean. The determinant clock from Section~\ref{sec:04_hidden_mediation} is not the affine-invariant geodesic distance once visible rank exceeds one. Its additivity comes from multiplicativity of determinant under hidden-load composition, not from metric geodesic length. The geometric parallels do not amount to an Einstein equation, a dynamical gravity theory, or a proved continuum limit. What the paper has established is kinematic geometric structure, not a gravitational dynamics.
\end{remark}

\appendix

\section{Bell and temporal frontier details}\label{app:bell}

This appendix provides the technical details underlying Theorem~\ref{thm:bell-collapse}, Proposition~\ref{prop:bell-strict-refinement}, Corollary~\ref{cor:bell-cube-body}, and the temporal-frontier discussion in Section~\ref{sec:06_bell_frontier}.

\subsection{Gaussian gluing on the Bell square}

We first isolate the Gaussian completion problem behind Theorem~\ref{thm:bell-collapse}. Write
\[
\Sigma_{xy}=\begin{pmatrix}\sigma_{A_x}^2 & c_{xy}\\ c_{xy} & \sigma_{B_y}^2\end{pmatrix},
\qquad
K_{xy}:=\frac{c_{xy}}{\sigma_{A_x}\sigma_{B_y}},
\]
where the one-site variances are assumed to satisfy the remote-setting consistency conditions \eqref{eq:variance-consistency-body}. Let $\mathcal E_{2,2}$ denote the bipartite elliptope, namely the set of all $2\times2$ matrices of the form
\[
L_{xy}=u_x\cdot v_y
\qquad
\text{with }u_x,v_y\text{ unit vectors in a real Hilbert space}.
\]

\begin{proposition}\label{prop:bell-elliptope}
Under the variance-consistency conditions \eqref{eq:variance-consistency-body}, the following are equivalent.
\begin{enumerate}
\item There exists a centred Gaussian law on $(A_0,A_1,B_0,B_1)$ whose $(A_x,B_y)$ marginals are exactly $\Sigma_{xy}$.
\item The normalised cross block $K=(K_{xy})$ belongs to $\mathcal E_{2,2}$.
\end{enumerate}
If the resulting full covariance is positive definite, this is equivalent to common Schur gluing of the visible pair precisions.
\end{proposition}

\begin{proof}
Assume first that a common centred Gaussian law exists. After standardising each coordinate by its one-site standard deviation, one obtains a correlation matrix on $(A_0,A_1,B_0,B_1)$ with diagonal entries equal to $1$ and cross block equal to $K$. Every correlation matrix is positive semidefinite and therefore a Gram matrix of unit vectors. Hence $K\in\mathcal E_{2,2}$.

Conversely, suppose $K_{xy}=u_x\cdot v_y$ for unit vectors $u_x,v_y$. Form the block Gram matrix
\[
R:=\begin{pmatrix}
(u_x\cdot u_{x'})_{x,x'} & (u_x\cdot v_y)_{x,y}\\
(v_y\cdot u_x)_{y,x} & (v_y\cdot v_{y'})_{y,y'}
\end{pmatrix}\succeq 0.
\]
Now scale $R$ by the diagonal matrix
\[
D:=\operatorname{diag}(\sigma_{A_0},\sigma_{A_1},\sigma_{B_0},\sigma_{B_1}).
\]
Then $\Sigma_\ast:=DRD\succeq 0$ is a covariance matrix whose $(A_x,B_y)$ marginal block is exactly $\Sigma_{xy}$. If $\Sigma_\ast\succ0$, inversion yields a visible precision $Q_\ast=\Sigma_\ast^{-1}$, and each pair precision is the corresponding Schur complement block. This is exactly common Schur gluing.
\end{proof}

\begin{proof}[Proof of Theorem~\ref{thm:bell-collapse}]
By Proposition~\ref{prop:bell-elliptope}, common Gaussian gluing is equivalent to variance consistency together with membership of the normalised cross block $K$ in $\mathcal E_{2,2}$. For pair precisions,
\[
K_{xy}=\frac{c_{xy}}{\sqrt{v^A_{xy}v^B_{xy}}}=\frac{-b_{xy}}{\sqrt{a_{xy}d_{xy}}}=\rho_{xy},
\]
which is \eqref{eq:bell-rho-def}.

Now define
\[
E_{xy}:=\frac{2}{\pi}\arcsin(K_{xy}).
\]
Standard Bell-square facts identify the frontier in three equivalent ways \cite{Fine1982}. Fine's theorem characterises Bell-locality of $E$ by the CHSH inequalities, the Gaussian sign-correlation formula shows that $E$ is the correlator table of the signs of a Gaussian family with correlation matrix having cross block $K$, and the same lifted correlator set is the relevant bipartite elliptope in the Bell square. Hence $K\in\mathcal E_{2,2}$ if and only if the lifted correlator table $E$ is Bell-local, equivalently if and only if the four arcsine CHSH inequalities \eqref{eq:arcsine-chsh-body-a} to \eqref{eq:arcsine-chsh-body-d} hold. Combining this with Proposition~\ref{prop:bell-elliptope} proves the theorem.
\end{proof}

\subsection{Strict refinement over Bell-locality}

Proposition~\ref{prop:bell-strict-refinement} already gives the minimal counterexample. The point is structural: Bell inequalities constrain only the binary-output correlator shadow, while common Gaussian gluing constrains the full pair law. The diagonal sector is therefore visible to the quotient frontier even when the Bell shadow is trivial.

A convenient quantitative summary is to separate the two obstructions. Define
\begin{equation}\label{eq:epsilon-var-app}
\varepsilon_{\mathrm{var}}:=\max\!\bigl\{|v^A_{00}-v^A_{01}|,|v^A_{10}-v^A_{11}|,|v^B_{00}-v^B_{10}|,|v^B_{01}-v^B_{11}|\bigr\},
\end{equation}
and let $\varepsilon_{\mathrm{CHSH}}$ be the excess of the largest left-hand side of \eqref{eq:arcsine-chsh-body-a} to \eqref{eq:arcsine-chsh-body-d} above $\pi$, truncated below at zero. Then common gluing on the Bell square is equivalent to the simultaneous vanishing of both obstructions,
\[
\varepsilon_{\mathrm{var}}=0,
\qquad
\varepsilon_{\mathrm{CHSH}}=0.
\]
This packages the Bell-square frontier as a two-coordinate compatibility problem: one coordinate is invisible to Bell correlators, while the other is the usual arcsine correlator obstruction.

\subsection{Low-rank mediated Gaussian sign shadows}

The Bell-local shadow can itself be realised by a very small mediated Gaussian model.

\begin{corollary}\label{cor:rank-two-sign-shadow}
Let $E=(E_{xy})$ be a Bell-local $2\times2$ correlator table, and set
\[
K_{xy}:=\sin\!\Bigl(\frac{\pi}{2}E_{xy}\Bigr).
\]
Then there exist unit vectors $u_0,u_1,v_0,v_1\in\mathbb R^2$ and a standard Gaussian $Z\sim N(0,I_2)$ such that
\[
U_x:=u_x\cdot Z,
\qquad
V_y:=v_y\cdot Z,
\qquad
\mathbb E[\operatorname{sgn}(U_x)\operatorname{sgn}(V_y)]=E_{xy}
\quad\text{for all }x,y.
\]
Thus every Bell-local Bell-square correlator table already sits inside a rank-two Gaussian sign-shadow.
\end{corollary}

\begin{proof}
By Theorem~\ref{thm:bell-collapse}, Bell-locality of $E$ is equivalent to $K\in\mathcal E_{2,2}$. Thus there exist unit vectors $u_x,v_y$ in some Euclidean space with $u_x\cdot v_y=K_{xy}$. Let $U:=\operatorname{span}\{u_0,u_1\}$, so $\dim U\le 2$, and let $p_y$ be the orthogonal projection of $v_y$ onto $U$. Then $u_x\cdot p_y=K_{xy}$ for all $x,y$. Choose $Z_U\sim N(0,I_U)$ and independent standard Gaussians $\xi_y$, and define
\[
U_x:=u_x\cdot Z_U,
\qquad
V_y:=p_y\cdot Z_U+\sqrt{1-\|p_y\|^2}\,\xi_y.
\]
These are centred Gaussian variables with unit variances and cross-correlations $\operatorname{Corr}(U_x,V_y)=K_{xy}$. The Gaussian sign-correlation formula gives
\[
\mathbb E[\operatorname{sgn}(U_x)\operatorname{sgn}(V_y)]
=\frac{2}{\pi}\arcsin(K_{xy})
=E_{xy},
\]
as required.
\end{proof}

\subsection{Bell cube normal form and exact repair}

The Bell square admits a global normal form in lifted coordinates. This globalises the one-face repair formulas and makes the Bell witness geometry completely explicit.

\begin{theorem}\label{thm:bell-cube-app}
Let $S$ be the $4\times 4$ matrix whose rows are the four CHSH sign vectors
\[
s_1=(1,1,1,-1),\qquad s_2=(1,1,-1,1),\qquad s_3=(1,-1,1,1),\qquad s_4=(-1,1,1,1).
\]
Then
\[
S^\top S=4I_4.
\]
For
\[
y:=\frac{1}{\pi}S\phi,
\qquad
\phi:=(\phi_{00},\phi_{01},\phi_{10},\phi_{11})^\top,
\]
the Bell-compatible region in arcsine coordinates becomes
\[
|y_i|\le 1,
\qquad i=1,\dots,4.
\]
Equivalently, the Bell square is globally the cube $[-1,1]^4$ in the lifted coordinates $y$. Moreover,
\[
\phi=\frac{\pi}{4}S^\top y,
\qquad
\|\delta\phi\|^2=\frac{\pi^2}{4}\,\|\delta y\|^2.
\]
Hence the Euclidean nearest-point map is coordinatewise clipping,
\[
y_i^\ast=\operatorname{clip}(y_i,-1,1),
\qquad
\phi^\ast=\frac{\pi}{4}S^\top y^\ast,
\]
and the squared Euclidean distance to the Bell frontier is
\[
\operatorname{dist}(\phi,\mathcal B)^2
=\frac{\pi^2}{4}\sum_{i=1}^4 (|y_i|-1)_+^2.
\]
If only one lifted CHSH coordinate violates the cube, with defect $e>0$, then the repair is
\[
\delta\phi^\ast=-\frac{e\pi}{4}s_i,
\qquad
\|\delta\phi^\ast\|_\infty=\frac{e\pi}{4}.
\]
\end{theorem}

\begin{proof}
The four inequalities \eqref{eq:arcsine-chsh-body-a} to \eqref{eq:arcsine-chsh-body-d} are precisely
\[
|s_i\cdot \phi|\le \pi,
\qquad i=1,\dots,4.
\]
Since the rows of $S$ are the vectors $s_i$, this is equivalent to $|y_i|\le 1$ for all $i$, so the compatible region is the cube $[-1,1]^4$ in the lifted coordinates. Direct computation gives $S^\top S=4I_4$, hence $S^{-1}=\frac14 S^\top$ and therefore
\[
\phi=\frac{\pi}{4}S^\top y.
\]
The Euclidean metric transforms accordingly,
\[
\|\delta\phi\|^2
=\frac{\pi^2}{16}\,\delta y^\top SS^\top \delta y
=\frac{\pi^2}{4}\,\|\delta y\|^2,
\]
because $SS^\top=4I_4$. Euclidean nearest-point repair in $y$ is therefore just coordinatewise clipping to the cube, and transforming back gives the stated projection formula and distance formula. If only one coordinate exceeds the cube by $e$, then clipping changes only that coordinate by $e$, so
\[
\delta\phi^\ast=\frac{\pi}{4}S^\top(-e e_i)=-\frac{e\pi}{4}s_i,
\]
which also gives $\|\delta\phi^\ast\|_\infty=e\pi/4$.
\end{proof}

\begin{corollary}\label{cor:bell-support-app}
For any covector $u\in\mathbb R^4$, the Bell support function in arcsine coordinates is
\[
h_{\mathcal B}(u):=\sup_{\phi\in\mathcal B}u\cdot\phi
=\frac{\pi}{4}\,\|Su\|_1.
\]
\end{corollary}

\begin{proof}
Using $\phi=(\pi/4)S^\top y$ and $|y_i|\le 1$, one has
\[
h_{\mathcal B}(u)=\sup_{|y_i|\le 1}\frac{\pi}{4}(Su)\cdot y
=\frac{\pi}{4}\sum_{i=1}^4 |(Su)_i|
=\frac{\pi}{4}\,\|Su\|_1.
\]
\end{proof}

\subsection{Observation graphs and the first temporal clique}

The Bell square is the first nonchordal compatibility graph treated in the main text. The same gluing question may be asked on any finite observation graph, with Bell and temporal contexts appearing as two graph-indexed instances of one Gaussian pair-law completion problem. The next proposition isolates the law-level object.

\begin{proposition}[Graph-indexed Gaussian gluing]\label{prop:graph-gluing}
Let $G=(V,E)$ be a finite graph. For each edge $\{i,j\}\in E$, let
\[
\Sigma_{ij}=\begin{pmatrix} v_i^{(ij)} & c_{ij}\\ c_{ij} & v_j^{(ij)}\end{pmatrix}\in\SPD(2),
\qquad
\rho_{ij}:=\frac{c_{ij}}{\sqrt{v_i^{(ij)}v_j^{(ij)}}}.
\]
Then there exists a centred Gaussian law on $(X_i)_{i\in V}$ whose $(i,j)$ marginal is $\Sigma_{ij}$ for every edge if and only if both of the following hold:
\begin{enumerate}[label=\textup{(\roman*)},leftmargin=2.3em]
\item vertex variance consistency, namely there exist numbers $v_i>0$ with
\[
v_i^{(ij)}=v_i
\qquad
\text{for every vertex $i$ and every incident edge $\{i,j\}$};
\]
\item correlation completion, namely there exists a correlation matrix $R\succeq0$ on $V$ such that
\[
R_{ii}=1,
\qquad
R_{ij}=\rho_{ij}
\quad
\text{for all }\{i,j\}\in E.
\]
\end{enumerate}
\end{proposition}

\begin{proof}
If a global centred Gaussian law exists, its one-site variances are independent of which edge marginal is inspected, so \textup{(i)} holds. Dividing the global covariance by the diagonal standard-deviation matrix gives a correlation matrix $R$ with the stated edge restrictions, so \textup{(ii)} holds.

Conversely, assume \textup{(i)} and \textup{(ii)}. Let $D=\operatorname{diag}(v_i)_{i\in V}$ and set $\Sigma:=D^{1/2}RD^{1/2}$. Then $\Sigma\succeq0$ and every edge marginal is
\[
\Sigma_{\{i,j\}}=\begin{pmatrix}
v_i & \sqrt{v_iv_j}\,R_{ij}\\
\sqrt{v_iv_j}\,R_{ij} & v_j
\end{pmatrix}
=\begin{pmatrix}
v_i^{(ij)} & c_{ij}\\ c_{ij} & v_j^{(ij)}
\end{pmatrix}
=\Sigma_{ij},
\]
because $R_{ij}=\rho_{ij}=c_{ij}/\sqrt{v_i^{(ij)}v_j^{(ij)}}$ and the visible variances agree by \textup{(i)}.
\end{proof}

The contextual object is therefore the edge-correlation completion set
\[
\mathcal E(G):=\{\rho\in(-1,1)^E:\rho\text{ extends to a correlation matrix on }V\}.
\]
Bell and temporal inequalities sit below $\mathcal E(G)$ as readout shadows. At law level the primary obstruction is gluing failure for the full pair laws.

The next theorem identifies the easy graph class.

\begin{theorem}[Chordal clique gluing]\label{thm:chordal-gluing}
Let $G$ be a chordal graph with maximal cliques $C_1,\dots,C_m$. For each clique $C_\alpha$, let $\nu_{C_\alpha}$ be a centred Gaussian law with positive definite clique covariance. Assume that whenever two neighbouring cliques in a clique tree meet in a separator $S=C_\alpha\cap C_\beta$, the induced $S$-marginals of $\nu_{C_\alpha}$ and $\nu_{C_\beta}$ agree. Then there exists a centred Gaussian law on all of $V$ whose $C_\alpha$-marginal is $\nu_{C_\alpha}$ for every $\alpha$.
\end{theorem}

\begin{proof}
The key separator-gluing step is explicit. If centred Gaussian laws on index sets $U$ and $W$ agree on the overlap $S=U\cap W$, then
\[
p(x_{U\cup W})=\frac{p_U(x_U)p_W(x_W)}{p_S(x_S)}
\]
defines a centred Gaussian law on $U\cup W$ with the prescribed $U$- and $W$-marginals. The quotient is well-defined because the overlap marginal agrees, and direct integration recovers the stated marginals.

Now choose a clique tree for $G$ and argue by induction on the number of cliques. For one clique there is nothing to prove. For more than one clique, remove a leaf clique $C_\ell$ with separator $S=C_\ell\cap C'$. By induction the remaining clique family admits a common Gaussian law. Gluing that law to $\nu_{C_\ell}$ across $S$ by the separator formula above produces the desired global Gaussian law.
\end{proof}

In particular, if $G$ is a tree, then the maximal cliques are just edges. Hence edgewise positive pair laws together with vertex variance consistency are sufficient for common Gaussian gluing. Trees therefore have no further law-level obstruction. The first genuinely nontrivial temporal obstruction appears when one passes from a tree to the first temporal clique.

The first fully observed temporal case is the triangle $K_3$. Because it is chordal, the correlation-completion condition reduces to positivity of the single $3\times3$ correlation matrix.

\begin{theorem}[Exact three-time triangle theorem]\label{thm:temporal-triangle}
Let $\Sigma_{01},\Sigma_{12},\Sigma_{02}\in\SPD(2)$ be centred Gaussian pair laws with entries
\[
\Sigma_{ij}=\begin{pmatrix} v_i^{(ij)} & c_{ij}\\ c_{ij} & v_j^{(ij)}\end{pmatrix},
\qquad
\rho_{ij}:=\frac{c_{ij}}{\sqrt{v_i^{(ij)}v_j^{(ij)}}}.
\]
There exists a centred Gaussian law on $(X_0,X_1,X_2)$ with these three pair marginals if and only if
\[
v_0^{(01)}=v_0^{(02)},\qquad
v_1^{(01)}=v_1^{(12)},\qquad
v_2^{(02)}=v_2^{(12)},
\]
and
\[
R=\begin{pmatrix}
1 & \rho_{01} & \rho_{02}\\
\rho_{01} & 1 & \rho_{12}\\
\rho_{02} & \rho_{12} & 1
\end{pmatrix}\succeq0.
\]
Equivalently,
\begin{equation}\label{eq:temporal-triangle-det}
1+2\rho_{01}\rho_{02}\rho_{12}-\rho_{01}^2-\rho_{02}^2-\rho_{12}^2\ge0.
\end{equation}
\end{theorem}

\begin{proof}
The gluing criterion is Proposition~\ref{prop:graph-gluing} specialised to $G=K_3$. For the determinant form, positivity of $R$ is equivalent to nonnegativity of its determinant because the $1\times1$ and $2\times2$ principal minors are already positive. A direct computation gives \eqref{eq:temporal-triangle-det}.
\end{proof}

\begin{corollary}[Temporal sign-shadow tetrahedron]\label{cor:temporal-sign-shadow}
Let $(X_0,X_1,X_2)$ be any centred Gaussian triple satisfying Theorem~\ref{thm:temporal-triangle}, define $S_i:=\operatorname{sgn}(X_i)$, and set
\[
y_{ij}:=\frac{2}{\pi}\arcsin(\rho_{ij}).
\]
Then
\[
\mathbb E[S_iS_j]=y_{ij},
\]
and the lifted correlator triple $(y_{01},y_{02},y_{12})$ lies in the classical three-variable correlation tetrahedron:
\begin{align}
1+y_{01}+y_{02}+y_{12}&\ge0,\label{eq:temporal-tetra-a}\\
1+y_{01}-y_{02}-y_{12}&\ge0,\label{eq:temporal-tetra-b}\\
1-y_{01}+y_{02}-y_{12}&\ge0,\label{eq:temporal-tetra-c}\\
1-y_{01}-y_{02}+y_{12}&\ge0.\label{eq:temporal-tetra-d}
\end{align}
In particular, the first temporal lifted shadow is the Leggett-Garg tetrahedron attached to the three-time clique \cite{LeggettGarg1985}.
\end{corollary}

\begin{proof}
For a centred Gaussian pair with correlation $\rho_{ij}$, Sheppard's formula gives
\[
\mathbb E[\operatorname{sgn}(X_i)\operatorname{sgn}(X_j)]=\frac{2}{\pi}\arcsin(\rho_{ij}).
\]
The tetrahedral inequalities are the classical three-variable correlation constraints for sign variables \cite{LeggettGarg1985}. Applying them to $(S_0,S_1,S_2)$ yields \eqref{eq:temporal-tetra-a}-\eqref{eq:temporal-tetra-d}.
\end{proof}

The temporal tree-versus-clique divide is therefore now visible. Trees glue automatically after repeated-marginal consistency, but they need not determine a unique global law.

\begin{proposition}[Path completion interval and nonuniqueness]\label{prop:path-completion-interval}
Fix unit variances on the path $0-1-2$ and edge correlations $\rho_{01},\rho_{12}\in(-1,1)$. A Gaussian completion exists for every
\[
\rho_{02}\in\Bigl[\rho_{01}\rho_{12}-\sqrt{(1-\rho_{01}^2)(1-\rho_{12}^2)},\,\rho_{01}\rho_{12}+\sqrt{(1-\rho_{01}^2)(1-\rho_{12}^2)}\Bigr].
\]
This interval is nondegenerate whenever $|\rho_{01}|<1$ and $|\rho_{12}|<1$. The distinguished tree-Markov construction selects the centre point $\rho_{02}=\rho_{01}\rho_{12}$.
\end{proposition}

\begin{proof}
A completion with correlation matrix
\[
R=\begin{pmatrix}1&\rho_{01}&\rho_{02}\\ \rho_{01}&1&\rho_{12}\\ \rho_{02}&\rho_{12}&1\end{pmatrix}
\]
exists if and only if $R\succeq0$. Since the $2\times2$ principal minors are positive, this is equivalent to
\[
1+2\rho_{01}\rho_{12}\rho_{02}-\rho_{01}^2-\rho_{12}^2-\rho_{02}^2\ge0.
\]
This is a quadratic inequality in $\rho_{02}$ with roots
\[
\rho_{01}\rho_{12}\pm\sqrt{(1-\rho_{01}^2)(1-\rho_{12}^2)},
\]
so the admissible interval is the one stated. Its midpoint is $\rho_{01}\rho_{12}$.
\end{proof}

\begin{proposition}[Fibre crossing at fixed correlation shadow]\label{prop:graph-fibre-crossing}
Let $G$ have a vertex of degree at least $2$, and let $\rho\in\mathcal E(G)$ lie in the interior of the correlation-completion set. Then there exist both gluable and non-gluable full edge-law families having the same edge-correlation data $\rho$.
\end{proposition}

\begin{proof}
Choose any centred Gaussian completion with unit variances, so each edge block is
\[
\Sigma_{ij}=\begin{pmatrix}1 & \rho_{ij}\\ \rho_{ij} & 1\end{pmatrix}
\]
and the family is gluable. Now pick a vertex $i$ used in at least two edges and choose one incident edge $\{i,j\}$. Replace only the variance of $i$ on that edge by a number $\lambda\neq1$, and replace the covariance on that edge by $\rho_{ij}\sqrt{\lambda}$. The normalised correlation on that edge remains $\rho_{ij}$, and all other edge correlations are untouched, so the edge-correlation shadow is unchanged. But the variance of vertex $i$ is now inconsistent across its incident edges, so Proposition~\ref{prop:graph-gluing} forbids common gluing.
\end{proof}

Thus Bell correlators do not exhaust the law-level frontier even before one passes to binary readout inequalities. They are shadows of a fuller pair-law compatibility problem, and the Bell square treated in the main text is the first nonchordal instance of that broader contextual branch.

\section{Scalar summaries and Schur-gap diagnostics}\label{app:scalar}

This appendix records the minimal scalar diagnostics supporting the Schur-gap discussion in Section~\ref{sec:quotient-precision}. No new structure is introduced here.

\subsection{Scalar mediation invariants}

Let $G\succeq0$ denote the Schur gap from \eqref{eq:schur-gap}. The basic scalar diagnostics are
\begin{equation}\label{eq:scalar-invariants}
\Gamma_{\mathrm{tr}}:=\operatorname{tr}(G),
\qquad
\Gamma_{\mathrm{op}}:=\|G\|_{\mathrm{op}},
\qquad
\Gamma_{\det}:=\log\det H_{\mathrm{comp}}-\log\det H_{\mathrm{quot}}.
\end{equation}
They measure total hidden contribution, maximal hidden contribution, and volumetric hidden contribution, respectively.

\begin{proposition}\label{prop:scalar-invariants}
The quantities in \eqref{eq:scalar-invariants} are all nonnegative.
\end{proposition}

\begin{proof}
Since $G\succeq0$, its trace and operator norm are nonnegative. Also $H_{\mathrm{comp}}\succeq H_{\mathrm{quot}}\succ0$ by \eqref{eq:schur-gap}, so monotonicity of the determinant on the positive cone gives $\det H_{\mathrm{comp}}\ge \det H_{\mathrm{quot}}$, hence $\Gamma_{\det}\ge0$.
\end{proof}

\begin{remark}
These diagnostics are only summaries. The object itself is the Schur gap operator $G$, and the main text uses the scalar summaries only as numerical readers of hidden contribution.
\end{remark}

\section{Adjacency, fill-in, and locality diagnostics}\label{app:adjacency}

This appendix provides the service definitions behind the locality discussion in Section~\ref{sec:05_hidden_locality}. It introduces quotient-induced adjacency and the associated fill-in diagnostics, but it does not add any new locality theorem beyond the main text.

\subsection{Fibre-induced visible adjacency}

To turn adjacency loss into an intrinsic quotient invariant, the visible mask must be induced by quotient geometry rather than chosen externally.

\begin{definition}[Fibre distance and quotient adjacency]\label{def:fibre-distance}
Let $\pi:X\to Y$ be a quotient of a metric graph or metric state space $(X,d_X)$. Define the lifted fibre distance by
\begin{equation}\label{eq:fibre-distance}
d_\pi(y,y'):=\inf\{d_X(u,v): u\in\pi^{-1}(y),\ v\in\pi^{-1}(y')\}.
\end{equation}
If upstairs locality has interaction range $r$, define visible adjacency by
\begin{equation}\label{eq:quotient-adjacency}
y\sim_\pi y'\quad\Longleftrightarrow\quad d_\pi(y,y')\le r.
\end{equation}
\end{definition}

Let $M_\pi$ denote the mask that keeps only matrix entries consistent with \eqref{eq:quotient-adjacency}.

\begin{definition}[Locality-defect operator]\label{def:locality-defect}
Let $H_{\mathrm{vis}}\in\SPD(Y)$ be a visible precision and let the visible Hilbert structure on $Y$ be fixed once and for all. The quotient locality-defect operator is
\begin{equation}\label{eq:locality-defect}
\mathcal N_\pi:=H_{\mathrm{vis}}-M_\pi(H_{\mathrm{vis}}).
\end{equation}
Its singular values, stable rank, and retained Frobenius mass are the quotient adjacency-loss diagnostics. These are numerical readers of Schur-complement fill-in and spectral complexity in sparse positive systems \cite{GeorgeLiu1981,HornJohnson2013,RudelsonVershynin2007}.
\end{definition}

\begin{proposition}[Intrinsic diagnostics]\label{prop:intrinsic-diagnostics}
Relative to a fixed visible Hilbert structure and interaction range $r$, the singular spectrum of $\mathcal N_\pi$ and the scalar mediation observables of Appendix~\ref{app:scalar} are intrinsic to the quotient construction.
\end{proposition}

\begin{proof}
Once $d_\pi$ and $r$ are fixed, the mask $M_\pi$ is determined by the quotient itself. The operator $\mathcal N_\pi$ is then determined by $H_{\mathrm{vis}}$ and the quotient geometry. The scalar observables of Appendix~\ref{app:scalar} depend only on the Schur gap \eqref{eq:schur-gap}, which is itself determined by the comparison between quotient observation and naive compression.
\end{proof}

\section{Renewal / transformed semi-coarse-graining benchmark}\label{app:renewal}

This appendix supports the transformed semi-coarse-graining test discussed in Section~\ref{sec:10_benchmark_synthesis}. The point of the benchmark is methodological. It sits just beyond the static quotient theorem domain of the main text, but it still yields a tractable visible law. The source model is the four-state benchmark introduced in the transformed semi-coarse-graining study of entropy-production estimation under partial observation by Kapustin, Ghosal, and Bisker \cite{KapustinGhosalBisker2024}. In spirit it belongs near coarse-graining comparisons such as \cite{TezaStella2020}, even though the transformed observer here is renewal-level rather than a static state quotient. What matters here is that the transformed observer is not itself a static quotient of the original continuous-time Markov chain, yet the visible law closes far enough to support a finite completion and a genuine benchmark comparison. For basic finite-state chain notation and path-block bookkeeping we follow standard Markov-chain conventions \cite{KemenySnell1976}.

\subsection{Four-state benchmark and transformed observer}

Let the original continuous-time Markov chain have states $\{1,2,3,4\}$ with hidden block $H=\{3,4\}$ and rates
\begin{equation}\label{eq:bisker-rates}
1\to2:2e^x,
\quad
1\to4:1,
\quad
2\to1:3e^{-x},
\quad
2\to3:2,
\quad
2\to4:35,
\end{equation}
\begin{equation}\label{eq:bisker-rates-hidden}
3\to2:50,
\quad
3\to4:0.7,
\quad
4\to1:8,
\quad
4\to2:0.2,
\quad
4\to3:75.
\end{equation}
Following \cite{KapustinGhosalBisker2024}, a transformed semi-coarse-grained trajectory groups each completed run of consecutive visits to the hidden macrostate $H$ into a symbol $H_n$, where $n$ is the number of hidden visits in that run and the waiting time of $H_n$ is the sum of the waiting times in those $n$ visits. To keep the previous visible state explicit, write
\[
A_n:=(1,H_n),
\qquad
B_n:=(2,H_n).
\]
Define
\begin{equation}\label{eq:bisker-r3r4}
r_3:=\frac{0.7}{50.7},
\qquad
r_4:=\frac{75}{83.2},
\qquad
r:=r_3r_4.
\end{equation}
Numerically,
\[
r\approx 0.0124459490,
\qquad
\Pr(n\le 2\mid\text{enter hidden})=1-r\approx 0.987554051,
\qquad
\Pr(n>4\mid\text{enter hidden})=r^2\approx 1.54901647\times10^{-4}.
\]
Thus the transformed observer is already sharply concentrated on very short hidden excursions.

\begin{proposition}[Exact parity collapse of the transformed sequence law]\label{prop:bisker-parity-collapse}
For the benchmark \eqref{eq:bisker-rates} to \eqref{eq:bisker-r3r4}, the countable augmented transformed sequence law on
\[
\{1,2\}\cup\{A_n,B_n:n\ge1\}
\]
collapses, at the visible-law level, to a four-symbol second-order parity observer with alphabet
\begin{equation}\label{eq:bisker-parity-alphabet}
\mathcal Y_{\mathrm{par}}=\{A_o,A_e,B_o,B_e\},
\end{equation}
where the previous visible state is retained as the second-order memory coordinate and parity records whether the completed hidden run has odd or even length. Equivalently, after unfolding the previous visible state, one obtains a six-symbol first-order Markov representation. More precisely, conditioned on entry from visible state $1$,
\begin{equation}\label{eq:bisker-parity-1}
P_o^{(1)}=\frac{1-r_4}{1-r},
\qquad
P_e^{(1)}=\frac{r_4(1-r_3)}{1-r},
\end{equation}
and conditioned on entry from visible state $2$,
\begin{equation}\label{eq:bisker-parity-2}
P_o^{(2)}=
\frac{\frac{2}{37}(1-r_3)+\frac{35}{37}(1-r_4)}{1-r},
\qquad
P_e^{(2)}=
\frac{\frac{2}{37}r_3(1-r_4)+\frac{35}{37}r_4(1-r_3)}{1-r}.
\end{equation}
The next visible state depends on the completed excursion only through the previous visible state and this odd or even parity label.
\end{proposition}

\begin{proof}
Inside the hidden block the chain alternates between states $3$ and $4$ until escape to the visible sector. For entry through state $4$ the completed excursion exits from hidden state $4$ when the number of hidden visits is odd and from hidden state $3$ when it is even. The probabilities \eqref{eq:bisker-parity-1} follow by summing the corresponding geometric series. For entry from visible state $2$, the first hidden state is a fixed mixture of $3$ and $4$, and summing the odd and even geometric series gives \eqref{eq:bisker-parity-2}. Since the visible return law depends only on the hidden exit microstate, it depends only on previous visible state and parity.
\end{proof}

The corresponding visible return probabilities are explicit. In particular,
\begin{equation}\label{eq:bisker-return-A}
A_o\to1 \text{ with probability } \frac{8}{8.2},
\qquad
A_o\to2 \text{ with probability } \frac{0.2}{8.2},
\qquad
A_e\to2 \text{ with probability }1,
\end{equation}
and for the $B$-symbols,
\begin{equation}\label{eq:bisker-return-B}
\Pr(1\mid B_o)\approx 0.6207014153,
\qquad
\Pr(2\mid B_o)\approx 0.3792985847,
\end{equation}
\begin{equation}\label{eq:bisker-return-Be}
\Pr(1\mid B_e)\approx 8.5325870\times 10^{-5},
\qquad
\Pr(2\mid B_e)\approx 0.9999146741.
\end{equation}
Thus the transformed observer exposes the hidden $3\leftrightarrow4$ mode as an almost visible parity channel.

\begin{theorem}[Two-kernel timing collapse]\label{thm:bisker-two-kernel}
Let
\[
L_i(s):=\frac{\lambda_i}{\lambda_i+s},
\qquad
\lambda_3=50.7,
\qquad
\lambda_4=83.2.
\]
Then the waiting-time transforms of the parity symbols lie exactly in the two-dimensional span generated by
\begin{equation}\label{eq:bisker-G4o}
G_{4,o}(s)=\frac{(1-r)L_4(s)}{1-rL_3(s)L_4(s)},
\qquad
G_{4,e}(s)=\frac{(1-r)L_3(s)L_4(s)}{1-rL_3(s)L_4(s)}.
\end{equation}
More precisely,
\begin{equation}\label{eq:bisker-G3o}
G_{3,o}(s)=\frac{\lambda_3}{\lambda_4}G_{4,o}(s)+\Bigl(1-\frac{\lambda_3}{\lambda_4}\Bigr)G_{4,e}(s)
=\frac{39}{64}G_{4,o}(s)+\frac{25}{64}G_{4,e}(s),
\end{equation}
\begin{equation}\label{eq:bisker-GBo}
G_{B_o}(s)=w\,G_{4,o}(s)+(1-w)\,G_{4,e}(s),
\qquad
w=\frac{15093}{17593},
\end{equation}
while
\begin{equation}\label{eq:bisker-GBe}
G_{B_e}(s)=G_{4,e}(s).
\end{equation}
All four transforms therefore share the same quadratic denominator
\begin{equation}\label{eq:bisker-common-den}
D(s)=(\lambda_3+s)(\lambda_4+s)-r\lambda_3\lambda_4=(\lambda_3+s)(\lambda_4+s)-52.5.
\end{equation}
\end{theorem}

\begin{proof}
The odd and even transforms for entry through state $4$ are summed geometric series in the product $L_3(s)L_4(s)$, which gives \eqref{eq:bisker-G4o}. Entry through state $3$ either exits immediately through $3$ or first moves to $4$, which yields \eqref{eq:bisker-G3o}. The $B$-symbol transforms are then the conditional mixtures induced by the two possible hidden entry microstates, giving \eqref{eq:bisker-GBo} and \eqref{eq:bisker-GBe}. The common denominator is immediate from \eqref{eq:bisker-G4o}.
\end{proof}

Theorem~\ref{thm:bisker-two-kernel} shows that the hidden timing sector of the transformed observer is not an uncontrolled memory tower. It is exactly two-dimensional at the visible-law level.

\begin{proposition}[Exact finite completion]\label{prop:bisker-finite-completion}
The parity-lifted transformed observer admits an explicit finite CTMC completion. One such completion has visible states $1,2$ together with four hidden phase blocks: a two-state block for $A_o$, a two-state block for $A_e$, a four-state block for $B_o$, and a four-state block for $B_e$. The two basic hidden excursion generators are
\begin{equation}\label{eq:bisker-S4}
S^{(4)}=
\begin{pmatrix}
-50.7 & 50.7\\
175/169 & -83.2
\end{pmatrix},
\end{equation}
\begin{equation}\label{eq:bisker-S3}
S^{(3)}=
\begin{pmatrix}
-50.7 & 525/832\\
83.2 & -83.2
\end{pmatrix},
\end{equation}
and the $B$-blocks are direct sums of these two generators with the conditional odd or even entry weights. Hence an finite completion exists with total state count $14$.
\end{proposition}

\begin{proof}
The parity transforms in Theorem~\ref{thm:bisker-two-kernel} are rational phase-type laws of order at most two \cite{Buchholz1994}. The $A$-symbols require one conditioned two-state excursion block each. The $B$-symbols are conditional mixtures over the two possible hidden entry microstates, so each is realised exactly by the corresponding direct-sum block. Adding the visible states $1,2$ gives $2+2+2+4+4=14$ states.
\end{proof}

\begin{remark}[One exact completion and the renewal frontier]\label{rem:bisker-renewal-frontier}
The transformed semi-coarse-graining benchmark works for a structural reason, not because of an accidental good observer choice. In the 4-state benchmark, the hidden 2-state block alternates deterministically, so the transformed sequence law depends only on the parity of the hidden run length, not on the full run length. Proposition~\ref{prop:bisker-parity-collapse} therefore gives a finite parity collapse in the sequence channel. The timing channel also closes: by Theorem~\ref{thm:bisker-two-kernel}, the parity-symbol waiting-time transforms lie in a two-dimensional phase-type span and share a common quadratic denominator, so the transformed timing sector is not an uncontrolled memory tower. Together these two facts show that the transformed visible law is already finite in disguise. The explicit 14-state CTMC construction is therefore not an ad hoc completion but a realisation of an intrinsically finite parity-lifted semi-Markov visible law.

On the explicit 14-state completion from Proposition~\ref{prop:bisker-finite-completion}, the transformed observer has active visible tangent rank $4$ and hidden-load rank $2$ at the visible stalling point. Numerically,
\[
\lambda(T_{\mathrm{tr}})\approx (0,0.81855,0.88615,9.26522,10.93096),
\]
\[
\lambda(X_{\mathrm{tr}})\approx (0,0.69542,0.85140,8.94220,10.77309),
\]
\[
\lambda(\Lambda_{\mathrm{tr}})\approx (0,0,0.07834,0.19437).
\]
By contrast, the original full-CG observer on the 4-state chain is only two-dimensional and is already a one-hard-direction quotient observer, with hidden-load spectrum approximately $(0,0.274678)$ at stall. This sharpens the true open problem: not whether one finite completion exists, but whether the completion-independent tangent ceiling and hidden-load structure can be fixed directly from the renewal visible law itself rather than from a chosen finite completion. This is the natural extension frontier beyond the static quotient geometry proved in the main text.
\end{remark}


\section{Blind benchmark diagnostics}\label{app:blind-benchmark}

This appendix records the auxiliary diagnostics behind the lock-then-reveal benchmark discussed in Section~\ref{sec:11_falsification}. They are included to document how the protocol behaved under finite observation and adaptive reveal, not to enlarge the theorem domain of the paper.

\subsection{Exact-ensemble finite-error spread}

On the 256-world ensemble, the retained-signal and detectability diagnostics show a nontrivial but controlled finite-error spread across worlds. The important methodological point is not that every world is reconstructed with the same accuracy, but that the ensemble error distribution remains concentrated enough for the finite tangent law to function as an operational discriminator at the level reported in the main text. This is the diagnostic content that underlay the ensemble error histograms in the archive.

\subsection{Adaptive reveal as a diagnostic frontier}

The adaptive-reveal experiment was run as a secondary check on whether theorem-targeted observation can improve the operational signal. On the 64-world diagnostic panel, the oracle theorem-target reveal improved the mean second-step revealed signal relative to the random baseline. The same comparison in the normalised detectability score was again positive, but smaller. This is why Section~\ref{sec:11_falsification} treats adaptive reveal as a diagnostic frontier rather than as a new theorem-bearing claim: the gain is real, but it belongs to the protocol layer, not to the intrinsic quotient-visible law.

\subsection{Exact versus sampled Bell verdicts}

The Bell-facing part of the blind benchmark was evaluated in both law-level and sampled modes. On the recorded hero families, law-level verdict recovery cleanly identified the glueable versus non-glueable cases at the law level. The sampled mode was materially weaker and missed the glueable families on the recorded benchmark, which is why the paper treats the Bell verdict as a law-level statement and uses the sampled track only as a robustness layer. This is the boundary that matters for interpretation: the lock-then-reveal protocol works as a test of the law-level compatibility machinery, but not yet as a theorem-safe sampled Bell pipeline.

\section{Methodological bridges and limits of the positive-survivor programme}\label{app:bridges}

This appendix retains one methodological bridge only, namely the distinction between structural and practical identifiability. Its role is to support the transfer boundary stated in Section~\ref{sec:12_scope}: neighbouring hidden-to-visible problems can exhibit the same survivor logic without thereby falling inside the theorem domain of the present paper. What is claimed here is limited and explicit. A regular structural quotient carries the descended Fisher form, enlarging the admissible experiment class sharpens that quotient, and the hidden-load geometry from the main text gives a finite-dimensional instance of the same split. What is not claimed is a universal theorem unifying all hidden-to-visible reductions.

\subsection{Bridge A: structural versus practical identifiability}

Let $\Theta$ be a smooth parameter manifold, let $\mathcal E$ be an admissible experiment class, and for each $E\in\mathcal E$ let
\[
p_E(z\mid\theta)
\]
be the observational law. Define structural equivalence relative to $\mathcal E$ by
\[
\theta\sim\theta'\quad\Longleftrightarrow\quad p_E(\cdot\mid\theta)=p_E(\cdot\mid\theta')\ \text{for all }E\in\mathcal E.
\]
Assume that near a regular point the equivalence classes form a smooth quotient manifold
\[
q:\Theta\to \bar\Theta:=\Theta/\!\sim.
\]
For a fixed experiment $E$, let $I_E(\theta)$ denote the Fisher information bilinear form on $T_\theta\Theta$.

\begin{theorem}\label{thm:bridge-identifiability}
Under the regular quotient hypothesis, each experiment $E\in\mathcal E$ factors uniquely through the quotient: there exists a reduced statistical model
\[
\bar p_E(z\mid \bar\theta),\qquad \bar\theta\in\bar\Theta,
\]
such that
\[
p_E(z\mid\theta)=\bar p_E(z\mid q(\theta)).
\]
Moreover:
\begin{enumerate}
\item for every vertical vector $v\in\ker dq_\theta$, the score annihilates,
\[
\partial_v \log p_E(z\mid\theta)=0;
\]
\item the Fisher form vanishes on vertical directions and descends uniquely to a bilinear form $\bar I_E$ on $T_{q(\theta)}\bar\Theta$ satisfying
\[
I_E(\theta)(u,w)=\bar I_E(q(\theta))(dq_\theta u,dq_\theta w);
\]
\item structural nonidentifiability is the vertical degeneracy of $q$, while practical nonidentifiability for a given experiment is degeneracy or strong anisotropy of the descended form $\bar I_E$ on the quotient tangent.
\end{enumerate}
\end{theorem}

\begin{proof}
By definition, each $p_E(\cdot\mid\theta)$ is constant on every structural fibre, so it depends only on $q(\theta)$; this gives the unique reduced model. Differentiating the factorisation along a vertical vector $v\in\ker dq_\theta$ gives the score annihilation. The Fisher identity then shows that $I_E(\theta)(v,\cdot)=0$, so $I_E(\theta)$ depends only on quotient tangent classes and therefore descends uniquely to $\bar I_E$.
\end{proof}

The theorem is local rather than global. It says that the correct canonical object is not a matrix written in some auxiliary horizontal complement, but the descended bilinear form on the structural quotient itself. A direct consequence is that enlarging the admissible experiment class can only refine the quotient and remove vertical directions; it cannot create a new structural fibre.

\begin{proposition}\label{prop:bridge-identifiability-refinement}
Let $\mathcal E_1\subseteq \mathcal E_2$ be two experiment classes, and assume both regular quotients exist near the point of interest:
\[
q_i:\Theta\to \bar\Theta_i=\Theta/\!\sim_i,\qquad i=1,2.
\]
Then there exists a unique local surjection $\pi:\bar\Theta_2\to \bar\Theta_1$ such that
\[
q_1=\pi\circ q_2.
\]
For every experiment $E\in\mathcal E_1$, the descended Fisher forms satisfy the pullback relation
\[
\bar I_E^{(2)}=\pi^*\bar I_E^{(1)}.
\]
Hence enlarging the experiment class can only sharpen the quotient geometry.
\end{proposition}

\begin{proof}
If two parameters are equivalent for the larger class $\mathcal E_2$, they are equivalent for the smaller class $\mathcal E_1$. Therefore every $\sim_2$-class sits inside a $\sim_1$-class, which induces the unique local map $\pi$ with $q_1=\pi\circ q_2$. The Fisher pullback identity then follows from Theorem~\ref{thm:bridge-identifiability} applied to the factorisation of the statistical model through the two quotients.
\end{proof}

A canonical toy model is the scalar decay law
\[
\dot x=-k x,
\qquad
Y_i=c x_0 e^{-k t_i}+\varepsilon_i,
\qquad
\varepsilon_i\sim\mathcal N(0,\sigma^2),
\]
with parameters $(c,x_0,k)$. The full observational law depends only on $a=cx_0$ and $k$, so the structural fibres are
\[
(c,x_0,k)\sim(\lambda c,x_0/\lambda,k),\qquad \lambda>0.
\]
The quotient is therefore two-dimensional with coordinates $(a,k)$, and the quotient Fisher form is the Fisher matrix of the reduced mean model $a e^{-k t_i}$. This shows in one line that a direction can be structurally invisible upstairs while the remaining quotient directions are still experiment-dependent in visibility.

In the present quotient setting this split becomes finite-dimensional rather than merely local.

\begin{theorem}[Exact inverse split for the hidden-load class]\label{thm:bridge-hidden-load-split}
Let $T\succeq0$ and let $S:=\operatorname{Ran}(T)$. Let $X$ be a visible law in the support-preserving interior beneath $T$, so that by Theorem~\ref{thm:hidden-load}
\[
X=T^{1/2}(I_S+\Lambda)^{-1}T^{1/2},\qquad \Lambda\succeq0\text{ on }S.
\]
Then:
\begin{enumerate}[label=\textup{(\roman*)},leftmargin=2.4em]
\item conditional on the tangent ceiling $T$, the hidden-load operator is uniquely identified by $X$ and equals
\begin{equation}\label{eq:bridge-load-recovery}
\Lambda=(T|_S)^{1/2}(X|_S)^{-1}(T|_S)^{1/2}-I_S;
\end{equation}
moreover,
\[
\rank(\Lambda)=\rank(T-X),
\qquad
\log\det(I_S+\Lambda)=\log\pdet(T)-\log\pdet(X);
\]
\item without an independently fixed ceiling, $X$ alone does not identify a unique hidden load: for every $T'\succeq X$ with $\operatorname{Ran}(T')=S$, there exists a unique $\Lambda_{T'}\succeq0$ on $S$ such that
\[
X=T'^{1/2}(I_S+\Lambda_{T'})^{-1}T'^{1/2},
\qquad
\Lambda_{T'}=(T'|_S)^{1/2}(X|_S)^{-1}(T'|_S)^{1/2}-I_S.
\]
Hence $X$ determines a full ceiling cone of compatible ceiling-load explanations.
\end{enumerate}
\end{theorem}

\begin{proof}
Part \textup{(i)} is the explicit recovery formula already built into Theorem~\ref{thm:hidden-load}. The rank and determinant identities are Corollary~\ref{cor:hidden-load-rank-volume}.

For part \textup{(ii)}, fix $T'\succeq X$ on the same support $S$ and define $\Lambda_{T'}$ by the displayed formula. Then
\[
T'^{-1/2}XT'^{-1/2}\preceq I_S.
\]
Since inversion reverses Loewner order on positive definite operators,
\[
(T'^{-1/2}XT'^{-1/2})^{-1}\succeq I_S,
\]
so $\Lambda_{T'}\succeq0$. Rearranging gives the stated representation of $X$, and uniqueness follows because $T'^{-1/2}XT'^{-1/2}=(I_S+\Lambda_{T'})^{-1}$.
\end{proof}

Combined with Corollary~\ref{cor:hidden-load-monotone}, this gives the order reversal at fixed ceiling:
\[
X_1\preceq X_0
\quad\Longleftrightarrow\quad
\Lambda_1\succeq\Lambda_0.
\]
It also gives the zero-load boundary $\Lambda=0\Longleftrightarrow X=T$.

\begin{proposition}[Gram rigidity and hidden-factor gauge]\label{prop:bridge-hidden-factor-gauge}
Fix $T$ and $X$ and let $\Lambda$ be the uniquely identified hidden-load operator from Theorem~\ref{thm:bridge-hidden-load-split}. Write $r:=\rank(\Lambda)$. Then every minimal hidden realisation is a Gram factorisation
\[
\Lambda=BB^\top,
\qquad B\in\mathbb R^{|S|\times r},
\]
and if $B,B'\in\mathbb R^{|S|\times r}$ both have full column rank and satisfy $BB^\top=B'B'^\top=\Lambda$, then
\[
B'=BQ
\]
for some orthogonal matrix $Q\in O(r)$. Consequently $(T,X)$ identifies the hidden Gram operator and its intrinsic invariants, but not a unique hidden basis, hidden factorisation, or hidden sparsity pattern.
\end{proposition}

\begin{proof}
The existence of a minimal factorisation follows from Theorem~\ref{thm:min-hidden-dim}. Assume $BB^\top=B'B'^\top=\Lambda$ with both factors of column rank $r$, and define
\[
Q:=B^+B',
\qquad
B^+=(B^\top B)^{-1}B^\top.
\]
Because $\operatorname{Ran}(B)=\operatorname{Ran}(\Lambda)=\operatorname{Ran}(B')$, one has
\[
BQ=BB^+B'=B'.
\]
Moreover,
\[
QQ^\top
=B^+B'B'^\top(B^+)^\top
=B^+\Lambda(B^+)^\top
=B^+BB^\top(B^+)^\top
=I_r.
\]
Since $Q$ is an $r\times r$ matrix, $QQ^\top=I_r$ implies $Q\in O(r)$. Hence $B'=BQ$.
\end{proof}


\begin{thebibliography}{58}\bibitem[{M.~D. Donsker} and {S.~R.~S.
  Varadhan}(1975{\natexlab{a}})]{DonskerVaradhan1975a}
{M.~D. Donsker} and {S.~R.~S. Varadhan}.
\newblock Asymptotic evaluation of certain {M}arkov process expectations for
  large time.~{I}, 1975{\natexlab{a}}.
\newblock \emph{Communications on Pure and Applied Mathematics},
  \textbf{28}:1-47, 1975.

\bibitem[{M.~D. Donsker} and {S.~R.~S.
  Varadhan}(1975{\natexlab{b}})]{DonskerVaradhan1975b}
{M.~D. Donsker} and {S.~R.~S. Varadhan}.
\newblock Asymptotic evaluation of certain {M}arkov process expectations for
  large time.~{II}, 1975{\natexlab{b}}.
\newblock \emph{Communications on Pure and Applied Mathematics},
  \textbf{28}:279-301, 1975.

\bibitem[{M.~D. Donsker} and {S.~R.~S. Varadhan}(1976)]{DonskerVaradhan1976}
{M.~D. Donsker} and {S.~R.~S. Varadhan}.
\newblock Asymptotic evaluation of certain {M}arkov process expectations for
  large time.~{III}, 1976.
\newblock \emph{Communications on Pure and Applied Mathematics},
  \textbf{29}(4):389-461, 1976.

\bibitem[{R.~S. Ellis}(2006)]{Ellis1985}
{R.~S. Ellis}.
\newblock \emph{Entropy, Large Deviations, and Statistical Mechanics}, 2006.
\newblock Grundlehren der mathematischen Wissenschaften, Vol.~271, Springer,
  1985. Reprinted as Classics in Mathematics, 2006.

\bibitem[{H.~Touchette}(2009)]{Touchette2009}
{H.~Touchette}.
\newblock The large deviation approach to statistical mechanics, 2009.
\newblock \emph{Physics Reports}, \textbf{478}:1-69, 2009.

\bibitem[{L.~Bertini} et~al.(2015){L.~Bertini}, {A.~Faggionato}, and
  {D.~Gabrielli}]{BertiniFaggionatoGabrielli2015}
{L.~Bertini}, {A.~Faggionato}, and {D.~Gabrielli}.
\newblock Large deviations of the empirical flow for continuous time {M}arkov
  chains, 2015.
\newblock \emph{Annales de l'I.H.P. Probabilit{\'e}s et Statistiques},
  \textbf{51}(3):867-900, 2015.

\bibitem[{J.~R. Dunkley}(2026)]{Dunkley2026FiniteObservation}
{J.~R. Dunkley}.
\newblock \emph{Finite Observation: A Spectral Theory of Nonequilibrium
  Visibility}, 2026.
\newblock Preprint, 2026.

\bibitem[{J.~Schnakenberg}(1976)]{Schnakenberg1976}
{J.~Schnakenberg}.
\newblock Network theory of microscopic and macroscopic behavior of master
  equation systems, 1976.
\newblock \emph{Reviews of Modern Physics}, \textbf{48}:571-585, 1976.

\bibitem[{J.~L. Lebowitz} and {H.~Spohn}(1999)]{LebowitzSpohn1999}
{J.~L. Lebowitz} and {H.~Spohn}.
\newblock A {G}allavotti-{C}ohen-type symmetry in the large deviation
  functional for stochastic dynamics, 1999.
\newblock \emph{Journal of Statistical Physics}, \textbf{95}:333-365, 1999.

\bibitem[{U.~Seifert}(2012)]{Seifert2012}
{U.~Seifert}.
\newblock Stochastic thermodynamics, fluctuation theorems and molecular
  machines, 2012.
\newblock \emph{Reports on Progress in Physics}, \textbf{75}:126001, 2012.

\bibitem[{M.~Esposito}(2012)]{Esposito2012}
{M.~Esposito}.
\newblock Stochastic thermodynamics under coarse graining, 2012.
\newblock \emph{Physical Review E}, \textbf{85}:041125, 2012.

\bibitem[{C.~Maes}(2020)]{Maes2020}
{C.~Maes}.
\newblock Frenesy: time-symmetric dynamical activity in nonequilibria, 2020.
\newblock \emph{Physics Reports}, \textbf{850}:1-33, 2020.

\bibitem[{I.~Schur}(1917)]{Schur1917}
{I.~Schur}.
\newblock {\"U}ber {P}otenzreihen, die im {I}nnern des {E}inheitskreises
  beschr{\"a}nkt sind, 1917.
\newblock \emph{Journal f{\"u}r die reine und angewandte Mathematik},
  \textbf{147}:205-232, 1917.

\bibitem[{W.~N. Anderson} and {Jr}(1971)]{Anderson1971}
{W.~N. Anderson} and {Jr}.
\newblock Shorted operators, 1971.
\newblock \emph{SIAM Journal on Applied Mathematics}, \textbf{20}(3):520-525,
  1971.

\bibitem[{W.~N. Anderson, Jr.} and {R.~J. Duffin}(1969)]{AndersonDuffin1969}
{W.~N. Anderson, Jr.} and {R.~J. Duffin}.
\newblock Series and parallel addition of matrices, 1969.
\newblock \emph{Journal of Mathematical Analysis and Applications},
  \textbf{26}(3):576-594, 1969.

\bibitem[{W.~N. Anderson, Jr.} and {G.~E. Trapp}(1975)]{AndersonTrapp1975}
{W.~N. Anderson, Jr.} and {G.~E. Trapp}.
\newblock Shorted operators.~{II}, 1975.
\newblock \emph{SIAM Journal on Applied Mathematics}, \textbf{28}(1):60-71,
  1975.

\bibitem[{R.~Bhatia}(2007)]{Bhatia2007}
{R.~Bhatia}.
\newblock \emph{Positive Definite Matrices}, 2007.
\newblock Princeton University Press, 2007.

\bibitem[{K.~L{\"o}wner}(1934)]{Loewner1934}
{K.~L{\"o}wner}.
\newblock {\"U}ber monotone {M}atrixfunktionen, 1934.
\newblock \emph{Mathematische Zeitschrift}, \textbf{38}:177-216, 1934.

\bibitem[{T.~Ando}(1979)]{Ando1979}
{T.~Ando}.
\newblock Concavity of certain maps on positive definite matrices and
  applications to {H}adamard products, 1979.
\newblock \emph{Linear Algebra and Its Applications}, \textbf{26}:203-241,
  1979.

\bibitem[{F.~Hansen} and {G.~K. Pedersen}(1982)]{HansenPedersen1982}
{F.~Hansen} and {G.~K. Pedersen}.
\newblock Jensen's inequality for operators and {L}{\"o}wner's theorem, 1982.
\newblock \emph{Mathematische Annalen}, \textbf{258}(3):229-241, 1982.

\bibitem[{R.~Bhatia}(1997)]{Bhatia1997}
{R.~Bhatia}.
\newblock \emph{Matrix Analysis}, 1997.
\newblock Graduate Texts in Mathematics, Vol.~169, Springer, 1997.

\bibitem[{F.~Kubo} and {T.~Ando}(1980)]{KuboAndo1980}
{F.~Kubo} and {T.~Ando}.
\newblock Means of positive linear operators, 1980.
\newblock \emph{Mathematische Annalen}, \textbf{246}:205-224, 1980.

\bibitem[{W.~Pusz} and {S.~L. Woronowicz}(1975)]{PuszWoronowicz1975}
{W.~Pusz} and {S.~L. Woronowicz}.
\newblock Functional calculus for sesquilinear forms and the purification map,
  1975.
\newblock \emph{Reports on Mathematical Physics}, \textbf{8}(2):159-170, 1975.

\bibitem[{L.~P. Kadanoff}(1966)]{Kadanoff1966}
{L.~P. Kadanoff}.
\newblock Scaling laws for {I}sing models near ${T}_c$, 1966.
\newblock \emph{Physics Physique Fizika}, \textbf{2}:263-272, 1966.

\bibitem[{K.~G. Wilson} and {J.~Kogut}(1974)]{WilsonKogut1974}
{K.~G. Wilson} and {J.~Kogut}.
\newblock The renormalization group and the $\varepsilon$ expansion, 1974.
\newblock \emph{Physics Reports}, \textbf{12}(2):75-200, 1974.

\bibitem[{J.~Polchinski}(1984)]{Polchinski1984}
{J.~Polchinski}.
\newblock Renormalization and effective {L}agrangians, 1984.
\newblock \emph{Nuclear Physics B}, \textbf{231}(2):269-295, 1984.

\bibitem[{K.~Fan}(1951)]{Fan1951}
{K.~Fan}.
\newblock Maximum properties and inequalities for the eigenvalues of completely
  continuous operators, 1951.
\newblock \emph{Proceedings of the National Academy of Sciences},
  \textbf{37}:760-766, 1951.

\bibitem[{R.~A. Horn} and {C.~R. Johnson}(2013)]{HornJohnson2013}
{R.~A. Horn} and {C.~R. Johnson}.
\newblock \emph{Matrix Analysis}, 2013.
\newblock Cambridge University Press, second edition, 2013.

\bibitem[{F.~Zhang} and {editor}(2005)]{Zhang2005}
{F.~Zhang} and {editor}.
\newblock \emph{The Schur Complement and Its Applications}, 2005.
\newblock Springer, New York, 2005.

\bibitem[{E.~V. Haynsworth}(1968)]{Haynsworth1968}
{E.~V. Haynsworth}.
\newblock On the {S}chur complement, 1968.
\newblock \emph{Basel Mathematical Notes}, BMN~20, 1968.

\bibitem[{D.~Carlson} et~al.(1974){D.~Carlson}, {E.~Haynsworth}, and
  {T.~Markham}]{CarlsonHaynsworthMarkham1974}
{D.~Carlson}, {E.~Haynsworth}, and {T.~Markham}.
\newblock A generalization of the {S}chur complement by means of the
  {M}oore-{P}enrose inverse, 1974.
\newblock \emph{SIAM Journal on Applied Mathematics}, \textbf{26}(1):169-175,
  1974.

\bibitem[{E.~B. Bogomol'nyi}(1976)]{Bogomolnyi1976}
{E.~B. Bogomol'nyi}.
\newblock Stability of classical solutions, 1976.
\newblock \emph{Soviet Journal of Nuclear Physics}, \textbf{24}(4):449-454,
  1976.

\bibitem[{A.~P. Dempster}(1972)]{Dempster1972}
{A.~P. Dempster}.
\newblock Covariance selection, 1972.
\newblock \emph{Biometrics}, \textbf{28}:157-175, 1972.

\bibitem[{S.~L. Lauritzen}(1996)]{Lauritzen1996}
{S.~L. Lauritzen}.
\newblock \emph{Graphical Models}, 1996.
\newblock Oxford Statistical Science Series, No.~17, Oxford University Press,
  1996.

\bibitem[{V.~Chandrasekaran} et~al.(2012){V.~Chandrasekaran}, {P.~A. Parrilo},
  and {A.~S. Willsky}]{ChandrasekaranParriloWillsky2012}
{V.~Chandrasekaran}, {P.~A. Parrilo}, and {A.~S. Willsky}.
\newblock Latent variable graphical model selection via convex optimization,
  2012.
\newblock \emph{The Annals of Statistics}, \textbf{40}(4):1935-1967, 2012.

\bibitem[{R.~Grone} et~al.(1984){R.~Grone}, {C.~R. Johnson}, {E.~M. S{\'a}},
  and {H.~Wolkowicz}]{GroneJohnsonSaWolkowicz1984}
{R.~Grone}, {C.~R. Johnson}, {E.~M. S{\'a}}, and {H.~Wolkowicz}.
\newblock Positive definite completions of partial {H}ermitian matrices, 1984.
\newblock \emph{Linear Algebra and Its Applications}, \textbf{58}:109-124,
  1984.

\bibitem[{J.~Williamson}(1936)]{Williamson1936}
{J.~Williamson}.
\newblock On the algebraic problem concerning the normal forms of linear
  dynamical systems, 1936.
\newblock \emph{American Journal of Mathematics}, \textbf{58}(1):141-163,
  1936.

\bibitem[{R.~A. Horn} and {C.~R. Johnson}(1991)]{HornJohnson1991}
{R.~A. Horn} and {C.~R. Johnson}.
\newblock \emph{Topics in Matrix Analysis}, 1991.
\newblock Cambridge University Press, 1991.

\bibitem[{A.~George} and {J.~W.~H. Liu}(1981)]{GeorgeLiu1981}
{A.~George} and {J.~W.~H. Liu}.
\newblock \emph{Computer Solution of Large Sparse Positive Definite Systems},
  1981.
\newblock Prentice-Hall, Englewood Cliffs, NJ, 1981.

\bibitem[{P.~G. Doyle} and {J.~L. Snell}(1984)]{DoyleSnell1984}
{P.~G. Doyle} and {J.~L. Snell}.
\newblock \emph{Random Walks and Electric Networks}, 1984.
\newblock Carus Mathematical Monographs, Vol.~22, Mathematical Association of
  America, 1984.

\bibitem[{D.~M. Malioutov} et~al.(2006){D.~M. Malioutov}, {J.~K. Johnson}, and
  {A.~S. Willsky}]{MalioutovJohnsonWillsky2006}
{D.~M. Malioutov}, {J.~K. Johnson}, and {A.~S. Willsky}.
\newblock Walk-sums and belief propagation in {G}aussian graphical models,
  2006.
\newblock \emph{Journal of Machine Learning Research}, \textbf{7}:2031-2064,
  2006.

\bibitem[{S.~Demko} et~al.(1984){S.~Demko}, {W.~F. Moss}, and {P.~W.
  Smith}]{DemkoMossSmith1984}
{S.~Demko}, {W.~F. Moss}, and {P.~W. Smith}.
\newblock Decay rates for inverses of band matrices, 1984.
\newblock \emph{Mathematics of Computation}, \textbf{43}(168):491-499, 1984.

\bibitem[{M.~Benzi} and {G.~H. Golub}(1999)]{BenziGolub1999}
{M.~Benzi} and {G.~H. Golub}.
\newblock Bounds for the entries of matrix functions with applications to
  preconditioning, 1999.
\newblock \emph{BIT Numerical Mathematics}, \textbf{39}(3):417-438, 1999.

\bibitem[{M.~Benzi} and {N.~Razouk}(2007)]{BenziRazouk2007}
{M.~Benzi} and {N.~Razouk}.
\newblock Decay bounds and {O}(n) algorithms for approximating functions of
  sparse matrices, 2007.
\newblock \emph{Electronic Transactions on Numerical Analysis},
  \textbf{28}:16-39, 2007.

\bibitem[{T.~J. Rivlin}(1990)]{Rivlin1974}
{T.~J. Rivlin}.
\newblock \emph{Chebyshev Polynomials: From Approximation Theory to Algebra and
  Number Theory}, 1990.
\newblock Wiley, 1974; second edition, 1990.

\bibitem[{L.~N. Trefethen}(2013)]{Trefethen2013}
{L.~N. Trefethen}.
\newblock \emph{Approximation Theory and Approximation Practice}, 2013.
\newblock SIAM, Philadelphia, 2013.

\bibitem[{U.~Kapustin} et~al.(2024){U.~Kapustin}, {A.~Ghosal}, and
  {G.~Bisker}]{KapustinGhosalBisker2024}
{U.~Kapustin}, {A.~Ghosal}, and {G.~Bisker}.
\newblock Utilizing time-series measurements for entropy-production estimation
  in partially observed systems, 2024.
\newblock \emph{Physical Review Research}, \textbf{6}:023039, 2024.

\bibitem[Liepelt and Lipowsky(2007)]{LiepeltLipowsky2007}
{A.~D. Liepelt} and {R.~Lipowsky}.
\newblock Kinesin's network of chemomechanical motor cycles.
\newblock \emph{Physical Review Letters}, \textbf{98}(25):258102, 2007.

\bibitem[Banerjee et~al.(2017)]{BanerjeeKolomeiskyIgoshin2017}
{S.~Banerjee}, {A.~B. Kolomeisky}, and {O.~A. Igoshin}.
\newblock Elucidating interplay of speed and accuracy in biological error correction, 2017.
\newblock \emph{Proceedings of the National Academy of Sciences}, \textbf{114}(20):5183-5188, 2017.

\bibitem[{C.~R. Rao}(1945)]{Rao1945}
{C.~R. Rao}.
\newblock Information and accuracy attainable in the estimation of statistical
  parameters, 1945.
\newblock \emph{Bulletin of the Calcutta Mathematical Society},
  \textbf{37}:81-91, 1945.

\bibitem[{S.~Amari} and {H.~Nagaoka}(2000)]{AmariNagaoka2000}
{S.~Amari} and {H.~Nagaoka}.
\newblock \emph{Methods of Information Geometry}, 2000.
\newblock Translations of Mathematical Monographs, Vol.~191, American
  Mathematical Society, 2000.

\bibitem[{N.~Ay} et~al.(2017){N.~Ay}, {J.~Jost}, {H.~V. L{\^e}}, and
  {L.~Schwachh{\"o}fer}]{AyJostLeSchwachhofer2017}
{N.~Ay}, {J.~Jost}, {H.~V. L{\^e}}, and {L.~Schwachh{\"o}fer}.
\newblock \emph{Information Geometry}, 2017.
\newblock Ergebnisse der Mathematik und ihrer Grenzgebiete, Vol.~64, Springer,
  2017.

\bibitem{Fine1982}
A.~Fine.
\newblock Hidden variables, joint probability, and the {B}ell inequalities.
\newblock \emph{Journal of Mathematical Physics}, \textbf{23}(7):1306-1310,
  1982.

\bibitem{LeggettGarg1985}
A.~J. Leggett and A.~Garg.
\newblock Quantum mechanics versus macroscopic realism: is the flux there when nobody looks?
\newblock \emph{Physical Review Letters}, \textbf{54}(9):857-860, 1985.

\bibitem[{M.~Rudelson} and {R.~Vershynin}(2007)]{RudelsonVershynin2007}
{M.~Rudelson} and {R.~Vershynin}.
\newblock Sampling from large matrices: an approach through geometric
  functional analysis, 2007.
\newblock \emph{Journal of the ACM}, \textbf{54}(4), Article~21, 2007.

\bibitem[{G.~Teza} and {A.~L. Stella}(2020)]{TezaStella2020}
{G.~Teza} and {A.~L. Stella}.
\newblock Exact coarse graining preserves entropy production out of
  equilibrium, 2020.
\newblock \emph{Physical Review Letters}, \textbf{125}:110601, 2020.

\bibitem[{J.~G. Kemeny} and {J.~L. Snell}(1960)]{KemenySnell1976}
{J.~G. Kemeny} and {J.~L. Snell}.
\newblock \emph{Finite Markov Chains}, 1960.
\newblock Springer-Verlag, 1976. Originally published by D.~Van Nostrand, 1960.

\bibitem[{P.~Buchholz}(1994)]{Buchholz1994}
{P.~Buchholz}.
\newblock Exact and ordinary lumpability in finite {M}arkov chains, 1994.
\newblock \emph{Journal of Applied Probability}, \textbf{31}:59-75, 1994.

\end{thebibliography}
\end{document}\right\}$}}%lass[11pt]{article}
\usepackage[T1]{fontenc}
\usepackage[utf8]{inputenc}
\usepackage{newtxtext,newtxmath}

\setcounter{tocdepth}{1}
\usepackage[a4paper,margin=35mm]{geometry}
\linespread{0.90}
\setlength{\parindent}{0pt}
\setlength{\parskip}{6pt}
\usepackage{float}
\usepackage{cmap}
\usepackage{amsfonts}
\let\openbox\relax
\usepackage{amsthm}
\usepackage{amsmath}
\usepackage{mathtools}
\allowdisplaybreaks
\usepackage{bm}
\usepackage{graphicx}
\usepackage{tikz}
\usepackage{pgfplots}
\pgfplotsset{compat=1.18}

\hyphenation{superselection}
\usepackage{microtype}
\usepackage[numbers,sort&compress]{natbib}
\usepackage{xcolor}
\usepackage{booktabs}
\usepackage{enumitem}
\usepackage{array}
\usepackage{tabularx}
\usepackage{setspace}
\usepackage[explicit]{titlesec}
\usepackage[labelfont=bf]{caption}
\DeclareCaptionFont{captionthirtysmaller}{\fontsize{7.7}{9.2}\selectfont}
\captionsetup[figure]{font=captionthirtysmaller}
\captionsetup[table]{font=captionthirtysmaller}
\usepackage{mdframed}
\usepackage{needspace}
\usepackage{tocloft}
\usepackage{placeins}
\newcolumntype{Y}{>{\raggedright\arraybackslash}X}
\setlength{\tabcolsep}{4.5pt}
\renewcommand{\arraystretch}{1.1}
\renewcommand{\qedsymbol}{}
\AtBeginEnvironment{tabular}{\small}
\AtBeginEnvironment{tabularx}{\small}

\numberwithin{equation}{section}
\mathtoolsset{showonlyrefs=false}
\DeclarePairedDelimiter{\abs}{\lvert}{\rvert}
\DeclarePairedDelimiter{\norm}{\lVert}{\rVert}
\DeclareMathOperator{\Tr}{Tr}
\DeclareMathOperator{\diag}{diag}
\DeclareMathOperator{\rank}{rank}
\DeclareMathOperator{\sr}{sr}
\DeclareMathOperator{\logdet}{logdet}
\DeclareMathOperator{\pdet}{pdet}
\DeclareMathOperator{\supp}{supp}
\newcommand{\Sym}{\mathrm{Sym}}
\newcommand{\dd}{\mathrm{d}}
\newcommand{\ii}{\mathrm{i}}
\newcommand{\SPD}{\mathrm{SPD}}
\newcommand{\eps}{\varepsilon}
\newcommand{\Loew}{\mathcal{L}}

\setlist{nosep,leftmargin=1.25em}

\titleformat{\section}
{\normalfont\normalsize\bfseries\raggedright}
{\thesection}{0.75em}{#1}
[\vspace{0.12em}\titlerule]

\newcommand{\subaccent}{\vspace{0.01em}\noindent\rule{7ex}{0.2pt}}
\titleformat{\subsection}
{\normalfont\normalsize\bfseries}{\thesubsection}{0.5em}{#1}
[\subaccent]
\titlespacing*{\subsection}{0pt}{2.2ex plus 0.6ex}{1.5ex}

\titleformat{\subsubsection}
{\normalfont\normalsize\bfseries}{\thesubsubsection}{0.5em}{#1}
\titlespacing*{\subsubsection}{0pt}{1.6ex plus 0.4ex}{0.6ex}

\usepackage[
  pdfusetitle,
  hidelinks,
  bookmarksopen=true,
  bookmarksnumbered=true
]{hyperref}

\newtheoremstyle{thmcompact}
{0.5ex}{0.5ex}
{\itshape}
{0pt}
{\bfseries}
{.}
{0.6em}
{\thmname{#1}\ \thmnumber{#2}\ \thmnote{(\textit{#3})}}

\theoremstyle{thmcompact}
\newtheorem{theorem}{Theorem}[section]
\newtheorem{proposition}[theorem]{Proposition}
\newtheorem{lemma}[theorem]{Lemma}
\newtheorem{corollary}[theorem]{Corollary}
\newtheorem{definition}[theorem]{Definition}

\theoremstyle{remark}
\newtheorem*{remarkplain}{Remark}

\mdfdefinestyle{thmframe}{%
	leftline=true, rightline=false, topline=false, bottomline=false,
	linewidth=0.8pt, linecolor=black,
	innerleftmargin=16pt, innerrightmargin=8pt,
	innertopmargin=6pt, innerbottommargin=6pt,
	skipabove=6pt, skipbelow=6pt,
	nobreak=true
}
\surroundwithmdframed[style=thmframe]{theorem}
\surroundwithmdframed[style=thmframe]{proposition}
\surroundwithmdframed[style=thmframe]{lemma}
\surroundwithmdframed[style=thmframe]{corollary}
\surroundwithmdframed[style=thmframe]{definition}

\newmdenv[
linewidth=0.6pt,
topline=false,
bottomline=false,
rightline=false,
skipabove=\baselineskip,
skipbelow=\baselineskip,
backgroundcolor=gray!3,
leftmargin=0pt,
innerleftmargin=8pt,
innerrightmargin=6pt,
frametitleaboveskip=0.4em,
frametitlebelowskip=0.2em,
frametitlefont=\bfseries\small,
]{remarkbox}

\mdfdefinestyle{eqbox}{%
	leftline=true, rightline=true, topline=true, bottomline=true,
	linewidth=0.8pt, linecolor=black,
	innerleftmargin=14pt, innerrightmargin=10pt,
	innertopmargin=8pt, innerbottommargin=8pt,
	skipabove=8pt, skipbelow=8pt,
	backgroundcolor=gray!3,
	nobreak=true
}
\usepackage{titling}
\newenvironment{remark}[1][]%
{\begin{remarkbox}\begin{remarkplain}[#1]\normalfont}
		{\end{remarkplain}\end{remarkbox}}

\renewcommand{\contentsname}{Contents}

\titleformat{name=\section,numberless}
{\normalfont\normalsize\bfseries\raggedright}
{}{0em}{#1}
[\vspace{0.12em}\titlerule]
\makeatletter
\renewcommand{\l@section}[2]{\@dottedtocline{1}{0em}{2.3em}{\raggedright #1}{#2}}
\renewcommand{\l@subsection}[2]{\@dottedtocline{2}{1.8em}{3.2em}{\raggedright #1}{#2}}
\makeatother

\emergencystretch=3em
\hfuzz=1pt
\tolerance=1800
\hbadness=2500

\newcommand{\thmneed}{\Needspace{6\baselineskip}}
\AtBeginEnvironment{theorem}{\thmneed}
\AtBeginEnvironment{proposition}{\thmneed}
\AtBeginEnvironment{lemma}{\thmneed}
\AtBeginEnvironment{corollary}{\thmneed}

\AtBeginEnvironment{thebibliography}{\small}

\widowpenalty=10000
\clubpenalty=10000

\title{\textbf{Quotient Observation}\\[0.2em]
{\large Hidden Locality, Schur Mediation, and Visible Precision Geometry}\\[0.2em]
{\large \textbf{DRAFT PREPRINT}}}
\author{J.\,R.~Dunkley}
\date{21st March 2026}

\hypersetup{
  pdftitle={Quotient Observation: Hidden Locality, Schur Mediation, and Visible Precision Geometry},
  pdfauthor={J. R. Dunkley},
  pdfsubject={Nonequilibrium Markov chains, quotient observation, visible precision geometry}
}

\begin{document}
\maketitle

\begin{abstract}
We study what can be recovered under partial observation and argue that the intrinsic visible object is quotient-visible precision, denoted Phi, rather than the Donsker-Varadhan Hessian itself. The paper develops a finite, compositional theory of this visible object on the positive cone of precision operators. We prove that quotient observation composes along towers of observation, identify the visible algebraic descendant of the bridge developed in Finite, and show that in the scalar visible case the resulting structure closes into a closed system of precision, clock, transport, and conservation laws. This yields a picture of hidden mediation, including support rigidity, a hidden-load parametrisation, a unique additive interior clock, and a sharp distinction between quotient observation and naive compression. We then turn to nonequilibrium systems and analyse how the visible object is realised in practice. At a symmetry point we obtain a canonical hidden response carrier and a scalar collapse. Away from that point we identify a fast-memory response layer with a two-state homogenised limit, a first-correction Schur carrier, and a source law with an explicit finite-epsilon obstruction. In a separate strong-bias corridor regime we derive a different asymptotic sector and show that these regimes are both real and not presently unified by a single higher operator architecture. The theory is tested against a sequence of benchmark systems used to keep the developing framework tied to a real signal while algebraic closure was being pinned down. These include kinesin, proofreading, transformed semi-Markov, and hidden-feedback ladder tests. Together they support the central claim that quotient-visible precision is the renormalisation-natural visible invariant, while also making the present transfer boundaries explicit. The result is a finished quotient-visible core theory together with a sharply delimited correction-layer package, a falsification battery, and a clear map of what remains open.
\end{abstract}

\newpage

\tableofcontents

\newpage

\section{Introduction}
\label{sec:01_introduction}

\subsection{The observation problem}

Given that you only see part of a system, what can you recover about its internal structure?

A Markov process runs on $n$ states; an observer sees only the image under a linear surjection $C : \mathbb{R}^n \to \mathbb{R}^m$ with $m < n$. The question is, what quadratic structure on the visible sector is canonically determined by the full precision and the observation map, with no further choices?

Here we push the full precision through the observation constraint and extract the unique visible quadratic form that records the minimum energy required to realise each visible displacement. The result is a definite positive operator on the visible sector, the quotient-visible precision $\Phi$, and it is this operator that is the natural visible invariant.

The setting is the Donsker-Varadhan large-deviation geometry of empirical flows in finite-state Markov chains, \cite{DonskerVaradhan1975a,DonskerVaradhan1975b,DonskerVaradhan1976,Ellis1985,Touchette2009,BertiniFaggionatoGabrielli2015} following our earlier work \cite{Dunkley2026FiniteObservation} which established a finite quadratic bridge between the stationary backbone and the nonequilibrium Hessian. It also sits naturally alongside the network and stochastic-thermodynamic description of nonequilibrium observables \cite{Schnakenberg1976,LebowitzSpohn1999,Seifert2012,Esposito2012,Maes2020}. Now we ask, what does that bridge produce once hidden degrees of freedom are eliminated by quotient observation?

\subsection{Main claim and conceptual picture}

The central claim is that $\Phi$ is the intrinsic visible variable. It is defined variationally, it composes along towers of observation, it admits a stable first-order calculus, and it is provably different from compression. No other candidate satisfies all properties simultaneously:

\[
\begingroup
\setlength{\fboxsep}{8pt}%
\setlength{\fboxrule}{1.15pt}%
\boxed{\scalebox{1.08}{$\displaystyle \Phi_T^{\circ}=\left\{X=T^{1/2}\Pi T^{1/2}:\; 0\prec \Pi\le I\ \text{ on }\ S=\operatorname{Ran}(T)\right\}$}}%
\endgroup
\]

\[
\begingroup
\setlength{\fboxsep}{8pt}%
\setlength{\fboxrule}{1.15pt}%
\boxed{\scalebox{1.08}{$\displaystyle \overline{\Phi}_T=\left\{X=T^{1/2}\Pi T^{1/2}:\; 0\le \Pi\le I\ \text{ on }\ S=\operatorname{Ran}(T)\right\}=T^{1/2}[0,I]_S T^{1/2}$}}%
\endgroup
\]

The Donsker-Varadhan bridge \cite{Dunkley2026FiniteObservation} identifies a finite visible algebraic object, the tangent descendant of $\Phi$, but it is $\Phi$ itself that closes under observation. The Hessian is a bridge object: it matters because it leads to $\Phi$. More precisely, once the tangent ceiling has been fixed on the part of visible space where it acts nontrivially, the exact observed class consists of all visible laws obtained by contracting within that same active sector, without ever creating new visible directions. Finite hidden sectors realise the strict interior of that class, while allowing limits adds the boundary cases. Section~\ref{sec:04_hidden_mediation} makes this precise.

The paper contains two species of results. The first is an algebraic theory of the visible object: its variational characterisation, composition law, conservation structure, hidden-load parametrisation, and graph-local decay. These hold for any quotient observation on any positive precision operator.

The second evaluates this theory in specific nonequilibrium systems: symmetry theorems, fast-memory homogenisation, corridor asymptotics, and thermodynamic constraints. The boundary between the two is maintained throughout.

\subsection{Main results}

The paper establishes the following results, listed in the order they appear.

\begin{enumerate}[label=(\roman*)]
\item Quotient observation on the positive cone is variational, composes along observation towers, and produces a coisometric tangent transport with monotone Ky Fan detectability envelopes.

\item The Donsker-Varadhan bridge produces a finite visible algebraic object via quotient reduction. The first genuinely new quotient effect beyond the tangent law is quartic and sign-definite; smooth divergence functionals built from the visible shadow lose only at sextic order.

\item In scalar visible rank, the visible precision admits a generator-level Schur simplification and does not depend on the hidden completion at the matched visible-law level.

\item The visible precision splits as $\Phi = F_0 + W$, with $F_0$ the reversible backbone contribution and $W$ the irreversible correction. This conservation identity holds without remainder.

\item Hidden mediation is parametrised by a positive hidden-load operator on the tangent support. Sequential transport, a unique additive determinant clock, and support rigidity follow.

\item Graph-local decay of hidden contributions: the visible effect of a hidden perturbation decays with the graph distance between the perturbation site and the visible sector.

\item On the Bell square, common-cause visible Gaussian gluing is a law-level compatibility problem whose frontier has two independent coordinates.

\item At the symmetric point of the two-state visible ladder, a symmetry theorem identifies a canonical hidden response carrier, and the visible algebra collapses to one free parameter: $F_0 = \alpha^2/(T+\alpha)$ and $W = \alpha T/(T+\alpha)$.

\item Fast-memory homogenisation yields a source law and an explicit obstruction to hidden-structure recovery.

\item Corridor asymptotics separate three regimes (fast memory, intermediate, slow memory) with sharp crossover diagnostics.

\item A thermodynamic ceiling constrains visible precision: $\Phi \leq F_0 + \sigma/2$, where $\sigma$ is the entropy production rate.

\item Timescale duality, the EPR-Fisher bound, and tightness at both timescale limits constrain the visible object within a closed region of parameter space.
\end{enumerate}

\subsection{Roadmap}

Section~\ref{sec:quotient-precision} formulates quotient observation on the positive cone, proves the variational characterisation and the Fr\'echet calculus, establishes tower composition, identifies the Schur gap, and names $\Phi$ as the intrinsic visible invariant. Section~\ref{sec:03_dv_bridge} connects the quotient theory to the Donsker-Varadhan bridge from \cite{Dunkley2026FiniteObservation}, establishes the visible algebraic object, and proves the quartic defect and sextic robustness theorems.

Sections~\ref{sec:04_hidden_mediation}-\ref{sec:06_bell_frontier} develop the internal structure of $\Phi$: hidden mediation and the conservation split, hidden locality and graph-local decay, and the Bell-square gluing frontier. Sections~\ref{sec:07_exact_symmetry}-\ref{sec:09_corridor_regime} evaluate the algebraic theory in nonequilibrium systems: symmetry at the symmetric point, fast-memory homogenisation with source law and obstruction, and corridor asymptotics with regime separation.

Section~\ref{sec:10_benchmark_synthesis} synthesises the benchmark evidence, Section~\ref{sec:11_falsification} presents falsification diagnostics, Section~\ref{sec:12_scope} states the proved boundary, and Section~\ref{sec:13_discussion} discusses what remains open and what geometric structure lies at the frontier.

\subsection{Notation and glossary}

The \emph{upstairs} space $\mathbb{R}^n$ carries the full precision operator $H$; the \emph{downstairs} space $\mathbb{R}^m$ carries the visible precision $\Phi$. An \emph{observation} is a surjective linear map $C : \mathbb{R}^n \to \mathbb{R}^m$. An \emph{observation tower} is a chain of surjections $\mathbb{R}^n \to \mathbb{R}^{m_1} \to \mathbb{R}^{m_2}$; the composition theorem says that observing once or in two steps gives the same visible precision.

The \emph{backbone} $H_0$ is the time-symmetric part of the full precision. The \emph{irreversible correction} $W$ is the contribution from detailed-balance-breaking currents. The \emph{conservation split} is $\Phi = F_0 + W$, where $F_0 := \Phi_C(H_0)$ is the quotient-visible backbone.

We write $\Sym_{++}(n)$ for the cone of $n \times n$ symmetric positive definite matrices, $\Phi_C(H)$ or $F_C(H)$ for the quotient map, and $\succeq$ for the Loewner partial order.


\section{Quotient observation and the intrinsic visible object}
\label{sec:quotient-precision}

For a positive precision operator $H \in \Sym_{++}(n)$ and a surjective observation map $C \in \mathbb{R}^{m \times n}$, the visible operator is
\begin{equation}
\Phi_C(H)
:=
F_C(H)
:=
\bigl(C H^{-1} C^{\mathsf T}\bigr)^{-1}.
\label{eq:quotient-observation-map}
\end{equation}
This is the visible precision canonically induced by $(H,C)$. It is the positive-definite shorted or Schur-complement construction associated with the observation map in the operator-theoretic sense of \cite{Schur1917,Anderson1971,AndersonDuffin1969,AndersonTrapp1975,Bhatia2007}. It is not chosen or assumed, instead forced by the observation problem itself: for each visible displacement, it records the least full-space energy needed to realise that displacement.

It has a variational meaning, it admits a clean differential calculus, it composes along towers of observation, and it differs in a controlled way from naive visible compression. Those facts identify $\Phi$ as the visible invariant carried by quotient observation.

\subsection{Setting and quotient observation on the positive cone}

Fix a surjective linear map $C : \mathbb{R}^n \to \mathbb{R}^m$. For each visible displacement $y \in \mathbb{R}^m$, consider the constrained quadratic energy
\begin{equation}
\mathcal E_H(y)
:=
\inf\!\left\{ x^{\mathsf T} H x : Cx = y \right\}.
\label{eq:full-space-energy}
\end{equation}
The positive cone is forced from the start: once the ambient precision is positive definite, the minimisation problem is strictly convex, the minimiser is unique, and the visible quadratic form is again positive definite. The quotient-visible precision is characterised by
\begin{equation}
\mathcal E_H(y)=y^{\mathsf T} \Phi_C(H) y.
\label{eq:variational-characterisation}
\end{equation}
So the visible operator is not an auxiliary construction, it is the quadratic form selected by the problem.

The same minimisation admits a matching dual formulation. The constrained primal energy upstairs and the dual visible energy downstairs agree, and that agreement is the algebra behind the quotient map. From it one obtains well posedness on $\Sym_{++}(n)$, positive homogeneity, monotonicity in Loewner order, and concavity, in the standard positive-operator framework of \cite{Loewner1934,Ando1979,HansenPedersen1982,Bhatia1997,Bhatia2007}. One also obtains a canonical $H_0$-horizontal chart, which separates genuine visible motion from gauge motion inside the hidden kernel and makes the later tangent calculus transparent.

\begin{theorem}[Variational characterisation]
\label{thm:variational}
For every $H \in \Sym_{++}(n)$ and every surjective $C$, the quotient-visible precision $\Phi_C(H)$ is the unique operator in $\Sym_{++}(m)$ satisfying \eqref{eq:variational-characterisation}. Equivalently,
\begin{equation}
y^{\mathsf T}\Phi_C(H)y
=
\inf_{Cx=y} x^{\mathsf T} H x
\qquad
(y \in \mathbb{R}^m).
\end{equation}
\end{theorem}

\begin{corollary}[Primal-dual decomposition]
The visible energy admits matching primal and dual representations. In particular, the minimising lift and the maximising visible multiplier encode the same quadratic form.
\end{corollary}

In this language quotient precision is the minimum-energy visible form induced by the upstairs law, and it sits naturally in the shorted-operator and parallel-addition side of positive-operator theory \cite{AndersonDuffin1969,Anderson1971,KuboAndo1980,PuszWoronowicz1975}.

\begin{theorem}[Basic positive-cone calculus]
\label{thm:basic-F}
The map $H \mapsto \Phi_C(H)$ is well posed on $\Sym_{++}(n)$, positively homogeneous, monotone, and concave. Relative to a fixed base point $H_0$, it admits a canonical horizontal chart adapted to quotient observation.
\end{theorem}

What has been gained already is decisive. Observation no longer looks like a loss of structure followed by arbitrary repair. It produces a definite positive operator with a canonical variational meaning.

\subsection{Fr\'echet calculus and the tangent ceiling}

Once the quotient map is fixed, its first-order calculus is immediate. For a symmetric perturbation $\dot H$, the derivative is the visible shadow of that perturbation transported through $H^{-1}$ and returned by $C$. The main point is not the formula itself, but what it bounds: positive perturbations upstairs induce a maximal first-order visible response, and later hidden-mediation effects can only sit beneath that ceiling.

\begin{theorem}[Fr\'echet derivative]
\label{thm:frechet-derivative}
The quotient map is Fr\'echet differentiable on $\Sym_{++}(n)$. Its derivative at $H$ in direction $\dot H$ is
\begin{equation}
D\Phi_C[H](\dot H)
=
\Phi_C(H)\, C H^{-1} \dot H H^{-1} C^{\mathsf T}\, \Phi_C(H).
\label{eq:derivative-schur}
\end{equation}
\end{theorem}

\begin{proposition}[Tangent ceiling]
If $\dot H \succeq 0$, then $D\Phi_C[H](\dot H) \succeq 0$. More sharply, $D\Phi_C[H](\dot H)$ is the maximal first-order visible response compatible with the positive perturbation $\dot H$.
\end{proposition}

The tangent ceiling is infrastructure rather than endpoint. It identifies the maximal first-order visible response compatible with a given positive perturbation. Later hidden-mediation results will show that hidden structure redistributes or attenuates visible strength beneath this ceiling, but does not create a new visible direction from nothing.

\subsection{Exact composition along observation towers}

The decisive structural fact is composition. Suppose observation occurs in stages,
\begin{equation}
\mathbb{R}^n \xrightarrow{\ C_1\ } \mathbb{R}^{m_1} \xrightarrow{\ C_2\ } \mathbb{R}^{m_2},
\qquad
C = C_2 C_1.
\end{equation}
Then one may observe once or observe in two steps. The result is the same:
\begin{equation}
\Phi_C(H)
=
\Phi_{C_2}\!\bigl(\Phi_{C_1}(H)\bigr).
\label{eq:tower-composition}
\end{equation}
So quotient observation closes on itself. The visible object keeps its type under repeated elimination of hidden structure.

\begin{theorem}[Composition law]
\label{thm:composition-law}
For surjective observation maps $C_1$ and $C_2$ with $C=C_2C_1$,
\begin{equation}
\Phi_C(H)
=
\Phi_{C_2}\!\bigl(\Phi_{C_1}(H)\bigr).
\end{equation}
\end{theorem}

This is what makes $\Phi$ a renormalisation-natural object rather than an ad hoc one, in the precise sense that repeated elimination stays inside one visible variable rather than generating a new effective object at each stage \cite{Kadanoff1966,WilsonKogut1974,Polchinski1984}. No correction term is needed. No auxiliary state has to be remembered. Repeated observation stays inside the same category and returns the same kind of operator.

The tangent theory composes as well. The visible derivative shadow transports down an observation tower by coisometric compression. At the spectral level this yields monotone Ky Fan detectability envelopes: as observation becomes coarser, the visible singular spectrum can only contract \cite{Fan1951,HornJohnson2013}.

\begin{theorem}[Tangent shadow transport]
\label{thm:shadow-transport}
Under tower composition, the derivative shadow of quotient observation transports by a coisometric map $W$ satisfying $WW^{\mathsf T} = I$. Specifically, if $C = C_2 C_1$ and $\Phi_1 := \Phi_{C_1}(H)$, the derivative at the coarser level factors through the derivative at the finer level via coisometric compression. In particular, the Ky Fan singular values of the visible derivative shadow contract monotonically along any observation tower.
\end{theorem}

With composition known, the role of $\Phi$ changes. It is no longer merely a good definition. It is the object that survives observation in the only way a visible invariant should.

\subsection{Compression versus quotient: the Schur gap}

Quotient observation is not naive visible compression. Fix a visible-hidden splitting and write
\begin{equation}
H
=
\begin{pmatrix}
H_{vv} & H_{vh} \\
H_{hv} & H_{hh}
\end{pmatrix}.
\end{equation}
Compression keeps the principal visible block $H_{vv}$. Quotient observation eliminates the hidden sector through the inverse and returns the Schur complement \cite{Zhang2005,Haynsworth1968,CarlsonHaynsworthMarkham1974}
\begin{equation}
\Phi(H)=H_{vv}-H_{vh}H_{hh}^{-1}H_{hv}.
\label{eq:schur-complement}
\end{equation}
The difference is the Schur gap
\begin{equation}
G(H)
:=
H_{vv}-\Phi(H)
=
H_{vh}H_{hh}^{-1}H_{hv}
\succeq 0.
\label{eq:schur-gap}
\end{equation}
So compression and quotient coincide only when the hidden sector does not feed back into the visible one.

\begin{proposition}[Schur gap formula]
For every block decomposition adapted to a visible-hidden splitting, the difference between naive compression and quotient observation is the positive semidefinite Schur gap.
\end{proposition}

\begin{remark}[Scalar diagnostics]
In scalar visible rank, useful summaries include $\operatorname{tr} G$, $\|G\|_{\mathrm{op}}$, and the corresponding log-determinant defect. These are diagnostics of hidden contribution.
\end{remark}

\subsection{The intrinsic visible invariant \texorpdfstring{$\Phi$}{Phi}}

The story is now fixed and quotient observation produces a definite visible precision. That operator has a variational characterisation, a stable first-order calculus, a composition law along observation towers, and a precise distinction from naive compression. Those properties identify $\Phi$ as the intrinsic visible invariant.

The Donsker-Varadhan Hessian has a different role. It is a bridge object upstairs, and it will matter in the next section because it leads to a finite visible algebraic descendant. But the object that closes under observation is $\Phi$, with the Donsker-Varadhan law surviving as its tangent ceiling.

That is the operator one can transport, compare, and compose without leaving the observation category. The paired-spectrum obstruction (\S\ref{sec:04_hidden_mediation}) makes this precise: at fixed visible dimension, the Donsker-Varadhan Hessian generically does not close, and the quotient precision is the correct replacement.

A second structural theme begins here. Later results will show that, in scalar visible rank, the visible precision splits as
\begin{equation}
\Phi = F_0 + W,
\label{eq:conservation-split-preview}
\end{equation}
with $F_0$ the reversible backbone and $W$ the irreversible correction. The point here is only conceptual: $\Phi$ is large enough to carry both contributions while remaining stable under observation.


\section{The Donsker-Varadhan bridge and the visible algebraic object}
\label{sec:03_dv_bridge}

The Donsker-Varadhan bridge from \cite{Dunkley2026FiniteObservation} produces a finite visible algebraic object, the tangent descendant of $\Phi$. The quartic and sextic defect theorems below measure how much structure survives the transition from tangent to full covariance level.

\subsection{The bridge and its quotient reduction}

Section~3 of \cite{Dunkley2026FiniteObservation} proved the exact quadratic Donsker-Varadhan bridge
\begin{equation}
H_{\mathrm{DV}} = H_0 + \Delta_{\mathrm{DV}},
\qquad
\Delta_{\mathrm{DV}} = \tfrac{1}{4}\, K\, H_0^{-1}\, K^{\mathsf T} \succeq 0,
\label{eq:dv-bridge}
\end{equation}
where $H_0$ is the time-symmetric backbone precision and $K$ is the skew current operator. We read the output of this bridge through quotient observation.

Introduce a perturbation parameter by $K = \eps K_1$, so that
\begin{equation}
H(\eps) = H_0 + \eps^2 \Delta_2 + O(\eps^4).
\label{eq:eps-bookkeeping}
\end{equation}
In particular,
\begin{equation}
\Delta_2 = \tfrac{1}{4}\, K_1\, H_0^{-1}\, K_1^{\mathsf T}.
\label{eq:Delta2-definition}
\end{equation}

\begin{definition}[$H_0$-adapted quotient sector]
\label{def:h0-adapted}
A decomposition $V = U \oplus U^{\perp}$ is called $H_0$-adapted if
\begin{equation}
H_0 =
\begin{pmatrix}
H_{0,S} & 0 \\
0 & H_{0,F}
\end{pmatrix}
\label{eq:h0-adapted}
\end{equation}
in the associated block coordinates.
\end{definition}

The algebraic cascade from bridge to visible object proceeds in three steps:
\begin{equation}
\begin{gathered}
H_{\mathrm{DV}} \;=\; H_0 + \tfrac{1}{4}\, K\, H_0^{-1}\, K^{\mathsf T}
\qquad [n \times n \text{ bridge}] \\[4pt]
\big\downarrow \quad \text{quotient observation } F_C \\[4pt]
\Phi_b \;=\; \operatorname{Schur}_v(H_b)
\qquad [m \times m \text{ visible algebraic object}] \\[4pt]
\big\downarrow \quad \text{composition (Theorem~\ref{thm:composition-law})} \\[4pt]
\Phi \;=\; F_C(H)
\qquad [\text{intrinsic visible invariant, closes}]
\end{gathered}
\label{eq:algebraic-cascade}
\end{equation}

\begin{remark}[Source, propagator, field structure]
The bridge \eqref{eq:dv-bridge} has a physical reading. The current operator $K$ is the source, the resolvent $H_0^{-1}$ is the propagator, and the product $K H_0^{-1} K^{\mathsf T}$ is the induced field correction. The visible algebraic object is the gauge-invariant visible projection of this chain.
\end{remark}

\subsection{The visible algebraic object}

Write the finite reduced block precision as
\begin{equation}
H_b
\;=\;
H_{0,b} + \tfrac{1}{4}\, K_b\, H_{0,b}^{-1}\, K_b^{\mathsf T}
\;=\;
\begin{pmatrix}
A & B \\
B^{\mathsf T} & D
\end{pmatrix},
\label{eq:hb-block}
\end{equation}
where the first coordinate is the visible direction and the remaining coordinates are hidden reduced directions.

\begin{theorem}[Visible algebraic object]
\label{thm:visible-algebraic-object}
The visible algebraic descendant of the bridge output is the Schur complement
\begin{equation}
\Phi_b \;:=\; \operatorname{Schur}_v(H_b) \;=\; A - B\, D^{-1}\, B^{\mathsf T}.
\label{eq:phi-schur}
\end{equation}
The Donsker-Varadhan correction is quadratic in the skew channel, and quotient elimination of the hidden reduced directions turns the bridge output into a visible Schur complement. By the composition theorem (Theorem~\ref{thm:composition-law}), $\Phi_b$ is a tangent descendant of the intrinsic visible precision $\Phi$.
\end{theorem}

\subsection{Scalar visible Schur identity and completion-independence package}
\label{subsec:scalar-schur}

In scalar visible rank ($m = 1$), the Schur complement reduces to a generator-level identity that does not require explicit diagonalisation of the full precision. The visible precision is a rational function of the generator entries, computable from the stationary distribution and transition rates without spectral decomposition.

\begin{theorem}[Completion independence at $M=1$]
\label{thm:completion-independence}
At the matched visible-law level, the quotient-visible precision $\Phi$ does not depend on the hidden completion. Two hidden completions that produce the same visible law yield the same visible precision.
\end{theorem}

The completion-independence hierarchy:

\begin{itemize}
\item Exact for $M = 1$ (proved).
\item Verified numerically for tested ladder families at $M = 2$.
\item Open for general $M$ and general graph topologies.
\end{itemize}

\subsection{Intrinsic quartic shadow defect and sextic divergence robustness}

The quartic order is where quotient observation first departs from the tangent law.

\begin{theorem}[Signed quartic bridge]
\label{thm:quartic-bridge}
Assume the Donsker-Varadhan expansion \eqref{eq:dv-bridge} with bookkeeping \eqref{eq:eps-bookkeeping}, and let $V = U \oplus U^{\perp}$ be $H_0$-adapted. Write $\Delta_{SF}^{(2)}$ for the $(S,F)$ block of $\Delta_2$ from \eqref{eq:Delta2-definition}. Then:
\begin{enumerate}
\item compression recovers the finite-observation visible quadratic law on $U$;
\item quotient and compression agree through order $\eps^2$;
\item the first genuinely new quotient effect is quartic and sign-definite:
\begin{equation}
H_U^{\mathrm{comp}} - H_U^{\mathrm{quot}}
=
\eps^4\, \Delta_{SF}^{(2)}\, H_{0,F}^{-1}\, \Delta_{FS}^{(2)} + O(\eps^6) \succeq 0.
\label{eq:quartic-gap}
\end{equation}
\end{enumerate}
\end{theorem}

\begin{proof}
Because the splitting is $H_0$-adapted, the mixed backbone blocks vanish. Hence
\[
H_{SF}(\eps) = \eps^2 \Delta_{SF}^{(2)} + O(\eps^4),
\qquad
H_{FF}(\eps) = H_{0,F} + O(\eps^2).
\]
Compression gives the visible block $H_{SS} = H_{0,S} + \eps^2 \Delta_{SS}^{(2)} + O(\eps^4)$, the finite-observation visible quadratic law from \cite{Dunkley2026FiniteObservation}. The Schur gap formula (\S\ref{sec:quotient-precision}) gives
\[
H_U^{\mathrm{comp}} - H_U^{\mathrm{quot}} = H_{SF}\, H_{FF}^{-1}\, H_{FS}.
\]
Substituting the block expansions yields \eqref{eq:quartic-gap}.
\end{proof}

\begin{corollary}[Tangent closure of visible Donsker-Varadhan geometry]
\label{cor:tangent-closure}
In the $H_0$-adapted setting,
\begin{equation}
H_U^{\mathrm{quot}} = H_{0,S} + (\Delta_{\mathrm{DV}})_{SS} + O(\norm{K}^4).
\label{eq:tangent-closure}
\end{equation}
The finite-observation visible law is the quadratic tangent law of quotient observation in the Donsker-Varadhan sector.
\end{corollary}

\begin{proof}
Combine Theorem~\ref{thm:quartic-bridge} with \eqref{eq:dv-bridge}.
\end{proof}

The signed quartic bridge identifies where quotient observation departs from compression. At the algebraic level it is a square-completion defect with definite sign, in the same broad spirit as energy-splitting identities familiar from other positive-definite variational settings \cite{Bogomolnyi1976}. The next theorem shows that this departure is intrinsic to the quotient geometry.

\begin{theorem}[Intrinsic quartic shadow defect and sextic divergence robustness]
\label{thm:shadow-sextic}
Fix a surjection $C : V \to Y$ and a backbone $H_0 \in \SPD(V)$. Let
\begin{equation}
H(\eps) = H_0 + \eps^2 \Delta_2 + O(\eps^4),
\qquad
\Delta_2 \succeq 0,
\label{eq:dv-general-family}
\end{equation}
and set $F_0 := F_C(H_0)$. Define the tangent visible correction
\begin{equation}
T_0 := DF_C(H_0)[\Delta_2],
\label{eq:tangent-visible-correction}
\end{equation}
the visible shadow
\begin{equation}
\Sigma_{\mathrm{exact}}(\eps)
:=
F_0^{-1/2}\bigl(F_C(H(\eps)) - F_0\bigr)F_0^{-1/2},
\label{eq:exact-shadow}
\end{equation}
and the tangent shadow
\begin{equation}
\Sigma_{\mathrm{tan}}(\eps)
:=
\eps^2\, F_0^{-1/2}\, T_0\, F_0^{-1/2}.
\label{eq:tangent-shadow}
\end{equation}
Then:
\begin{enumerate}
\item the quotient-visible expansion has the intrinsic form
\begin{equation}
F_C(H(\eps))
=
F_0 + \eps^2 T_0 - \eps^4 Q_0(\Delta_2) + O(\eps^6),
\label{eq:hessian-defect-expansion}
\end{equation}
where
\begin{equation}
Q_0(\Delta_2) := -\tfrac{1}{2}\, D^2 F_C(H_0)[\Delta_2, \Delta_2] \succeq 0;
\label{eq:hessian-defect}
\end{equation}
\item the visible shadow loses first at quartic order,
\begin{equation}
\Sigma_{\mathrm{exact}}(\eps)
=
\Sigma_{\mathrm{tan}}(\eps) - \eps^4 M_0 + O(\eps^6),
\qquad
M_0 := F_0^{-1/2}\, Q_0(\Delta_2)\, F_0^{-1/2} \succeq 0;
\label{eq:quartic-shadow-defect}
\end{equation}
\item for every smooth spectral functional $\Psi$ on symmetric matrices with $\Psi(0) = 0$ and $D\Psi(0) = 0$,
\begin{equation}
\Psi\!\bigl(\Sigma_{\mathrm{tan}}(\eps)\bigr) - \Psi\!\bigl(\Sigma_{\mathrm{exact}}(\eps)\bigr) = O(\eps^6).
\label{eq:sextic-functional-gap}
\end{equation}
In particular, reverse Gaussian KL, forward Gaussian KL, and squared Hellinger first lose only at sextic order.
\end{enumerate}
\end{theorem}

\begin{proof}
The quotient map $F_C$ is smooth on the positive cone, so Taylor expansion at $H_0$ along the curve $H(\eps)=H_0+\eps^2\Delta_2$ gives
\[
F_C(H(\eps))
=
F_C(H_0)+\eps^2 DF_C(H_0)[\Delta_2]+\tfrac12\eps^4 D^2F_C(H_0)[\Delta_2,\Delta_2]+O(\eps^6).
\]
Using the definitions of $F_0$, $T_0$, and $Q_0(\Delta_2):=-\tfrac12 D^2F_C(H_0)[\Delta_2,\Delta_2]$ yields \eqref{eq:hessian-defect-expansion}. Conjugating by $F_0^{-1/2}$ gives \eqref{eq:quartic-shadow-defect}.

For the spectral-functional claim, both $\Sigma_{\mathrm{tan}}(\eps)$ and $\Sigma_{\mathrm{exact}}(\eps)$ are $O(\eps^2)$, and their difference is $O(\eps^4)$ by \eqref{eq:quartic-shadow-defect}. Since $\Psi(0)=0$ and $D\Psi(0)=0$, the Taylor expansion of $\Psi$ at the origin begins quadratically. Applying that expansion to the pair $(\Sigma_{\mathrm{tan}}(\eps),\Sigma_{\mathrm{exact}}(\eps))$ shows that the first possible difference appears at order $\eps^6$, which proves \eqref{eq:sextic-functional-gap}.
\end{proof}

\begin{remark}
In the canonical $H_0$-horizontal chart of \S\ref{sec:quotient-precision}, the quartic Hessian defect $Q_0(\Delta_2)$ reduces to the signed quartic bridge of Theorem~\ref{thm:quartic-bridge}. The adapted block formula is the coordinate representation of an intrinsic quotient-geometric defect.
\end{remark}

\subsection{What the bridge still contributes conceptually}

\begin{remark}[Computation versus interpretation]
In the scalar visible case ($m = 1$), the bridge may be algebraically bypassed: the visible precision can be computed directly from the generator. The bridge nonetheless matters as the systematic path from the Donsker-Varadhan Hessian to the quotient-visible precision. The distinction is between computation (which the scalar case streamlines) and interpretation (which the bridge provides at any rank).
\end{remark}

\subsection{Boundaries of the scalar simplification}

The scalar simplification of \S\ref{subsec:scalar-schur} does not extend to general visible rank. When $m > 1$, the visible precision $\Phi$ is a matrix whose internal structure (eigenvalues, eigenvectors, parametric dependence) is not determined by scalar diagnostics.

Visible-law encoding of $\Phi$ is $M = 1$-specific: one can read $\Phi$ from the visible switching rates. For general $M$, the visible law constrains but does not determine $\Phi$. What additional structure is needed at higher rank remains open.

The algebraic theory (\S\ref{sec:quotient-precision}) and the bridge connection are in hand. The next sections develop the internal structure of $\Phi$: hidden mediation, conservation, locality, and the gluing frontier.


% Sections 4 onwards intentionally reduced to headers only

\section{Hidden mediation, scalar clocks, and sequential transport}
\label{sec:04_hidden_mediation}

Section~\ref{sec:03_dv_bridge} identified the visible algebraic object. The next question is internal: what hidden mechanism produces a visible law beneath the tangent ceiling, how does that mechanism compose, and which scalar quantity measures its cumulative action? The answer is hidden mediation. Beneath a fixed tangent ceiling, every support-preserving visible law is classified by a positive hidden load, sequential composition transports that load by a rigid conjugation law, and the only continuous orthogonally invariant additive clock on the interior is the determinant clock. The section closes by isolating the obstruction that prevents same-dimension Donsker-Varadhan closure from being the correct observable class.

\subsection{Exact latent-contraction formula and the tangent ceiling}

Work in an $H_0$-adapted quotient chart and whiten by the backbone so that the upstairs precision takes the form
\begin{equation}\label{eq:whitened-LLt}
\widetilde H = I + L L^{\mathsf T},
\qquad
L=
\begin{pmatrix}
B\\
M
\end{pmatrix},
\end{equation}
where $B$ contains the visible rows and $M$ the hidden rows. In this representation the tangent ceiling is explicit: quotient observation keeps the same visible support ceiling as compression, but contracts beneath it by a positive hidden contribution. At the level of Gaussian graphical models this is hidden-variable precision geometry, not a heuristic latent-variable ansatz \cite{Dempster1972,Lauritzen1996,ChandrasekaranParriloWillsky2012}.

\begin{theorem}[Exact latent-contraction formula]\label{thm:latent-contraction}
In the setting \eqref{eq:whitened-LLt}, the visible quotient correction is
\begin{equation}\label{eq:exact-visible-correction}
X = B(I+M^{\mathsf T} M)^{-1} B^{\mathsf T}.
\end{equation}
The tangent ceiling is
\begin{equation}\label{eq:tangent-ceiling}
T:=BB^{\mathsf T},
\end{equation}
and the gap is
\begin{equation}\label{eq:exact-gap}
T-X = B\,M^{\mathsf T}(I+MM^{\mathsf T})^{-1}M\,B^{\mathsf T} \succeq 0.
\end{equation}
Equivalently,
\begin{equation}\label{eq:X-as-BRBt}
X = B R B^{\mathsf T},
\qquad
R:=(I+M^{\mathsf T} M)^{-1},
\qquad 0\prec R\le I.
\end{equation}
\end{theorem}

\begin{proof}
In block form,
\[
\widetilde H=
\begin{pmatrix}
I+BB^{\mathsf T} & BM^{\mathsf T}\\
MB^{\mathsf T} & I+MM^{\mathsf T}
\end{pmatrix}.
\]
The visible quotient precision is its Schur complement:
\[
I+BB^{\mathsf T} - BM^{\mathsf T}(I+MM^{\mathsf T})^{-1}MB^{\mathsf T}.
\]
Using the identity
\[
I - M^{\mathsf T}(I+MM^{\mathsf T})^{-1}M = (I+M^{\mathsf T}M)^{-1},
\]
we obtain
\[
I + B(I+M^{\mathsf T}M)^{-1}B^{\mathsf T}.
\]
Subtracting the whitened backbone $I$ gives \eqref{eq:exact-visible-correction}. The formulae \eqref{eq:exact-gap} and \eqref{eq:X-as-BRBt} follow immediately.
\end{proof}

\begin{corollary}[Exact support theorem]\label{cor:support-theorem}
For every finite hidden sector,
\begin{equation}\label{eq:support-theorem}
\operatorname{Ran}(X)=\operatorname{Ran}(T),
\qquad
\rank(X)=\rank(T).
\end{equation}
Thus finite hidden mediation preserves the visible correction support. In particular, the hidden sector attenuates inside the existing visible support rather than creating new visible directions, which is also the natural completion geometry for latent Gaussian models \cite{GroneJohnsonSaWolkowicz1984,ChandrasekaranParriloWillsky2012}.
\end{corollary}

\begin{proof}
Since $R\succ 0$ on the latent space, $BRB^{\mathsf T}$ and $BB^{\mathsf T}$ have the same range and rank.
\end{proof}

\begin{remark}[Blind directions remain blind]
Corollary~\ref{cor:support-theorem} is the support rigidity statement in ordinary language. If a visible direction receives no signal from the tangent ceiling $T$, then no finite hidden architecture can make that direction informative later, because the actual visible correction $X$ has the same support. Hidden complexity can attenuate or redistribute strength inside the tangent support, but it cannot create a new visible direction from nothing.
\end{remark}

\subsection{Support-preserving visible cone, its closure, and hidden-load parametrisation}

The latent-contraction formula says that hidden mediation acts by attenuation inside the support of the tangent ceiling. The next result identifies the full image of that action.

\begin{proposition}[Support-preserving interval and closure theorem]\label{prop:interval}
Let $T=BB^{\mathsf T}$ and let $S:=\operatorname{Ran}(T)$. A matrix $X$ arises from a finite hidden sector if and only if
\begin{equation}\label{eq:interval-form}
X=T^{1/2}\Pi T^{1/2},
\qquad 0\prec \Pi\le I \text{ on } S.
\end{equation}
Taking closure adds the boundary $0\le \Pi\le I$ on $S$. Equivalently, the closure of the visible image is the Loewner interval beneath $T$, whose interior is support-preserving.
\end{proposition}

\begin{proof}
From Theorem~\ref{thm:latent-contraction},
\[
X = BRB^{\mathsf T},
\qquad 0\prec R\le I.
\]
By polar decomposition, $B=T^{1/2}U$ for a partial isometry $U$ with initial space the latent support of $B$ and final space $S$. Therefore
\[
X=T^{1/2}(URU^{\mathsf T})T^{1/2},
\]
so \eqref{eq:interval-form} holds with $\Pi:=URU^{\mathsf T}$. Since $R\succ 0$ and $R\le I$, one has $0\prec\Pi\le I$ on $S$.

Conversely, let $0\prec\Pi\le I$ on $S$. Choose the latent space to be $S$, set $B:=T^{1/2}$, and define $R:=\Pi$. Since $R\succ 0$ and $R\le I$, the matrix $G:=R^{-1}-I$ is positive semidefinite. Pick $M$ with $M^{\mathsf T}M=G$. Then $(I+M^{\mathsf T}M)^{-1}=R$, so $X=BRB^{\mathsf T}$. The boundary case follows by limit.
\end{proof}

\noindent\textbf{Exact visible class and its closure.} With $S:=\operatorname{Ran}(T)$,

\[
\begingroup
\setlength{\fboxsep}{8pt}%
\setlength{\fboxrule}{1.05pt}%
\boxed{\displaystyle \Phi_T^{\circ}=\left\{X=T^{1/2}\Pi T^{1/2}:\; 0\prec \Pi\le I\ \text{ on }\ S=\operatorname{Ran}(T)\right\}}%
\endgroup
\]

\[
\begingroup
\setlength{\fboxsep}{8pt}%
\setlength{\fboxrule}{1.05pt}%
\boxed{\displaystyle \overline{\Phi}_T=\left\{X=T^{1/2}\Pi T^{1/2}:\; 0\le \Pi\le I\ \text{ on }\ S=\operatorname{Ran}(T)\right\}=T^{1/2}[0,I]_S T^{1/2}}%
\endgroup
\]

Here $[0,I]_S:=\{\Pi\in\Sym(S):0\le \Pi\le I_S\}$. Thus finite hidden sectors realise the support-preserving interior $\Phi_T^{\circ}$, while taking closure adds the boundary $\overline{\Phi}_T$. The hidden-load coordinate of Theorem~\ref{thm:hidden-load} is defined on $\Phi_T^{\circ}$, and the next results identify its rank and determinant invariants.

\begin{theorem}[Intrinsic hidden-load parametrisation]\label{thm:hidden-load}
Let $T\succeq 0$ and let $S:=\operatorname{Ran}(T)$. A visible law $X$ in the support-preserving interior beneath $T$ is equivalently an operator of the form
\[
X=T^{1/2}\Pi T^{1/2},
\qquad
0\prec \Pi\le I \text{ on } S.
\]
For such $X$, define the intrinsic hidden-load operator on $S$ by
\begin{equation}\label{eq:hidden-load}
\Lambda
:=
(T|_S)^{1/2}(X|_S)^{-1}(T|_S)^{1/2}-I_S.
\end{equation}
Then $\Lambda\succeq 0$ on $S$, and
\begin{equation}\label{eq:X-from-load}
X=T^{1/2}(I_S+\Lambda)^{-1}T^{1/2}.
\end{equation}
Conversely, every $\Lambda\succeq 0$ on $S$ determines a visible law by \eqref{eq:X-from-load}. Hence the support-preserving visible interior beneath $T$ is in bijection with the positive cone on $S$.
\end{theorem}

\begin{proof}
Write
\[
\Pi:=(T|_S)^{-1/2}(X|_S)(T|_S)^{-1/2}.
\]
Then $0\prec \Pi\le I_S$ on $S$. Therefore
\[
\Lambda=\Pi^{-1}-I_S \succeq 0,
\]
and
\[
\Pi=(I_S+\Lambda)^{-1}.
\]
Substituting this into $X=T^{1/2}\Pi T^{1/2}$ gives \eqref{eq:X-from-load}.

Conversely, if $\Lambda\succeq 0$ on $S$, then $\Pi:=(I_S+\Lambda)^{-1}$ satisfies
\[
0\prec \Pi\le I_S,
\]
so Proposition~\ref{prop:interval} yields a visible realisation
\[
X=T^{1/2}\Pi T^{1/2}=T^{1/2}(I_S+\Lambda)^{-1}T^{1/2}.
\]
This proves the bijection.
\end{proof}

\begin{corollary}[Hidden-load monotonicity]\label{cor:hidden-load-monotone}
Let
\[
X_i=T^{1/2}(I_S+\Lambda_i)^{-1}T^{1/2},
\qquad \Lambda_i\succeq 0 \text{ on } S,
\]
for $i=0,1$. Then
\begin{equation}\label{eq:hidden-load-order}
\Lambda_1\succeq \Lambda_0
\quad\Longleftrightarrow\quad
X_1\preceq X_0.
\end{equation}
Thus increasing hidden load can only attenuate the visible correction.
\end{corollary}

\begin{proof}
Restrict to $S$, where $T|_S\succ 0$. Congruence by $(T|_S)^{-1/2}$ gives
\[
X_1\preceq X_0
\quad\Longleftrightarrow\quad
(I_S+\Lambda_1)^{-1}\preceq (I_S+\Lambda_0)^{-1}.
\]
Since inversion reverses Loewner order on positive definite operators,
\[
(I_S+\Lambda_1)^{-1}\preceq (I_S+\Lambda_0)^{-1}
\quad\Longleftrightarrow\quad
I_S+\Lambda_1 \succeq I_S+\Lambda_0,
\]
which is \eqref{eq:hidden-load-order}.
\end{proof}

\begin{remark}[One positive-cone geometry seen from two sides]
The quotient programme now contains two positive-cone geometries. First, the global quotient map
\[
H \mapsto F_C(H)=(CH^{-1}C^{\mathsf T})^{-1}
\]
is monotone and concave on the upstairs positive cone. Second, beneath a fixed tangent ceiling $T$ on the active support $S=\operatorname{Ran}(T)$, every support-preserving visible law in the interior has the form
\[
X=T^{1/2}(I_S+\Lambda)^{-1}T^{1/2},
\qquad \Lambda\succeq 0.
\]
The first is the global quotient-side cone, the second is the intrinsic hidden-load cone after the ceiling has been fixed. They are the global and support-adapted faces of the same attenuation geometry.
\end{remark}

\begin{theorem}[Minimal hidden latent dimension]\label{thm:min-hidden-dim}
Let $X=T^{1/2}\Pi T^{1/2}$ with $0\prec\Pi\le I$ on $S=\operatorname{Ran}(T)$. Then the minimal hidden latent dimension needed to realise $X$ is
\begin{equation}\label{eq:min-hidden-dim}
d_{\min} = \rank(I-\Pi) = \rank(T-X).
\end{equation}
\end{theorem}

\begin{proof}
By Theorem~\ref{thm:hidden-load}, write
\[
X=T^{1/2}(I_S+\Lambda)^{-1}T^{1/2},
\qquad
\Lambda\succeq 0.
\]
Any realisation requires a Gram factorisation $\Lambda=M^{\mathsf T}M$, so the hidden latent dimension must be at least $\rank(\Lambda)$. Choosing a Gram factorisation of $\Lambda$ shows that $\rank(\Lambda)$ hidden dimensions suffice. Since
\[
T-X = T^{1/2}\Bigl(I_S-(I_S+\Lambda)^{-1}\Bigr)T^{1/2},
\]
and $I_S-(I_S+\Lambda)^{-1}=\Lambda(I_S+\Lambda)^{-1}$ has the same rank as $\Lambda$, we obtain
\[
\rank(\Lambda)=\rank(T-X)=\rank(I-\Pi).
\]
This proves the claim.
\end{proof}

\begin{corollary}[Rank and volume of hidden load]\label{cor:hidden-load-rank-volume}
With $\Lambda$ as in Theorem~\ref{thm:hidden-load},
\begin{equation}\label{eq:hidden-load-rank}
\rank(\Lambda)=\rank(T-X),
\end{equation}
and
\begin{equation}\label{eq:hidden-load-volume}
\log\pdet(T)-\log\pdet(X)=\log\det(I_S+\Lambda).
\end{equation}
In particular, the minimal hidden latent dimension is $\rank(\Lambda)=\rank(T-X)$.
\end{corollary}

\begin{proof}
From \eqref{eq:X-from-load},
\[
T-X
=
T^{1/2}\Bigl(I_S-(I_S+\Lambda)^{-1}\Bigr)T^{1/2}
=
T^{1/2}\Lambda(I_S+\Lambda)^{-1}T^{1/2}.
\]
Since $(I_S+\Lambda)^{-1}$ is invertible on $S$, this implies
\[
\rank(T-X)=\rank(\Lambda).
\]
Also, on $S$,
\[
X|_S=(T|_S)^{1/2}(I_S+\Lambda)^{-1}(T|_S)^{1/2},
\]
hence
\[
\det(X|_S)=\det(T|_S)\det(I_S+\Lambda)^{-1}.
\]
Taking logarithms gives \eqref{eq:hidden-load-volume}.
\end{proof}

\begin{remark}[Operator-interval picture]
Beneath a fixed tangent ceiling $T$, finite hidden sectors realise the support-preserving interior
\[
X=T^{1/2}\Pi T^{1/2},
\qquad 0\prec \Pi\le I \text{ on } S,
\]
and taking closure adds the boundary $0\le \Pi\le I$ on $S$. On the interior, the hidden-load coordinate $\Lambda=\Pi^{-1}-I$ turns the support-preserving class into the positive cone on the active support. The scalar quantity
\[
\log\det(I_S+\Lambda)=\log\pdet(T)-\log\pdet(X)
\]
measures latent contraction volume. Same-dimension skew-square closure will appear below as a special paired-spectrum slice inside this larger cone.
\end{remark}

\subsection{Sequential composition, hidden-load transport, and scalar clocks}
\label{sec:support_towers}

All statements in this subsection are made on one fixed active support stratum $S$. No claim is made here about a finite additive scalar clock across genuine rank-loss events.

\begin{proposition}[Support-preserving sequential effect law]\label{prop:sequential-effect-law}
Let
\[
\mathcal E^\circ(S):=\{\Pi\in\mathrm{Sym}(S):0\prec \Pi\le I_S\}.
\]
For $\Pi,\Xi\in\mathcal E^\circ(S)$ define
\begin{equation}\label{eq:sequential-effect-law}
\Pi\circ\Xi:=\Pi^{1/2}\Xi\Pi^{1/2}.
\end{equation}
Then $\Pi\circ\Xi\in\mathcal E^\circ(S)$.
\end{proposition}

\begin{proof}
Since $\Xi\succ 0$, congruence by $\Pi^{1/2}$ preserves positivity and gives $\Pi\circ\Xi\succ 0$. Also,
\[
0\prec \Xi\le I_S
\quad\Longrightarrow\quad
0\prec \Pi^{1/2}\Xi\Pi^{1/2}\le \Pi\le I_S.
\]
Hence $\Pi\circ\Xi\in\mathcal E^\circ(S)$.
\end{proof}

\begin{remark}[The visible law is not the associative object]
The visible product $\circ$ is not associative in general. The true associative object is the underlying contraction composition downstairs. The visible interval is the Gram-shadow image of that contraction composition.
\end{remark}

\begin{theorem}[Exact hidden-load transport law]\label{thm:hidden-load-transport}
Let
\[
\Pi=(I_S+\Lambda)^{-1},
\qquad
\Xi=(I_S+M)^{-1},
\qquad
\Lambda,M\succeq 0 \text{ on } S.
\]
Then
\begin{equation}\label{eq:load-transport-product}
\Pi\circ\Xi=(I_S+\Lambda_{\mathrm{tot}})^{-1},
\end{equation}
with
\begin{equation}\label{eq:load-transport-law}
\Lambda_{\mathrm{tot}}
=
\Lambda + (I_S+\Lambda)^{1/2} M (I_S+\Lambda)^{1/2}.
\end{equation}
\end{theorem}

\begin{proof}
Because
\[
\Pi^{-1/2}=(I_S+\Lambda)^{1/2},
\qquad
\Xi^{-1}=I_S+M,
\]
we have
\[
(\Pi\circ\Xi)^{-1}
=
\Pi^{-1/2}\Xi^{-1}\Pi^{-1/2}
=
(I_S+\Lambda)^{1/2}(I_S+M)(I_S+\Lambda)^{1/2}.
\]
Expanding the right-hand side gives
\[
(\Pi\circ\Xi)^{-1}
=
I_S + \Lambda + (I_S+\Lambda)^{1/2} M (I_S+\Lambda)^{1/2}.
\]
This is \eqref{eq:load-transport-law}.
\end{proof}

\begin{corollary}[Survivor transport law]\label{cor:survivor-transport}
Define the survivor complement
\[
\Sigma(\Pi):=I_S-\Pi.
\]
Then
\begin{equation}\label{eq:survivor-transport}
\Sigma(\Pi\circ\Xi)
=
\Sigma(\Pi)+\Pi^{1/2}\Sigma(\Xi)\Pi^{1/2}.
\end{equation}
\end{corollary}

\begin{proof}
Use
\[
I_S-\Pi^{1/2}\Xi\Pi^{1/2}
=
(I_S-\Pi)+\Pi^{1/2}(I_S-\Xi)\Pi^{1/2}.
\]
\end{proof}

The transport law isolates the correct associative structure, and the scalar clock theorem identifies the only continuous additive quantity compatible with it.

\begin{corollary}[Additive determinant clock and barrier identity]\label{cor:det-clock}
Define
\begin{equation}\label{eq:det-clock}
\tau(\Pi):=-\log\det\Pi
\end{equation}
on $\mathcal E^\circ(S)$. Then
\begin{equation}\label{eq:det-clock-additive}
\tau(\Pi\circ\Xi)=\tau(\Pi)+\tau(\Xi).
\end{equation}
In hidden-load coordinates,
\begin{equation}\label{eq:clock-load-barrier}
\tau(\Pi)=\log\det(I_S+\Lambda).
\end{equation}
If the unwhitened visible law is
\[
X=T^{1/2}\Pi T^{1/2},
\]
then on the active support,
\begin{equation}\label{eq:pdet-barrier}
\log\pdet(T)-\log\pdet(X)=\log\det(I_S+\Lambda).
\end{equation}
\end{corollary}

\begin{proof}
Since
\[
\det(\Pi^{1/2}\Xi\Pi^{1/2})=\det(\Pi)\det(\Xi),
\]
we obtain \eqref{eq:det-clock-additive}. Because $\Pi=(I_S+\Lambda)^{-1}$, this is equivalent to \eqref{eq:clock-load-barrier}. The pseudo-determinant identity is Corollary~\ref{cor:hidden-load-rank-volume}.
\end{proof}

\begin{theorem}[Uniqueness of the additive interior clock]\label{thm:clock-uniqueness}
Let $F:\mathcal E^\circ(S)\to\mathbb R$ be continuous, orthogonally invariant, and additive under $\circ$:
\[
F(\Pi\circ\Xi)=F(\Pi)+F(\Xi),
\qquad
F(I_S)=0.
\]
Then there exists $c\in\mathbb R$ such that
\begin{equation}\label{eq:clock-uniqueness}
F(\Pi)=c\,\log\det\Pi
\end{equation}
for all $\Pi\in\mathcal E^\circ(S)$.
\end{theorem}

\begin{proof}
By orthogonal invariance, it suffices to evaluate $F$ on diagonal effects. For diagonal matrices,
\[
\operatorname{diag}(a_1,\dots,a_r)\circ \operatorname{diag}(b_1,\dots,b_r)
=
\operatorname{diag}(a_1b_1,\dots,a_rb_r),
\]
with each $a_i,b_i\in(0,1]$. Hence continuity and additivity under coordinatewise multiplication imply
\[
F(\operatorname{diag}(a_1,\dots,a_r)) = c\sum_i \log a_i
\]
for some constant $c$. Orthogonal invariance extends this to all $\Pi\in\mathcal E^\circ(S)$.
\end{proof}

\begin{corollary}[No finite continuous additive clock on the closed interval]\label{cor:no-global-finite-clock}
Let $F:[0,I_S]\to\mathbb R$ be continuous, orthogonally invariant, additive under $\circ$, and satisfy $F(I_S)=0$. Then
\[
F\equiv 0.
\]
\end{corollary}

\begin{proof}
Restrict to the interior $0\prec \Pi\le I_S$. By Theorem~\ref{thm:clock-uniqueness},
\[
F(\Pi)=c\,\log\det\Pi
\]
there. Consider
\[
\Pi_\varepsilon=\operatorname{diag}(1,\dots,1,\varepsilon),
\qquad \varepsilon\downarrow 0.
\]
If $c\neq 0$, then $\log\det\Pi_\varepsilon=\log\varepsilon\to -\infty$, contradicting finite continuity on $[0,I_S]$. Thus $c=0$, hence $F\equiv 0$.
\end{proof}

\begin{corollary}[Extended-real completion of the clock]\label{cor:extended-clock}
Define
\[
\tau_{\mathrm{ext}}(\Pi)
=
\begin{cases}
-\log\det\Pi, & 0\prec \Pi\le I_S,\\
+\infty, & \Pi \text{ singular.}
\end{cases}
\]
Then $\tau_{\mathrm{ext}}$ is additive under $\circ$ in the extended-real sense.
\end{corollary}

\begin{proof}
Interior additivity is Corollary~\ref{cor:det-clock}. If either factor is singular, then $\Pi\circ\Xi$ is singular, so $+\infty$ is stable under composition.
\end{proof}

\begin{proposition}[Scalar mediation invariants]\label{prop:scalar-mediation}
Let
\[
G:=I_S-\Pi=\Lambda(I_S+\Lambda)^{-1}.
\]
If $\lambda_i$ are the eigenvalues of $\Lambda$, then the eigenvalues of $G$ are
\[
g_i=\frac{\lambda_i}{1+\lambda_i}.
\]
Hence
\begin{equation}\label{eq:trace-gap-load}
\operatorname{tr}(G)=\sum_i \frac{\lambda_i}{1+\lambda_i},
\end{equation}
\begin{equation}\label{eq:op-gap-load}
\|G\|_{\mathrm{op}}=\max_i \frac{\lambda_i}{1+\lambda_i},
\end{equation}
and
\begin{equation}\label{eq:logdet-gap-load}
\tau(\Pi)=\log\det(I_S+\Lambda)=\sum_i \log(1+\lambda_i).
\end{equation}
Thus trace, operator norm, and log-determinant provide complementary scalar diagnostics of total, peak, and volumetric hidden mediation.
\end{proposition}

\begin{proof}
Since $\Pi=(I_S+\Lambda)^{-1}$ is a spectral function of $\Lambda$, the operators $\Pi$, $G$, and $\Lambda$ commute and are simultaneously diagonalisable. On each eigenvector of $\Lambda$ with eigenvalue $\lambda_i$, the corresponding eigenvalue of $G$ is $\lambda_i/(1+\lambda_i)$. The displayed trace, norm, and log-determinant formulas follow immediately.
\end{proof}

\begin{proposition}[Exact superadditive ceiling-KL identities]\label{prop:ceiling-kl-superadditive}
Let
\[
A:=I_S+\Lambda,
\qquad
B:=I_S+M,
\qquad
A_{\mathrm{tot}}:=A^{1/2}BA^{1/2},
\]
and write
\[
P_0:=\mathcal N(0,I_S),
\qquad
P_\Lambda:=\mathcal N(0,A),
\qquad
P_M:=\mathcal N(0,B),
\qquad
P_{\Lambda_{\mathrm{tot}}}:=\mathcal N(0,A_{\mathrm{tot}}).
\]
Then
\begin{equation}\label{eq:reverse-kl-superadditive}
D_{KL}(P_{\Lambda_{\mathrm{tot}}}\|P_0)
=
D_{KL}(P_\Lambda\|P_0)+D_{KL}(P_M\|P_0)+\frac12\operatorname{tr}(\Lambda M),
\end{equation}
and, with
\[
G_\Lambda:=I_S-(I_S+\Lambda)^{-1},
\qquad
G_M:=I_S-(I_S+M)^{-1},
\]
one also has
\begin{equation}\label{eq:forward-kl-superadditive}
D_{KL}(P_0\|P_{\Lambda_{\mathrm{tot}}})
=
D_{KL}(P_0\|P_\Lambda)+D_{KL}(P_0\|P_M)+\frac12\operatorname{tr}(G_\Lambda G_M).
\end{equation}
In particular, both Gaussian KL directions relative to the ceiling are superadditive under support-preserving sequential quotienting.
\end{proposition}

\begin{proof}
For the reverse direction,
\[
D_{KL}(P_\Lambda\|P_0)=\frac12\bigl(\operatorname{tr}(A)-r-\log\det A\bigr),
\qquad r:=\dim S.
\]
Because
\[
\operatorname{tr}(A_{\mathrm{tot}})=\operatorname{tr}(AB)=r+\operatorname{tr}\Lambda+\operatorname{tr}M+\operatorname{tr}(\Lambda M)
\]
and
\[
\log\det A_{\mathrm{tot}}=\log\det A+\log\det B,
\]
substitution gives \eqref{eq:reverse-kl-superadditive}. For the forward direction,
\[
D_{KL}(P_0\|P_\Lambda)=\frac12\bigl(\log\det A+\operatorname{tr}(A^{-1})-r\bigr).
\]
Now
\[
A_{\mathrm{tot}}^{-1}=A^{-1/2}B^{-1}A^{-1/2},
\]
so cyclicity of trace gives
\[
\operatorname{tr}(A_{\mathrm{tot}}^{-1})=\operatorname{tr}(A^{-1}B^{-1}).
\]
Since
\[
A^{-1}=I_S-G_\Lambda,
\qquad
B^{-1}=I_S-G_M,
\]
we have
\[
\operatorname{tr}(A^{-1}B^{-1})=r-\operatorname{tr}G_\Lambda-\operatorname{tr}G_M+\operatorname{tr}(G_\Lambda G_M).
\]
Substituting into the forward Gaussian KL formula yields \eqref{eq:forward-kl-superadditive}.
\end{proof}

\begin{remark}[Hidden-load interaction measure]
The cross-term $\tfrac12\operatorname{tr}(\Lambda M)$ in \eqref{eq:reverse-kl-superadditive} is not bookkeeping noise. It is the interaction measure for sequential hidden mediations. It vanishes only when the two mediations are orthogonal in hidden-load space, and otherwise quantifies how much the second mediation is transported by the first before the visible law is read out.
\end{remark}

\subsection{Conservation split: $F_0 + W = \Phi$}

The determinant clock identifies the only additive scalar content of sequential mediation. The next statement isolates the additive operator content of the visible precision itself.

\begin{theorem}[Conservation split]\label{thm:conservation-split}
Fix a completion and write
\begin{equation}\label{eq:conservation-defs}
F_0 := F_{\mathrm{vis}}(H_0),
\qquad
\Phi := F_{\mathrm{vis}}(H_{\mathrm{alg}}),
\qquad
W := \Phi - F_0,
\end{equation}
where $H_0$ is the reversible backbone and $H_{\mathrm{alg}}$ is the algebraic precision carrying the full visible response. Then
\begin{equation}\label{eq:conservation-split}
\Phi = F_0 + W.
\end{equation}
In scalar visible rank, and more generally on every class of completions for which $\Phi$ is known to be completion-independent, the completion-dependent variations of $F_0$ and $W$ compensate so that the total visible precision remains fixed.
\end{theorem}

\begin{proof}
The identity \eqref{eq:conservation-split} is the definition of $W$ rewritten. The nontrivial statement is the compensation claim. In scalar visible rank this follows from the completion-independence package of Section~\ref{subsec:scalar-schur}: when two completions produce the same visible law, they produce the same $\Phi$, hence any change in $F_0$ must be cancelled by the opposite change in $W$. The same implication holds on any broader class of completions for which completion-independence of $\Phi$ has been established.
\end{proof}

\begin{remark}[Ward-identity reading]
The split \eqref{eq:conservation-split} has the structure of a Ward identity. The total visible precision $\Phi$ is the invariant object, while $F_0$ and $W$ depend on the chosen completion but are locked in sum. The paper does not need gauge language for the mathematics, but it is the right interpretation of the compensation mechanism.
\end{remark}

\begin{remark}[Reversible and irreversible sectors]
In the Donsker-Varadhan bridge, $F_0$ is the visible contribution of the time-symmetric backbone $H_0$ and $W$ is the visible increment carried by the quadratic skew correction $\Delta$. The split therefore separates the reversible backbone from the irreversible correction at the level of the intrinsic visible object. It is this separation, not the fixed-dimension Hessian itself, that survives quotient observation.
\end{remark}

\subsection{Same-dimension closure criterion and its failure boundary}

The hidden-load cone is the class closed by quotient elimination. The same observed-dimension Donsker-Varadhan skew-square class is smaller, and its failure boundary is sharp.

\begin{proposition}[Paired-spectrum criterion]\label{prop:paired-spectrum}
Let $H_{0,\mathrm{vis}}\succ0$ and let $\Delta_{\mathrm{vis}}\succeq0$ be a visible correction. Set
\begin{equation}\label{eq:whitened-visible}
W:=H_{0,\mathrm{vis}}^{-1/2}\,\Delta_{\mathrm{vis}}\,H_{0,\mathrm{vis}}^{-1/2}.
\end{equation}
There exists a visible skew matrix $J_{\mathrm{vis}}$ of the same observed dimension such that
\[
\Delta_{\mathrm{vis}}=\tfrac14 J_{\mathrm{vis}}H_{0,\mathrm{vis}}^{-1}J_{\mathrm{vis}}^{\mathsf T}
\]
if and only if the positive spectrum of $W$ is pairwise degenerate. In particular, odd visible correction rank is impossible for same-dimension skew-square closure.
\end{proposition}

This is the same real skew-symmetric pairing obstruction that underlies Williamson-type normal forms and the canonical even-multiplicity structure of skew squares \cite{Williamson1936,HornJohnson1991}.

\begin{proof}
If such $J_{\mathrm{vis}}$ exists, define
\[
K:=H_{0,\mathrm{vis}}^{-1/2}J_{\mathrm{vis}}H_{0,\mathrm{vis}}^{-1/2}.
\]
Then $K$ is real skew-symmetric and
\[
W=\tfrac14 K K^{\mathsf T}.
\]
The nonzero singular values of a real skew matrix come in equal pairs, hence the positive eigenvalues of $W$ do as well.

Conversely, if the positive spectrum of $W$ is pairwise degenerate, choose an orthogonal basis in which
\[
W=Q\,\diag(\mu_1,\mu_1,\mu_2,\mu_2,\dots,0,\dots,0)\,Q^{\mathsf T},
\qquad \mu_i>0.
\]
Let
\[
K := 2Q\,\diag\!\bigl(\sqrt{\mu_1}J_2,\sqrt{\mu_2}J_2,\dots,0\bigr)Q^{\mathsf T},
\]
where $J_2=\begin{psmallmatrix}0&1\\-1&0\end{psmallmatrix}$. Then $K$ is real skew-symmetric and $W=\tfrac14 K K^{\mathsf T}$. Setting
\[
J_{\mathrm{vis}}:=H_{0,\mathrm{vis}}^{1/2} K H_{0,\mathrm{vis}}^{1/2}
\]
gives the required representation.
\end{proof}

\begin{corollary}[Paired-spectrum defect as exact distance]\label{cor:paired-defect}
Let $W$ be as in \eqref{eq:whitened-visible}, with ordered eigenvalues
\[
\lambda_1\ge \lambda_2\ge \cdots \ge \lambda_n\ge 0.
\]
Write $m:=\lfloor n/2\rfloor$ and define
\[
\bar\lambda_i:=\frac{\lambda_{2i-1}+\lambda_{2i}}{2},
\qquad i=1,\dots,m.
\]
Let
\begin{equation}\label{eq:pair-projection}
P_{\mathrm{pair}}(W)
:=
Q\,\diag(\bar\lambda_1,\bar\lambda_1,\dots,\bar\lambda_m,\bar\lambda_m,\,0\text{ if }n\text{ is odd})\,Q^{\mathsf T},
\end{equation}
where $W=Q\diag(\lambda_1,\dots,\lambda_n)Q^{\mathsf T}$ is an orthogonal spectral decomposition. Then $P_{\mathrm{pair}}(W)$ belongs to the same observed-dimension skew-square slice from Proposition~\ref{prop:paired-spectrum}, and it is the unique Frobenius-nearest point in that slice. Equivalently,
\begin{equation}\label{eq:paired-defect-distance}
\inf\Bigl\{\norm{W-Z}_F: Z=\tfrac14 K K^{\mathsf T},\ K^{\mathsf T}=-K,\ K\in\mathbb R^{n\times n}\Bigr\}
=
\norm{W-P_{\mathrm{pair}}(W)}_F.
\end{equation}
Moreover,
\begin{equation}\label{eq:paired-defect-formula}
\norm{W-P_{\mathrm{pair}}(W)}_F^2
=
\frac12\sum_{i=1}^{m}(\lambda_{2i-1}-\lambda_{2i})^2
\,+\,
\begin{cases}
0, & n\text{ even},\\
\lambda_n^2, & n\text{ odd}.
\end{cases}
\end{equation}
In particular, the defect vanishes if and only if same observed-dimension skew-square closure holds.
\end{corollary}

\begin{proof}
By Proposition~\ref{prop:paired-spectrum}, the target class consists of matrices of the form $\tfrac14 K K^{\mathsf T}$ with $K^{\mathsf T}=-K$, equivalently positive semidefinite matrices whose positive eigenvalues occur in equal adjacent pairs, with one forced zero when $n$ is odd. Fix such a target matrix $Z$. For fixed eigenvalues of $Z$, the Frobenius distance
\[
\norm{W-Z}_F^2=\operatorname{tr}(W^2)+\operatorname{tr}(Z^2)-2\operatorname{tr}(WZ)
\]
is minimised when $W$ and $Z$ commute, by the symmetric form of von Neumann's trace inequality. Hence the nearest point may be taken in the eigenbasis of $W$, and the matrix problem reduces to Euclidean projection of $(\lambda_1,\dots,\lambda_n)$ onto the cone of vectors with equal adjacent pairs and final zero when $n$ is odd.

The orthogonal projection onto that linear constraint set is obtained by averaging each adjacent pair, which gives \eqref{eq:pair-projection}. Because the eigenvalues are ordered, these pair averages remain ordered and nonnegative, so the projected vector already lies in the admissible cone. This proves \eqref{eq:paired-defect-distance}. Expanding the squared norm pair by pair gives
\[
(\lambda_{2i-1}-\bar\lambda_i)^2+(\lambda_{2i}-\bar\lambda_i)^2
=\frac12(\lambda_{2i-1}-\lambda_{2i})^2,
\]
and when $n$ is odd the final coordinate is projected to zero, contributing $\lambda_n^2$. This yields \eqref{eq:paired-defect-formula}. The final claim is immediate from Proposition~\ref{prop:paired-spectrum}.\qedhere
\end{proof}

\begin{remark}[Robust empirical certificate]
Define the paired-spectrum defect by
\[
D_{\mathrm{pair}}(W):=\norm{W-P_{\mathrm{pair}}(W)}_F.
\]
Distance to a closed set is $1$-Lipschitz in Frobenius norm, so for any estimate $\widehat W$ one has
\[
\abs{D_{\mathrm{pair}}(\widehat W)-D_{\mathrm{pair}}(W)}\le \norm{\widehat W-W}_F.
\]
Hence any certified lower bound on $D_{\mathrm{pair}}(\widehat W)$ beyond estimation error is a robust witness that the visible correction does not lie in the same observed-dimension skew-square class. The defect is therefore the quantitative version of the binary obstruction in Proposition~\ref{prop:paired-spectrum}.
\end{remark}

\begin{remark}[What is and is not closed]
Theorems~\ref{thm:latent-contraction} and \ref{thm:hidden-load}, together with Proposition~\ref{prop:paired-spectrum}, show the correct closure picture. Exact quotienting preserves Gaussian precision geometry and the support-preserving attenuation class beneath a fixed tangent ceiling. It does not generically preserve the same observed-dimension Donsker-Varadhan skew-square class. The microscopic Donsker-Varadhan Hessian should therefore be regarded as the tangent ceiling and microscopic cone inside the larger quotient-visible class, not as the class closed by Schur elimination itself.
\end{remark}

\section{Hidden locality and graph-local structure}
\label{sec:05_hidden_locality}

Section~\ref{sec:04_hidden_mediation} showed that the quotient-visible object is produced by hidden mediation beneath a fixed tangent ceiling. The next question is spatial: if the upstairs precision is graph-local on the hidden sector, does the induced visible law respect that locality, or can Schur elimination generate arbitrary long-range coupling? The answer is no. Visible fill-in is organised by hidden connected components and is exponentially suppressed by hidden graph distance, with a universal rate determined by the hidden spectral window. The section closes by identifying corridor extremisers for this long-range transfer problem. Those extremisers belong to the abstract positive-cone locality theory developed here. Section~\ref{sec:09_corridor_regime} later studies a Markov corridor regime in the hidden-feedback ladder; it is inspired by the same geometry but is not the same statement.

\subsection{Componentwise hidden mediation}
\label{subsec:componentwise-hidden-mediation}

Let $V=V_{\mathrm{vis}}\oplus V_{\mathrm{hid}}$ and let
\begin{equation}\label{eq:local-block}
K_X=
\begin{pmatrix}
K_{VV} & K_{VH}\\
K_{HV} & K_{HH}
\end{pmatrix}
\in\SPD(V)
\end{equation}
be a graph-local precision operator. In this Gaussian graphical-model setting, precision sparsity encodes conditional independence and Schur elimination creates visible fill-in \cite{Dempster1972,Lauritzen1996,GeorgeLiu1981}. The point is that this fill-in is not arbitrary. It is first split by hidden connected component, and only then can one ask how large the resulting visible coupling can be.

\begin{theorem}[Hidden-component mediation]\label{thm:hidden-components}
The visible effective precision after hidden elimination is
\begin{equation}\label{eq:keff}
K_{\mathrm{eff}}=K_{VV}-K_{VH}K_{HH}^{-1}K_{HV}.
\end{equation}
If the hidden graph decomposes into connected components $H=\bigsqcup_{\alpha} H_\alpha$ respected by $K_{HH}$, then
\begin{equation}\label{eq:component-sum}
K_{VH}K_{HH}^{-1}K_{HV}
=
\sum_{\alpha} K_{V H_\alpha}K_{H_\alpha H_\alpha}^{-1}K_{H_\alpha V}.
\end{equation}
Consequently a visible pair can acquire a new coupling only if both visible coordinates connect to the same hidden component.
\end{theorem}

The decomposition in \eqref{eq:component-sum} is the Green-kernel or walk-sum mediation picture for hidden elimination \cite{DoyleSnell1984,MalioutovJohnsonWillsky2006}. It is the componentwise version of the hidden-mediation mechanism from Section~\ref{sec:04_hidden_mediation}: locality is not lost at the level of support, even though fill-in becomes dense inside the mediated visible block.

\begin{proof}
Equation~\eqref{eq:keff} is the Schur complement formula. If the hidden graph splits into connected components respected by $K_{HH}$, then after permuting indices one has
\[
K_{HH}=\diag(K_{H_1H_1},\dots,K_{H_mH_m}),
\]
so the inverse is block diagonal with the same decomposition. Substituting this into the Schur term yields \eqref{eq:component-sum}. The last claim is immediate from the support structure of the sum.
\end{proof}

\subsection{Intrinsic graph-local decay}
\label{subsec:intrinsic-graph-local-decay}

Componentwise mediation is only the first locality statement. The sharper question is quantitative: once two visible channels are forced to communicate through a distant hidden region, how strongly can they couple? The next theorem gives the universal answer. It extends classical decay-of-inverse results for sparse positive matrices to the hidden-mediation setting relevant for quotient observation \cite{DemkoMossSmith1984,BenziGolub1999,BenziRazouk2007}.

\begin{proposition}[Graph-distance decay via polynomial approximation]\label{thm:chebyshev-decay}
Let $D\succ0$ be a hidden precision matrix supported on a hidden graph $G_H$, and suppose
\[
\operatorname{spec}(D)\subset[\alpha,\beta],
\qquad 0<\alpha<\beta<\infty.
\]
Let $S,T$ be hidden vertex sets with graph distance at least $\delta$, and define
\begin{equation}\label{eq:best-poly-error}
E_\delta(\alpha,\beta):=
\inf_{\deg p\le \delta-1}
\max_{x\in[\alpha,\beta]}\abs{\frac1x-p(x)}.
\end{equation}
Then
\begin{equation}\label{eq:chebyshev-decay}
\|P_S D^{-1} P_T\|_{\mathrm{op}}\le E_\delta(\alpha,\beta).
\end{equation}
Moreover, classical Chebyshev approximation yields constants $C_{\mathrm{dec}}(\alpha,\beta)>0$ and
\[
\rho:=\frac{\sqrt{\kappa}-1}{\sqrt{\kappa}+1},
\qquad \kappa:=\beta/\alpha,
\]
for which
\begin{equation}\label{eq:chebyshev-decay-exp}
E_\delta(\alpha,\beta)\le C_{\mathrm{dec}}(\alpha,\beta)\,\rho^\delta.
\end{equation}
\end{proposition}

\begin{proof}
Let $p$ be any polynomial of degree at most $\delta-1$. Since $D^j$ propagates only along hidden walks of length at most $j$, one has
\[
P_S p(D) P_T = 0.
\]
Hence
\[
P_S D^{-1} P_T = P_S\bigl(D^{-1}-p(D)\bigr)P_T,
\]
so
\[
\|P_S D^{-1} P_T\|_{\mathrm{op}}
\le
\|D^{-1}-p(D)\|_{\mathrm{op}}
\le
\max_{x\in[\alpha,\beta]} \abs{\frac1x-p(x)}.
\]
Taking the infimum over all such $p$ gives \eqref{eq:chebyshev-decay}. The exponential estimate \eqref{eq:chebyshev-decay-exp} is the standard Chebyshev-approximation consequence for $1/x$ on $[\alpha,\beta]$ \cite{DemkoMossSmith1984,Rivlin1974,Trefethen2013}.
\end{proof}

\begin{corollary}[Visible coupling decay]\label{cor:visible-decay}
Let $b_u,b_v$ be visible-to-hidden coupling vectors supported on hidden sets at graph distance at least $\delta$. Then
\begin{equation}\label{eq:visible-decay}
\abs{b_u^T D^{-1} b_v}
\le
\norm{b_u}_2\,\norm{b_v}_2\,E_\delta(\alpha,\beta)
\le
\norm{b_u}_2\,\norm{b_v}_2\,C_{\mathrm{dec}}(\alpha,\beta)\,\rho^\delta.
\end{equation}
\end{corollary}

\begin{proof}
Apply Proposition~\ref{thm:chebyshev-decay} and Cauchy-Schwarz.
\end{proof}

\begin{remark}[Locality as a visible prediction]
Corollary~\ref{cor:visible-decay} is the headline locality statement of the programme-core theory. Once the hidden spectral window is controlled, long-range visible coupling is exponentially suppressed by hidden graph distance. This is a prediction about the visible object, not merely a structural statement about the upstairs operator.
\end{remark}

\subsection{Corridor extremisers and bundle structure}
\label{subsec:corridor-extremisers}

The Chebyshev bound gives a universal decay rate. The next theorem identifies the asymptotic graph shape that realises that rate. This is still an abstract statement about graph-local hidden precision geometry. The strong-driving corridor analysis in Section~\ref{sec:09_corridor_regime} concerns a specific hidden-feedback ladder and should not be conflated with the present extremiser theorem.

\begin{proposition}[Corridor realisation of the universal exponential rate]\label{thm:corridor}
Let $P_{\delta+1}$ be the path on $\delta+1$ hidden vertices and define the constant-coefficient corridor precision
\begin{equation}\label{eq:corridor-precision}
D_\delta^{\mathrm{corr}} := cI-dA_{P_{\delta+1}},
\qquad
c:=\frac{\alpha+\beta}{2},
\qquad
d:=\frac{\beta-\alpha}{4}.
\end{equation}
Then $\operatorname{spec}(D_\delta^{\mathrm{corr}})\subset[\alpha,\beta]$, the endpoint distance is $\delta$, and the endpoint transfer obeys
\begin{equation}\label{eq:corridor-rate}
\bigl(D_\delta^{\mathrm{corr}}\bigr)^{-1}_{1,\delta+1}
=
C_{\mathrm{corr}}(\alpha,\beta)\,\rho^\delta\,(1+o(1)),
\qquad \delta\to\infty,
\end{equation}
for an explicit positive constant $C_{\mathrm{corr}}(\alpha,\beta)$ and the same decay rate $\rho=(\sqrt\kappa-1)/(\sqrt\kappa+1)$ as in Proposition~\ref{thm:chebyshev-decay}. Consequently the one-dimensional corridor realises the universal exponential decay rate at the level of the exponent.
\end{proposition}

\begin{proof}
The spectral inclusion for \eqref{eq:corridor-precision} follows from the path adjacency spectrum. The inverse of a constant-coefficient Jacobi chain is explicit, and its endpoint entry has the stated asymptotic form \eqref{eq:corridor-rate}. Proposition~\ref{thm:chebyshev-decay} shows that no graph-local hidden precision with spectral window $[\alpha,\beta]$ can decay more slowly than exponentially at rate $\rho^\delta$. Hence the corridor family realises the universal exponent.
\end{proof}

\begin{remark}[Rate optimality versus extremiser questions]
Proposition~\ref{thm:corridor} is a rate-level statement. It identifies the exponential decay exponent and a corridor family that realises it. It does not prove a sharp finite-$\delta$ minimax constant, and it does not prove the stronger many-channel threshold-count extremiser statement. Those questions remain open in the abstract locality problem. The corridor regime in Section~\ref{sec:09_corridor_regime} should therefore be read as a separate dynamical investigation built on the same rate-level geometry, not as a proof of full finite-distance optimality.
\end{remark}

\subsection{Scope of the current locality theory}
\label{subsec:scope-locality-theory}

The locality theory proved here has two parts. First, hidden mediation splits by connected component. Second, once the hidden spectral window is fixed, off-mask visible coupling is exponentially suppressed with graph distance, and the corridor family realises the universal decay exponent asymptotically. What is not proved here is a sharp finite-distance minimax constant or a full classification of graph extremisers beyond the corridor bundle. Adjacency-loss diagnostics and fill-in summaries are deferred to Appendix~C.

\section{Bell and temporal frontiers as law-level evidence}
\label{sec:06_bell_frontier}

Sections~\ref{sec:04_hidden_mediation} and \ref{sec:05_hidden_locality} developed the internal geometry of $\Phi$: hidden mediation, transport, conservation, and graph-local structure. The next question is evidential. Does the quotient-visible viewpoint detect genuine law-level structure that simpler summaries erase? On the Bell square the answer is yes. The visible datum is not merely a correlator table. It is a family of visible pair laws, and the right compatibility question is whether those pair laws arise as quotient marginals of one common visible Gaussian law. On that square the resulting frontier has two independent coordinates: variance consistency and the arcsine CHSH obstruction. Correlator-only summaries retain only the second.

\subsection{Common visible gluing as a law-level compatibility problem}

The Bell square is the first nontrivial test of the programme-core viewpoint because it asks for compatibility of several visible quotient outputs at once. The question is not whether each pair law is individually admissible. It is whether the four pair laws can be realised simultaneously as quotient marginals of one common law upstairs on the visible variables. That is the kind of compatibility problem the quotient-visible formalism is designed to express.

\begin{definition}[Common Schur gluing of context families]\label{def:common-schur-gluing}
Let
\[
Q_{xy}\in\SPD(2),
\qquad x,y\in\{0,1\},
\]
be four visible pair precisions, and let
\[
P_{xy}:\mathbb R^4\to\mathbb R^2
\]
select the coordinate pair $(A_x,B_y)$ from $(A_0,A_1,B_0,B_1)$. We say that the family $(Q_{xy})$ is \emph{commonly Schur-glued} if there exists $Q_\ast\in\SPD(4)$ such that
\begin{equation}\label{eq:common-schur-gluing}
Q_{xy}=F_{P_{xy}}(Q_\ast)
\qquad\text{for all }x,y\in\{0,1\}.
\end{equation}
\end{definition}

Definition~\ref{def:common-schur-gluing} is the Bell-square instance of quotient compatibility. It asks whether four visible pair laws sit on one common point of the positive cone after quotienting by the four coordinate projections. The next theorem shows that, on the Bell square, this compatibility problem collapses to two law-level conditions.

\begin{theorem}[Bell-square collapse theorem]\label{thm:bell-collapse}
Let
\[
Q_{xy}=\begin{pmatrix} a_{xy} & b_{xy}\\ b_{xy} & d_{xy}\end{pmatrix}\in\SPD(2),
\qquad
\Sigma_{xy}:=Q_{xy}^{-1}=\begin{pmatrix} v^A_{xy} & c_{xy}\\ c_{xy} & v^B_{xy}\end{pmatrix},
\]
and define
\begin{equation}\label{eq:bell-rho-def}
\rho_{xy}:=\frac{c_{xy}}{\sqrt{v^A_{xy}v^B_{xy}}}=\frac{-b_{xy}}{\sqrt{a_{xy}d_{xy}}},
\qquad
\phi_{xy}:=\arcsin(\rho_{xy}).
\end{equation}
Then the following are equivalent.
\begin{enumerate}
\item There exists a centred Gaussian law on $(A_0,A_1,B_0,B_1)$ whose $(A_x,B_y)$ marginals are exactly the pair laws with covariances $\Sigma_{xy}$. In the positive definite case this is equivalent to common Schur gluing in the sense of Definition~\ref{def:common-schur-gluing}.
\item The family satisfies remote-setting variance consistency,
\begin{equation}\label{eq:variance-consistency-body}
v^A_{x0}=v^A_{x1}\quad (x=0,1),
\qquad
v^B_{0y}=v^B_{1y}\quad (y=0,1),
\end{equation}
and the four arcsine CHSH inequalities
\begin{equation}\label{eq:arcsine-chsh-body-a}
|\phi_{00}+\phi_{01}+\phi_{10}-\phi_{11}|\le \pi,
\end{equation}
\begin{equation}\label{eq:arcsine-chsh-body-b}
|\phi_{00}+\phi_{01}-\phi_{10}+\phi_{11}|\le \pi,
\end{equation}
\begin{equation}\label{eq:arcsine-chsh-body-c}
|\phi_{00}-\phi_{01}+\phi_{10}+\phi_{11}|\le \pi,
\end{equation}
\begin{equation}\label{eq:arcsine-chsh-body-d}
|-\phi_{00}+\phi_{01}+\phi_{10}+\phi_{11}|\le \pi.
\end{equation}
Equivalently, with
\begin{equation}\label{eq:bell-lifted-correlators}
E_{xy}:=\frac{2}{\pi}\phi_{xy}=\frac{2}{\pi}\arcsin\!\Bigl(\frac{-b_{xy}}{\sqrt{a_{xy}d_{xy}}}\Bigr),
\end{equation}
common gluing is equivalent to variance consistency together with Bell-locality of the lifted correlator table $E=(E_{xy})$.
\end{enumerate}
\end{theorem}

\begin{proof}
Appendix~A isolates the Gaussian gluing criterion and then specialises the Bell square using the standard arcsine lift. The new point for the present paper is structural: the quotient-visible compatibility problem is not exhausted by correlators. On the Bell square the full visible law contributes the diagonal variance-gluing constraints \eqref{eq:variance-consistency-body} and the lifted CHSH constraints \eqref{eq:arcsine-chsh-body-a}-\eqref{eq:arcsine-chsh-body-d}.
\end{proof}

The theorem should be read as support for the programme-core theory, not as an isolated Bell observation. Sections~\ref{sec:quotient-precision}-\ref{sec:04_hidden_mediation} say that visible objects must be compared at law level, because quotient observation lives on the positive cone and composes by Schur transport. The Bell square tests that claim in the first genuinely constrained multi-context setting. The answer is rigid: the compatibility frontier is two-coordinate from the outset.

\subsection{Bell-square frontier and strict refinement beyond Bell-locality}

The Bell-local shadow detects only one part of the frontier from Theorem~\ref{thm:bell-collapse}. The variance sector is a distinct obstruction, and it is invisible to correlator-only summaries.

\begin{proposition}[Strict refinement over Bell-locality]\label{prop:bell-strict-refinement}
Bell-locality of the lifted correlator table $E$ is necessary but not sufficient for common visible Gaussian gluing of the full pair laws.
\end{proposition}

\begin{proof}
Consider
\[
Q_{00}=\begin{pmatrix}1&0\\0&1\end{pmatrix},
\qquad
Q_{01}=\begin{pmatrix}1/4&0\\0&1\end{pmatrix},
\qquad
Q_{10}=Q_{11}=\begin{pmatrix}1&0\\0&1\end{pmatrix}.
\]
Then every off-diagonal entry vanishes, so \eqref{eq:bell-lifted-correlators} gives $E_{xy}=0$ for all $x,y$, hence the lifted correlator table is Bell-local. But
\[
Q_{00}^{-1}=\begin{pmatrix}1&0\\0&1\end{pmatrix},
\qquad
Q_{01}^{-1}=\begin{pmatrix}4&0\\0&1\end{pmatrix},
\]
so the variance of $A_0$ depends on Bob's setting. The variance-consistency conditions \eqref{eq:variance-consistency-body} therefore fail, and no common visible Gaussian gluing exists.
\end{proof}

\begin{remark}[Two Bell-square obstructions]\label{rem:two-bell-obstructions}
For a Bell-square family, define
\begin{equation}\label{eq:epsilon-var}
\varepsilon_{\mathrm{var}}:=\max\!\bigl\{|v^A_{00}-v^A_{01}|,|v^A_{10}-v^A_{11}|,|v^B_{00}-v^B_{10}|,|v^B_{01}-v^B_{11}|\bigr\},
\end{equation}
and let $\varepsilon_{\mathrm{CHSH}}$ be the excess of the largest left-hand side of \eqref{eq:arcsine-chsh-body-a}-\eqref{eq:arcsine-chsh-body-d} above $\pi$, truncated below at zero. Then common gluing on the Bell square is equivalent to the simultaneous vanishing of both obstructions,
\[
\varepsilon_{\mathrm{var}}=0,
\qquad
\varepsilon_{\mathrm{CHSH}}=0.
\]
This packages the Bell-square frontier as a two-coordinate compatibility problem: one coordinate is invisible to Bell correlators, while the other is the usual arcsine CHSH obstruction.
\end{remark}

\begin{corollary}[Bell cube normal form]\label{cor:bell-cube-body}
Let $S$ be the $4\times 4$ matrix whose rows are the four CHSH sign vectors
\[
(1,1,1,-1),\qquad (1,1,-1,1),\qquad (1,-1,1,1),\qquad (-1,1,1,1),
\]
and define lifted Bell coordinates
\[
y:=\frac{1}{\pi}S\phi,
\qquad
\phi:=(\phi_{00},\phi_{01},\phi_{10},\phi_{11})^\top.
\]
Then the Bell-compatible region in arcsine coordinates is exactly the cube
\[
|y_i|\le 1,
\qquad i=1,\dots,4.
\]
Equivalently, the four CHSH inequalities are the coordinate inequalities of one global hypercubic normal form. In particular, Euclidean nearest-point repair in $y$ is coordinatewise clipping, and the Bell support function is
\[
h_{\mathcal B}(u)=\frac{\pi}{4}\,\|Su\|_1.
\]
Appendix~A records the derivation and the exact repair formulas.
\end{corollary}

\begin{corollary}[Single-readout CHSH locality]\label{cor:chsh-guardrail}
Suppose a Bell-context family $(Q_{xy})$ is commonly Schur-glued, and let
\[
Z=(Z_{A_0},Z_{A_1},Z_{B_0},Z_{B_1})\sim N(0,Q_\ast^{-1})
\]
for some $Q_\ast$ satisfying \eqref{eq:common-schur-gluing}. Let
\[
A_x=f_x(Z_{A_x}),\qquad B_y=g_y(Z_{B_y}),
\qquad |f_x|\le 1,\quad |g_y|\le 1.
\]
Then the correlators
\[
\widetilde E_{xy}:=\mathbb E[A_xB_y]
\]
obey the CHSH bound
\begin{equation}\label{eq:chsh-guardrail}
\bigl|\widetilde E_{00}+\widetilde E_{01}+\widetilde E_{10}-\widetilde E_{11}\bigr|\le 2,
\end{equation}
and likewise for the other three sign choices with one minus sign.
\end{corollary}

\begin{proof}
All four outputs are measurable functions of the single common random vector $Z$. Hence for every realisation one has
\[
\bigl|A_0(B_0+B_1)+A_1(B_0-B_1)\bigr|
\le |B_0+B_1|+|B_0-B_1|\le 2,
\]
because $|A_x|\le 1$ and $|B_y|\le 1$. Taking expectations gives \eqref{eq:chsh-guardrail}, and the other sign choices follow by relabelling.
\end{proof}

\subsection{Why the quotient-visible perspective sees more than correlators}

The Bell-facing conclusion is now sharper than a CHSH guardrail. Common visible Gaussian gluing on the Bell square is a law-level compatibility problem. Bell correlators detect only the arcsine shadow of that problem, whereas the quotient-visible formalism keeps the full pair law and therefore keeps the variance sector as well. Proposition~\ref{prop:bell-strict-refinement} proves that this extra sector is not cosmetic. Two families can have identical Bell-local correlator shadow and still differ on the existence of a common quotient-visible law.

This is the kind of evidence the paper needs at this stage. Theorem~\ref{thm:bell-collapse} does not merely add another reformulation of Bell-locality. It shows that the intrinsic visible variable from Section~\ref{sec:quotient-precision} resolves frontier structure that correlator compression destroys. The Bell square is therefore a genuine law-level stress test for $\Phi$. The passage from pair laws to correlators is a further readout layer, and it loses one of the two compatibility coordinates. The quotient-visible perspective sees more because it works before that readout collapse.

Corollary~\ref{cor:chsh-guardrail} then recovers the operational statement inside the fuller law-level picture. Once common gluing exists, every single bounded local readout sits under the classical CHSH bound. The usual guardrail is therefore recovered as a downstream shadow of common visible gluing, not as the primary object.

\subsection{Temporal frontier}

The same compatibility question can be asked on an observation graph of times rather than Bell settings. One specifies visible pair laws on the edges of a time-indexed graph and asks whether they admit one common Gaussian law whose edge marginals realise the prescribed family. Trees have only repeated-marginal consistency, whereas the first genuinely temporal clique already introduces a nontrivial law-level obstruction. Appendix~A records the first fully observed temporal case and its sign-shadow interpretation. We do not develop a full temporal theory in the main text. The point here is only structural: Bell and temporal frontiers are two graph-indexed instances of one quotient-visible gluing problem.
\section{Exact symmetry and the hidden response carrier}
\label{sec:07_exact_symmetry}

We now turn from the abstract quotient-visible theory to the controlled hidden-feedback ladder family. The visible variable is the activity $a\in\{0,1\}$, the hidden variable is the memory level $m\in\{0,\dots,M\}$, and the observer sees only the activity projection. The question is the first one worth asking once the abstract theory is in place: does the intrinsic visible object $\Phi$ ever close in a completely rigid way inside a genuine nonequilibrium model? At $h=0$ and $\lambda=0$ the answer is yes. A discrete symmetry forces an exact eigenmode, the visible line becomes an exact eigenline of the reduced conjugated operator, and the scalar visible algebra collapses to one parameter.

\subsection{The exact eigenmode at $\lambda = 0$}

Write the ladder generator in the split form
\begin{equation}\label{eq:ladder-generator-split}
Q = \alpha\bigl(Q_{\mathrm{vis}} + \eps Q_{\mathrm{mem}}\bigr),
\end{equation}
where $Q_{\mathrm{vis}}$ changes activity and $Q_{\mathrm{mem}}$ changes memory only. At $h=0$ and $\lambda=0$ the visible flip rates are symmetric, so the lumped visible chain has rates $\beta_0=\beta_1=\alpha$. Define the activity sign vector
\begin{equation}\label{eq:activity-sign-vector}
u(a,m) := \mathbf 1_{\{a=0\}} - \mathbf 1_{\{a=1\}}.
\end{equation}
This vector is constant along each memory fibre and records only the visible activity parity.

\begin{theorem}[Exact symmetry at the symmetric point]\label{thm:exact-symmetry-point}
Consider the hidden-feedback ladder at $h=0$ and $\lambda=0$, with scalar visible observation given by activity. Then the following statements hold exactly for every finite $M$ and every finite $\eps$.
\begin{enumerate}[label=(\roman*)]
\item The activity sign vector $u$ satisfies
\begin{equation}\label{eq:Qu-eigenmode}
Q u = -2\alpha u.
\end{equation}
\item In the visible-first reduced tangent basis, the visible line is an exact right eigenline of the reduced conjugated generator with eigenvalue $-2\alpha$.
\item The quotient-visible algebraic precision satisfies
\begin{equation}\label{eq:phi-equals-alpha}
\Phi = \alpha.
\end{equation}
\item Writing $F_0$ for the reversible visible backbone, $W:=\Phi-F_0$ for the irreversible correction, and $T$ for the scalar tangent load on the visible line, one has
\begin{equation}\label{eq:exact-symmetry-scalars}
v^{\mathsf T} H_0 v = \alpha,
\qquad
W\alpha = F_0 T.
\end{equation}
\end{enumerate}
\end{theorem}

\begin{proof}
Memory moves preserve activity and therefore preserve the sign vector, so $Q_{\mathrm{mem}}u=0$. At $h=0$ and $\lambda=0$, every visible flip changes the sign and occurs at rate $\alpha$, hence $Q_{\mathrm{vis}}u=-2u$. Multiplying by the prefactor in \eqref{eq:ladder-generator-split} gives \eqref{eq:Qu-eigenmode}. The reduction to the visible-first tangent basis preserves the visible line generated by the activity mode, so the induced reduced operator has the same visible eigenvalue. In scalar visible rank, the generator-level Schur identity from Section~\ref{subsec:scalar-schur} identifies $2\Phi$ with the visible Schur complement of $-L_{\mathrm{red}}$. Because the visible line is exact and carries eigenvalue $-2\alpha$, that Schur complement equals $2\alpha$, which gives \eqref{eq:phi-equals-alpha}. The identities in \eqref{eq:exact-symmetry-scalars} are then the scalar Schur relations for the decomposition of the reduced conjugated operator into reversible backbone plus skew correction.
\end{proof}

The theorem is the centre of Part III. Nothing asymptotic has happened yet. The visible clock is already pinned to a constant, and the remaining freedom is only in how that fixed total precision is split between reversible and irreversible sectors.

\begin{corollary}[Single-parameter reduction and frozen clock]\label{cor:single-parameter-reduction}
On the slice $h=0$ and $\lambda=0$, the scalar visible algebra is determined by the single parameter $T(\eps,M)$. More precisely,
\begin{equation}\label{eq:single-parameter-reduction}
F_0 = \frac{\alpha^2}{T+\alpha},
\qquad
W = \frac{\alpha T}{T+\alpha},
\qquad
F_0 + W = \alpha.
\end{equation}
Consequently $\Phi=\alpha$ is frozen across all $\eps$ and $M$, while $F_0$ decreases and $W$ increases monotonically as functions of $T$.
\end{corollary}

\begin{proof}
By Theorem~\ref{thm:exact-symmetry-point}, one has $\Phi=\alpha$ and $W\alpha=F_0T$. Together with the exact conservation split from Theorem~\ref{thm:conservation-split}, namely $F_0+W=\Phi$, this gives
\[
(\alpha-F_0)\alpha = F_0T.
\]
Solving for $F_0$ and then substituting into $W=\alpha-F_0$ yields \eqref{eq:single-parameter-reduction}. The monotonicity statements are immediate from
\[
\frac{dF_0}{dT} = -\frac{\alpha^2}{(T+\alpha)^2} < 0,
\qquad
\frac{dW}{dT} = \frac{\alpha^2}{(T+\alpha)^2} > 0.
\]
\end{proof}

\begin{remarkplain}[Trading segment and corridor endpoint]
Corollary~\ref{cor:single-parameter-reduction} puts every point on the symmetric $\lambda=0$ slice onto the line segment $F_0+W=\alpha$. The equilibrium endpoint is $(F_0,W)=(\alpha,0)$. Section~\ref{sec:09_corridor_regime} shows that the strong-corridor endpoint is $(\alpha/(M+1),\alpha M/(M+1))$. The total visible precision does not move. Only its internal partition moves.
\end{remarkplain}

\begin{remarkplain}[Why the eigenmode exists]
The mechanism behind Theorem~\ref{thm:exact-symmetry-point} is a $\mathbb Z_2$ activity-exchange symmetry. At $h=0$ and $\lambda=0$, exchanging the two activity classes leaves the generator invariant and makes the sign representation one-dimensional. The mode $u$ is therefore forced, not guessed. This is also the reason the phenomenon transfers to any binary-visible model whose generator has the same activity-exchange symmetry.
\end{remarkplain}

\subsection{The visible-first Schur reduction}

The previous theorem identifies the correct visible line. Once that line is fixed, the correct reduced basis is forced: first the true visible tangent vector, then an orthonormal basis of its hidden complement. The point is structural, not cosmetic. The first-correction analysis lives in this basis because only in this basis does the visible eigenline become a block constraint.

Let $L_{\mathrm{red}}(\lambda)$ denote the reduced conjugated tangent operator in a visible-first orthonormal basis, and write its block form as
\begin{equation}\label{eq:Lred-block-general}
L_{\mathrm{red}}(\lambda)=
\begin{pmatrix}
a(\lambda) & r(\lambda)^{\mathsf T}\\
c(\lambda) & \mathcal H(\lambda)
\end{pmatrix},
\end{equation}
where the first coordinate is the visible line and the remaining coordinates form the hidden complement.

\begin{proposition}[Visible-first Schur reduction at $\lambda=0$]\label{prop:visible-first-schur}
For arbitrary field $h$, provided $\lambda=0$, the visible line is the exact tangent line of the lumped two-state visible chain and is an exact right eigenline of the reduced conjugated generator. Equivalently,
\begin{equation}\label{eq:Lred-block-triangular}
L_{\mathrm{red}}(0)=
\begin{pmatrix}
a_0 & r_0^{\mathsf T}\\
0 & \mathcal H_0
\end{pmatrix},
\qquad
a_0 = -(\beta_0+\beta_1).
\end{equation}
In particular, the lower-left block vanishes exactly, so all hidden-to-visible correction passes through the single Schur channel determined by $r_0$, $\mathcal H_0$, and the derivative of the lower-left block.
\end{proposition}

\begin{proof}
Exact lumpability at $\lambda=0$ identifies the true visible tangent direction with the activity mode of the two-state visible chain. Writing the reduced conjugated operator in a basis whose first vector spans that line, the eigenline condition means exactly that the image of the first basis vector has no hidden-complement component. This is equivalent to $c(0)=0$, which yields the block upper-triangular form \eqref{eq:Lred-block-triangular}. The visible scalar block is the negative total visible escape rate of the lumped chain, namely $-(\beta_0+\beta_1)$.
\end{proof}

Proposition~\ref{prop:visible-first-schur} is the basis theorem for the correction layer. It explains why the naive odd-pair basis is not the right global coordinate system for the fast-memory expansion. The visible object responds through the true visible tangent line, not through an arbitrary parity chart. Section~\ref{sec:08_fast_memory} uses this exact block form as the starting point for the first-correction expansion.

\subsection{The canonical hidden response carrier}

Once the visible-first block form is exact, the first hidden derivative that can influence the visible Schur coefficient is forced. There is only one scalar contraction left.

\begin{theorem}[Canonical hidden response carrier]\label{thm:canonical-hidden-carrier}
Assume $\lambda=0$ and write the visible-first reduced operator as in \eqref{eq:Lred-block-general}, with $c(0)=0$ as in Proposition~\ref{prop:visible-first-schur}. Then the odd hidden contribution to the first derivative of the scalar visible Schur coefficient is carried by the single scalar
\begin{equation}\label{eq:Jhid-definition}
J_{\mathrm{hid}} := -\frac{\eps}{2}\, r_0^{\mathsf T}\mathcal H_0^{-1} c'(0).
\end{equation}
Equivalently, all competing hidden derivative channels are killed by the exact upper-triangular reduction at $\lambda=0$.

Moreover, once the visible line is fixed, $J_{\mathrm{hid}}$ is invariant under arbitrary orthogonal changes of hidden complement. Thus $J_{\mathrm{hid}}$ is canonical inside the visible-first Schur setting.
\end{theorem}

\begin{proof}
For scalar visible rank, the visible Schur coefficient of \eqref{eq:Lred-block-general} is
\begin{equation}\label{eq:scalar-schur-lambda}
S(\lambda) = a(\lambda) - r(\lambda)^{\mathsf T}\mathcal H(\lambda)^{-1}c(\lambda).
\end{equation}
Differentiating at $\lambda=0$ and using $c(0)=0$ gives
\begin{equation}\label{eq:scalar-schur-derivative}
S'(0) = a'(0) - r_0^{\mathsf T}\mathcal H_0^{-1}c'(0).
\end{equation}
Every term involving $r'(0)$, $\mathcal H'(0)$, or derivatives of the inverse multiplied by $c(0)$ vanishes. The only hidden contraction that survives is therefore $r_0^{\mathsf T}\mathcal H_0^{-1}c'(0)$. The normalization in \eqref{eq:Jhid-definition} is the one natural for the ladder scaling. For canonicality, let $O$ be an orthogonal change of basis on the hidden complement. Then
\[
r_0^{\mathsf T}\mapsto r_0^{\mathsf T}O^{\mathsf T},
\qquad
\mathcal H_0\mapsto O\mathcal H_0 O^{\mathsf T},
\qquad
c'(0)\mapsto O c'(0),
\]
so the scalar contraction is unchanged:
\[
r_0^{\mathsf T}O^{\mathsf T}(O\mathcal H_0 O^{\mathsf T})^{-1}Oc'(0)
= r_0^{\mathsf T}\mathcal H_0^{-1}c'(0).
\]
This proves both claims.
\end{proof}

\begin{remarkplain}[Meaning of canonical]
The word ``canonical'' is used here in a precise and limited sense. The scalar $J_{\mathrm{hid}}$ is not a basis-free universal observable on every reduced model. It is canonical once the visible line has been fixed and the reduction is written in the visible-first Schur basis forced by Proposition~\ref{prop:visible-first-schur}. Inside that setting there is exactly one hidden derivative contraction, and it is invariant under all residual orthogonal freedom.
\end{remarkplain}

\subsection{What this does and does not solve}

Section~\ref{sec:07_exact_symmetry} solves the exact algebraic centre of the correction analysis at $\lambda=0$. It identifies the correct reduction basis, proves that the visible precision freezes to $\Phi=\alpha$ on the symmetric slice, and collapses the scalar visible algebra to the single function $T(\eps,M)$. What it does not solve is equally important. It does not show that the eigenmode or the identity $\Phi=\alpha$ survive when $h\neq0$ or $\lambda\neq0$, it does not determine the fast-memory obstruction, and it does not explain the corridor regime. Those are the jobs of Sections~\ref{sec:08_fast_memory} and \ref{sec:09_corridor_regime}. The mechanistic Woodbury termination behind the multiplicative identity remains open even though the algebraic identity itself is proved. Finally, the canonical carrier identified here is shown in Section~\ref{sec:08_fast_memory} to be the unique survivor among the plausible first-correction reduction routes.

\section{Fast-memory response and thermodynamic constraints}
\label{sec:08_fast_memory}

The exact carrier from Section~\ref{sec:07_exact_symmetry} answers the symmetry question at the algebraic centre. The next question is asymptotic. When hidden memory equilibrates much faster than the visible activity flips, what visible law survives, and how does hidden nonequilibrium re-enter at the first nontrivial order? The answer has two parts. First, the visible law homogenises to a two-state chain, so the leading visible precision is completely explicit. Second, the first correction already remembers hidden geometry through one Schur channel, and that channel comes with both a source law and a sharp list of failures of the obvious simplifications.

\subsection{Fast-memory homogenisation}

Write $k_{a\to 1-a}(m)$ for the visible flip rate out of activity $a$ at memory level $m$, and let $\nu_a$ denote the stationary law of the pure memory dynamics on the fibre with activity fixed to $a$. The corresponding homogenised visible exit rates are
\begin{equation}\label{eq:beta-hom-def}
\beta_a:=\sum_{m=0}^M \nu_a(m)\,k_{a\to 1-a}(m),
\qquad a\in\{0,1\}.
\end{equation}
These are the rates seen after the memory sector has fully mixed before the next visible jump.

\begin{proposition}[Fast-memory homogenisation]\label{prop:fast-memory-homogenisation}
Fix finite $M$ and fixed parameters $(\mu,\lambda,h)$. As $\eps\to\infty$, the visible activity process homogenises to a two-state continuous-time Markov chain with rates $\beta_0$ and $\beta_1$ from \eqref{eq:beta-hom-def}. In particular, the leading visible dwell law is exponential in each activity sector.
\end{proposition}

\begin{proof}
This is the standard fast-memory timescale-separation mechanism, used here as a verified asymptotic input rather than as a standalone homogenisation proof. At timescale $\eps^{-1}$ the memory generator relaxes inside each activity fibre before a visible flip typically occurs. The visible jump hazard therefore averages against the fibre equilibrium $\nu_a$, which yields the effective exit rates $\beta_a$. Once those rates are fixed, the homogenised visible law is the two-state chain with generator
\[
Q_{\mathrm{hom}}=
\begin{pmatrix}
-\beta_0 & \beta_0\\
\beta_1 & -\beta_1
\end{pmatrix}.
\]
The exponential visible dwell law is then immediate.
\end{proof}

The first point of contact with the quotient-visible theory is that the homogenised limit is already scalar-visible, so the visible precision can be read off directly.

\begin{corollary}[Homogenised visible precision]\label{cor:phi-hom}
For the homogenised two-state chain from Proposition~\ref{prop:fast-memory-homogenisation},
\begin{equation}\label{eq:phi-hom}
\Phi_{\mathrm{hom}}=\frac{\beta_0+\beta_1}{2}.
\end{equation}
Moreover, along the ladder family,
\begin{equation}\label{eq:phi-fast-limit}
\Phi\to \Phi_{\mathrm{hom}},
\qquad
W\to 0,
\qquad \eps\to\infty.
\end{equation}
Thus the irreversible contribution is subleading in the fast-memory limit.
\end{corollary}

\begin{proof}
For a two-state visible chain with rates $a$ and $b$, the scalar visible Schur computation gives $\Phi=(a+b)/2$. Applying this to $Q_{\mathrm{hom}}$ proves \eqref{eq:phi-hom}. The convergence statement is the visible consequence of Proposition~\ref{prop:fast-memory-homogenisation}: once the visible law converges to the homogenised two-state chain, the scalar visible precision converges to its exact two-state value, while the hidden irreversible correction disappears at leading order.
\end{proof}

The basis theorem from Section~\ref{sec:07_exact_symmetry} is already active here. The correct asymptotic basis is not an arbitrary parity chart. It is the visible-first basis forced by Proposition~\ref{prop:visible-first-schur}, because that basis keeps the true visible tangent line fixed while the hidden block diverges.

\subsection{First-correction structure}

The homogenised limit is too coarse to explain how hidden nonequilibrium first re-enters the visible precision. That information appears one order later. The visible-first basis from Section~\ref{sec:07_exact_symmetry} is the right coordinate system because it separates scalar visible drift from hidden-mediated feedback before any asymptotic truncation is taken.

The first object that enters is the correct entry law for a visible dwell. If the process is about to enter activity $a$, it does not enter according to the within-activity equilibrium $\nu_a$. It enters according to jump-weighted mass coming from the opposite activity.

\begin{proposition}[Leading fast-memory entry law]\label{prop:fast-memory-entry-law}
Define
\begin{equation}\label{eq:eta-entry-law}
\eta_a(m):=\frac{\nu_{1-a}(m)\,k_{1-a\to a}(m)}{\beta_{1-a}},
\qquad a\in\{0,1\}.
\end{equation}
Then $\eta_a$ is the leading fast-memory entry law for dwells in activity $a$ in the fast-memory asymptotic regime. In particular, the exact entry law differs from $\eta_a$ by $O(\eps^{-1})$, and at fixed $(M,\mu,\lambda)$ the law $\eta_a$ is independent of the uniform field $h$.
\end{proposition}

\begin{proof}
A dwell in activity $a$ begins with a jump from the opposite activity, so the entering memory distribution is weighted by the source equilibrium $\nu_{1-a}$ multiplied by the jump hazard into $a$. Normalisation by the total inflow rate $\beta_a$ gives \eqref{eq:eta-entry-law}. The $O(\eps^{-1})$ correction is the usual first transient defect before complete fibre mixing. The field-independence is a normalisation cancellation: the uniform field multiplies both numerator and denominator in the same way and therefore drops out of the leading ratio.
\end{proof}

The next input is spectral. Let $T_{a,\eps}=\eps G_a-\diag(k_{a\to 1-a})$ denote the killed memory operator in activity sector $a$.

\begin{proposition}[Principal eigenvalue shift]\label{prop:fast-memory-eigenvalue}
For each activity sector,
\begin{equation}\label{eq:lambda-max-fast-memory}
\lambda_{\max}(T_{a,\eps})=-\beta_a+\frac{c_{\mathrm{eig},a}}{\eps}+o(\eps^{-1}),
\qquad \eps\to\infty.
\end{equation}
The first correction to the visible mean dwell is therefore not purely an effective-rate effect: the leading eigenvalue shift is only one contribution to the full first-order visible correction.
\end{proposition}

\begin{proof}
At leading order the operator $\eps G_a$ mixes to the fibre equilibrium, so the killed generator sees only the averaged hazard $\beta_a$. Standard first-order spectral perturbation around that averaged hazard gives \eqref{eq:lambda-max-fast-memory}. The final sentence records that the entry-law and transient-relaxation terms survive at the same order, so the eigenvalue shift alone does not close the visible correction.
\end{proof}

The visible-first Schur basis now exposes the full first coefficient. Write the reduced conjugated operator of $-L$ in that basis as
\begin{equation}\label{eq:fast-memory-block-expansion}
-L_{\mathrm{red},\eps}=
\begin{pmatrix}
2\Phi_{\mathrm{hom}}+a_1\eps^{-1}+o(\eps^{-1}) & r_0^{\mathsf T}+o(1)\\
c_0+o(1) & \eps\mathcal H_1+o(\eps)
\end{pmatrix}.
\end{equation}
The hidden block diverges like $\eps$, but the visible line stays fixed. This is precisely why the Schur feedback survives at order $\eps^{-1}$.

\begin{theorem}[First fast-memory correction in visible-first form]\label{thm:fast-memory-first-correction}
Assume the expansion \eqref{eq:fast-memory-block-expansion}. Then the scalar visible precision satisfies
\begin{equation}\label{eq:phi-first-correction}
\Phi_{\mathrm{hom}}-\Phi
=
\frac{1}{2\eps}\Bigl[-a_1+r_0^{\mathsf T}\mathcal H_1^{-1}c_0\Bigr]+o(\eps^{-1}).
\end{equation}
Thus the first visible correction splits into a scalar drift term and a hidden-mediated Schur term. At $\lambda=0$, the second term is exactly the fast-memory manifestation of the canonical hidden carrier from Theorem~\ref{thm:canonical-hidden-carrier}.
\end{theorem}

\begin{proof}
For scalar visible rank, Section~\ref{subsec:scalar-schur} identifies $2\Phi$ with the visible Schur complement of $-L_{\mathrm{red},\eps}$. Applying the scalar Schur formula to \eqref{eq:fast-memory-block-expansion} gives
\[
2\Phi
=
2\Phi_{\mathrm{hom}}+a_1\eps^{-1}-r_0^{\mathsf T}(\eps\mathcal H_1)^{-1}c_0+o(\eps^{-1}).
\]
Since $(\eps\mathcal H_1)^{-1}=\eps^{-1}\mathcal H_1^{-1}$, rearranging yields \eqref{eq:phi-first-correction}. The identification with the carrier from Section~\ref{sec:07_exact_symmetry} is exactly the statement that the only hidden-mediated channel surviving in visible-first coordinates is the canonical Schur contraction.
\end{proof}

Equation~\eqref{eq:phi-first-correction} is the technical centre of the section. It shows why the first correction cannot be read from homogenised rates alone. The visible coefficient remembers hidden resolvent geometry through the single Schur channel $r_0^{\mathsf T}\mathcal H_1^{-1}c_0$.

\subsection{Source law and obstruction}

We now specialise to the symmetric coupling slice $\lambda=0$, where the activity-swap involution is exact. Section~\ref{sec:07_exact_symmetry} isolated the canonical hidden carrier $J_{\mathrm{hid}}$. The present question is how that carrier varies with the field $h$ in the fast-memory regime.

\begin{theorem}[Parity and exact perturbation decomposition]\label{thm:field-parity-decomposition}
Assume $\lambda=0$ and fix the visible-first basis from Proposition~\ref{prop:visible-first-schur}. Write
\[
r(h)=r_0+\Delta r(h),
\qquad
c'_h=c'(0)+\Delta c(h),
\qquad
\mathcal H(h)=\mathcal H_0+\Delta\mathcal H(h),
\]
and let
\[
R(h):=\mathcal H(h)^{-1}-\mathcal H_0^{-1}.
\]
Then the canonical hidden carrier
\[
J_{\mathrm{hid}}(h):=-\frac{\eps}{2}\,r(h)^{\mathsf T}\mathcal H(h)^{-1}c'_h
\]
is an even function of $h$. The normalised datum
\begin{equation}\label{eq:Jstar-def}
J_*(h):=\frac{J_{\mathrm{hid}}(h)}{1+\cosh(h)}
\end{equation}
is also even. Moreover, the field defect
\[
\Delta J_{\mathrm{hid}}:=J_{\mathrm{hid}}(h)-J_{\mathrm{hid}}(0)
\]
admits the exact decomposition
\begin{equation}\label{eq:Jh-defect-decomp}
\Delta J_{\mathrm{hid}}=T_r+T_c+T_H+T_{\mathrm{mix}},
\end{equation}
where
\begin{align}
T_r&:=-\frac{\eps}{2}\,(\Delta r)^{\mathsf T}\mathcal H_0^{-1}c'(0),\label{eq:Tr-def}\\
T_c&:=-\frac{\eps}{2}\,r_0^{\mathsf T}\mathcal H_0^{-1}\Delta c,\label{eq:Tc-def}\\
T_H&:=-\frac{\eps}{2}\,r_0^{\mathsf T}R(h)c'(0),\label{eq:TH-def}
\end{align}
and
\begin{align}\label{eq:Tmix-def}
T_{\mathrm{mix}}:={}&-\frac{\eps}{2}\Bigl[(\Delta r)^{\mathsf T}\mathcal H_0^{-1}\Delta c
+(\Delta r)^{\mathsf T}R(h)c'(0)
+r_0^{\mathsf T}R(h)\Delta c
+(\Delta r)^{\mathsf T}R(h)\Delta c\Bigr].
\end{align}
Here $T_r$ is row-tail drift, $T_c$ is input-channel drift, $T_H$ is hidden-resolvent drift, and $T_{\mathrm{mix}}$ collects the mixed interaction terms.
\end{theorem}

\begin{proof}
The evenness is the activity-swap symmetry from Section~\ref{sec:07_exact_symmetry}: replacing $h$ by $-h$ exchanges the two activity classes while leaving the visible-first scalar contraction invariant, so $J_{\mathrm{hid}}(h)=J_{\mathrm{hid}}(-h)$. Since $1+\cosh(h)$ is itself even, the same holds for $J_*(h)$.

For the decomposition, expand
\[
r(h)^{\mathsf T}\mathcal H(h)^{-1}c'_h
=
\bigl(r_0+\Delta r\bigr)^{\mathsf T}\bigl(\mathcal H_0^{-1}+R(h)\bigr)\bigl(c'(0)+\Delta c\bigr)
\]
and subtract the $h=0$ contribution $r_0^{\mathsf T}\mathcal H_0^{-1}c'(0)$. Grouping the surviving terms by whether the perturbation enters through the row, the input channel, the hidden resolvent, or a mixed product yields \eqref{eq:Jh-defect-decomp} with \eqref{eq:Tr-def}-\eqref{eq:Tmix-def}.
\end{proof}

The point of Theorem~\ref{thm:field-parity-decomposition} is structural. There is a clean source law, but there is also a genuine obstruction. After normalisation by $1+\cosh(h)$, the carrier becomes asymptotically field-blind at leading fast-memory order, yet not exactly field-blind at finite $\eps$.

\begin{proposition}[Asymptotic source law and obstruction class]\label{prop:source-law-obstruction}
Assume $\lambda=0$. Then in the fast-memory regime,
\begin{equation}\label{eq:Jstar-field-blind}
J_*(h)=J_*(0)+O\!\left(\frac{\cosh(h)-1}{\eps}\right),
\end{equation}
and hence
\begin{equation}\label{eq:Jstar-obstruction-h2}
J_*(h)-J_*(0)=O\!\left(\frac{h^2}{\eps}\right)
\qquad (h\to 0).
\end{equation}
In the same regime the odd first-correction coefficients obey the asymptotic redistribution laws
\begin{equation}\label{eq:odd-redistribution-scalar}
-\frac{c_{\mathrm{scalar}}(h)}{\cosh(h)}\sim J_*(0),
\qquad
\frac{c_{\mathrm{schur}}(h)}{1+\cosh(h)}\sim J_*(0),
\qquad \eps\to\infty.
\end{equation}
Thus the source law is asymptotically clean, while the residual defect sits in an even obstruction class of order $O(h^2\eps^{-1})$.
\end{proposition}

\begin{proof}
Because $J_*$ is even, its small-field defect must begin at quadratic order in $h$. Fast-memory scaling pushes the coefficient of that even defect down by one power of $\eps^{-1}$, which yields \eqref{eq:Jstar-field-blind} and \eqref{eq:Jstar-obstruction-h2}. The redistribution laws follow from the visible-first Schur decomposition of the odd first-correction slope: the scalar part carries a $\cosh(h)$ factor, while the hidden Schur part carries a $(1+\cosh(h))$ factor. After those factors are divided out, both coefficients approach the same source datum $J_*(0)$.
\end{proof}

At larger depth the same mechanism persists in a parity-resolved form. The fast-memory correction separates into a dominant odd part and a smaller negative even defect; the odd response is linear to first order in small $\lambda$; and, after parity decomposition, both parity sectors are extensive in depth. That is the sense in which the hidden carrier from Section~\ref{sec:07_exact_symmetry} really does organise the fast-memory layer. It does not solve the whole correction problem, but it isolates the unique dominant odd channel.

\subsection{What the first correction is not}

The first-correction formula is informative partly because of what it excludes. Every obvious scalar shortcut was tested against the ladder family, and each one fails for a structural reason: the visible coefficient remembers hidden resolvent geometry, not just a few coarse summaries.

First, the carrier itself is not recovered by replacing the exact stationary law with the naive product-limit stationary law. The product-limit picture gets the scaling intuition right, but it does not reproduce the exact hidden contraction at finite $\eps$.

Second, the coefficient
\begin{equation}\label{eq:Cphi-def}
C_\Phi:=\eps\bigl(\Phi_{\mathrm{hom}}-\Phi\bigr)
\end{equation}
is not determined by the homogenised rate pair $(\beta_0,\beta_1)$ alone. Nor is it captured by simple summaries of the entry laws $\eta_a$ together with those rates. The reason is visible in Theorem~\ref{thm:fast-memory-first-correction}: the term $r_0^{\mathsf T}\mathcal H_1^{-1}c_0$ depends on hidden resolvent data that no low-dimensional rate summary can see.

Third, the first-order visible dwell correction is not determined solely by the principal eigenvalue shift \eqref{eq:lambda-max-fast-memory}. Entry-law and transient-relaxation contributions survive at the same order. Likewise, in the large-$M$ fast-memory regime the odd correction density does not collapse to simple Poisson-energy scalars such as $\nu(u_1^2)$, $\nu(k u_1^2)$, $\eta(u_1)$, or $\eta(u_1^2)$. Those quantities do not encode the hidden resolvent geometry of the Schur term.

Fourth, the field defect from Proposition~\ref{prop:source-law-obstruction} does not admit a one-slot reduction. The exact decomposition \eqref{eq:Jh-defect-decomp} proves that row-tail drift alone, input-channel drift alone, and hidden-resolvent drift alone are each insufficient. The mixed channel is not optional. For the same reason, the negative even defect is not generated by continuing the odd carrier $J_{\mathrm{hid}}$ to second order in $\lambda$.

Finally, the finite-forcing odd redistribution law is not exact on the full parameter set. The asymptotic source law is real, but away from the fast-memory limit higher-order finite-forcing terms spoil exact factorisation. This is why the source law is a leading-order theorem rather than a full finite-parameter identity.

Taken together, these exclusions are not side remarks. They are the rigidity part of the correction theory. They show that the first visible correction cannot be reduced to homogenised rates, entry-law moments, product-limit stationarity, one-slot field drift, or a continuation of the odd carrier. The visible object is already genuinely higher-order at first subleading order.

\subsection{Thermodynamic constraints on the visible object}

The correction layer is not unconstrained. The same ladder family that reveals the source law also puts a thermodynamic ceiling on how large the visible object can become. The cleanest summary is one line.

\begin{corollary}[Thermodynamic ceiling]\label{cor:thermodynamic-ceiling}
Assume the verified half-bound $W\le \sigma/2$, where $\sigma$ is the total entropy production rate. Then
\begin{equation}\label{eq:thermo-ceiling}
\Phi\le F_0+\frac{\sigma}{2}.
\end{equation}
\end{corollary}

\begin{proof}
By the exact conservation split from Theorem~\ref{thm:conservation-split},
\[
\Phi=F_0+W.
\]
Combining this with the verified half-bound $W\le \sigma/2$ yields \eqref{eq:thermo-ceiling}.
\end{proof}

The ceiling inherits the status of the bound behind it. Within the ladder family the stronger verified statement is
\begin{equation}\label{eq:epr-fisher-sum}
W+T\le \sigma,
\end{equation}
with the individual half-bounds
\begin{equation}\label{eq:epr-fisher-half}
W\le \frac{\sigma}{2},
\qquad
T\le \frac{\sigma}{2}.
\end{equation}
These inequalities were verified across the full tested parameter grid, including nonzero $\lambda$. They are programme-core constraints, not a peculiarity of one asymptotic slice.

There is also an exact symmetry statement on the simplest slice. At $M=1$, $\lambda=0$, and $h=0$, all visible Fisher quantities are invariant under the timescale duality
\begin{equation}\label{eq:eps-duality}
\eps\longmapsto \eps^{-1}.
\end{equation}
In particular,
\begin{equation}\label{eq:eps-duality-components}
F_0(\eps)=F_0(\eps^{-1}),
\qquad
T(\eps)=T(\eps^{-1}),
\qquad
W(\eps)=W(\eps^{-1}),
\qquad
\Phi(\eps)=\Phi(\eps^{-1}).
\end{equation}
Since the unique fixed point of the duality is $\eps=1$, the irreversible correction is maximised there on that slice.

\begin{proposition}[Complete $\eps$-profile at $M=1$]\label{prop:eps-profile-M1}
At $M=1$, $\lambda=0$, and $h=0$, the quantities $F_0$, $T$, $W$, and $\Phi$ are symmetric in $\log\eps$ around $\eps=1$. Moreover,
\begin{equation}\label{eq:T-peak-eps1}
T(\eps)\le T(1)
\end{equation}
for all $\eps>0$, and the maximum of $W$ occurs at $\eps=1$.
\end{proposition}

\begin{proof}
The symmetry in $\log\eps$ is exactly the duality \eqref{eq:eps-duality-components}. The fixed point of the duality is $\eps=1$, so any unique extremum must occur there. The complete profile follows once this is combined with the verified tightness at both timescale extremes recorded below.
\end{proof}

The duality is exceptional. It fails cleanly once either of its structural hypotheses is removed: at $M\ge 2$ the odd sector is no longer two-dimensional, and at $\lambda\neq0$ the activity-exchange symmetry is broken. Those failure modes matter because they show that the duality is a rigid symmetry, not a numerical accident.

The asymptotic tightness statements are the two endpoint companions to the ceiling. In the slow-memory regime,
\begin{equation}\label{eq:slow-memory-tightness}
1-\frac{W}{T}=O(\eps),
\qquad \eps\to 0,
\end{equation}
and in the fast-memory regime,
\begin{equation}\label{eq:fast-memory-tightness}
1-\frac{W}{T}=O(\eps^{-1}),
\qquad \eps\to \infty.
\end{equation}
So the tangent envelope becomes tight in both timescale extremes, from opposite sides.

Finally, the correction layer has a clean near-equilibrium and small-driving structure. Near detailed balance, with plaquette affinity $A_\square\to 0$,
\begin{equation}\label{eq:near-equilibrium-scaling}
W\asymp (A_\square^2)^2,
\qquad
T\asymp (A_\square^2)^2,
\qquad
\sigma\asymp (A_\square^2)^2.
\end{equation}
Thus the ratios $W/\sigma$ and $T/\sigma$ remain finite in the linear-response regime. But $\Phi$ is not determined by $A_\square$ alone. The full Schur object depends on more than the scalar thermodynamic force. Likewise, at small driving with $\lambda=0$ one has
\begin{equation}\label{eq:small-mu-symmetry}
W=T+O(\mu^4),
\qquad \mu\to 0,
\end{equation}
so the tangent and irreversible corrections coincide to leading order before separating at quartic order.

\subsection{Rigor boundary}

Section~\ref{sec:08_fast_memory} contains four different epistemic layers, and they should be kept distinct. The homogenised two-state visible precision \eqref{eq:phi-hom} and the visible-first Schur expansion \eqref{eq:phi-first-correction} are exact once the asymptotic block expansion is granted. The parity statement for $J_{\mathrm{hid}}$ and the perturbation decomposition \eqref{eq:Jh-defect-decomp} are exact on the slice $\lambda=0$. The field-obstruction class, the EPR-Fisher bounds, the $M=1$ duality, the endpoint tightness laws, and the scaling claims are verified asymptotic or verified structural statements in the frozen ladder programme. The open pieces are the sharp obstruction constants, the operator that generates the even defect, and analytic proofs of the EPR-Fisher bound and the $M=1$ timescale duality.

Two proof routes are visible already. First, an analytic proof of $W+T\le \sigma$ would likely follow if the tangent quantity $T$ can be identified as a current-type precision in the sense required by the thermodynamic uncertainty relation. Second, the corridor decay transition discussed in Section~\ref{sec:09_corridor_regime} looks accessible by a direct spectral-gap analysis of the path-kernel restricted generator. Both are genuine research directions rather than present results.

\section{Corridor regime and the regime boundary}
\label{sec:09_corridor_regime}

Section~\ref{sec:08_fast_memory} analysed the fast-memory direction $\eps\to\infty$. We now turn a different knob. Fix $h=0$ and $\lambda=0$, keep the exact symmetry backbone from Section~\ref{sec:07_exact_symmetry}, and drive the bias $\mu$ large. A second asymptotic mechanism appears. Its natural carrier is not the visible-first Schur channel of the fast-memory layer, but an odd-sector path kernel with an endpoint Schur complement. The section has two jobs. First, it records the corridor asymptotics cleanly. Second, it proves that this corridor sector is genuinely different from the fast-memory response sector. That separation is a result, not a matter of taste. \par

\subsection{The strong-driving limit}
\label{subsec:strong-driving-limit}

At $h=0$ and $\lambda=0$, Section~\ref{sec:07_exact_symmetry} gave the exact scalar backbone $\Phi=\alpha$. The corridor regime asks how the decomposition of that fixed visible precision changes when $\mu\to\infty$. The answer is that the reversible part collapses to an endpoint law of order $O(M^{-1})$, while the tangent load grows linearly in depth.

\begin{proposition}[Corridor endpoint law and refined asymptotics]\label{thm:corridor-path-kernel}
Fix $h=0$, $\lambda=0$, and finite depth $M$. In the strong-driving limit $\mu\to\infty$, the odd sector is governed by a path-kernel reduction whose visible output is an endpoint Schur complement. The endpoint law is
\begin{equation}\label{eq:corridor-F0-limit}
F_0(M,\mu)\to \frac{\alpha}{M+1}.
\end{equation}
Consequently, using the exact scalar identity $\Phi=\alpha$ on the same slice,
\begin{equation}\label{eq:corridor-T-limit}
T(M,\mu)\to \alpha M,
\end{equation}
\begin{equation}\label{eq:corridor-WT-limit}
\frac{W}{T}(M,\mu)\to \frac{1}{M+1},
\end{equation}
and
\begin{equation}\label{eq:corridor-W-limit}
W(M,\mu)\to \frac{\alpha M}{M+1}.
\end{equation}
In the archived ladder computations, this convergence is consistent with the refined asymptotic form
\begin{equation}\label{eq:corridor-T-refined}
T(M,\mu)=\alpha M\tanh^2\!\left(\frac{\mu}{4}\right)+O\!\left(e^{-\gamma_M\mu}\right),
\qquad
\gamma_M=
\begin{cases}
1, & M=2,\\
\tfrac12, & M\ge 3.
\end{cases}
\end{equation}
\end{proposition}

\begin{proof}
The endpoint law \eqref{eq:corridor-F0-limit} is the corridor endpoint Schur computation on the odd-sector path kernel. Section~\ref{sec:07_exact_symmetry} gives the exact identity $\Phi=\alpha$ on the present slice, so
\[
W=\Phi-F_0=\alpha-F_0.
\]
Taking the limit in \eqref{eq:corridor-F0-limit} gives \eqref{eq:corridor-W-limit}. The ratio limit \eqref{eq:corridor-WT-limit} and the tangent limit \eqref{eq:corridor-T-limit} are then equivalent through \eqref{eq:corridor-W-limit}. The refined form \eqref{eq:corridor-T-refined} is the verified corridor spectral pattern from the archived ladder computations.
\end{proof}

\begin{corollary}[Trading relation and corridor endpoint]\label{cor:corridor-trading-endpoint}
On the slice $h=0$ and $\lambda=0$, the pair $(F_0,W)$ lies on the exact line
\begin{equation}\label{eq:corridor-trading-segment}
F_0+W=\alpha.
\end{equation}
Its equilibrium endpoint is $(\alpha,0)$ and its corridor endpoint is
\begin{equation}\label{eq:corridor-endpoint-pair}
\left(\frac{\alpha}{M+1},\frac{\alpha M}{M+1}\right).
\end{equation}
No monotonicity or no-overshoot statement in $\mu$ is claimed here.
\end{corollary}

\begin{proof}
The exact identity $F_0+W=\Phi=\alpha$ holds on the whole $h=0$, $\lambda=0$ slice by Section~\ref{sec:07_exact_symmetry}. The equilibrium endpoint is the small-driving limit $\mu\to 0$, and the corridor endpoint is given by Proposition~\ref{thm:corridor-path-kernel}.
\end{proof}

\begin{remark}[Different asymptotic direction, different carrier]
The fast-memory limit and the corridor limit are independent asymptotic directions in parameter space. One sends $\eps\to\infty$ at fixed $(M,\mu,h,\lambda)$ and keeps the visible-first Schur carrier. The other sends $\mu\to\infty$ on the exact-symmetry slice and replaces that carrier by an endpoint Schur law on the odd-sector path kernel. The asymptotic formulas may both involve the same scalar visible objects $(F_0,T,W,\Phi)$, but they are produced by different internal reductions.
\end{remark}

\subsection{Why this is not the same regime}
\label{subsec:not-the-same-regime}

The point of the corridor analysis is not merely that one gets a second formula. The point is that the second formula comes from a different operator sector. The next statement is the load-bearing boundary theorem for Part III.

\begin{proposition}[Verified regime separation in the current ladder family]\label{thm:regime-separation}
Within the current hidden-feedback ladder analysis, the fast-memory response sector and the large-$\mu$ corridor sector do not behave as one operator regime in disguise. The separation has three independent components.

First, the two sectors use different natural bases: the fast-memory expansion is carried by the true visible tangent vector completed to a visible-first orthonormal basis, whereas the corridor reduction is carried by the odd pair basis.

Second, the two sectors use different carrier geometries: the fast-memory odd response is organised by the canonical hidden Schur carrier
\[
J_{\mathrm{hid}}=-\frac{\eps}{2}r_0^{\mathsf T}\mathcal H_0^{-1}c'(0),
\]
while the corridor sector is organised by an endpoint path-kernel carrier on the odd chain.

Third, the two sectors have different scaling classes on strong-bias slices: the fast-memory odd hidden density remains $O(1)$ in $M$, whereas the corridor visible reversible law is $O(M^{-1})$.

In addition, for $M\ge 2$ at $\lambda=0$ and $\eps=1$, the tangent correction does not factorise as
\begin{equation}\label{eq:no-separable-cMg}
T(M,\mu)=c(M)g(\mu).
\end{equation}
So even inside the corridor sector, depth and driving remain genuinely entangled.
\end{proposition}

\begin{proof}
The basis statement is the contrast between Section~\ref{sec:08_fast_memory}, where the visible-first Schur basis is forced by the true visible tangent line, and Theorem~\ref{thm:corridor-path-kernel}, where the strong-driving reduction lives on the odd pair sector. The carrier statement is immediate from the constructions: Section~\ref{sec:08_fast_memory} is driven by the canonical hidden Schur scalar $J_{\mathrm{hid}}$, while the corridor reduction is an endpoint Schur complement of a path Green kernel. The scaling-class statement follows by comparing the verified large-$M$ odd-density class from the fast-memory package with \eqref{eq:corridor-F0-limit}, which is of order $O(M^{-1})$ at fixed large $\mu$. The non-factorisation claim is exactly the observed corridor negative boundary: the dependence of $T(M,\mu)/\tanh^2(\mu/4)$ on $\mu$ varies with $M$ for every tested $M\ge2$, so no decomposition of the form \eqref{eq:no-separable-cMg} can hold on that domain.
\end{proof}

\begin{remark}[What the theorem does not say]
Theorem~\ref{thm:regime-separation} does not deny that a higher containing architecture may exist. It says something narrower and stronger: the present fast-memory and corridor sectors are already provably distinct at the level of basis, carrier, and scaling class. Any future programme would therefore have to contain both sectors without collapsing these differences.
\end{remark}

\subsection{The regime boundary}
\label{subsec:regime-boundary}

The present theory is therefore genuinely multi-regime. Section~\ref{sec:07_exact_symmetry} gives the exact scalar backbone at $h=0$ and $\lambda=0$. Section~\ref{sec:08_fast_memory} gives the visible-first Schur carrier in the fast-memory direction. The present section gives the endpoint path-kernel carrier in the strong-driving direction. These are all real pieces of the same correction programme, but they are not yet one unified operator sector.

\begin{proposition}[Current regime boundary]\label{prop:current-regime-boundary}
The current correction theory consists of at least two distinct asymptotic sectors beyond the exact-symmetry backbone: the fast-memory sector of Section~\ref{sec:08_fast_memory} and the corridor sector of Theorem~\ref{thm:corridor-path-kernel}. A higher operator architecture containing both sectors without collapsing their distinctions remains open.
\end{proposition}

\begin{proof}
The existence of the two sectors follows from Section~\ref{sec:08_fast_memory} and Theorem~\ref{thm:corridor-path-kernel}. Their non-collapse is exactly Theorem~\ref{thm:regime-separation}. The remaining statement is therefore the open containing-architecture problem.
\end{proof}

The strength of Proposition~\ref{prop:current-regime-boundary} is that it says exactly enough. It does not pretend that the present paper has a single master kernel. It proves the boundary instead. That is the honest state of the theory.

\subsection{Large-$M$ scaling}
\label{subsec:large-M-scaling}

The strong-driving corridor endpoint is only one part of the depth story. The fast-memory package also revealed a depth-normalised odd hidden density whose behaviour across $M$ is remarkably stable. That object is not the corridor carrier, but it is the main large-$M$ datum that survives the present correction analysis, so it belongs here at the regime boundary.

\begin{proposition}[Verified large-$M$ odd-density class]\label{prop:large-M-odd-density}
At $h=0$ and on tested $\mu$ slices, the normalised odd hidden density is consistent with the asymptotic class
\begin{equation}\label{eq:largeM-density-class}
\frac{J_*(M,\mu)}{M-1}=L(\mu)+\frac{A(\mu)}{M}+O\!\left(\frac{1}{M^2}\right).
\end{equation}
This is finite-limit evidence in a stable asymptotic class, not a proved limit theorem.
\end{proposition}

\begin{proof}
The statement summarises the verified large-depth fits: after the extensive factor $(M-1)$ is removed, the data are consistent with a stable linear-in-$1/M$ correction and a smaller residual. The point here is not to claim a proof of limit existence, but to record the stable asymptotic form supported by the current computations.
\end{proof}

\begin{proposition}[Crossover at $\mu=1$ is not a singular threshold]\label{prop:mu-one-crossover}
Across the tested large-$M$ window, the approach direction in \eqref{eq:largeM-density-class} changes across $\mu=1$: it is from below for $\mu<1$, from above for $\mu>1$, and nearly flat near $\mu=1$. At present there is no numerical evidence that $\mu=1$ is a singular threshold in the odd hidden density itself.
\end{proposition}

\begin{proof}
The sign of the leading $1/M$ correction changes across $\mu=1$, which forces the approach direction to switch. At the same time, the underlying scans remain smooth in a neighbourhood of $\mu=1$, so the available evidence supports a crossover centre rather than a singular carrier threshold.
\end{proof}

Two open problems remain attached to this large-$M$ sector. First, the actual limit function $L(\mu)$ is not yet identified analytically. Second, the special role of $\mu=1$ is only phenomenological. We know it marks the centre of the observed crossover in approach direction, but we do not yet know whether it carries a deeper operator meaning. Those are the natural questions left by the present data.


\section{Benchmark synthesis and evidence}
\label{sec:10_benchmark_synthesis}

Parts~I-III built the quotient-visible theory and then opened the hidden-feedback ladder family in enough detail to expose its multi-regime correction architecture. The remaining question is methodological as well as external: while algebraic closure was still being pinned down, did the theory stay tied to a real signal under deliberately varied tests? The four archived studies were the sequence of tests used to answer that question. They were not run after the final closure layer was in hand, and the paper should say so openly. Their role was to keep the developing theory sharply grounded while the algebraic picture was still being fixed. That is also why the main text does not pretend they already use the final closure method. Read at their true historical strength, they do exactly what Part~IV needs. One study supplies an exact transformed semi-Markov anchor, one supplies same-visible-law completion tests, one supplies a decoupling benchmark in which visible precision survives while observed entropy production vanishes, and one supplies the exact visible-law family that eventually made the completion picture and closure push possible.

\subsection{Why these four tests matter}

These four tests probe four different structural predictions, and they do so in roughly the order the theory itself was being stabilised. The Kapustin-Ghosal-Bisker transformed semi-Markov test asks whether exact transformed visible structure can exist in a finite partially observed system and whether the exact compression mechanism is genuinely structural rather than numerical \citep{KapustinGhosalBisker2024}. The matched-completion tests ask whether the total algebraic visible quantity survives when the hidden completion is changed without changing the visible law; in the archive this includes both abstract matched completions and a real-model sanity check built from the Liepelt-Lipowsky six-state kinesin chemomechanical network \citep{LiepeltLipowsky2007}. The Banerjee-Kolomeisky-Igoshin proofreading test asks whether visible precision can remain nonzero after coarse observation kills the observed entropy-production signal on a biologically motivated kinetic proofreading model \citep{BanerjeeKolomeiskyIgoshin2017}. The hidden-feedback ladder test family asks whether visible memory, hidden dissipation, and visible algebraic correction can be separated inside one exact visible-law family with genuine two-way hidden feedback. Together they test exact structure, invariance, decoupling, and frontier behaviour while also documenting the route by which the theory was kept empirically disciplined before final algebraic closure was achieved.


\subsection{Exact transformed semi-Markov anchor}

The Kapustin-Ghosal-Bisker transformed semi-Markov benchmark \citep{KapustinGhosalBisker2024} remains the cleanest exact transformed anchor in the archive. On the canonical four-state benchmark, the transformed observer has an exact sequence-level gain, the transformed waiting-time contribution vanishes exactly, and the exact compact observer is the four-symbol second-order parity lift recorded in the archived theorem spine. That exact compression mechanism is not merely approximate: it retains the transformed quantity at machine precision on the benchmark, while the threshold-2 observer remains near-exact but not exact. The same archive also carries the right negative contrast. In the flashing ratchet test there is no comparable compact exact lift, and the best compact retentions remain far below the exact four-state parity benchmark.

The algebraic layer extracted from the same test is also scientifically useful. On the benchmark four-state system it reports
\[
\begin{aligned}
F_0 &= 5.4104119541, & \Phi &= 5.4286720632,\\
W &= 0.0182601092, & T &= 4.7306689027,\\
\frac{W}{T} &= 0.003860.
\end{aligned}
\]
So the exact transformed anchor sits in an overwhelmingly backbone-dominated regime. This test therefore does two jobs at once. At the semi-Markov level it provides the exact transformed benchmark promised in Appendix~D. At the algebraic level it gives a small but nonzero visible correction on a benchmark whose transformed structure is known exactly.

\subsection{Same-visible-law completion tests}

These matched-completion tests probe the claim that the total scalar visible quantity survives same-visible-law completion changes even when the hidden architecture is radically altered. Its strongest archived result is the four-state versus fourteen-state same-law comparison:
\[
\Phi_{4\text{-state}}=5.4286720632,
\qquad
\Phi_{14\text{-state}}=5.4286720632,
\qquad
|\Delta\Phi|=2.7\times 10^{-15}.
\]
This is the exact kind of test Section~3 asks for. The visible law is held fixed while the hidden phase representation is enlarged, and the total visible algebraic precision survives to machine precision. The decomposition beneath that total does \emph{not} stay fixed. In the archived outputs the four-state system has
\[
F_0=5.4104119541,
\qquad W=0.0182601092,
\qquad T=4.7306689027,
\]
whereas the fourteen-state completion has
\[
F_0=4.1993218800,
\qquad W=1.2293501833,
\qquad T=5.1336332970.
\]
So this test demonstrates the right invariant and the right non-invariants in the same comparison.

The archive also contains a second independent same-law check at the renewal level. Duplicating one hidden state to produce an augmented five-state system preserves the visible semi-Markov law exactly and leaves the visible correction unchanged to machine precision,
\[
W_{\mathrm{orig}}=0.0230962835,
\qquad
W_{\mathrm{dup}}=0.0230962835,
\qquad
|\Delta W|=8.88\times 10^{-16}.
\]
That second check matters because it rules out the possibility that the fourteen-state match is a one-off feature of a particular phase-type construction.


The same archive also contains a useful real-model sanity check based on the six-state Liepelt-Lipowsky kinesin chemomechanical network, observed through a two-class D-versus-T readout \citep{LiepeltLipowsky2007}. For that test the archived values are
\[
W=7.3411775678,
\qquad
T=9.9969400245,
\qquad
W/T=0.734342.
\]
This is not a same-law completion test. It is a reminder that once the hidden topology really carries a cycle current visible to the chosen partition, the quotient-visible correction can be large.

\subsection{Banerjee-Kolomeisky-Igoshin proofreading test}

The Banerjee-Kolomeisky-Igoshin proofreading test probes a different prediction, namely that visible precision need not track observed entropy production after coarse observation on a biological error-correction network \citep{BanerjeeKolomeiskyIgoshin2017}. Under the default closure, all three coarse observers in the archived Banerjee benchmark retain exactly zero observed entropy production, but the quotient analysis reports nonzero visible precision for all four presets. On the scalar E-only observer the archived values are
\[
17.5036,\quad 16.2444,\quad 15.7350,\quad 387.9618,
\]
for the WT, HYP, ERR, and T7 presets respectively. On the correct-versus-wrong observer the corresponding traces are
\[
31.3115,\quad 22.1071,\quad 37.3374,\quad 811.0236.
\]
So in this test zero observed entropy production emphatically does not imply zero visible precision.

The test also tells us what kind of decoupling this is. The archived irreversibility fractions remain modest,
\[
W/T\in[0.070,0.258],
\]
across the four presets. The visible algebraic object is therefore still backbone-dominated. What survives the coarse observation is not a large visible irreversible sector but a finite local precision object supported by hidden structure that the coarse entropy-production observer cannot see.

This test archive also contains the right negative control. When the closure is changed in a way that changes the visible semi-Markov law, completion independence fails exactly as it should. The archive compares a five-state turnover closure to a seven-state product closure and finds visibly different $F_{\mathrm{vis}}$ values for the same coarse observer. That failure is a feature, not a bug. It shows the test is probing the correct invariance statement: same visible law, not arbitrary closure change.


\subsection{Hidden-feedback ladder test family}

The hidden-feedback ladder test family has a dual status in the paper. It is a benchmark family in its own right, and it is also the system through which the completion picture and later closure push became possible. It therefore has to be framed more carefully than the other three archives. The archived ladder runs predate the final algebraic-closure stage, but they already give an exact visible-law family of the right kind: under the activity-only observer the visible process is an exact alternating semi-Markov process with phase-type dwell laws. The promoted atlas then separates baseline Markovity, recorder-only hidden dissipation, modulator-only visible memory, two-way adaptive feedback, and equilibrium-feedback reversible controls. Historically, this is the family that kept the theory anchored while the algebraic closure mechanism was still being isolated.

The equilibrium-feedback line is the cleanest control in the archive. The points labelled \texttt{eq\_feedback\_weak} and \texttt{eq\_feedback\_strong} both have $A_\square=0$ and entropy production at numerical zero, but their visible dwell coefficients of variation are $1.094$ and $1.294$ respectively. So the archive contains memory-rich visible laws with effectively zero plaquette-level dissipation. At the algebraic layer both points satisfy
\[
W=0,
\qquad
F_0=\Phi,
\]
and therefore act as reversible memory controls for the correction theory.

The recorder and modulator points then separate hidden cost from visible memory. The archived \texttt{pure\_recorder} point has
\[
\begin{aligned}
\mathrm{EPR} &= 0.3618, & cv &= 1,\\
F_0 &= 0.86768, & \Phi &= 1.00000,\\
W &= 0.13232, & \frac{W}{T} &= 0.86768,
\end{aligned}
\]
so it carries substantial algebraic visible correction without visible non-Markovity. By contrast the archived \texttt{pure\_modulator} point has
\[
\begin{aligned}
\mathrm{EPR} &= 0.0210, & cv &= 1.0754,\\
F_0 &= 0.49433, & \Phi &= 0.49925,\\
W &= 0.00492, & \frac{W}{T} &= 0.54210,
\end{aligned}
\]
so it carries visible memory with only a very small algebraic correction. That single pair is already enough to show that visible memory, hidden cost, and visible correction are three different coordinates in this family.

The adaptive and depth-sensitive points push the family to the current boundary. The archived \texttt{adaptive\_strong} point has $cv\approx 1.280$ and $W/T\approx 0.216$; \texttt{adaptive\_slow} has smaller dissipation but remains visibly non-Markov; and the $M=8$ \texttt{depth\_sensitive\_hidden\_resolution} point changes the visible law only moderately while remaining noticeably nonlinear and numerically stiff. The frozen nonlinear report is already scientifically informative, but it explicitly classifies several promoted points as optimisation-sensitive rather than fully bankable closure constants. That is the correct boundary to state in the main text, especially because these runs were part of the route toward closure rather than post-closure victory laps. The exact backbone and the promoted algebraic layer are bankable now; the nonlinear layer is qualitatively supportive and quantitatively selective.


\subsection{Cross-test synthesis}

\begin{enumerate}[label=(\roman*)]
\item The Kapustin-Ghosal-Bisker transformed semi-Markov test shows that exact transformed visible structure can exist in a finite benchmark and can admit an exact compact observer, but the ratchet contrast shows that this is a structural mechanism, not a generic compression miracle.
\item The matched-completion tests show that at scalar visible rank the total visible algebraic precision can survive drastic same-law completion changes even while the backbone-correction split moves substantially underneath it.
\item The Banerjee-Kolomeisky-Igoshin proofreading test shows that visible precision and observed entropy production can decouple sharply under coarse observation, exactly in the direction predicted by the quotient-visible theory.
\item The hidden-feedback ladder test family shows that visible memory, hidden dissipation, and visible algebraic correction are genuinely distinct coordinates and can be separated inside one exact visible-law family.
\item Across all four archives, the honest reading is methodological as well as mathematical: these were the tests used to keep the developing theory tied to a real signal while algebraic closure was still being isolated. That is why they matter historically, and it is why the main text reads them at their archived strength rather than pretending they were post-closure reruns.
\item The current boundary is not absence of structure but controlled selectivity: the exact and algebraic layers are already strong, whereas the nonlinear branch is informative only where the archived optimisation diagnostics say it is stable enough to trust.
\end{enumerate}



\section{Synthetic falsification battery and blind benchmark}
\label{sec:11_falsification}

Section~\ref{sec:10_benchmark_synthesis} showed that the theory remains coherent across independent benchmark tests grounded in transformed semi-Markov inference, biological proofreading, and molecular-motor coarse observation \cite{KapustinGhosalBisker2024,BanerjeeKolomeiskyIgoshin2017,LiepeltLipowsky2007}. This section applies a stricter standard. The point is not further illustration, but deliberate adversarial pressure on the sharpest structural claims. We therefore separate two forms of stress. First, a synthetic falsification battery targets the exposed theorem-level predictions directly. Second, a blind digital benchmark asks whether locked predictions made from restricted visible data survive reveal of the hidden world.

\subsection{Falsification principles and synthetic battery}

The synthetic battery was designed to attack the claims that would most clearly fail if the theory were tracking numerical artefacts rather than genuine structure: quartic bridge scaling, support preservation beneath quotient observation, the paired-spectrum obstruction to fixed-dimension closure, graph-local decay, and the corridor-versus-bundle extremiser picture. The tests were not tuning exercises. Each family was built to try to break a definite prediction.

\begin{table}[H]
\centering
\begin{tabularx}{\textwidth}{>{\raggedright\arraybackslash}p{3.1cm} >{\raggedright\arraybackslash}p{4.0cm} Y Y}
\toprule
Diagnostic & Tested regime & Quantitative outcome & Interpretation\\
\midrule
Quartic bridge slope & $H_0$-adapted bridge sweeps with visible and hidden dimensions $2$ to $4$ & Grand mean slope $3.9957$; nine dimension-pair means from $3.9936$ to $3.9977$ & Consistent with the quartic hidden-mediation bridge and stable across the tested dimension range.\\
Worst bridge-gap eigenvalue & Same adapted bridge sweeps & Worst reported bridge-gap eigenvalue $-1.17\times 10^{-15}$ & Numerically nonnegative up to floating-point tolerance, as required by the sign-definite quartic bridge law.\\
Support preservation & Same adapted bridge sweeps & Rank-preservation rate $100\%$; maximum reported principal-angle sine $5.58\times 10^{-8}$ & Finite hidden sectors preserve tangent support to numerical precision, matching the support-rigidity theorem.\\
Same-dimension closure & Generic quotient Gaussian ensembles with visible dimensions $3$ to $8$ & Failure fraction $1.0$ throughout; pairable fraction $0.0$ throughout & Generic same-dimension skew-square closure fails completely in the tested ensemble, matching the paired-spectrum obstruction.\\
Hidden-local decay bound & Random local hidden graphs of size $24$, $36$, and $48$ & Bound-compliance rate $100\%$ across all sampled size-distance cells & Visible long-range transfer behaves as a controlled hidden-local tail, consistent with the graph-local decay theory.\\
Topology and bundle tests & Paths, cycles, ladders, trees, and disjoint corridor bundles & Corridor geometry realises the sharp decay class; disjoint bundles give flat singular spectra; under the tested finite-distance normalisation the best one-channel performer was the cycle & Supports the rate-level corridor picture and the flat-spectrum many-channel extremiser law while refining the finite-distance prefactor statement.\\
\bottomrule
\end{tabularx}
\caption{Synthetic falsification highlights. Each row is a deliberate attempt to break a sharp structural prediction.}
\label{tab:falsification_battery}
\end{table}

The outcomes are clear. Across the adapted bridge sweeps the quartic slope stays pinned near $4$, the bridge-gap eigenvalue remains at floating-point scale, and support preservation holds to numerical precision. Across visible dimensions $3$ to $8$, generic same-dimension skew-square closure fails in every sampled case. Across the hidden-local graph families, every tested transfer entry obeys the decay bound. The one substantive refinement forced by the synthetic battery is finite-distance rather than structural: the one-channel topology statement should be made at the level of universal exponential rate, not universal finite-distance constant, because looped topologies can improve prefactors under the tested spectral-window normalisation.

\subsection{Blind digital benchmark}

The blind benchmark was designed as a lock-then-reveal protocol. The hidden world was fixed first, only a restricted visible sector was exposed, predictions were committed before reveal, and only then was the concealed world released for scoring. The point of this protocol was not to enlarge the theorem domain of the paper. It was to check whether the existing finite tangent law, quotient-visible attenuation geometry, and Bell-frontier compatibility theory can function as an instrument when the hidden world is genuinely withheld.

Each digital world coupled three exact layers. First, a finite-state nonequilibrium backbone fixed the reduced detailed-balance geometry $H_0$, the skew object $\widehat J$, and the finite Donsker-Varadhan signal $\Delta_{DV}=\tfrac14\widehat J H_0^{-1}\widehat J^\top$ on the observer space. Second, an adapted visible sector $U$ fixed the tangent ceiling $T=P_U\Delta_{DV}P_U|_U$, and the actual quotient-visible correction was generated inside the proved attenuation class
\[
X=T^{1/2}(I+\Lambda)^{-1}T^{1/2},\qquad \Lambda\succeq 0.
\]
Third, Bell-context families were attached in three exact types: commonly glueable families, Bell-local but variance-inconsistent families, and variance-consistent families with arcsine-CHSH violation. Hidden truth and visible challenge artefacts were separated on disk, and the analysis stage read only the visible manifest before reveal.

On the exact-mode hero world, the relative errors were $2.61\times 10^{-2}$ for the retained-signal curve $K_d$ and $1.67\times 10^{-3}$ for the detectability curve $\Theta_d$. On the quotient side, the minimum hidden dimension and contraction-volume scores were recovered exactly in the recorded score files, with identity residuals at machine scale. The stronger test is the frozen 256-world exact ensemble. Across that ensemble the mean relative error was $1.461\times 10^{-1}$ for the finite $K$-curve and $2.541\times 10^{-2}$ for the detectability curve $\Theta_d$. Quotient scores again remained exact on the recorded ensemble, Bell verdict agreement was $1.0$, locality compatibility matched in $92.58\%$ of worlds, and motif agreement was $45.70\%$. We therefore treat locality only as a class-level diagnostic reader, not as an exact hidden-graph reconstruction claim.

Two cautions are essential. First, the sampled track is a robustness layer, not the primary claim. On the sampled hero world the finite module remained stable and quotient scores again remained exact, but the sampled Bell layer missed glueable families, so the sampled benchmark should be read as interface support rather than as the cleanest statement of the law-level frontier. Second, adaptive reveal is already informative but should still be described cautiously. On a 64-world exact panel, oracle theorem-target reveal improved the mean second-step raw revealed signal over a random baseline by $6.39\times 10^{-2}$ and the corresponding detectability score by $1.64\times 10^{-2}$. We therefore treat adaptive reveal as a diagnostic of what the observation ladder can target next, not yet as a deployed blind acquisition policy. Detailed benchmark diagnostics are collected in Appendix~\ref{app:blind-benchmark}.

\subsection{What survived and what did not}

The falsification pressure was productive because it did not leave the theory unchanged. What survived is the structural core. The quartic bridge law, support preservation, the paired-spectrum obstruction, and the hidden-local decay bound all survived direct adversarial testing. The exact-track blind benchmark also survived in the sense that the finite tangent law, quotient-visible attenuation geometry, and Bell-frontier verdicts remained operational under a real lock-then-reveal protocol.

What did not survive unchanged are the stronger informal readings that were never entitled to theorem status. The one-channel topology claim had to be weakened from a universal finite-distance statement to a universal rate statement. Motif reconstruction did not become an exact hidden-graph reader. The sampled Bell layer is not yet a clean benchmark instrument. And adaptive reveal remains a promising diagnostic policy rather than a finished acquisition procedure. That is the right outcome for this section. The theory survives where it claimed structural necessity, and it bends precisely where the data says only a weaker statement is justified.

\section{Scope, guardrails, and transfer boundaries}
\label{sec:12_scope}

This section states the status boundaries of the paper in one place. Its purpose is not rhetorical. It is to tell the reader exactly which statements are proved, which are asymptotic theorems, which are verified but not yet proved, and which questions remain open. The division below is the contract under which the rest of the paper should be read.

\subsection{What is proved: the abstract algebraic theory}

The abstract quotient-visible theory developed in Sections~\ref{sec:quotient-precision}-\ref{sec:06_bell_frontier} is proved at the level stated in the theorem bodies. For any quotient observation on the positive cone, the paper proves the variational characterisation of quotient precision, its homogeneity, monotonicity, and concavity, and its exact composition along towers of surjections. In that sense the central visible object $\Phi$ is not an ad hoc diagnostic but the intrinsic visible variable of quotient observation.

The same abstract theory also proves the internal algebraic structure of the visible object once a tangent ceiling has been fixed. Hidden mediation preserves support, every support-preserving visible law in the interior beneath the ceiling is parametrised by a positive hidden-load operator, sequential mediation transports that load by an exact Schur law, the determinant clock is the unique continuous orthogonally invariant additive interior clock, and the visible precision splits exactly into the reversible contribution and the irreversible correction through the conservation identity $\Phi=F_0+W$. These statements are exact structural results, not asymptotic approximations.

The locality and compatibility results proved in Sections~\ref{sec:05_hidden_locality} and \ref{sec:06_bell_frontier} belong to the same exact quotient-visible theory. Graph-local hidden precision yields componentwise mediation and exponential graph-distance suppression of long-range visible coupling, with corridor architectures realising the universal asymptotic rate class. On the Bell square, common visible Gaussian gluing is proved to be equivalent to remote-setting variance consistency together with Bell-locality of the arcsine lift, and this yields the single-readout CHSH guardrail as a downstream consequence of common gluing. These results hold at the level of quotient-visible law geometry itself; they do not depend on the hidden-feedback ladder analysis of Sections~\ref{sec:07_exact_symmetry}-\ref{sec:09_corridor_regime}.

\subsection{What is proved and verified: the nonequilibrium analysis}

The nonequilibrium analysis in Sections~\ref{sec:07_exact_symmetry}-\ref{sec:09_corridor_regime} is carried out for the hidden-feedback ladder family. Within that family, several statements are exact. At $h=0$ and $\lambda=0$, the activity sign mode is an exact generator eigenmode, the visible tangent line is an exact eigenline of the reduced conjugated generator, the scalar visible precision satisfies $\Phi=\alpha$ exactly, and the symmetric-point identities linking $\Phi$, $F_0$, $W$, and $T$ hold exactly. At $\lambda=0$ for arbitrary field, the true visible line remains exact, the visible-first reduced basis is exact, and the odd first-correction slope is carried by the canonical hidden scalar $J_{\mathrm{hid}}$, which is invariant under orthogonal changes of hidden complement.

The same ladder analysis also contains proved asymptotic statements. In the strong-bias corridor regime at $h=0$ and $\lambda=0$, the reversible visible law, the tangent correction, and the ratio $W/T$ obey the explicit corridor asymptotics stated in Section~\ref{sec:09_corridor_regime}. Those are genuine asymptotic theorems for that regime. They are not merely numerical trends, but they do not imply that the corridor carrier and the fast-memory carrier are one and the same object.

A larger part of the ladder analysis is verified rather than proved. This includes fast-memory homogenisation to a two-state visible chain, convergence of $\Phi$ to the homogenised value, decay of $W$ in the fast-memory tail, the first-correction structure beyond the leading homogenised rates, the field-normalised source law and its obstruction class, the large-$M$ normalised odd sector, the timescale-duality profile on the $M=1$, $\lambda=0$, $h=0$ slice, the EPR-Fisher bounds, endpoint tightness in the slow-memory and fast-memory limits, near-equilibrium scaling, and the small-driving symmetry of $W$ and $T$. Throughout the paper these statements are presented as verified asymptotic or verified structural facts, not as already-promoted theorems.

The benchmark studies of Sections~\ref{sec:10_benchmark_synthesis} and \ref{sec:11_falsification} do not enlarge the theorem boundary. Their role is methodological. They show that the developing algebraic picture was kept tied to a real signal while closure was being pinned down, and they apply adversarial pressure to the visible object and its diagnostics. They support transfer and robustness, but they are not used to convert verified statements into proved ones.

\subsection{What remains open}

Completion independence of $\Phi$ beyond the proved $M=1$ case remains open, and this matters because a full visible-law descent theorem would identify the exact minimal visible carrier of the algebraic quantity beyond the scalar benchmark. The multiplicative mechanism behind the $\lambda=0$ identities is identified but not yet proved analytically in full Schur-termination form. The field-law package does not yet have sharp operator constants beyond the present obstruction class $O((\cosh h-1)/\eps)$.

The large-$M$ odd sector has a stable asymptotic class in the tested range, but a proved limit theorem and an explicit identification of the limit function remain open. The status of the crossover near $\mu=1$ is also open at the operator level: the numerical evidence identifies a change in approach direction, but not yet a deeper structural interpretation. The operator generating the negative even defect has not yet been identified, and its relation, if any, to the odd hidden carrier is not known.

At the level of regime architecture, it remains open whether there exists a higher operator framework that contains both the fast-memory and corridor sectors without erasing their distinction. The EPR-Fisher bound and the $M=1$ timescale duality are currently verified rather than analytically proved. The full nonlinear Markov re-embedding problem after quotient observation also remains open. Finally, the transfer from the present discrete quotient-visible architecture to any genuine continuum or dynamical gravity theory is open; the paper proves structural parallels only, not a continuum field theory.

\subsection{Transfer question}

The abstract quotient-visible theory of Sections~\ref{sec:quotient-precision}-\ref{sec:06_bell_frontier} is independent of the hidden-feedback ladder. Its statements are formulated for quotient observation on positive precision operators and therefore transfer at that level without reference to any one benchmark family. The ladder enters only when the paper asks what the visible object looks like inside a nonequilibrium system and how the first correction organises itself.

The analysis of Sections~\ref{sec:07_exact_symmetry}-\ref{sec:09_corridor_regime} is ladder-specific in its present form. The exact symmetry results, the fast-memory source law, the obstruction package, the corridor asymptotics, and the regime-separation statements are established for that family. Sections~\ref{sec:10_benchmark_synthesis} and \ref{sec:11_falsification} then test partial transfer: the archived kinesin and proofreading studies show that the algebraic observables remain meaningful under controlled biochemical observers, the transformed semi-Markov benchmark shows that the visible programme survives a nontrivial observer change, and the falsification battery shows that the diagnostics are not tuned to one toy topology. What those tests do not yet prove is that the ladder-specific correction architecture transfers unchanged to every binary-visible nonequilibrium system.

The paper therefore draws a sharp line. The quotient-visible theory transfers as abstract geometry. The correction architecture is presently established in one benchmark family and partially pressure-tested beyond it. The transfer question is real, and the paper treats it as open where it has not been proved.

\subsection{What the present paper does not claim}

The paper does \textit{not} claim a full ontological unification of the fast-memory and corridor regimes.

It does not claim a continuum field theory, a dynamical gravity theory, or an Einstein-equation analogue.

It does not claim a nonlinear Markov re-embedding theorem after quotient observation.

It does not claim a general visible-law descent theorem beyond the proved $M=1$ scalar case.

\section{Discussion}
\label{sec:13_discussion}

\subsection{Why $\Phi$ is the intrinsic visible variable}

The central interpretive claim of the paper is now sharp. The intrinsic visible variable is quotient precision $\Phi$, not the Donsker-Varadhan Hessian taken at fixed visible dimension. The reason is structural, not stylistic. Section~\ref{sec:quotient-precision} proves that quotient observation composes exactly along towers, and Theorem~\ref{thm:composition-law} makes that closure explicit. Section~\ref{sec:03_dv_bridge} then shows that the Donsker-Varadhan bridge remains essential, but as a bridge variable: it is the route by which the microscopic skew-square structure descends to a visible algebraic object. What closes under observation is the quotient-visible precision itself.

That distinction is not cosmetic. The tangent closure statement in Corollary~\ref{cor:tangent-closure} identifies the Donsker-Varadhan skew-square sector as the tangent microscopic cone inside a larger quotient-visible attenuation class. Section~\ref{sec:04_hidden_mediation} explains the geometry of that larger class by hidden load, support rigidity, additive clock, and conservation. Once those structures are in place, the right reading is forced: the Donsker-Varadhan Hessian is the bridge variable that generates the correct visible object, while $\Phi$ is the visible invariant that survives composition, hidden elimination, and law-level comparison.

At scalar visible rank the conclusion is even stronger. Within that domain, ceiling, clock, precision, clock additivity, and conservation form a closed scalar system. The tangent ceiling fixes where visible correction can live, the determinant clock measures hidden contraction volume, and the conservation split from Section~\ref{sec:04_hidden_mediation} separates reversible and irreversible visible contributions without leaving the quotient-visible framework. In that precise sense, the scalar visible theory is complete: it is not merely a collection of formulas, but a closed algebraic system for the visible object. Conceptually this places the visible side of the theory closer to an information-geometric coarse description than to a fixed-dimension microscopic parametrisation \cite{Rao1945,AmariNagaoka2000,AyJostLeSchwachhofer2017}.

\subsection{What the DV bridge means after quotient reduction}

The bridge from \emph{Finite} is not demoted by the quotient theory. It is repositioned. Before quotient reduction, the bridge organises the microscopic correction as a skew-square law built from the reversible backbone. After quotient reduction, that same bridge is understood as the doorway to the right visible object. It tells us how the microscopic correction enters the visible sector, but it does not itself define the renormalisation-natural invariant. The visible algebraic object of Section~\ref{sec:03_dv_bridge} is therefore best read as the bridge's finite descendant, while $\Phi$ is the quotient-theoretic object that remains well defined after observation towers, hidden mediation, and law-level comparison have been taken seriously.

\subsection{Apparent nonlinearity as coordinate artefact}

Within the paper's proved domain, apparent nonlinearity repeatedly resolves into one of two things. Either it disappears after the correct coordinate choice, or it reappears as a finite-dimensional algebraic feature with explicit operational meaning. Section~\ref{sec:04_hidden_mediation} is the clearest example. Hidden-load composition looks nonlinear in the raw visible variable $X$, but becomes additive transport in the determinant clock and linear order on the hidden-load cone. The same phenomenon appears in Section~\ref{sec:06_bell_frontier}: Bell-square compatibility looks nonlinear in correlator language, but after the arcsine lift the compatible region becomes a hypercubic coordinate constraint with exact repair by clipping. In both cases the apparent nonlinearity is not false, but misplaced. It belongs to the choice of coordinates rather than to the underlying compatibility law.

This claim must be stated with its boundary. The paper does not assert that all nonlinearity in nonequilibrium observation theory is a coordinate artefact. It asserts a narrower and stronger statement: within the quotient-visible theory developed here, every nonlinear structure that has been brought under exact control either linearises in canonical coordinates or reduces to a finite-dimensional algebraic feature with explicit physical meaning. Hidden-load composition, Bell-square compatibility, and the scalar visible Schur laws are the proved examples. Beyond that domain, the paper makes no general ontological claim.

\subsection{The multi-regime architecture}

One of the paper's mature conclusions is that the nonequilibrium visible theory is genuinely multi-regime. The abstract quotient-visible theory of Sections~\ref{sec:quotient-precision}-\ref{sec:06_bell_frontier} is exact and regime-independent. The hidden-feedback analysis of Sections~\ref{sec:07_exact_symmetry}-\ref{sec:09_corridor_regime} is not. It splits into an exact symmetry backbone, a fast-memory response sector, and a strong-driving corridor sector. The paper does not hide that separation. It proves it.

That separation is a strength. Section~\ref{sec:08_fast_memory} shows that fast-memory response is organised by the visible-first Schur carrier, the canonical odd hidden response scalar, and a source law with explicit obstruction. Section~\ref{sec:09_corridor_regime} shows that the strong-driving corridor regime is controlled instead by path-kernel endpoint Schur laws and a distinct scaling architecture. The regime-separation theorem, Theorem~\ref{thm:regime-separation}, and the current boundary statement in Proposition~\ref{prop:current-regime-boundary} make the point sharply: these are not two parametrisations of one already unified operator sector. They are two real asymptotic sectors with different carriers, different dominant mechanisms, and a currently open containing architecture.

The hidden-feedback ladder therefore teaches two lessons at once. First, the visible object is rigid enough to support exact symmetry laws, asymptotic source laws, regime boundaries, and thermodynamic constraints. Second, those laws do not collapse into one prematurely simplified picture. The theory has become strong enough to prove the separation, and that is a more valuable outcome than a forced unity would have been.

\subsection{Geometric structure of quotient observation}

The quotient map is also geometric in a precise sense.

\begin{remark}[Affine-invariant geometric viewpoint]
Equip $\Sym_{++}(n)$ with the affine-invariant metric
\[
g_H(U,V)=\operatorname{tr}(H^{-1}UH^{-1}V).
\]
In the metric-normalised shadow coordinates of Section~\ref{sec:quotient-precision}, the differential of
\[
F_C(H)=(CH^{-1}C^{\mathsf T})^{-1}
\]
acts by the coisometric compression of Theorem~\ref{thm:shadow-transport}. This is the tangent-level isometry condition one would need for a Riemannian-submersion interpretation of quotient observation in the affine-invariant geometry of the positive cone.

That geometric reading is enough for the present paper. A full Riemannian-submersion theorem would additionally require an explicit horizontal distribution, surjectivity of the differential onto the visible tangent space, and a complete metric-preservation argument on horizontal directions. We do not supply that full proof here. What the paper uses is the tangent coisometry and the resulting structural analogy.

It is equally important to state what this does \emph{not} mean. The determinant clock from Section~\ref{sec:04_hidden_mediation} is not the affine-invariant geodesic distance once visible rank exceeds one. Its additivity comes from multiplicativity of determinant under hidden-load composition, not from metric geodesic length. The geometric parallels do not amount to an Einstein equation, a dynamical gravity theory, or a proved continuum limit. What the paper has established is kinematic geometric structure, not a gravitational dynamics.
\end{remark}

\appendix

\section{Bell and temporal frontier details}\label{app:bell}

This appendix provides the technical details underlying Theorem~\ref{thm:bell-collapse}, Proposition~\ref{prop:bell-strict-refinement}, Corollary~\ref{cor:bell-cube-body}, and the temporal-frontier discussion in Section~\ref{sec:06_bell_frontier}.

\subsection{Gaussian gluing on the Bell square}

We first isolate the Gaussian completion problem behind Theorem~\ref{thm:bell-collapse}. Write
\[
\Sigma_{xy}=\begin{pmatrix}\sigma_{A_x}^2 & c_{xy}\\ c_{xy} & \sigma_{B_y}^2\end{pmatrix},
\qquad
K_{xy}:=\frac{c_{xy}}{\sigma_{A_x}\sigma_{B_y}},
\]
where the one-site variances are assumed to satisfy the remote-setting consistency conditions \eqref{eq:variance-consistency-body}. Let $\mathcal E_{2,2}$ denote the bipartite elliptope, namely the set of all $2\times2$ matrices of the form
\[
L_{xy}=u_x\cdot v_y
\qquad
\text{with }u_x,v_y\text{ unit vectors in a real Hilbert space}.
\]

\begin{proposition}\label{prop:bell-elliptope}
Under the variance-consistency conditions \eqref{eq:variance-consistency-body}, the following are equivalent.
\begin{enumerate}
\item There exists a centred Gaussian law on $(A_0,A_1,B_0,B_1)$ whose $(A_x,B_y)$ marginals are exactly $\Sigma_{xy}$.
\item The normalised cross block $K=(K_{xy})$ belongs to $\mathcal E_{2,2}$.
\end{enumerate}
If the resulting full covariance is positive definite, this is equivalent to common Schur gluing of the visible pair precisions.
\end{proposition}

\begin{proof}
Assume first that a common centred Gaussian law exists. After standardising each coordinate by its one-site standard deviation, one obtains a correlation matrix on $(A_0,A_1,B_0,B_1)$ with diagonal entries equal to $1$ and cross block equal to $K$. Every correlation matrix is positive semidefinite and therefore a Gram matrix of unit vectors. Hence $K\in\mathcal E_{2,2}$.

Conversely, suppose $K_{xy}=u_x\cdot v_y$ for unit vectors $u_x,v_y$. Form the block Gram matrix
\[
R:=\begin{pmatrix}
(u_x\cdot u_{x'})_{x,x'} & (u_x\cdot v_y)_{x,y}\\
(v_y\cdot u_x)_{y,x} & (v_y\cdot v_{y'})_{y,y'}
\end{pmatrix}\succeq 0.
\]
Now scale $R$ by the diagonal matrix
\[
D:=\operatorname{diag}(\sigma_{A_0},\sigma_{A_1},\sigma_{B_0},\sigma_{B_1}).
\]
Then $\Sigma_\ast:=DRD\succeq 0$ is a covariance matrix whose $(A_x,B_y)$ marginal block is exactly $\Sigma_{xy}$. If $\Sigma_\ast\succ0$, inversion yields a visible precision $Q_\ast=\Sigma_\ast^{-1}$, and each pair precision is the corresponding Schur complement block. This is exactly common Schur gluing.
\end{proof}

\begin{proof}[Proof of Theorem~\ref{thm:bell-collapse}]
By Proposition~\ref{prop:bell-elliptope}, common Gaussian gluing is equivalent to variance consistency together with membership of the normalised cross block $K$ in $\mathcal E_{2,2}$. For pair precisions,
\[
K_{xy}=\frac{c_{xy}}{\sqrt{v^A_{xy}v^B_{xy}}}=\frac{-b_{xy}}{\sqrt{a_{xy}d_{xy}}}=\rho_{xy},
\]
which is \eqref{eq:bell-rho-def}.

Now define
\[
E_{xy}:=\frac{2}{\pi}\arcsin(K_{xy}).
\]
Standard Bell-square facts identify the frontier in three equivalent ways \cite{Fine1982}. Fine's theorem characterises Bell-locality of $E$ by the CHSH inequalities, the Gaussian sign-correlation formula shows that $E$ is the correlator table of the signs of a Gaussian family with correlation matrix having cross block $K$, and the same lifted correlator set is the relevant bipartite elliptope in the Bell square. Hence $K\in\mathcal E_{2,2}$ if and only if the lifted correlator table $E$ is Bell-local, equivalently if and only if the four arcsine CHSH inequalities \eqref{eq:arcsine-chsh-body-a} to \eqref{eq:arcsine-chsh-body-d} hold. Combining this with Proposition~\ref{prop:bell-elliptope} proves the theorem.
\end{proof}

\subsection{Strict refinement over Bell-locality}

Proposition~\ref{prop:bell-strict-refinement} already gives the minimal counterexample. The point is structural: Bell inequalities constrain only the binary-output correlator shadow, while common Gaussian gluing constrains the full pair law. The diagonal sector is therefore visible to the quotient frontier even when the Bell shadow is trivial.

A convenient quantitative summary is to separate the two obstructions. Define
\begin{equation}\label{eq:epsilon-var-app}
\varepsilon_{\mathrm{var}}:=\max\!\bigl\{|v^A_{00}-v^A_{01}|,|v^A_{10}-v^A_{11}|,|v^B_{00}-v^B_{10}|,|v^B_{01}-v^B_{11}|\bigr\},
\end{equation}
and let $\varepsilon_{\mathrm{CHSH}}$ be the excess of the largest left-hand side of \eqref{eq:arcsine-chsh-body-a} to \eqref{eq:arcsine-chsh-body-d} above $\pi$, truncated below at zero. Then common gluing on the Bell square is equivalent to the simultaneous vanishing of both obstructions,
\[
\varepsilon_{\mathrm{var}}=0,
\qquad
\varepsilon_{\mathrm{CHSH}}=0.
\]
This packages the Bell-square frontier as a two-coordinate compatibility problem: one coordinate is invisible to Bell correlators, while the other is the usual arcsine correlator obstruction.

\subsection{Low-rank mediated Gaussian sign shadows}

The Bell-local shadow can itself be realised by a very small mediated Gaussian model.

\begin{corollary}\label{cor:rank-two-sign-shadow}
Let $E=(E_{xy})$ be a Bell-local $2\times2$ correlator table, and set
\[
K_{xy}:=\sin\!\Bigl(\frac{\pi}{2}E_{xy}\Bigr).
\]
Then there exist unit vectors $u_0,u_1,v_0,v_1\in\mathbb R^2$ and a standard Gaussian $Z\sim N(0,I_2)$ such that
\[
U_x:=u_x\cdot Z,
\qquad
V_y:=v_y\cdot Z,
\qquad
\mathbb E[\operatorname{sgn}(U_x)\operatorname{sgn}(V_y)]=E_{xy}
\quad\text{for all }x,y.
\]
Thus every Bell-local Bell-square correlator table already sits inside a rank-two Gaussian sign-shadow.
\end{corollary}

\begin{proof}
By Theorem~\ref{thm:bell-collapse}, Bell-locality of $E$ is equivalent to $K\in\mathcal E_{2,2}$. Thus there exist unit vectors $u_x,v_y$ in some Euclidean space with $u_x\cdot v_y=K_{xy}$. Let $U:=\operatorname{span}\{u_0,u_1\}$, so $\dim U\le 2$, and let $p_y$ be the orthogonal projection of $v_y$ onto $U$. Then $u_x\cdot p_y=K_{xy}$ for all $x,y$. Choose $Z_U\sim N(0,I_U)$ and independent standard Gaussians $\xi_y$, and define
\[
U_x:=u_x\cdot Z_U,
\qquad
V_y:=p_y\cdot Z_U+\sqrt{1-\|p_y\|^2}\,\xi_y.
\]
These are centred Gaussian variables with unit variances and cross-correlations $\operatorname{Corr}(U_x,V_y)=K_{xy}$. The Gaussian sign-correlation formula gives
\[
\mathbb E[\operatorname{sgn}(U_x)\operatorname{sgn}(V_y)]
=\frac{2}{\pi}\arcsin(K_{xy})
=E_{xy},
\]
as required.
\end{proof}

\subsection{Bell cube normal form and exact repair}

The Bell square admits a global normal form in lifted coordinates. This globalises the one-face repair formulas and makes the Bell witness geometry completely explicit.

\begin{theorem}\label{thm:bell-cube-app}
Let $S$ be the $4\times 4$ matrix whose rows are the four CHSH sign vectors
\[
s_1=(1,1,1,-1),\qquad s_2=(1,1,-1,1),\qquad s_3=(1,-1,1,1),\qquad s_4=(-1,1,1,1).
\]
Then
\[
S^\top S=4I_4.
\]
For
\[
y:=\frac{1}{\pi}S\phi,
\qquad
\phi:=(\phi_{00},\phi_{01},\phi_{10},\phi_{11})^\top,
\]
the Bell-compatible region in arcsine coordinates becomes
\[
|y_i|\le 1,
\qquad i=1,\dots,4.
\]
Equivalently, the Bell square is globally the cube $[-1,1]^4$ in the lifted coordinates $y$. Moreover,
\[
\phi=\frac{\pi}{4}S^\top y,
\qquad
\|\delta\phi\|^2=\frac{\pi^2}{4}\,\|\delta y\|^2.
\]
Hence the Euclidean nearest-point map is coordinatewise clipping,
\[
y_i^\ast=\operatorname{clip}(y_i,-1,1),
\qquad
\phi^\ast=\frac{\pi}{4}S^\top y^\ast,
\]
and the squared Euclidean distance to the Bell frontier is
\[
\operatorname{dist}(\phi,\mathcal B)^2
=\frac{\pi^2}{4}\sum_{i=1}^4 (|y_i|-1)_+^2.
\]
If only one lifted CHSH coordinate violates the cube, with defect $e>0$, then the repair is
\[
\delta\phi^\ast=-\frac{e\pi}{4}s_i,
\qquad
\|\delta\phi^\ast\|_\infty=\frac{e\pi}{4}.
\]
\end{theorem}

\begin{proof}
The four inequalities \eqref{eq:arcsine-chsh-body-a} to \eqref{eq:arcsine-chsh-body-d} are precisely
\[
|s_i\cdot \phi|\le \pi,
\qquad i=1,\dots,4.
\]
Since the rows of $S$ are the vectors $s_i$, this is equivalent to $|y_i|\le 1$ for all $i$, so the compatible region is the cube $[-1,1]^4$ in the lifted coordinates. Direct computation gives $S^\top S=4I_4$, hence $S^{-1}=\frac14 S^\top$ and therefore
\[
\phi=\frac{\pi}{4}S^\top y.
\]
The Euclidean metric transforms accordingly,
\[
\|\delta\phi\|^2
=\frac{\pi^2}{16}\,\delta y^\top SS^\top \delta y
=\frac{\pi^2}{4}\,\|\delta y\|^2,
\]
because $SS^\top=4I_4$. Euclidean nearest-point repair in $y$ is therefore just coordinatewise clipping to the cube, and transforming back gives the stated projection formula and distance formula. If only one coordinate exceeds the cube by $e$, then clipping changes only that coordinate by $e$, so
\[
\delta\phi^\ast=\frac{\pi}{4}S^\top(-e e_i)=-\frac{e\pi}{4}s_i,
\]
which also gives $\|\delta\phi^\ast\|_\infty=e\pi/4$.
\end{proof}

\begin{corollary}\label{cor:bell-support-app}
For any covector $u\in\mathbb R^4$, the Bell support function in arcsine coordinates is
\[
h_{\mathcal B}(u):=\sup_{\phi\in\mathcal B}u\cdot\phi
=\frac{\pi}{4}\,\|Su\|_1.
\]
\end{corollary}

\begin{proof}
Using $\phi=(\pi/4)S^\top y$ and $|y_i|\le 1$, one has
\[
h_{\mathcal B}(u)=\sup_{|y_i|\le 1}\frac{\pi}{4}(Su)\cdot y
=\frac{\pi}{4}\sum_{i=1}^4 |(Su)_i|
=\frac{\pi}{4}\,\|Su\|_1.
\]
\end{proof}

\subsection{Observation graphs and the first temporal clique}

The Bell square is the first nonchordal compatibility graph treated in the main text. The same gluing question may be asked on any finite observation graph, with Bell and temporal contexts appearing as two graph-indexed instances of one Gaussian pair-law completion problem. The next proposition isolates the law-level object.

\begin{proposition}[Graph-indexed Gaussian gluing]\label{prop:graph-gluing}
Let $G=(V,E)$ be a finite graph. For each edge $\{i,j\}\in E$, let
\[
\Sigma_{ij}=\begin{pmatrix} v_i^{(ij)} & c_{ij}\\ c_{ij} & v_j^{(ij)}\end{pmatrix}\in\SPD(2),
\qquad
\rho_{ij}:=\frac{c_{ij}}{\sqrt{v_i^{(ij)}v_j^{(ij)}}}.
\]
Then there exists a centred Gaussian law on $(X_i)_{i\in V}$ whose $(i,j)$ marginal is $\Sigma_{ij}$ for every edge if and only if both of the following hold:
\begin{enumerate}[label=\textup{(\roman*)},leftmargin=2.3em]
\item vertex variance consistency, namely there exist numbers $v_i>0$ with
\[
v_i^{(ij)}=v_i
\qquad
\text{for every vertex $i$ and every incident edge $\{i,j\}$};
\]
\item correlation completion, namely there exists a correlation matrix $R\succeq0$ on $V$ such that
\[
R_{ii}=1,
\qquad
R_{ij}=\rho_{ij}
\quad
\text{for all }\{i,j\}\in E.
\]
\end{enumerate}
\end{proposition}

\begin{proof}
If a global centred Gaussian law exists, its one-site variances are independent of which edge marginal is inspected, so \textup{(i)} holds. Dividing the global covariance by the diagonal standard-deviation matrix gives a correlation matrix $R$ with the stated edge restrictions, so \textup{(ii)} holds.

Conversely, assume \textup{(i)} and \textup{(ii)}. Let $D=\operatorname{diag}(v_i)_{i\in V}$ and set $\Sigma:=D^{1/2}RD^{1/2}$. Then $\Sigma\succeq0$ and every edge marginal is
\[
\Sigma_{\{i,j\}}=\begin{pmatrix}
v_i & \sqrt{v_iv_j}\,R_{ij}\\
\sqrt{v_iv_j}\,R_{ij} & v_j
\end{pmatrix}
=\begin{pmatrix}
v_i^{(ij)} & c_{ij}\\ c_{ij} & v_j^{(ij)}
\end{pmatrix}
=\Sigma_{ij},
\]
because $R_{ij}=\rho_{ij}=c_{ij}/\sqrt{v_i^{(ij)}v_j^{(ij)}}$ and the visible variances agree by \textup{(i)}.
\end{proof}

The contextual object is therefore the edge-correlation completion set
\[
\mathcal E(G):=\{\rho\in(-1,1)^E:\rho\text{ extends to a correlation matrix on }V\}.
\]
Bell and temporal inequalities sit below $\mathcal E(G)$ as readout shadows. At law level the primary obstruction is gluing failure for the full pair laws.

The next theorem identifies the easy graph class.

\begin{theorem}[Chordal clique gluing]\label{thm:chordal-gluing}
Let $G$ be a chordal graph with maximal cliques $C_1,\dots,C_m$. For each clique $C_\alpha$, let $\nu_{C_\alpha}$ be a centred Gaussian law with positive definite clique covariance. Assume that whenever two neighbouring cliques in a clique tree meet in a separator $S=C_\alpha\cap C_\beta$, the induced $S$-marginals of $\nu_{C_\alpha}$ and $\nu_{C_\beta}$ agree. Then there exists a centred Gaussian law on all of $V$ whose $C_\alpha$-marginal is $\nu_{C_\alpha}$ for every $\alpha$.
\end{theorem}

\begin{proof}
The key separator-gluing step is explicit. If centred Gaussian laws on index sets $U$ and $W$ agree on the overlap $S=U\cap W$, then
\[
p(x_{U\cup W})=\frac{p_U(x_U)p_W(x_W)}{p_S(x_S)}
\]
defines a centred Gaussian law on $U\cup W$ with the prescribed $U$- and $W$-marginals. The quotient is well-defined because the overlap marginal agrees, and direct integration recovers the stated marginals.

Now choose a clique tree for $G$ and argue by induction on the number of cliques. For one clique there is nothing to prove. For more than one clique, remove a leaf clique $C_\ell$ with separator $S=C_\ell\cap C'$. By induction the remaining clique family admits a common Gaussian law. Gluing that law to $\nu_{C_\ell}$ across $S$ by the separator formula above produces the desired global Gaussian law.
\end{proof}

In particular, if $G$ is a tree, then the maximal cliques are just edges. Hence edgewise positive pair laws together with vertex variance consistency are sufficient for common Gaussian gluing. Trees therefore have no further law-level obstruction. The first genuinely nontrivial temporal obstruction appears when one passes from a tree to the first temporal clique.

The first fully observed temporal case is the triangle $K_3$. Because it is chordal, the correlation-completion condition reduces to positivity of the single $3\times3$ correlation matrix.

\begin{theorem}[Exact three-time triangle theorem]\label{thm:temporal-triangle}
Let $\Sigma_{01},\Sigma_{12},\Sigma_{02}\in\SPD(2)$ be centred Gaussian pair laws with entries
\[
\Sigma_{ij}=\begin{pmatrix} v_i^{(ij)} & c_{ij}\\ c_{ij} & v_j^{(ij)}\end{pmatrix},
\qquad
\rho_{ij}:=\frac{c_{ij}}{\sqrt{v_i^{(ij)}v_j^{(ij)}}}.
\]
There exists a centred Gaussian law on $(X_0,X_1,X_2)$ with these three pair marginals if and only if
\[
v_0^{(01)}=v_0^{(02)},\qquad
v_1^{(01)}=v_1^{(12)},\qquad
v_2^{(02)}=v_2^{(12)},
\]
and
\[
R=\begin{pmatrix}
1 & \rho_{01} & \rho_{02}\\
\rho_{01} & 1 & \rho_{12}\\
\rho_{02} & \rho_{12} & 1
\end{pmatrix}\succeq0.
\]
Equivalently,
\begin{equation}\label{eq:temporal-triangle-det}
1+2\rho_{01}\rho_{02}\rho_{12}-\rho_{01}^2-\rho_{02}^2-\rho_{12}^2\ge0.
\end{equation}
\end{theorem}

\begin{proof}
The gluing criterion is Proposition~\ref{prop:graph-gluing} specialised to $G=K_3$. For the determinant form, positivity of $R$ is equivalent to nonnegativity of its determinant because the $1\times1$ and $2\times2$ principal minors are already positive. A direct computation gives \eqref{eq:temporal-triangle-det}.
\end{proof}

\begin{corollary}[Temporal sign-shadow tetrahedron]\label{cor:temporal-sign-shadow}
Let $(X_0,X_1,X_2)$ be any centred Gaussian triple satisfying Theorem~\ref{thm:temporal-triangle}, define $S_i:=\operatorname{sgn}(X_i)$, and set
\[
y_{ij}:=\frac{2}{\pi}\arcsin(\rho_{ij}).
\]
Then
\[
\mathbb E[S_iS_j]=y_{ij},
\]
and the lifted correlator triple $(y_{01},y_{02},y_{12})$ lies in the classical three-variable correlation tetrahedron:
\begin{align}
1+y_{01}+y_{02}+y_{12}&\ge0,\label{eq:temporal-tetra-a}\\
1+y_{01}-y_{02}-y_{12}&\ge0,\label{eq:temporal-tetra-b}\\
1-y_{01}+y_{02}-y_{12}&\ge0,\label{eq:temporal-tetra-c}\\
1-y_{01}-y_{02}+y_{12}&\ge0.\label{eq:temporal-tetra-d}
\end{align}
In particular, the first temporal lifted shadow is the Leggett-Garg tetrahedron attached to the three-time clique \cite{LeggettGarg1985}.
\end{corollary}

\begin{proof}
For a centred Gaussian pair with correlation $\rho_{ij}$, Sheppard's formula gives
\[
\mathbb E[\operatorname{sgn}(X_i)\operatorname{sgn}(X_j)]=\frac{2}{\pi}\arcsin(\rho_{ij}).
\]
The tetrahedral inequalities are the classical three-variable correlation constraints for sign variables \cite{LeggettGarg1985}. Applying them to $(S_0,S_1,S_2)$ yields \eqref{eq:temporal-tetra-a}-\eqref{eq:temporal-tetra-d}.
\end{proof}

The temporal tree-versus-clique divide is therefore now visible. Trees glue automatically after repeated-marginal consistency, but they need not determine a unique global law.

\begin{proposition}[Path completion interval and nonuniqueness]\label{prop:path-completion-interval}
Fix unit variances on the path $0-1-2$ and edge correlations $\rho_{01},\rho_{12}\in(-1,1)$. A Gaussian completion exists for every
\[
\rho_{02}\in\Bigl[\rho_{01}\rho_{12}-\sqrt{(1-\rho_{01}^2)(1-\rho_{12}^2)},\,\rho_{01}\rho_{12}+\sqrt{(1-\rho_{01}^2)(1-\rho_{12}^2)}\Bigr].
\]
This interval is nondegenerate whenever $|\rho_{01}|<1$ and $|\rho_{12}|<1$. The distinguished tree-Markov construction selects the centre point $\rho_{02}=\rho_{01}\rho_{12}$.
\end{proposition}

\begin{proof}
A completion with correlation matrix
\[
R=\begin{pmatrix}1&\rho_{01}&\rho_{02}\\ \rho_{01}&1&\rho_{12}\\ \rho_{02}&\rho_{12}&1\end{pmatrix}
\]
exists if and only if $R\succeq0$. Since the $2\times2$ principal minors are positive, this is equivalent to
\[
1+2\rho_{01}\rho_{12}\rho_{02}-\rho_{01}^2-\rho_{12}^2-\rho_{02}^2\ge0.
\]
This is a quadratic inequality in $\rho_{02}$ with roots
\[
\rho_{01}\rho_{12}\pm\sqrt{(1-\rho_{01}^2)(1-\rho_{12}^2)},
\]
so the admissible interval is the one stated. Its midpoint is $\rho_{01}\rho_{12}$.
\end{proof}

\begin{proposition}[Fibre crossing at fixed correlation shadow]\label{prop:graph-fibre-crossing}
Let $G$ have a vertex of degree at least $2$, and let $\rho\in\mathcal E(G)$ lie in the interior of the correlation-completion set. Then there exist both gluable and non-gluable full edge-law families having the same edge-correlation data $\rho$.
\end{proposition}

\begin{proof}
Choose any centred Gaussian completion with unit variances, so each edge block is
\[
\Sigma_{ij}=\begin{pmatrix}1 & \rho_{ij}\\ \rho_{ij} & 1\end{pmatrix}
\]
and the family is gluable. Now pick a vertex $i$ used in at least two edges and choose one incident edge $\{i,j\}$. Replace only the variance of $i$ on that edge by a number $\lambda\neq1$, and replace the covariance on that edge by $\rho_{ij}\sqrt{\lambda}$. The normalised correlation on that edge remains $\rho_{ij}$, and all other edge correlations are untouched, so the edge-correlation shadow is unchanged. But the variance of vertex $i$ is now inconsistent across its incident edges, so Proposition~\ref{prop:graph-gluing} forbids common gluing.
\end{proof}

Thus Bell correlators do not exhaust the law-level frontier even before one passes to binary readout inequalities. They are shadows of a fuller pair-law compatibility problem, and the Bell square treated in the main text is the first nonchordal instance of that broader contextual branch.

\section{Scalar summaries and Schur-gap diagnostics}\label{app:scalar}

This appendix records the minimal scalar diagnostics supporting the Schur-gap discussion in Section~\ref{sec:quotient-precision}. No new structure is introduced here.

\subsection{Scalar mediation invariants}

Let $G\succeq0$ denote the Schur gap from \eqref{eq:schur-gap}. The basic scalar diagnostics are
\begin{equation}\label{eq:scalar-invariants}
\Gamma_{\mathrm{tr}}:=\operatorname{tr}(G),
\qquad
\Gamma_{\mathrm{op}}:=\|G\|_{\mathrm{op}},
\qquad
\Gamma_{\det}:=\log\det H_{\mathrm{comp}}-\log\det H_{\mathrm{quot}}.
\end{equation}
They measure total hidden contribution, maximal hidden contribution, and volumetric hidden contribution, respectively.

\begin{proposition}\label{prop:scalar-invariants}
The quantities in \eqref{eq:scalar-invariants} are all nonnegative.
\end{proposition}

\begin{proof}
Since $G\succeq0$, its trace and operator norm are nonnegative. Also $H_{\mathrm{comp}}\succeq H_{\mathrm{quot}}\succ0$ by \eqref{eq:schur-gap}, so monotonicity of the determinant on the positive cone gives $\det H_{\mathrm{comp}}\ge \det H_{\mathrm{quot}}$, hence $\Gamma_{\det}\ge0$.
\end{proof}

\begin{remark}
These diagnostics are only summaries. The object itself is the Schur gap operator $G$, and the main text uses the scalar summaries only as numerical readers of hidden contribution.
\end{remark}

\section{Adjacency, fill-in, and locality diagnostics}\label{app:adjacency}

This appendix provides the service definitions behind the locality discussion in Section~\ref{sec:05_hidden_locality}. It introduces quotient-induced adjacency and the associated fill-in diagnostics, but it does not add any new locality theorem beyond the main text.

\subsection{Fibre-induced visible adjacency}

To turn adjacency loss into an intrinsic quotient invariant, the visible mask must be induced by quotient geometry rather than chosen externally.

\begin{definition}[Fibre distance and quotient adjacency]\label{def:fibre-distance}
Let $\pi:X\to Y$ be a quotient of a metric graph or metric state space $(X,d_X)$. Define the lifted fibre distance by
\begin{equation}\label{eq:fibre-distance}
d_\pi(y,y'):=\inf\{d_X(u,v): u\in\pi^{-1}(y),\ v\in\pi^{-1}(y')\}.
\end{equation}
If upstairs locality has interaction range $r$, define visible adjacency by
\begin{equation}\label{eq:quotient-adjacency}
y\sim_\pi y'\quad\Longleftrightarrow\quad d_\pi(y,y')\le r.
\end{equation}
\end{definition}

Let $M_\pi$ denote the mask that keeps only matrix entries consistent with \eqref{eq:quotient-adjacency}.

\begin{definition}[Locality-defect operator]\label{def:locality-defect}
Let $H_{\mathrm{vis}}\in\SPD(Y)$ be a visible precision and let the visible Hilbert structure on $Y$ be fixed once and for all. The quotient locality-defect operator is
\begin{equation}\label{eq:locality-defect}
\mathcal N_\pi:=H_{\mathrm{vis}}-M_\pi(H_{\mathrm{vis}}).
\end{equation}
Its singular values, stable rank, and retained Frobenius mass are the quotient adjacency-loss diagnostics. These are numerical readers of Schur-complement fill-in and spectral complexity in sparse positive systems \cite{GeorgeLiu1981,HornJohnson2013,RudelsonVershynin2007}.
\end{definition}

\begin{proposition}[Intrinsic diagnostics]\label{prop:intrinsic-diagnostics}
Relative to a fixed visible Hilbert structure and interaction range $r$, the singular spectrum of $\mathcal N_\pi$ and the scalar mediation observables of Appendix~\ref{app:scalar} are intrinsic to the quotient construction.
\end{proposition}

\begin{proof}
Once $d_\pi$ and $r$ are fixed, the mask $M_\pi$ is determined by the quotient itself. The operator $\mathcal N_\pi$ is then determined by $H_{\mathrm{vis}}$ and the quotient geometry. The scalar observables of Appendix~\ref{app:scalar} depend only on the Schur gap \eqref{eq:schur-gap}, which is itself determined by the comparison between quotient observation and naive compression.
\end{proof}

\section{Renewal / transformed semi-coarse-graining benchmark}\label{app:renewal}

This appendix supports the transformed semi-coarse-graining test discussed in Section~\ref{sec:10_benchmark_synthesis}. The point of the benchmark is methodological. It sits just beyond the static quotient theorem domain of the main text, but it still yields a tractable visible law. The source model is the four-state benchmark introduced in the transformed semi-coarse-graining study of entropy-production estimation under partial observation by Kapustin, Ghosal, and Bisker \cite{KapustinGhosalBisker2024}. In spirit it belongs near coarse-graining comparisons such as \cite{TezaStella2020}, even though the transformed observer here is renewal-level rather than a static state quotient. What matters here is that the transformed observer is not itself a static quotient of the original continuous-time Markov chain, yet the visible law closes far enough to support a finite completion and a genuine benchmark comparison. For basic finite-state chain notation and path-block bookkeeping we follow standard Markov-chain conventions \cite{KemenySnell1976}.

\subsection{Four-state benchmark and transformed observer}

Let the original continuous-time Markov chain have states $\{1,2,3,4\}$ with hidden block $H=\{3,4\}$ and rates
\begin{equation}\label{eq:bisker-rates}
1\to2:2e^x,
\quad
1\to4:1,
\quad
2\to1:3e^{-x},
\quad
2\to3:2,
\quad
2\to4:35,
\end{equation}
\begin{equation}\label{eq:bisker-rates-hidden}
3\to2:50,
\quad
3\to4:0.7,
\quad
4\to1:8,
\quad
4\to2:0.2,
\quad
4\to3:75.
\end{equation}
Following \cite{KapustinGhosalBisker2024}, a transformed semi-coarse-grained trajectory groups each completed run of consecutive visits to the hidden macrostate $H$ into a symbol $H_n$, where $n$ is the number of hidden visits in that run and the waiting time of $H_n$ is the sum of the waiting times in those $n$ visits. To keep the previous visible state explicit, write
\[
A_n:=(1,H_n),
\qquad
B_n:=(2,H_n).
\]
Define
\begin{equation}\label{eq:bisker-r3r4}
r_3:=\frac{0.7}{50.7},
\qquad
r_4:=\frac{75}{83.2},
\qquad
r:=r_3r_4.
\end{equation}
Numerically,
\[
r\approx 0.0124459490,
\qquad
\Pr(n\le 2\mid\text{enter hidden})=1-r\approx 0.987554051,
\qquad
\Pr(n>4\mid\text{enter hidden})=r^2\approx 1.54901647\times10^{-4}.
\]
Thus the transformed observer is already sharply concentrated on very short hidden excursions.

\begin{proposition}[Exact parity collapse of the transformed sequence law]\label{prop:bisker-parity-collapse}
For the benchmark \eqref{eq:bisker-rates} to \eqref{eq:bisker-r3r4}, the countable augmented transformed sequence law on
\[
\{1,2\}\cup\{A_n,B_n:n\ge1\}
\]
collapses, at the visible-law level, to a four-symbol second-order parity observer with alphabet
\begin{equation}\label{eq:bisker-parity-alphabet}
\mathcal Y_{\mathrm{par}}=\{A_o,A_e,B_o,B_e\},
\end{equation}
where the previous visible state is retained as the second-order memory coordinate and parity records whether the completed hidden run has odd or even length. Equivalently, after unfolding the previous visible state, one obtains a six-symbol first-order Markov representation. More precisely, conditioned on entry from visible state $1$,
\begin{equation}\label{eq:bisker-parity-1}
P_o^{(1)}=\frac{1-r_4}{1-r},
\qquad
P_e^{(1)}=\frac{r_4(1-r_3)}{1-r},
\end{equation}
and conditioned on entry from visible state $2$,
\begin{equation}\label{eq:bisker-parity-2}
P_o^{(2)}=
\frac{\frac{2}{37}(1-r_3)+\frac{35}{37}(1-r_4)}{1-r},
\qquad
P_e^{(2)}=
\frac{\frac{2}{37}r_3(1-r_4)+\frac{35}{37}r_4(1-r_3)}{1-r}.
\end{equation}
The next visible state depends on the completed excursion only through the previous visible state and this odd or even parity label.
\end{proposition}

\begin{proof}
Inside the hidden block the chain alternates between states $3$ and $4$ until escape to the visible sector. For entry through state $4$ the completed excursion exits from hidden state $4$ when the number of hidden visits is odd and from hidden state $3$ when it is even. The probabilities \eqref{eq:bisker-parity-1} follow by summing the corresponding geometric series. For entry from visible state $2$, the first hidden state is a fixed mixture of $3$ and $4$, and summing the odd and even geometric series gives \eqref{eq:bisker-parity-2}. Since the visible return law depends only on the hidden exit microstate, it depends only on previous visible state and parity.
\end{proof}

The corresponding visible return probabilities are explicit. In particular,
\begin{equation}\label{eq:bisker-return-A}
A_o\to1 \text{ with probability } \frac{8}{8.2},
\qquad
A_o\to2 \text{ with probability } \frac{0.2}{8.2},
\qquad
A_e\to2 \text{ with probability }1,
\end{equation}
and for the $B$-symbols,
\begin{equation}\label{eq:bisker-return-B}
\Pr(1\mid B_o)\approx 0.6207014153,
\qquad
\Pr(2\mid B_o)\approx 0.3792985847,
\end{equation}
\begin{equation}\label{eq:bisker-return-Be}
\Pr(1\mid B_e)\approx 8.5325870\times 10^{-5},
\qquad
\Pr(2\mid B_e)\approx 0.9999146741.
\end{equation}
Thus the transformed observer exposes the hidden $3\leftrightarrow4$ mode as an almost visible parity channel.

\begin{theorem}[Two-kernel timing collapse]\label{thm:bisker-two-kernel}
Let
\[
L_i(s):=\frac{\lambda_i}{\lambda_i+s},
\qquad
\lambda_3=50.7,
\qquad
\lambda_4=83.2.
\]
Then the waiting-time transforms of the parity symbols lie exactly in the two-dimensional span generated by
\begin{equation}\label{eq:bisker-G4o}
G_{4,o}(s)=\frac{(1-r)L_4(s)}{1-rL_3(s)L_4(s)},
\qquad
G_{4,e}(s)=\frac{(1-r)L_3(s)L_4(s)}{1-rL_3(s)L_4(s)}.
\end{equation}
More precisely,
\begin{equation}\label{eq:bisker-G3o}
G_{3,o}(s)=\frac{\lambda_3}{\lambda_4}G_{4,o}(s)+\Bigl(1-\frac{\lambda_3}{\lambda_4}\Bigr)G_{4,e}(s)
=\frac{39}{64}G_{4,o}(s)+\frac{25}{64}G_{4,e}(s),
\end{equation}
\begin{equation}\label{eq:bisker-GBo}
G_{B_o}(s)=w\,G_{4,o}(s)+(1-w)\,G_{4,e}(s),
\qquad
w=\frac{15093}{17593},
\end{equation}
while
\begin{equation}\label{eq:bisker-GBe}
G_{B_e}(s)=G_{4,e}(s).
\end{equation}
All four transforms therefore share the same quadratic denominator
\begin{equation}\label{eq:bisker-common-den}
D(s)=(\lambda_3+s)(\lambda_4+s)-r\lambda_3\lambda_4=(\lambda_3+s)(\lambda_4+s)-52.5.
\end{equation}
\end{theorem}

\begin{proof}
The odd and even transforms for entry through state $4$ are summed geometric series in the product $L_3(s)L_4(s)$, which gives \eqref{eq:bisker-G4o}. Entry through state $3$ either exits immediately through $3$ or first moves to $4$, which yields \eqref{eq:bisker-G3o}. The $B$-symbol transforms are then the conditional mixtures induced by the two possible hidden entry microstates, giving \eqref{eq:bisker-GBo} and \eqref{eq:bisker-GBe}. The common denominator is immediate from \eqref{eq:bisker-G4o}.
\end{proof}

Theorem~\ref{thm:bisker-two-kernel} shows that the hidden timing sector of the transformed observer is not an uncontrolled memory tower. It is exactly two-dimensional at the visible-law level.

\begin{proposition}[Exact finite completion]\label{prop:bisker-finite-completion}
The parity-lifted transformed observer admits an explicit finite CTMC completion. One such completion has visible states $1,2$ together with four hidden phase blocks: a two-state block for $A_o$, a two-state block for $A_e$, a four-state block for $B_o$, and a four-state block for $B_e$. The two basic hidden excursion generators are
\begin{equation}\label{eq:bisker-S4}
S^{(4)}=
\begin{pmatrix}
-50.7 & 50.7\\
175/169 & -83.2
\end{pmatrix},
\end{equation}
\begin{equation}\label{eq:bisker-S3}
S^{(3)}=
\begin{pmatrix}
-50.7 & 525/832\\
83.2 & -83.2
\end{pmatrix},
\end{equation}
and the $B$-blocks are direct sums of these two generators with the conditional odd or even entry weights. Hence an finite completion exists with total state count $14$.
\end{proposition}

\begin{proof}
The parity transforms in Theorem~\ref{thm:bisker-two-kernel} are rational phase-type laws of order at most two \cite{Buchholz1994}. The $A$-symbols require one conditioned two-state excursion block each. The $B$-symbols are conditional mixtures over the two possible hidden entry microstates, so each is realised exactly by the corresponding direct-sum block. Adding the visible states $1,2$ gives $2+2+2+4+4=14$ states.
\end{proof}

\begin{remark}[One exact completion and the renewal frontier]\label{rem:bisker-renewal-frontier}
The transformed semi-coarse-graining benchmark works for a structural reason, not because of an accidental good observer choice. In the 4-state benchmark, the hidden 2-state block alternates deterministically, so the transformed sequence law depends only on the parity of the hidden run length, not on the full run length. Proposition~\ref{prop:bisker-parity-collapse} therefore gives a finite parity collapse in the sequence channel. The timing channel also closes: by Theorem~\ref{thm:bisker-two-kernel}, the parity-symbol waiting-time transforms lie in a two-dimensional phase-type span and share a common quadratic denominator, so the transformed timing sector is not an uncontrolled memory tower. Together these two facts show that the transformed visible law is already finite in disguise. The explicit 14-state CTMC construction is therefore not an ad hoc completion but a realisation of an intrinsically finite parity-lifted semi-Markov visible law.

On the explicit 14-state completion from Proposition~\ref{prop:bisker-finite-completion}, the transformed observer has active visible tangent rank $4$ and hidden-load rank $2$ at the visible stalling point. Numerically,
\[
\lambda(T_{\mathrm{tr}})\approx (0,0.81855,0.88615,9.26522,10.93096),
\]
\[
\lambda(X_{\mathrm{tr}})\approx (0,0.69542,0.85140,8.94220,10.77309),
\]
\[
\lambda(\Lambda_{\mathrm{tr}})\approx (0,0,0.07834,0.19437).
\]
By contrast, the original full-CG observer on the 4-state chain is only two-dimensional and is already a one-hard-direction quotient observer, with hidden-load spectrum approximately $(0,0.274678)$ at stall. This sharpens the true open problem: not whether one finite completion exists, but whether the completion-independent tangent ceiling and hidden-load structure can be fixed directly from the renewal visible law itself rather than from a chosen finite completion. This is the natural extension frontier beyond the static quotient geometry proved in the main text.
\end{remark}


\section{Blind benchmark diagnostics}\label{app:blind-benchmark}

This appendix records the auxiliary diagnostics behind the lock-then-reveal benchmark discussed in Section~\ref{sec:11_falsification}. They are included to document how the protocol behaved under finite observation and adaptive reveal, not to enlarge the theorem domain of the paper.

\subsection{Exact-ensemble finite-error spread}

On the 256-world ensemble, the retained-signal and detectability diagnostics show a nontrivial but controlled finite-error spread across worlds. The important methodological point is not that every world is reconstructed with the same accuracy, but that the ensemble error distribution remains concentrated enough for the finite tangent law to function as an operational discriminator at the level reported in the main text. This is the diagnostic content that underlay the ensemble error histograms in the archive.

\subsection{Adaptive reveal as a diagnostic frontier}

The adaptive-reveal experiment was run as a secondary check on whether theorem-targeted observation can improve the operational signal. On the 64-world diagnostic panel, the oracle theorem-target reveal improved the mean second-step revealed signal relative to the random baseline. The same comparison in the normalised detectability score was again positive, but smaller. This is why Section~\ref{sec:11_falsification} treats adaptive reveal as a diagnostic frontier rather than as a new theorem-bearing claim: the gain is real, but it belongs to the protocol layer, not to the intrinsic quotient-visible law.

\subsection{Exact versus sampled Bell verdicts}

The Bell-facing part of the blind benchmark was evaluated in both law-level and sampled modes. On the recorded hero families, law-level verdict recovery cleanly identified the glueable versus non-glueable cases at the law level. The sampled mode was materially weaker and missed the glueable families on the recorded benchmark, which is why the paper treats the Bell verdict as a law-level statement and uses the sampled track only as a robustness layer. This is the boundary that matters for interpretation: the lock-then-reveal protocol works as a test of the law-level compatibility machinery, but not yet as a theorem-safe sampled Bell pipeline.

\section{Methodological bridges and limits of the positive-survivor programme}\label{app:bridges}

This appendix retains one methodological bridge only, namely the distinction between structural and practical identifiability. Its role is to support the transfer boundary stated in Section~\ref{sec:12_scope}: neighbouring hidden-to-visible problems can exhibit the same survivor logic without thereby falling inside the theorem domain of the present paper. What is claimed here is limited and explicit. A regular structural quotient carries the descended Fisher form, enlarging the admissible experiment class sharpens that quotient, and the hidden-load geometry from the main text gives a finite-dimensional instance of the same split. What is not claimed is a universal theorem unifying all hidden-to-visible reductions.

\subsection{Bridge A: structural versus practical identifiability}

Let $\Theta$ be a smooth parameter manifold, let $\mathcal E$ be an admissible experiment class, and for each $E\in\mathcal E$ let
\[
p_E(z\mid\theta)
\]
be the observational law. Define structural equivalence relative to $\mathcal E$ by
\[
\theta\sim\theta'\quad\Longleftrightarrow\quad p_E(\cdot\mid\theta)=p_E(\cdot\mid\theta')\ \text{for all }E\in\mathcal E.
\]
Assume that near a regular point the equivalence classes form a smooth quotient manifold
\[
q:\Theta\to \bar\Theta:=\Theta/\!\sim.
\]
For a fixed experiment $E$, let $I_E(\theta)$ denote the Fisher information bilinear form on $T_\theta\Theta$.

\begin{theorem}\label{thm:bridge-identifiability}
Under the regular quotient hypothesis, each experiment $E\in\mathcal E$ factors uniquely through the quotient: there exists a reduced statistical model
\[
\bar p_E(z\mid \bar\theta),\qquad \bar\theta\in\bar\Theta,
\]
such that
\[
p_E(z\mid\theta)=\bar p_E(z\mid q(\theta)).
\]
Moreover:
\begin{enumerate}
\item for every vertical vector $v\in\ker dq_\theta$, the score annihilates,
\[
\partial_v \log p_E(z\mid\theta)=0;
\]
\item the Fisher form vanishes on vertical directions and descends uniquely to a bilinear form $\bar I_E$ on $T_{q(\theta)}\bar\Theta$ satisfying
\[
I_E(\theta)(u,w)=\bar I_E(q(\theta))(dq_\theta u,dq_\theta w);
\]
\item structural nonidentifiability is the vertical degeneracy of $q$, while practical nonidentifiability for a given experiment is degeneracy or strong anisotropy of the descended form $\bar I_E$ on the quotient tangent.
\end{enumerate}
\end{theorem}

\begin{proof}
By definition, each $p_E(\cdot\mid\theta)$ is constant on every structural fibre, so it depends only on $q(\theta)$; this gives the unique reduced model. Differentiating the factorisation along a vertical vector $v\in\ker dq_\theta$ gives the score annihilation. The Fisher identity then shows that $I_E(\theta)(v,\cdot)=0$, so $I_E(\theta)$ depends only on quotient tangent classes and therefore descends uniquely to $\bar I_E$.
\end{proof}

The theorem is local rather than global. It says that the correct canonical object is not a matrix written in some auxiliary horizontal complement, but the descended bilinear form on the structural quotient itself. A direct consequence is that enlarging the admissible experiment class can only refine the quotient and remove vertical directions; it cannot create a new structural fibre.

\begin{proposition}\label{prop:bridge-identifiability-refinement}
Let $\mathcal E_1\subseteq \mathcal E_2$ be two experiment classes, and assume both regular quotients exist near the point of interest:
\[
q_i:\Theta\to \bar\Theta_i=\Theta/\!\sim_i,\qquad i=1,2.
\]
Then there exists a unique local surjection $\pi:\bar\Theta_2\to \bar\Theta_1$ such that
\[
q_1=\pi\circ q_2.
\]
For every experiment $E\in\mathcal E_1$, the descended Fisher forms satisfy the pullback relation
\[
\bar I_E^{(2)}=\pi^*\bar I_E^{(1)}.
\]
Hence enlarging the experiment class can only sharpen the quotient geometry.
\end{proposition}

\begin{proof}
If two parameters are equivalent for the larger class $\mathcal E_2$, they are equivalent for the smaller class $\mathcal E_1$. Therefore every $\sim_2$-class sits inside a $\sim_1$-class, which induces the unique local map $\pi$ with $q_1=\pi\circ q_2$. The Fisher pullback identity then follows from Theorem~\ref{thm:bridge-identifiability} applied to the factorisation of the statistical model through the two quotients.
\end{proof}

A canonical toy model is the scalar decay law
\[
\dot x=-k x,
\qquad
Y_i=c x_0 e^{-k t_i}+\varepsilon_i,
\qquad
\varepsilon_i\sim\mathcal N(0,\sigma^2),
\]
with parameters $(c,x_0,k)$. The full observational law depends only on $a=cx_0$ and $k$, so the structural fibres are
\[
(c,x_0,k)\sim(\lambda c,x_0/\lambda,k),\qquad \lambda>0.
\]
The quotient is therefore two-dimensional with coordinates $(a,k)$, and the quotient Fisher form is the Fisher matrix of the reduced mean model $a e^{-k t_i}$. This shows in one line that a direction can be structurally invisible upstairs while the remaining quotient directions are still experiment-dependent in visibility.

In the present quotient setting this split becomes finite-dimensional rather than merely local.

\begin{theorem}[Exact inverse split for the hidden-load class]\label{thm:bridge-hidden-load-split}
Let $T\succeq0$ and let $S:=\operatorname{Ran}(T)$. Let $X$ be a visible law in the support-preserving interior beneath $T$, so that by Theorem~\ref{thm:hidden-load}
\[
X=T^{1/2}(I_S+\Lambda)^{-1}T^{1/2},\qquad \Lambda\succeq0\text{ on }S.
\]
Then:
\begin{enumerate}[label=\textup{(\roman*)},leftmargin=2.4em]
\item conditional on the tangent ceiling $T$, the hidden-load operator is uniquely identified by $X$ and equals
\begin{equation}\label{eq:bridge-load-recovery}
\Lambda=(T|_S)^{1/2}(X|_S)^{-1}(T|_S)^{1/2}-I_S;
\end{equation}
moreover,
\[
\rank(\Lambda)=\rank(T-X),
\qquad
\log\det(I_S+\Lambda)=\log\pdet(T)-\log\pdet(X);
\]
\item without an independently fixed ceiling, $X$ alone does not identify a unique hidden load: for every $T'\succeq X$ with $\operatorname{Ran}(T')=S$, there exists a unique $\Lambda_{T'}\succeq0$ on $S$ such that
\[
X=T'^{1/2}(I_S+\Lambda_{T'})^{-1}T'^{1/2},
\qquad
\Lambda_{T'}=(T'|_S)^{1/2}(X|_S)^{-1}(T'|_S)^{1/2}-I_S.
\]
Hence $X$ determines a full ceiling cone of compatible ceiling-load explanations.
\end{enumerate}
\end{theorem}

\begin{proof}
Part \textup{(i)} is the explicit recovery formula already built into Theorem~\ref{thm:hidden-load}. The rank and determinant identities are Corollary~\ref{cor:hidden-load-rank-volume}.

For part \textup{(ii)}, fix $T'\succeq X$ on the same support $S$ and define $\Lambda_{T'}$ by the displayed formula. Then
\[
T'^{-1/2}XT'^{-1/2}\preceq I_S.
\]
Since inversion reverses Loewner order on positive definite operators,
\[
(T'^{-1/2}XT'^{-1/2})^{-1}\succeq I_S,
\]
so $\Lambda_{T'}\succeq0$. Rearranging gives the stated representation of $X$, and uniqueness follows because $T'^{-1/2}XT'^{-1/2}=(I_S+\Lambda_{T'})^{-1}$.
\end{proof}

Combined with Corollary~\ref{cor:hidden-load-monotone}, this gives the order reversal at fixed ceiling:
\[
X_1\preceq X_0
\quad\Longleftrightarrow\quad
\Lambda_1\succeq\Lambda_0.
\]
It also gives the zero-load boundary $\Lambda=0\Longleftrightarrow X=T$.

\begin{proposition}[Gram rigidity and hidden-factor gauge]\label{prop:bridge-hidden-factor-gauge}
Fix $T$ and $X$ and let $\Lambda$ be the uniquely identified hidden-load operator from Theorem~\ref{thm:bridge-hidden-load-split}. Write $r:=\rank(\Lambda)$. Then every minimal hidden realisation is a Gram factorisation
\[
\Lambda=BB^\top,
\qquad B\in\mathbb R^{|S|\times r},
\]
and if $B,B'\in\mathbb R^{|S|\times r}$ both have full column rank and satisfy $BB^\top=B'B'^\top=\Lambda$, then
\[
B'=BQ
\]
for some orthogonal matrix $Q\in O(r)$. Consequently $(T,X)$ identifies the hidden Gram operator and its intrinsic invariants, but not a unique hidden basis, hidden factorisation, or hidden sparsity pattern.
\end{proposition}

\begin{proof}
The existence of a minimal factorisation follows from Theorem~\ref{thm:min-hidden-dim}. Assume $BB^\top=B'B'^\top=\Lambda$ with both factors of column rank $r$, and define
\[
Q:=B^+B',
\qquad
B^+=(B^\top B)^{-1}B^\top.
\]
Because $\operatorname{Ran}(B)=\operatorname{Ran}(\Lambda)=\operatorname{Ran}(B')$, one has
\[
BQ=BB^+B'=B'.
\]
Moreover,
\[
QQ^\top
=B^+B'B'^\top(B^+)^\top
=B^+\Lambda(B^+)^\top
=B^+BB^\top(B^+)^\top
=I_r.
\]
Since $Q$ is an $r\times r$ matrix, $QQ^\top=I_r$ implies $Q\in O(r)$. Hence $B'=BQ$.
\end{proof}


\begin{thebibliography}{58}\bibitem[{M.~D. Donsker} and {S.~R.~S.
  Varadhan}(1975{\natexlab{a}})]{DonskerVaradhan1975a}
{M.~D. Donsker} and {S.~R.~S. Varadhan}.
\newblock Asymptotic evaluation of certain {M}arkov process expectations for
  large time.~{I}, 1975{\natexlab{a}}.
\newblock \emph{Communications on Pure and Applied Mathematics},
  \textbf{28}:1-47, 1975.

\bibitem[{M.~D. Donsker} and {S.~R.~S.
  Varadhan}(1975{\natexlab{b}})]{DonskerVaradhan1975b}
{M.~D. Donsker} and {S.~R.~S. Varadhan}.
\newblock Asymptotic evaluation of certain {M}arkov process expectations for
  large time.~{II}, 1975{\natexlab{b}}.
\newblock \emph{Communications on Pure and Applied Mathematics},
  \textbf{28}:279-301, 1975.

\bibitem[{M.~D. Donsker} and {S.~R.~S. Varadhan}(1976)]{DonskerVaradhan1976}
{M.~D. Donsker} and {S.~R.~S. Varadhan}.
\newblock Asymptotic evaluation of certain {M}arkov process expectations for
  large time.~{III}, 1976.
\newblock \emph{Communications on Pure and Applied Mathematics},
  \textbf{29}(4):389-461, 1976.

\bibitem[{R.~S. Ellis}(2006)]{Ellis1985}
{R.~S. Ellis}.
\newblock \emph{Entropy, Large Deviations, and Statistical Mechanics}, 2006.
\newblock Grundlehren der mathematischen Wissenschaften, Vol.~271, Springer,
  1985. Reprinted as Classics in Mathematics, 2006.

\bibitem[{H.~Touchette}(2009)]{Touchette2009}
{H.~Touchette}.
\newblock The large deviation approach to statistical mechanics, 2009.
\newblock \emph{Physics Reports}, \textbf{478}:1-69, 2009.

\bibitem[{L.~Bertini} et~al.(2015){L.~Bertini}, {A.~Faggionato}, and
  {D.~Gabrielli}]{BertiniFaggionatoGabrielli2015}
{L.~Bertini}, {A.~Faggionato}, and {D.~Gabrielli}.
\newblock Large deviations of the empirical flow for continuous time {M}arkov
  chains, 2015.
\newblock \emph{Annales de l'I.H.P. Probabilit{\'e}s et Statistiques},
  \textbf{51}(3):867-900, 2015.

\bibitem[{J.~R. Dunkley}(2026)]{Dunkley2026FiniteObservation}
{J.~R. Dunkley}.
\newblock \emph{Finite Observation: A Spectral Theory of Nonequilibrium
  Visibility}, 2026.
\newblock Preprint, 2026.

\bibitem[{J.~Schnakenberg}(1976)]{Schnakenberg1976}
{J.~Schnakenberg}.
\newblock Network theory of microscopic and macroscopic behavior of master
  equation systems, 1976.
\newblock \emph{Reviews of Modern Physics}, \textbf{48}:571-585, 1976.

\bibitem[{J.~L. Lebowitz} and {H.~Spohn}(1999)]{LebowitzSpohn1999}
{J.~L. Lebowitz} and {H.~Spohn}.
\newblock A {G}allavotti-{C}ohen-type symmetry in the large deviation
  functional for stochastic dynamics, 1999.
\newblock \emph{Journal of Statistical Physics}, \textbf{95}:333-365, 1999.

\bibitem[{U.~Seifert}(2012)]{Seifert2012}
{U.~Seifert}.
\newblock Stochastic thermodynamics, fluctuation theorems and molecular
  machines, 2012.
\newblock \emph{Reports on Progress in Physics}, \textbf{75}:126001, 2012.

\bibitem[{M.~Esposito}(2012)]{Esposito2012}
{M.~Esposito}.
\newblock Stochastic thermodynamics under coarse graining, 2012.
\newblock \emph{Physical Review E}, \textbf{85}:041125, 2012.

\bibitem[{C.~Maes}(2020)]{Maes2020}
{C.~Maes}.
\newblock Frenesy: time-symmetric dynamical activity in nonequilibria, 2020.
\newblock \emph{Physics Reports}, \textbf{850}:1-33, 2020.

\bibitem[{I.~Schur}(1917)]{Schur1917}
{I.~Schur}.
\newblock {\"U}ber {P}otenzreihen, die im {I}nnern des {E}inheitskreises
  beschr{\"a}nkt sind, 1917.
\newblock \emph{Journal f{\"u}r die reine und angewandte Mathematik},
  \textbf{147}:205-232, 1917.

\bibitem[{W.~N. Anderson} and {Jr}(1971)]{Anderson1971}
{W.~N. Anderson} and {Jr}.
\newblock Shorted operators, 1971.
\newblock \emph{SIAM Journal on Applied Mathematics}, \textbf{20}(3):520-525,
  1971.

\bibitem[{W.~N. Anderson, Jr.} and {R.~J. Duffin}(1969)]{AndersonDuffin1969}
{W.~N. Anderson, Jr.} and {R.~J. Duffin}.
\newblock Series and parallel addition of matrices, 1969.
\newblock \emph{Journal of Mathematical Analysis and Applications},
  \textbf{26}(3):576-594, 1969.

\bibitem[{W.~N. Anderson, Jr.} and {G.~E. Trapp}(1975)]{AndersonTrapp1975}
{W.~N. Anderson, Jr.} and {G.~E. Trapp}.
\newblock Shorted operators.~{II}, 1975.
\newblock \emph{SIAM Journal on Applied Mathematics}, \textbf{28}(1):60-71,
  1975.

\bibitem[{R.~Bhatia}(2007)]{Bhatia2007}
{R.~Bhatia}.
\newblock \emph{Positive Definite Matrices}, 2007.
\newblock Princeton University Press, 2007.

\bibitem[{K.~L{\"o}wner}(1934)]{Loewner1934}
{K.~L{\"o}wner}.
\newblock {\"U}ber monotone {M}atrixfunktionen, 1934.
\newblock \emph{Mathematische Zeitschrift}, \textbf{38}:177-216, 1934.

\bibitem[{T.~Ando}(1979)]{Ando1979}
{T.~Ando}.
\newblock Concavity of certain maps on positive definite matrices and
  applications to {H}adamard products, 1979.
\newblock \emph{Linear Algebra and Its Applications}, \textbf{26}:203-241,
  1979.

\bibitem[{F.~Hansen} and {G.~K. Pedersen}(1982)]{HansenPedersen1982}
{F.~Hansen} and {G.~K. Pedersen}.
\newblock Jensen's inequality for operators and {L}{\"o}wner's theorem, 1982.
\newblock \emph{Mathematische Annalen}, \textbf{258}(3):229-241, 1982.

\bibitem[{R.~Bhatia}(1997)]{Bhatia1997}
{R.~Bhatia}.
\newblock \emph{Matrix Analysis}, 1997.
\newblock Graduate Texts in Mathematics, Vol.~169, Springer, 1997.

\bibitem[{F.~Kubo} and {T.~Ando}(1980)]{KuboAndo1980}
{F.~Kubo} and {T.~Ando}.
\newblock Means of positive linear operators, 1980.
\newblock \emph{Mathematische Annalen}, \textbf{246}:205-224, 1980.

\bibitem[{W.~Pusz} and {S.~L. Woronowicz}(1975)]{PuszWoronowicz1975}
{W.~Pusz} and {S.~L. Woronowicz}.
\newblock Functional calculus for sesquilinear forms and the purification map,
  1975.
\newblock \emph{Reports on Mathematical Physics}, \textbf{8}(2):159-170, 1975.

\bibitem[{L.~P. Kadanoff}(1966)]{Kadanoff1966}
{L.~P. Kadanoff}.
\newblock Scaling laws for {I}sing models near ${T}_c$, 1966.
\newblock \emph{Physics Physique Fizika}, \textbf{2}:263-272, 1966.

\bibitem[{K.~G. Wilson} and {J.~Kogut}(1974)]{WilsonKogut1974}
{K.~G. Wilson} and {J.~Kogut}.
\newblock The renormalization group and the $\varepsilon$ expansion, 1974.
\newblock \emph{Physics Reports}, \textbf{12}(2):75-200, 1974.

\bibitem[{J.~Polchinski}(1984)]{Polchinski1984}
{J.~Polchinski}.
\newblock Renormalization and effective {L}agrangians, 1984.
\newblock \emph{Nuclear Physics B}, \textbf{231}(2):269-295, 1984.

\bibitem[{K.~Fan}(1951)]{Fan1951}
{K.~Fan}.
\newblock Maximum properties and inequalities for the eigenvalues of completely
  continuous operators, 1951.
\newblock \emph{Proceedings of the National Academy of Sciences},
  \textbf{37}:760-766, 1951.

\bibitem[{R.~A. Horn} and {C.~R. Johnson}(2013)]{HornJohnson2013}
{R.~A. Horn} and {C.~R. Johnson}.
\newblock \emph{Matrix Analysis}, 2013.
\newblock Cambridge University Press, second edition, 2013.

\bibitem[{F.~Zhang} and {editor}(2005)]{Zhang2005}
{F.~Zhang} and {editor}.
\newblock \emph{The Schur Complement and Its Applications}, 2005.
\newblock Springer, New York, 2005.

\bibitem[{E.~V. Haynsworth}(1968)]{Haynsworth1968}
{E.~V. Haynsworth}.
\newblock On the {S}chur complement, 1968.
\newblock \emph{Basel Mathematical Notes}, BMN~20, 1968.

\bibitem[{D.~Carlson} et~al.(1974){D.~Carlson}, {E.~Haynsworth}, and
  {T.~Markham}]{CarlsonHaynsworthMarkham1974}
{D.~Carlson}, {E.~Haynsworth}, and {T.~Markham}.
\newblock A generalization of the {S}chur complement by means of the
  {M}oore-{P}enrose inverse, 1974.
\newblock \emph{SIAM Journal on Applied Mathematics}, \textbf{26}(1):169-175,
  1974.

\bibitem[{E.~B. Bogomol'nyi}(1976)]{Bogomolnyi1976}
{E.~B. Bogomol'nyi}.
\newblock Stability of classical solutions, 1976.
\newblock \emph{Soviet Journal of Nuclear Physics}, \textbf{24}(4):449-454,
  1976.

\bibitem[{A.~P. Dempster}(1972)]{Dempster1972}
{A.~P. Dempster}.
\newblock Covariance selection, 1972.
\newblock \emph{Biometrics}, \textbf{28}:157-175, 1972.

\bibitem[{S.~L. Lauritzen}(1996)]{Lauritzen1996}
{S.~L. Lauritzen}.
\newblock \emph{Graphical Models}, 1996.
\newblock Oxford Statistical Science Series, No.~17, Oxford University Press,
  1996.

\bibitem[{V.~Chandrasekaran} et~al.(2012){V.~Chandrasekaran}, {P.~A. Parrilo},
  and {A.~S. Willsky}]{ChandrasekaranParriloWillsky2012}
{V.~Chandrasekaran}, {P.~A. Parrilo}, and {A.~S. Willsky}.
\newblock Latent variable graphical model selection via convex optimization,
  2012.
\newblock \emph{The Annals of Statistics}, \textbf{40}(4):1935-1967, 2012.

\bibitem[{R.~Grone} et~al.(1984){R.~Grone}, {C.~R. Johnson}, {E.~M. S{\'a}},
  and {H.~Wolkowicz}]{GroneJohnsonSaWolkowicz1984}
{R.~Grone}, {C.~R. Johnson}, {E.~M. S{\'a}}, and {H.~Wolkowicz}.
\newblock Positive definite completions of partial {H}ermitian matrices, 1984.
\newblock \emph{Linear Algebra and Its Applications}, \textbf{58}:109-124,
  1984.

\bibitem[{J.~Williamson}(1936)]{Williamson1936}
{J.~Williamson}.
\newblock On the algebraic problem concerning the normal forms of linear
  dynamical systems, 1936.
\newblock \emph{American Journal of Mathematics}, \textbf{58}(1):141-163,
  1936.

\bibitem[{R.~A. Horn} and {C.~R. Johnson}(1991)]{HornJohnson1991}
{R.~A. Horn} and {C.~R. Johnson}.
\newblock \emph{Topics in Matrix Analysis}, 1991.
\newblock Cambridge University Press, 1991.

\bibitem[{A.~George} and {J.~W.~H. Liu}(1981)]{GeorgeLiu1981}
{A.~George} and {J.~W.~H. Liu}.
\newblock \emph{Computer Solution of Large Sparse Positive Definite Systems},
  1981.
\newblock Prentice-Hall, Englewood Cliffs, NJ, 1981.

\bibitem[{P.~G. Doyle} and {J.~L. Snell}(1984)]{DoyleSnell1984}
{P.~G. Doyle} and {J.~L. Snell}.
\newblock \emph{Random Walks and Electric Networks}, 1984.
\newblock Carus Mathematical Monographs, Vol.~22, Mathematical Association of
  America, 1984.

\bibitem[{D.~M. Malioutov} et~al.(2006){D.~M. Malioutov}, {J.~K. Johnson}, and
  {A.~S. Willsky}]{MalioutovJohnsonWillsky2006}
{D.~M. Malioutov}, {J.~K. Johnson}, and {A.~S. Willsky}.
\newblock Walk-sums and belief propagation in {G}aussian graphical models,
  2006.
\newblock \emph{Journal of Machine Learning Research}, \textbf{7}:2031-2064,
  2006.

\bibitem[{S.~Demko} et~al.(1984){S.~Demko}, {W.~F. Moss}, and {P.~W.
  Smith}]{DemkoMossSmith1984}
{S.~Demko}, {W.~F. Moss}, and {P.~W. Smith}.
\newblock Decay rates for inverses of band matrices, 1984.
\newblock \emph{Mathematics of Computation}, \textbf{43}(168):491-499, 1984.

\bibitem[{M.~Benzi} and {G.~H. Golub}(1999)]{BenziGolub1999}
{M.~Benzi} and {G.~H. Golub}.
\newblock Bounds for the entries of matrix functions with applications to
  preconditioning, 1999.
\newblock \emph{BIT Numerical Mathematics}, \textbf{39}(3):417-438, 1999.

\bibitem[{M.~Benzi} and {N.~Razouk}(2007)]{BenziRazouk2007}
{M.~Benzi} and {N.~Razouk}.
\newblock Decay bounds and {O}(n) algorithms for approximating functions of
  sparse matrices, 2007.
\newblock \emph{Electronic Transactions on Numerical Analysis},
  \textbf{28}:16-39, 2007.

\bibitem[{T.~J. Rivlin}(1990)]{Rivlin1974}
{T.~J. Rivlin}.
\newblock \emph{Chebyshev Polynomials: From Approximation Theory to Algebra and
  Number Theory}, 1990.
\newblock Wiley, 1974; second edition, 1990.

\bibitem[{L.~N. Trefethen}(2013)]{Trefethen2013}
{L.~N. Trefethen}.
\newblock \emph{Approximation Theory and Approximation Practice}, 2013.
\newblock SIAM, Philadelphia, 2013.

\bibitem[{U.~Kapustin} et~al.(2024){U.~Kapustin}, {A.~Ghosal}, and
  {G.~Bisker}]{KapustinGhosalBisker2024}
{U.~Kapustin}, {A.~Ghosal}, and {G.~Bisker}.
\newblock Utilizing time-series measurements for entropy-production estimation
  in partially observed systems, 2024.
\newblock \emph{Physical Review Research}, \textbf{6}:023039, 2024.

\bibitem[Liepelt and Lipowsky(2007)]{LiepeltLipowsky2007}
{A.~D. Liepelt} and {R.~Lipowsky}.
\newblock Kinesin's network of chemomechanical motor cycles.
\newblock \emph{Physical Review Letters}, \textbf{98}(25):258102, 2007.

\bibitem[Banerjee et~al.(2017)]{BanerjeeKolomeiskyIgoshin2017}
{S.~Banerjee}, {A.~B. Kolomeisky}, and {O.~A. Igoshin}.
\newblock Elucidating interplay of speed and accuracy in biological error correction, 2017.
\newblock \emph{Proceedings of the National Academy of Sciences}, \textbf{114}(20):5183-5188, 2017.

\bibitem[{C.~R. Rao}(1945)]{Rao1945}
{C.~R. Rao}.
\newblock Information and accuracy attainable in the estimation of statistical
  parameters, 1945.
\newblock \emph{Bulletin of the Calcutta Mathematical Society},
  \textbf{37}:81-91, 1945.

\bibitem[{S.~Amari} and {H.~Nagaoka}(2000)]{AmariNagaoka2000}
{S.~Amari} and {H.~Nagaoka}.
\newblock \emph{Methods of Information Geometry}, 2000.
\newblock Translations of Mathematical Monographs, Vol.~191, American
  Mathematical Society, 2000.

\bibitem[{N.~Ay} et~al.(2017){N.~Ay}, {J.~Jost}, {H.~V. L{\^e}}, and
  {L.~Schwachh{\"o}fer}]{AyJostLeSchwachhofer2017}
{N.~Ay}, {J.~Jost}, {H.~V. L{\^e}}, and {L.~Schwachh{\"o}fer}.
\newblock \emph{Information Geometry}, 2017.
\newblock Ergebnisse der Mathematik und ihrer Grenzgebiete, Vol.~64, Springer,
  2017.

\bibitem{Fine1982}
A.~Fine.
\newblock Hidden variables, joint probability, and the {B}ell inequalities.
\newblock \emph{Journal of Mathematical Physics}, \textbf{23}(7):1306-1310,
  1982.

\bibitem{LeggettGarg1985}
A.~J. Leggett and A.~Garg.
\newblock Quantum mechanics versus macroscopic realism: is the flux there when nobody looks?
\newblock \emph{Physical Review Letters}, \textbf{54}(9):857-860, 1985.

\bibitem[{M.~Rudelson} and {R.~Vershynin}(2007)]{RudelsonVershynin2007}
{M.~Rudelson} and {R.~Vershynin}.
\newblock Sampling from large matrices: an approach through geometric
  functional analysis, 2007.
\newblock \emph{Journal of the ACM}, \textbf{54}(4), Article~21, 2007.

\bibitem[{G.~Teza} and {A.~L. Stella}(2020)]{TezaStella2020}
{G.~Teza} and {A.~L. Stella}.
\newblock Exact coarse graining preserves entropy production out of
  equilibrium, 2020.
\newblock \emph{Physical Review Letters}, \textbf{125}:110601, 2020.

\bibitem[{J.~G. Kemeny} and {J.~L. Snell}(1960)]{KemenySnell1976}
{J.~G. Kemeny} and {J.~L. Snell}.
\newblock \emph{Finite Markov Chains}, 1960.
\newblock Springer-Verlag, 1976. Originally published by D.~Van Nostrand, 1960.

\bibitem[{P.~Buchholz}(1994)]{Buchholz1994}
{P.~Buchholz}.
\newblock Exact and ordinary lumpability in finite {M}arkov chains, 1994.
\newblock \emph{Journal of Applied Probability}, \textbf{31}:59-75, 1994.

\end{thebibliography}
\end{document}
